% 本文件为账本事件的首次读者呈现：按并行政权、来源与材料限度编排。
\section{元朝治理与区域分化的逐年入口}

本节把同年不同政权的可见行动并置，而不把某一朝廷年号当作唯一时间轴。每一锚点都回到账本的核心材料、前后压力与可说范围；材料稀薄处保留为年序坐标，不补写未见的情节。

\begin{sourcenote}{行省条格能说明什么，不能说明什么}
\eventref{YUAN-1281-EVENT-179}{一二八一年对日远征}与行省、站赤、河运的记录可由《元史》、条格、碑刻和区域材料互读。制度语言说明机构与规范，不能单独证明每一路执行一致，更不能替沿海、江南或东北的不同经验作证。
\end{sourcenote}

\subsection{1271年（元至元八年；干支、月日待核；公元换算据《元史》本纪年号对照）}

\noindent\eventanchor{YUAN-1271-EVENT-173}{忽必烈定国号元}\textbf{元。}
元的1271年（元至元八年；干支、月日待核；公元换算据《元史》本纪年号对照）记录写到“忽必烈定国号元”。我们可以据此辨认一个可见的选择或压力，而不能把零散的官方记载、碑刻或后出叙事拼成无缝的社会全景。

账本把这一节点放在““忽必烈定国号元”发生在忽必烈政权、征宋和元朝建制的具体压力与机会之中，相关行动者的选择不能脱离此前事件链解释。”与““忽必烈定国号元”为后续政治、边防或资源配置留下条件；其社会影响与实际覆盖范围仍须以异类材料限定。”之间；两者提示前后压力，并不把时间相邻写成单一因果。

本条首先回到《元史》相关本纪；《元朝秘史》、波斯文史籍或《高丽史》对应叙事；《元典章》《通制条格》、多语碑刻与文书（据事核）；《元史》明初速修，纪年与人名有异同；蒙古大汗政权不等同所有汗国，征服、赋役和身份分类须按地区及材料语言分列。 因而蒙古帝国、元朝行省财政、多语文书与区域社会研究不是装饰性的书目，而是检验这一叙述边界的路径。

\subsection{1272年（元至元9年；公元换算据《元史》本纪纪年）}

\noindent\eventanchor{SYD-1272-YUAN-001}{:今遣汝等，水陸並進，布告遐邇，使咸知之}\textbf{元。}
1272年（元至元9年；公元换算据《元史》本纪纪年），“:今遣汝等，水陸並進，布告遐邇，使咸知之”使元与相邻行动者的关系进入可核的外交时点。名号、贡赐和盟约是各方政治语言的一部分；它们能够确认交涉，却不能单凭措辞判定实际隶属或边地生活。

其前提记为“本纪保留灾害或工程的报告地点与朝廷处置；灾情范围需另证。”；其后续是“可与地方文书、河工和环境材料比较，不将单点记录外推为全域因果。”。这是一条待后续材料检验的事件链，而不是对所有行动者意图的替写。

可核的材料入口是《元史》卷008〈本纪第8卷本年条〉。《元史》为明初仓促官修，本纪保存大量实录条目但有编次与纪年问题；行省执行、数字、族称及对外关系须以条格、多语文书和对方史料互校。 就此而言，元代行省、赋役、河工与多语文书研究有助于把年序同区域、制度或物质材料并读。

\noindent\eventanchor{SYD-1272-YUAN-073}{:爰自太祖皇帝以來，與宋使介交通}\textbf{元与南宋。}
1272年（元至元9年；公元换算据《元史》本纪纪年），“:爰自太祖皇帝以來，與宋使介交通”使元与南宋与相邻行动者的关系进入可核的外交时点。名号、贡赐和盟约是各方政治语言的一部分；它们能够确认交涉，却不能单凭措辞判定实际隶属或边地生活。

此前的条件是“本纪年条提供日期和中枢可见行动；前因不据单条补定。”，此后的可见影响是“可回查同年及后续条目，地方影响仍须另证。”。这里的“影响”须按地点、机构和材料种类分别衡量。

这一叙述依据《元史》卷008〈本纪第8卷本年条〉，但《元史》为明初仓促官修，本纪保存大量实录条目但有编次与纪年问题；行省执行、数字、族称及对外关系须以条格、多语文书和对方史料互校。。研究元代行省、赋役、河工与多语文书研究可以补问材料未直接说出的空间与社会差异。

\noindent\eventanchor{SYD-1272-YUAN-145}{敕杖合丹，斥無入宿衞，謫往西川效死軍中，餘定罪有差}\textbf{元。}
在1272年（元至元9年；公元换算据《元史》本纪纪年）的可见材料中，元留下“敕杖合丹，斥無入宿衞，謫往西川效死軍中，餘定罪有差”。这一动作把具体人群、机构或地方条件带进年序，但一项记录不能替未被记载者概括整体反应。

账本把这一节点放在“本纪年条提供日期和中枢可见行动；前因不据单条补定。”与“可回查同年及后续条目，地方影响仍须另证。”之间；两者提示前后压力，并不把时间相邻写成单一因果。

本条首先回到《元史》卷008〈本纪第8卷本年条〉；《元史》为明初仓促官修，本纪保存大量实录条目但有编次与纪年问题；行省执行、数字、族称及对外关系须以条格、多语文书和对方史料互校。 因而元代行省、赋役、河工与多语文书研究不是装饰性的书目，而是检验这一叙述边界的路径。

\noindent\eventanchor{SYD-1272-YUAN-217}{從其請，仍降璽書，遣使諭江陵府制置司及高達已下官吏軍民}\textbf{元。}
1272年（元至元9年；公元换算据《元史》本纪纪年），元的行动落在“從其請，仍降璽書，遣使諭江陵府制置司及高達已下官吏軍民”所指的战事与通道上。可见的不是抽象的推进：攻守双方必须让人员、消息与补给穿过具体的城防、道路或水路；账本所列的结果只说明这一节点改变了下一步选择。

此前的条件是“本纪年条记录政令或机构处置；实施背景须参照制度材料。”，此后的可见影响是“可在同年及后续政令中追踪；条文不单独证明地方执行。”。这里的“影响”须按地点、机构和材料种类分别衡量。

这一叙述依据《元史》卷008〈本纪第8卷本年条〉，但《元史》为明初仓促官修，本纪保存大量实录条目但有编次与纪年问题；行省执行、数字、族称及对外关系须以条格、多语文书和对方史料互校。。研究元代行省、赋役、河工与多语文书研究可以补问材料未直接说出的空间与社会差异。

\noindent\eventanchor{SYD-1272-YUAN-289}{帝是其言，復改淮西行中書省為行樞密院}\textbf{元。}
在1272年（元至元9年；公元换算据《元史》本纪纪年）的可见材料中，元留下“帝是其言，復改淮西行中書省為行樞密院”。这一动作把具体人群、机构或地方条件带进年序，但一项记录不能替未被记载者概括整体反应。

从账本的前后项看，“本纪年条记录政令或机构处置；实施背景须参照制度材料。”提供了背景压力；“可在同年及后续政令中追踪；条文不单独证明地方执行。”标出后来需要追踪的变化，二者之间仍存在未被材料填满的环节。

核心线索为《元史》卷008〈本纪第8卷本年条〉。《元史》为明初仓促官修，本纪保存大量实录条目但有编次与纪年问题；行省执行、数字、族称及对外关系须以条格、多语文书和对方史料互校。 现代研究可从元代行省、赋役、河工与多语文书研究继续追问。

\subsection{1273年（元至元10年；公元换算据《元史》本纪纪年）}

\noindent\eventanchor{SYD-1273-YUAN-002}{安南使者還，言陳光昞受詔不拜}\textbf{元与安南。}
元与安南于1273年（元至元10年；公元换算据《元史》本纪纪年）处理“安南使者還，言陳光昞受詔不拜”时，面对的是道路、翻译、礼仪和资源交换共同构成的关系网。书面称谓可见，地方执行和持续性则仍须由后续材料检验。

其前提记为“本纪年条记录政令或机构处置；实施背景须参照制度材料。”；其后续是“可在同年及后续政令中追踪；条文不单独证明地方执行。”。这是一条待后续材料检验的事件链，而不是对所有行动者意图的替写。

可核的材料入口是《元史》卷008〈至元十年〉；《安南志略》及当地碑刻材料（互校）。《元史》为明初仓促官修，本纪保存大量实录条目但有编次与纪年问题；行省执行、数字、族称及对外关系须以条格、多语文书和对方史料互校。 就此而言，元代行省、赋役、河工与多语文书研究有助于把年序同区域、制度或物质材料并读。

\noindent\eventanchor{SYD-1273-YUAN-074}{帝以農事有益，詔勿禁}\textbf{元。}
元的1273年（元至元10年；公元换算据《元史》本纪纪年）记录写到“帝以農事有益，詔勿禁”。我们可以据此辨认一个可见的选择或压力，而不能把零散的官方记载、碑刻或后出叙事拼成无缝的社会全景。

此前的条件是“本纪年条记录政令或机构处置；实施背景须参照制度材料。”，此后的可见影响是“可在同年及后续政令中追踪；条文不单独证明地方执行。”。这里的“影响”须按地点、机构和材料种类分别衡量。

这一叙述依据《元史》卷008〈至元十年〉，但《元史》为明初仓促官修，本纪保存大量实录条目但有编次与纪年问题；行省执行、数字、族称及对外关系须以条格、多语文书和对方史料互校。。研究元代行省、赋役、河工与多语文书研究可以补问材料未直接说出的空间与社会差异。

\noindent\eventanchor{SYD-1273-YUAN-146}{帝曰：「汝言雖是，若坐視宋人戍之，罪亦不免也}\textbf{元与南宋。}
元与南宋的1273年（元至元10年；公元换算据《元史》本纪纪年）记录写到“帝曰：「汝言雖是，若坐視宋人戍之，罪亦不免也”。我们可以据此辨认一个可见的选择或压力，而不能把零散的官方记载、碑刻或后出叙事拼成无缝的社会全景。

账本把这一节点放在“本纪年条提供日期和中枢可见行动；前因不据单条补定。”与“可回查同年及后续条目，地方影响仍须另证。”之间；两者提示前后压力，并不把时间相邻写成单一因果。

本条首先回到《元史》卷008〈至元十年〉；《元史》为明初仓促官修，本纪保存大量实录条目但有编次与纪年问题；行省执行、数字、族称及对外关系须以条格、多语文书和对方史料互校。 因而元代行省、赋役、河工与多语文书研究不是装饰性的书目，而是检验这一叙述边界的路径。

\noindent\eventanchor{SYD-1273-YUAN-218}{庚寅，河南水，發粟賑民饑，仍免今年田租}\textbf{元。}
元在1273年（元至元10年；公元换算据《元史》本纪纪年）留下的这条记录是“庚寅，河南水，發粟賑民饑，仍免今年田租”。它照亮资源如何被命名和调度，也暴露材料的盲点：数字、额度或禁令若无地方文书和物质材料互证，不能被写成全境同速的现实。

此前的条件是“本纪所记财用、赈济或征发属于特定报告场景。”，此后的可见影响是“可与食货志、文书和物质材料比较，不外推为全境总量。”。这里的“影响”须按地点、机构和材料种类分别衡量。

这一叙述依据《元史》卷008〈至元十年〉，但《元史》为明初仓促官修，本纪保存大量实录条目但有编次与纪年问题；行省执行、数字、族称及对外关系须以条格、多语文书和对方史料互校。。研究元代行省、赋役、河工与多语文书研究可以补问材料未直接说出的空间与社会差异。

\noindent\eventanchor{SYD-1273-YUAN-290}{乙亥，詔：「免民代輸簽軍戶絲銀，及伐木夫戶賦稅}\textbf{元。}
在1273年（元至元10年；公元换算据《元史》本纪纪年）的材料中，元面对的是乙亥，詔：「免民代輸簽軍戶絲銀，及伐木夫戶賦稅。这里应先看行动所依赖的关隘、城镇、河道或粮道，再讨论政治后果；军队抵达某处，并不等于它已能稳定征收、安置或控制周边。

此前的条件是“本纪年条记录政令或机构处置；实施背景须参照制度材料。”，此后的可见影响是“可在同年及后续政令中追踪；条文不单独证明地方执行。”。这里的“影响”须按地点、机构和材料种类分别衡量。

这一叙述依据《元史》卷008〈至元十年〉，但《元史》为明初仓促官修，本纪保存大量实录条目但有编次与纪年问题；行省执行、数字、族称及对外关系须以条格、多语文书和对方史料互校。。研究元代行省、赋役、河工与多语文书研究可以补问材料未直接说出的空间与社会差异。

\subsection{1273年}

\noindent\eventanchor{YUAN-1273-EVENT-174}{襄阳、樊城陷落}\textbf{元与南宋。}
在1273年的材料中，元与南宋面对的是襄阳、樊城陷落。这里应先看行动所依赖的关隘、城镇、河道或粮道，再讨论政治后果；军队抵达某处，并不等于它已能稳定征收、安置或控制周边。

其前提记为“忽必烈政权、征宋和元朝建制中既有的联盟、交通与补给压力，为“襄阳、樊城陷落”提供了直接的军事或动员背景。”；其后续是““襄阳、樊城陷落”改变了相关城镇、道路或政权间的军事选择；伤亡、俘获和控制范围仅可按各方材料分别确认。”。这是一条待后续材料检验的事件链，而不是对所有行动者意图的替写。

可核的材料入口是《元史》相关本纪；《元朝秘史》、波斯文史籍或《高丽史》对应叙事；《元典章》《通制条格》、多语碑刻与文书（据事核）。《元史》明初速修，纪年与人名有异同；蒙古大汗政权不等同所有汗国，征服、赋役和身份分类须按地区及材料语言分列。 就此而言，蒙古帝国、元朝行省财政、多语文书与区域社会研究有助于把年序同区域、制度或物质材料并读。

\subsection{1274年（元至元11年；公元换算据《元史》本纪纪年）}

\noindent\eventanchor{SYD-1274-YUAN-003}{丁亥，罷各路洞冶總管府，以轉運司兼領}\textbf{元。}
“丁亥，罷各路洞冶總管府，以轉運司兼領”是元在1274年（元至元11年；公元换算据《元史》本纪纪年）可追查的节点。它提示命令、环境、宗教、迁徙或地方组织如何进入政治过程；影响的范围和速度仍要避免由一条材料外推。

其前提记为“本纪年条记录政令或机构处置；实施背景须参照制度材料。”；其后续是“可在同年及后续政令中追踪；条文不单独证明地方执行。”。这是一条待后续材料检验的事件链，而不是对所有行动者意图的替写。

可核的材料入口是《元史》卷007〈本纪第7卷本年条〉。《元史》为明初仓促官修，本纪保存大量实录条目但有编次与纪年问题；行省执行、数字、族称及对外关系须以条格、多语文书和对方史料互校。 就此而言，元代行省、赋役、河工与多语文书研究有助于把年序同区域、制度或物质材料并读。

\noindent\eventanchor{SYD-1274-YUAN-075}{改御史臺典事為都事}\textbf{元。}
“改御史臺典事為都事”是元在1274年（元至元11年；公元换算据《元史》本纪纪年）可追查的节点。它提示命令、环境、宗教、迁徙或地方组织如何进入政治过程；影响的范围和速度仍要避免由一条材料外推。

此前的条件是“本纪年条记录政令或机构处置；实施背景须参照制度材料。”，此后的可见影响是“可在同年及后续政令中追踪；条文不单独证明地方执行。”。这里的“影响”须按地点、机构和材料种类分别衡量。

这一叙述依据《元史》卷007〈本纪第7卷本年条〉，但《元史》为明初仓促官修，本纪保存大量实录条目但有编次与纪年问题；行省执行、数字、族称及对外关系须以条格、多语文书和对方史料互校。。研究元代行省、赋役、河工与多语文书研究可以补问材料未直接说出的空间与社会差异。

\noindent\eventanchor{SYD-1274-YUAN-147}{命陝西等路宣撫使趙良弼為祕書監，充國信使，使日本}\textbf{元。}
在“命陝西等路宣撫使趙良弼為祕書監，充國信使，使日本”中，元的选择经由使节、礼物、边市或协商被记录下来。跨境文本的观察位置不同，因而应把可确认的往来与对动机、服从程度的推断分开。

账本把这一节点放在“本纪年条记录政令或机构处置；实施背景须参照制度材料。”与“可在同年及后续政令中追踪；条文不单独证明地方执行。”之间；两者提示前后压力，并不把时间相邻写成单一因果。

本条首先回到《元史》卷007〈本纪第7卷本年条〉；《元史》为明初仓促官修，本纪保存大量实录条目但有编次与纪年问题；行省执行、数字、族称及对外关系须以条格、多语文书和对方史料互校。 因而元代行省、赋役、河工与多语文书研究不是装饰性的书目，而是检验这一叙述边界的路径。

\noindent\eventanchor{SYD-1274-YUAN-219}{世子愖奏乞隨朝及尚主，不許，命隨其父還國}\textbf{元。}
在“世子愖奏乞隨朝及尚主，不許，命隨其父還國”中，元的选择经由使节、礼物、边市或协商被记录下来。跨境文本的观察位置不同，因而应把可确认的往来与对动机、服从程度的推断分开。

此前的条件是“本纪年条记录政令或机构处置；实施背景须参照制度材料。”，此后的可见影响是“可在同年及后续政令中追踪；条文不单独证明地方执行。”。这里的“影响”须按地点、机构和材料种类分别衡量。

这一叙述依据《元史》卷007〈本纪第7卷本年条〉，但《元史》为明初仓促官修，本纪保存大量实录条目但有编次与纪年问题；行省执行、数字、族称及对外关系须以条格、多语文书和对方史料互校。。研究元代行省、赋役、河工与多语文书研究可以补问材料未直接说出的空间与社会差异。

\noindent\eventanchor{SYD-1274-YUAN-291}{又詔令國王頭輦哥等舉軍入高麗舊京，以脫朵兒、焦天翼為其國達魯花赤，護送禃還國，仍下詔：「林衍廢立，罪不可赦}\textbf{元与高丽。}
元与高丽的“又詔令國王頭輦哥等舉軍入高麗舊京，以脫朵兒、焦天翼為其國達魯花赤，護送禃還國，仍下詔：「林衍廢立，罪不可赦”并非只改换一个统治者。任官、赏赐、军权与信息通路要在随后重新接合；材料较能确认先后，较难替当事人写出唯一动机。

此前的条件是“本纪年条记录政令或机构处置；实施背景须参照制度材料。”，此后的可见影响是“可在同年及后续政令中追踪；条文不单独证明地方执行。”。这里的“影响”须按地点、机构和材料种类分别衡量。

这一叙述依据《元史》卷007〈本纪第7卷本年条〉；《高丽史》同年条（互校），但《元史》为明初仓促官修，本纪保存大量实录条目但有编次与纪年问题；行省执行、数字、族称及对外关系须以条格、多语文书和对方史料互校。。研究元代行省、赋役、河工与多语文书研究可以补问材料未直接说出的空间与社会差异。

\subsection{1274年}

\noindent\eventanchor{YUAN-1274-EVENT-175}{元军分路南下进攻南宋}\textbf{元与南宋。}
在1274年的材料中，元与南宋面对的是元军分路南下进攻南宋。这里应先看行动所依赖的关隘、城镇、河道或粮道，再讨论政治后果；军队抵达某处，并不等于它已能稳定征收、安置或控制周边。

从账本的前后项看，“忽必烈政权、征宋和元朝建制中既有的联盟、交通与补给压力，为“元军分路南下进攻南宋”提供了直接的军事或动员背景。”提供了背景压力；““元军分路南下进攻南宋”改变了相关城镇、道路或政权间的军事选择；伤亡、俘获和控制范围仅可按各方材料分别确认。”标出后来需要追踪的变化，二者之间仍存在未被材料填满的环节。

核心线索为《元史》相关本纪；《元朝秘史》、波斯文史籍或《高丽史》对应叙事；《元典章》《通制条格》、多语碑刻与文书（据事核）。《元史》明初速修，纪年与人名有异同；蒙古大汗政权不等同所有汗国，征服、赋役和身份分类须按地区及材料语言分列。 现代研究可从蒙古帝国、元朝行省财政、多语文书与区域社会研究继续追问。

\subsection{1275年（元至元12年；公元换算据《元史》本纪纪年）}

\noindent\eventanchor{SYD-1275-YUAN-004}{:誕膺景命，奄四海以宅尊}\textbf{元。}
在1275年（元至元12年；公元换算据《元史》本纪纪年）的可见材料中，元留下“:誕膺景命，奄四海以宅尊”。这一动作把具体人群、机构或地方条件带进年序，但一项记录不能替未被记载者概括整体反应。

其前提记为“本纪年条记录政令或机构处置；实施背景须参照制度材料。”；其后续是“可在同年及后续政令中追踪；条文不单独证明地方执行。”。这是一条待后续材料检验的事件链，而不是对所有行动者意图的替写。

可核的材料入口是《元史》卷007〈本纪第7卷本年条〉。《元史》为明初仓促官修，本纪保存大量实录条目但有编次与纪年问题；行省执行、数字、族称及对外关系须以条格、多语文书和对方史料互校。 就此而言，元代行省、赋役、河工与多语文书研究有助于把年序同区域、制度或物质材料并读。

\noindent\eventanchor{SYD-1275-YUAN-076}{:我太祖聖武皇帝，握乾符而起朔土，以神武而膺帝圖，四震天聲，大恢土宇，輿圖之廣，歷古所無}\textbf{元。}
在1275年（元至元12年；公元换算据《元史》本纪纪年）的可见材料中，元留下“:我太祖聖武皇帝，握乾符而起朔土，以神武而膺帝圖，四震天聲，大恢土宇，輿圖之廣，歷古所無”。这一动作把具体人群、机构或地方条件带进年序，但一项记录不能替未被记载者概括整体反应。

此前的条件是“本纪保留灾害或工程的报告地点与朝廷处置；灾情范围需另证。”，此后的可见影响是“可与地方文书、河工和环境材料比较，不将单点记录外推为全域因果。”。这里的“影响”须按地点、机构和材料种类分别衡量。

这一叙述依据《元史》卷007〈本纪第7卷本年条〉，但《元史》为明初仓促官修，本纪保存大量实录条目但有编次与纪年问题；行省执行、数字、族称及对外关系须以条格、多语文书和对方史料互校。。研究元代行省、赋役、河工与多语文书研究可以补问材料未直接说出的空间与社会差异。

\noindent\eventanchor{SYD-1275-YUAN-148}{丙子，改山東、河間、陝西三路鹽課都轉運司為都轉運鹽使司}\textbf{元。}
1275年（元至元12年；公元换算据《元史》本纪纪年），“丙子，改山東、河間、陝西三路鹽課都轉運司為都轉運鹽使司”使元与相邻行动者的关系进入可核的外交时点。名号、贡赐和盟约是各方政治语言的一部分；它们能够确认交涉，却不能单凭措辞判定实际隶属或边地生活。

账本把这一节点放在“本纪年条记录政令或机构处置；实施背景须参照制度材料。”与“可在同年及后续政令中追踪；条文不单独证明地方执行。”之间；两者提示前后压力，并不把时间相邻写成单一因果。

本条首先回到《元史》卷007〈本纪第7卷本年条〉；《元史》为明初仓促官修，本纪保存大量实录条目但有编次与纪年问题；行省执行、数字、族称及对外关系须以条格、多语文书和对方史料互校。 因而元代行省、赋役、河工与多语文书研究不是装饰性的书目，而是检验这一叙述边界的路径。

\noindent\eventanchor{SYD-1275-YUAN-220}{高麗國王王禃遣使奉表，為世子愖請昏}\textbf{元与高丽。}
1275年（元至元12年；公元换算据《元史》本纪纪年），“高麗國王王禃遣使奉表，為世子愖請昏”使元与高丽与相邻行动者的关系进入可核的外交时点。名号、贡赐和盟约是各方政治语言的一部分；它们能够确认交涉，却不能单凭措辞判定实际隶属或边地生活。

其前提记为“本纪记录使节、贡赐或往来，不等于稳定统属关系。”；其后续是“可与对方编年和交通、文书材料核对其后续落实。”。这是一条待后续材料检验的事件链，而不是对所有行动者意图的替写。

可核的材料入口是《元史》卷007〈本纪第7卷本年条〉；《高丽史》同年条（互校）。《元史》为明初仓促官修，本纪保存大量实录条目但有编次与纪年问题；行省执行、数字、族称及对外关系须以条格、多语文书和对方史料互校。 就此而言，元代行省、赋役、河工与多语文书研究有助于把年序同区域、制度或物质材料并读。

\noindent\eventanchor{SYD-1275-YUAN-292}{己亥，詔招諭宋襄陽守臣呂文煥}\textbf{元与南宋。}
1275年（元至元12年；公元换算据《元史》本纪纪年），元与南宋的行动落在“己亥，詔招諭宋襄陽守臣呂文煥”所指的战事与通道上。可见的不是抽象的推进：攻守双方必须让人员、消息与补给穿过具体的城防、道路或水路；账本所列的结果只说明这一节点改变了下一步选择。

此前的条件是“本纪年条记录政令或机构处置；实施背景须参照制度材料。”，此后的可见影响是“可在同年及后续政令中追踪；条文不单独证明地方执行。”。这里的“影响”须按地点、机构和材料种类分别衡量。

这一叙述依据《元史》卷007〈本纪第7卷本年条〉，但《元史》为明初仓促官修，本纪保存大量实录条目但有编次与纪年问题；行省执行、数字、族称及对外关系须以条格、多语文书和对方史料互校。。研究元代行省、赋役、河工与多语文书研究可以补问材料未直接说出的空间与社会差异。

\subsection{1276年（元至元13年；公元换算据《元史》本纪纪年）}

\noindent\eventanchor{SYD-1276-YUAN-005}{敕令整為書復之，賞整，使還軍中，誅永寧僧及其黨友}\textbf{元。}
元于1276年（元至元13年；公元换算据《元史》本纪纪年）处理“敕令整為書復之，賞整，使還軍中，誅永寧僧及其黨友”时，面对的是道路、翻译、礼仪和资源交换共同构成的关系网。书面称谓可见，地方执行和持续性则仍须由后续材料检验。

其前提记为“本纪年条记录政令或机构处置；实施背景须参照制度材料。”；其后续是“可在同年及后续政令中追踪；条文不单独证明地方执行。”。这是一条待后续材料检验的事件链，而不是对所有行动者意图的替写。

可核的材料入口是《元史》卷007〈本纪第7卷本年条〉。《元史》为明初仓促官修，本纪保存大量实录条目但有编次与纪年问题；行省执行、数字、族称及对外关系须以条格、多语文书和对方史料互校。 就此而言，元代行省、赋役、河工与多语文书研究有助于把年序同区域、制度或物质材料并读。

\noindent\eventanchor{SYD-1276-YUAN-077}{敕燕王遣使持香旛，祠岳瀆、后土、五臺興國寺}\textbf{元。}
元于1276年（元至元13年；公元换算据《元史》本纪纪年）处理“敕燕王遣使持香旛，祠岳瀆、后土、五臺興國寺”时，面对的是道路、翻译、礼仪和资源交换共同构成的关系网。书面称谓可见，地方执行和持续性则仍须由后续材料检验。

此前的条件是“本纪记录使节、贡赐或往来，不等于稳定统属关系。”，此后的可见影响是“可与对方编年和交通、文书材料核对其后续落实。”。这里的“影响”须按地点、机构和材料种类分别衡量。

这一叙述依据《元史》卷007〈本纪第7卷本年条〉，但《元史》为明初仓促官修，本纪保存大量实录条目但有编次与纪年问题；行省执行、数字、族称及对外关系须以条格、多语文书和对方史料互校。。研究元代行省、赋役、河工与多语文书研究可以补问材料未直接说出的空间与社会差异。

\noindent\eventanchor{SYD-1276-YUAN-149}{癸酉，同僉河南省事崔斌訟右丞阿里妄奏軍數二萬，敕杖而罷之}\textbf{元。}
元的1276年（元至元13年；公元换算据《元史》本纪纪年）记录写到“癸酉，同僉河南省事崔斌訟右丞阿里妄奏軍數二萬，敕杖而罷之”。我们可以据此辨认一个可见的选择或压力，而不能把零散的官方记载、碑刻或后出叙事拼成无缝的社会全景。

账本把这一节点放在“本纪年条记录政令或机构处置；实施背景须参照制度材料。”与“可在同年及后续政令中追踪；条文不单独证明地方执行。”之间；两者提示前后压力，并不把时间相邻写成单一因果。

本条首先回到《元史》卷007〈本纪第7卷本年条〉；《元史》为明初仓促官修，本纪保存大量实录条目但有编次与纪年问题；行省执行、数字、族称及对外关系须以条格、多语文书和对方史料互校。 因而元代行省、赋役、河工与多语文书研究不是装饰性的书目，而是检验这一叙述边界的路径。

\noindent\eventanchor{SYD-1276-YUAN-221}{壬辰，高麗國王王禃遣其臣齊安侯王淑來賀改國號}\textbf{元与高丽。}
元与高丽的1276年（元至元13年；公元换算据《元史》本纪纪年）记录写到“壬辰，高麗國王王禃遣其臣齊安侯王淑來賀改國號”。我们可以据此辨认一个可见的选择或压力，而不能把零散的官方记载、碑刻或后出叙事拼成无缝的社会全景。

其前提记为“本纪年条记录政令或机构处置；实施背景须参照制度材料。”；其后续是“可在同年及后续政令中追踪；条文不单独证明地方执行。”。这是一条待后续材料检验的事件链，而不是对所有行动者意图的替写。

可核的材料入口是《元史》卷007〈本纪第7卷本年条〉；《高丽史》同年条（互校）。《元史》为明初仓促官修，本纪保存大量实录条目但有编次与纪年问题；行省执行、数字、族称及对外关系须以条格、多语文书和对方史料互校。 就此而言，元代行省、赋役、河工与多语文书研究有助于把年序同区域、制度或物质材料并读。

\noindent\eventanchor{SYD-1276-YUAN-293}{詔自今凡詔令並以蒙古字行，仍遣百官子弟入學}\textbf{元。}
元的1276年（元至元13年；公元换算据《元史》本纪纪年）记录写到“詔自今凡詔令並以蒙古字行，仍遣百官子弟入學”。我们可以据此辨认一个可见的选择或压力，而不能把零散的官方记载、碑刻或后出叙事拼成无缝的社会全景。

此前的条件是“本纪年条记录政令或机构处置；实施背景须参照制度材料。”，此后的可见影响是“可在同年及后续政令中追踪；条文不单独证明地方执行。”。这里的“影响”须按地点、机构和材料种类分别衡量。

这一叙述依据《元史》卷007〈本纪第7卷本年条〉，但《元史》为明初仓促官修，本纪保存大量实录条目但有编次与纪年问题；行省执行、数字、族称及对外关系须以条格、多语文书和对方史料互校。。研究元代行省、赋役、河工与多语文书研究可以补问材料未直接说出的空间与社会差异。

\subsection{1276年（元至元十三年、宋德祐二年；干支、月日待核；公元换算据双方本纪年号对照）}

\noindent\eventanchor{YUAN-1276-EVENT-176}{元军进入临安，宋恭帝降元}\textbf{元与南宋。}
在1276年（元至元十三年、宋德祐二年；干支、月日待核；公元换算据双方本纪年号对照）的材料中，元与南宋面对的是元军进入临安，宋恭帝降元。这里应先看行动所依赖的关隘、城镇、河道或粮道，再讨论政治后果；军队抵达某处，并不等于它已能稳定征收、安置或控制周边。

账本把这一节点放在““元军进入临安，宋恭帝降元”发生在忽必烈政权、征宋和元朝建制的具体压力与机会之中，相关行动者的选择不能脱离此前事件链解释。”与““元军进入临安，宋恭帝降元”为后续政治、边防或资源配置留下条件；其社会影响与实际覆盖范围仍须以异类材料限定。”之间；两者提示前后压力，并不把时间相邻写成单一因果。

本条首先回到《元史》相关本纪；《元朝秘史》、波斯文史籍或《高丽史》对应叙事；《元典章》《通制条格》、多语碑刻与文书（据事核）；《元史》明初速修，纪年与人名有异同；蒙古大汗政权不等同所有汗国，征服、赋役和身份分类须按地区及材料语言分列。 因而蒙古帝国、元朝行省财政、多语文书与区域社会研究不是装饰性的书目，而是检验这一叙述边界的路径。

\subsection{1277年}

\noindent\eventanchor{YUAN-1277-EVENT-177}{元在江南建立行省、路府与财赋接收安排}\textbf{元。}
1277年的“元在江南建立行省、路府与财赋接收安排”将元的财政或生产安排推到可见处。制度文本能说明官府打算如何收取、转运或调节，却不能直接等同于不同地区的实收、流通或家庭负担。

此前的条件是“忽必烈政权、征宋和元朝建制对中枢资源、交通或行政协同的要求，推动了“元在江南建立行省、路府与财赋接收安排”这一制度或空间安排。”，此后的可见影响是““元在江南建立行省、路府与财赋接收安排”提供新的治理框架或资源通道，但颁行、复申与不同地区的实际执行须分开核验。”。这里的“影响”须按地点、机构和材料种类分别衡量。

这一叙述依据《元史》相关本纪；《元朝秘史》、波斯文史籍或《高丽史》对应叙事；《元典章》《通制条格》、多语碑刻与文书（据事核），但《元史》明初速修，纪年与人名有异同；蒙古大汗政权不等同所有汗国，征服、赋役和身份分类须按地区及材料语言分列。。研究蒙古帝国、元朝行省财政、多语文书与区域社会研究可以补问材料未直接说出的空间与社会差异。

\subsection{1278年（元至元15年；公元换算据《元史》本纪纪年）}

\noindent\eventanchor{SYD-1278-YUAN-006}{近侍劉鐵木兒因言：「阿里海牙屬吏張鼎，今亦參知攻事}\textbf{元。}
1278年（元至元15年；公元换算据《元史》本纪纪年），元的行动落在“近侍劉鐵木兒因言：「阿里海牙屬吏張鼎，今亦參知攻事”所指的战事与通道上。可见的不是抽象的推进：攻守双方必须让人员、消息与补给穿过具体的城防、道路或水路；账本所列的结果只说明这一节点改变了下一步选择。

此前的条件是“本纪从朝廷视角记录军政行动；出兵与战果需分辨。”，此后的可见影响是“后续路线、控制与损失应与对方记录和遗址材料复核。”。这里的“影响”须按地点、机构和材料种类分别衡量。

这一叙述依据《元史》卷010〈至元十五年〉，但《元史》为明初仓促官修，本纪保存大量实录条目但有编次与纪年问题；行省执行、数字、族称及对外关系须以条格、多语文书和对方史料互校。。研究元代行省、赋役、河工与多语文书研究可以补问材料未直接说出的空间与社会差异。

\noindent\eventanchor{SYD-1278-YUAN-078}{申嚴無籍軍虜掠及傭奴代軍之禁}\textbf{元。}
在1278年（元至元15年；公元换算据《元史》本纪纪年）的可见材料中，元留下“申嚴無籍軍虜掠及傭奴代軍之禁”。这一动作把具体人群、机构或地方条件带进年序，但一项记录不能替未被记载者概括整体反应。

从账本的前后项看，“本纪年条记录政令或机构处置；实施背景须参照制度材料。”提供了背景压力；“可在同年及后续政令中追踪；条文不单独证明地方执行。”标出后来需要追踪的变化，二者之间仍存在未被材料填满的环节。

核心线索为《元史》卷010〈至元十五年〉。《元史》为明初仓促官修，本纪保存大量实录条目但有编次与纪年问题；行省执行、数字、族称及对外关系须以条格、多语文书和对方史料互校。 现代研究可从元代行省、赋役、河工与多语文书研究继续追问。

\noindent\eventanchor{SYD-1278-YUAN-150}{戊申，從阿合馬請，自今御史臺非白于省，毋擅召倉庫吏，亦毋究錢穀數，及集議中書不至者罪之}\textbf{元。}
在1278年（元至元15年；公元换算据《元史》本纪纪年）的可见材料中，元留下“戊申，從阿合馬請，自今御史臺非白于省，毋擅召倉庫吏，亦毋究錢穀數，及集議中書不至者罪之”。这一动作把具体人群、机构或地方条件带进年序，但一项记录不能替未被记载者概括整体反应。

账本把这一节点放在“本纪年条记录政令或机构处置；实施背景须参照制度材料。”与“可在同年及后续政令中追踪；条文不单独证明地方执行。”之间；两者提示前后压力，并不把时间相邻写成单一因果。

本条首先回到《元史》卷010〈至元十五年〉；《元史》为明初仓促官修，本纪保存大量实录条目但有编次与纪年问题；行省执行、数字、族称及对外关系须以条格、多语文书和对方史料互校。 因而元代行省、赋役、河工与多语文书研究不是装饰性的书目，而是检验这一叙述边界的路径。

\noindent\eventanchor{SYD-1278-YUAN-222}{乙亥，敕省、院、臺諸司應聞奏事，必由起居注}\textbf{元。}
在1278年（元至元15年；公元换算据《元史》本纪纪年）的可见材料中，元留下“乙亥，敕省、院、臺諸司應聞奏事，必由起居注”。这一动作把具体人群、机构或地方条件带进年序，但一项记录不能替未被记载者概括整体反应。

此前的条件是“本纪年条记录政令或机构处置；实施背景须参照制度材料。”，此后的可见影响是“可在同年及后续政令中追踪；条文不单独证明地方执行。”。这里的“影响”须按地点、机构和材料种类分别衡量。

这一叙述依据《元史》卷010〈至元十五年〉，但《元史》为明初仓促官修，本纪保存大量实录条目但有编次与纪年问题；行省执行、数字、族称及对外关系须以条格、多语文书和对方史料互校。。研究元代行省、赋役、河工与多语文书研究可以补问材料未直接说出的空间与社会差异。

\noindent\eventanchor{SYD-1278-YUAN-294}{詔以也速海牙總制之}\textbf{元。}
在1278年（元至元15年；公元换算据《元史》本纪纪年）的可见材料中，元留下“詔以也速海牙總制之”。这一动作把具体人群、机构或地方条件带进年序，但一项记录不能替未被记载者概括整体反应。

从账本的前后项看，“本纪年条记录政令或机构处置；实施背景须参照制度材料。”提供了背景压力；“可在同年及后续政令中追踪；条文不单独证明地方执行。”标出后来需要追踪的变化，二者之间仍存在未被材料填满的环节。

核心线索为《元史》卷010〈至元十五年〉。《元史》为明初仓促官修，本纪保存大量实录条目但有编次与纪年问题；行省执行、数字、族称及对外关系须以条格、多语文书和对方史料互校。 现代研究可从元代行省、赋役、河工与多语文书研究继续追问。

\subsection{1279年（元至元16年；公元换算据《元史》本纪纪年）}

\noindent\eventanchor{SYD-1279-YUAN-007}{西川軍官父死子繼，勤勞四十年，乞顯加爵秩}\textbf{元。}
元的1279年（元至元16年；公元换算据《元史》本纪纪年）记录写到“西川軍官父死子繼，勤勞四十年，乞顯加爵秩”。我们可以据此辨认一个可见的选择或压力，而不能把零散的官方记载、碑刻或后出叙事拼成无缝的社会全景。

此前的条件是“本纪年条提供日期和中枢可见行动；前因不据单条补定。”，此后的可见影响是“可回查同年及后续条目，地方影响仍须另证。”。这里的“影响”须按地点、机构和材料种类分别衡量。

这一叙述依据《元史》卷010〈至元十六年〉，但《元史》为明初仓促官修，本纪保存大量实录条目但有编次与纪年问题；行省执行、数字、族称及对外关系须以条格、多语文书和对方史料互校。。研究元代行省、赋役、河工与多语文书研究可以补问材料未直接说出的空间与社会差异。

\noindent\eventanchor{SYD-1279-YUAN-079}{甲戌，給要束合所領工匠牛二千，就令運米二千石供軍}\textbf{元。}
元的1279年（元至元16年；公元换算据《元史》本纪纪年）记录写到“甲戌，給要束合所領工匠牛二千，就令運米二千石供軍”。我们可以据此辨认一个可见的选择或压力，而不能把零散的官方记载、碑刻或后出叙事拼成无缝的社会全景。

从账本的前后项看，“本纪所记财用、赈济或征发属于特定报告场景。”提供了背景压力；“可与食货志、文书和物质材料比较，不外推为全境总量。”标出后来需要追踪的变化，二者之间仍存在未被材料填满的环节。

核心线索为《元史》卷010〈至元十六年〉。《元史》为明初仓促官修，本纪保存大量实录条目但有编次与纪年问题；行省执行、数字、族称及对外关系须以条格、多语文书和对方史料互校。 现代研究可从元代行省、赋役、河工与多语文书研究继续追问。

\noindent\eventanchor{SYD-1279-YUAN-151}{明日，下詔皇太子燕王參決朝政，凡中書省、樞密院、御史臺及百司之事，皆先啟後聞}\textbf{元。}
元于1279年（元至元16年；公元换算据《元史》本纪纪年）处理“明日，下詔皇太子燕王參決朝政，凡中書省、樞密院、御史臺及百司之事，皆先啟後聞”时，面对的是道路、翻译、礼仪和资源交换共同构成的关系网。书面称谓可见，地方执行和持续性则仍须由后续材料检验。

账本把这一节点放在“本纪年条记录政令或机构处置；实施背景须参照制度材料。”与“可在同年及后续政令中追踪；条文不单独证明地方执行。”之间；两者提示前后压力，并不把时间相邻写成单一因果。

本条首先回到《元史》卷010〈至元十六年〉；《元史》为明初仓促官修，本纪保存大量实录条目但有编次与纪年问题；行省执行、数字、族称及对外关系须以条格、多语文书和对方史料互校。 因而元代行省、赋役、河工与多语文书研究不是装饰性的书目，而是检验这一叙述边界的路径。

\noindent\eventanchor{SYD-1279-YUAN-223}{女直、水達達軍不出征者，令隸民籍輸賦}\textbf{元。}
元的1279年（元至元16年；公元换算据《元史》本纪纪年）记录写到“女直、水達達軍不出征者，令隸民籍輸賦”。我们可以据此辨认一个可见的选择或压力，而不能把零散的官方记载、碑刻或后出叙事拼成无缝的社会全景。

此前的条件是“本纪所记财用、赈济或征发属于特定报告场景。”，此后的可见影响是“可与食货志、文书和物质材料比较，不外推为全境总量。”。这里的“影响”须按地点、机构和材料种类分别衡量。

这一叙述依据《元史》卷010〈至元十六年〉，但《元史》为明初仓促官修，本纪保存大量实录条目但有编次与纪年问题；行省执行、数字、族称及对外关系须以条格、多语文书和对方史料互校。。研究元代行省、赋役、河工与多语文书研究可以补问材料未直接说出的空间与社会差异。

\noindent\eventanchor{SYD-1279-YUAN-295}{數斬唐兀帶、武伯不花，餘減死論，以所掠者還其民}\textbf{元。}
元的1279年（元至元16年；公元换算据《元史》本纪纪年）记录写到“數斬唐兀帶、武伯不花，餘減死論，以所掠者還其民”。我们可以据此辨认一个可见的选择或压力，而不能把零散的官方记载、碑刻或后出叙事拼成无缝的社会全景。

从账本的前后项看，“本纪年条提供日期和中枢可见行动；前因不据单条补定。”提供了背景压力；“可回查同年及后续条目，地方影响仍须另证。”标出后来需要追踪的变化，二者之间仍存在未被材料填满的环节。

核心线索为《元史》卷010〈至元十六年〉。《元史》为明初仓促官修，本纪保存大量实录条目但有编次与纪年问题；行省执行、数字、族称及对外关系须以条格、多语文书和对方史料互校。 现代研究可从元代行省、赋役、河工与多语文书研究继续追问。

\subsection{1279年（元至元十六年、宋祥兴二年；干支、月日待核；公元换算据双方本纪年号对照）}

\noindent\eventanchor{YUAN-1279-EVENT-178}{崖山海战后南宋政权终结}\textbf{元与南宋。}
1279年（元至元十六年、宋祥兴二年；干支、月日待核；公元换算据双方本纪年号对照），元与南宋的行动落在“崖山海战后南宋政权终结”所指的战事与通道上。可见的不是抽象的推进：攻守双方必须让人员、消息与补给穿过具体的城防、道路或水路；账本所列的结果只说明这一节点改变了下一步选择。

从账本的前后项看，“忽必烈政权、征宋和元朝建制中既有的联盟、交通与补给压力，为“崖山海战后南宋政权终结”提供了直接的军事或动员背景。”提供了背景压力；““崖山海战后南宋政权终结”改变了相关城镇、道路或政权间的军事选择；伤亡、俘获和控制范围仅可按各方材料分别确认。”标出后来需要追踪的变化，二者之间仍存在未被材料填满的环节。

核心线索为《元史》相关本纪；《元朝秘史》、波斯文史籍或《高丽史》对应叙事；《元典章》《通制条格》、多语碑刻与文书（据事核）。《元史》明初速修，纪年与人名有异同；蒙古大汗政权不等同所有汗国，征服、赋役和身份分类须按地区及材料语言分列。 现代研究可从蒙古帝国、元朝行省财政、多语文书与区域社会研究继续追问。

\subsection{1280年（元至元17年；公元换算据《元史》本纪纪年）}

\noindent\eventanchor{SYD-1280-YUAN-008}{甲戌，給要束合所領工匠牛二千，就令運米二千石供軍}\textbf{元。}
元的1280年（元至元17年；公元换算据《元史》本纪纪年）记录写到“甲戌，給要束合所領工匠牛二千，就令運米二千石供軍”。我们可以据此辨认一个可见的选择或压力，而不能把零散的官方记载、碑刻或后出叙事拼成无缝的社会全景。

从账本的前后项看，“本纪所记财用、赈济或征发属于特定报告场景。”提供了背景压力；“可与食货志、文书和物质材料比较，不外推为全境总量。”标出后来需要追踪的变化，二者之间仍存在未被材料填满的环节。

核心线索为《元史》卷010〈本纪第10卷本年条〉。《元史》为明初仓促官修，本纪保存大量实录条目但有编次与纪年问题；行省执行、数字、族称及对外关系须以条格、多语文书和对方史料互校。 现代研究可从元代行省、赋役、河工与多语文书研究继续追问。

\noindent\eventanchor{SYD-1280-YUAN-080}{近侍劉鐵木兒因言：「阿里海牙屬吏張鼎，今亦參知攻事}\textbf{元。}
在1280年（元至元17年；公元换算据《元史》本纪纪年）的材料中，元面对的是近侍劉鐵木兒因言：「阿里海牙屬吏張鼎，今亦參知攻事。这里应先看行动所依赖的关隘、城镇、河道或粮道，再讨论政治后果；军队抵达某处，并不等于它已能稳定征收、安置或控制周边。

从账本的前后项看，“本纪从朝廷视角记录军政行动；出兵与战果需分辨。”提供了背景压力；“后续路线、控制与损失应与对方记录和遗址材料复核。”标出后来需要追踪的变化，二者之间仍存在未被材料填满的环节。

核心线索为《元史》卷010〈本纪第10卷本年条〉。《元史》为明初仓促官修，本纪保存大量实录条目但有编次与纪年问题；行省执行、数字、族称及对外关系须以条格、多语文书和对方史料互校。 现代研究可从元代行省、赋役、河工与多语文书研究继续追问。

\noindent\eventanchor{SYD-1280-YUAN-152}{明日，下詔皇太子燕王參決朝政，凡中書省、樞密院、御史臺及百司之事，皆先啟後聞}\textbf{元。}
元于1280年（元至元17年；公元换算据《元史》本纪纪年）处理“明日，下詔皇太子燕王參決朝政，凡中書省、樞密院、御史臺及百司之事，皆先啟後聞”时，面对的是道路、翻译、礼仪和资源交换共同构成的关系网。书面称谓可见，地方执行和持续性则仍须由后续材料检验。

其前提记为“本纪年条记录政令或机构处置；实施背景须参照制度材料。”；其后续是“可在同年及后续政令中追踪；条文不单独证明地方执行。”。这是一条待后续材料检验的事件链，而不是对所有行动者意图的替写。

可核的材料入口是《元史》卷010〈本纪第10卷本年条〉。《元史》为明初仓促官修，本纪保存大量实录条目但有编次与纪年问题；行省执行、数字、族称及对外关系须以条格、多语文书和对方史料互校。 就此而言，元代行省、赋役、河工与多语文书研究有助于把年序同区域、制度或物质材料并读。

\noindent\eventanchor{SYD-1280-YUAN-224}{女直、水達達軍不出征者，令隸民籍輸賦}\textbf{元。}
元的1280年（元至元17年；公元换算据《元史》本纪纪年）记录写到“女直、水達達軍不出征者，令隸民籍輸賦”。我们可以据此辨认一个可见的选择或压力，而不能把零散的官方记载、碑刻或后出叙事拼成无缝的社会全景。

从账本的前后项看，“本纪所记财用、赈济或征发属于特定报告场景。”提供了背景压力；“可与食货志、文书和物质材料比较，不外推为全境总量。”标出后来需要追踪的变化，二者之间仍存在未被材料填满的环节。

核心线索为《元史》卷010〈本纪第10卷本年条〉。《元史》为明初仓促官修，本纪保存大量实录条目但有编次与纪年问题；行省执行、数字、族称及对外关系须以条格、多语文书和对方史料互校。 现代研究可从元代行省、赋役、河工与多语文书研究继续追问。

\noindent\eventanchor{SYD-1280-YUAN-296}{申嚴無籍軍虜掠及傭奴代軍之禁}\textbf{元。}
元的1280年（元至元17年；公元换算据《元史》本纪纪年）记录写到“申嚴無籍軍虜掠及傭奴代軍之禁”。我们可以据此辨认一个可见的选择或压力，而不能把零散的官方记载、碑刻或后出叙事拼成无缝的社会全景。

账本把这一节点放在“本纪年条记录政令或机构处置；实施背景须参照制度材料。”与“可在同年及后续政令中追踪；条文不单独证明地方执行。”之间；两者提示前后压力，并不把时间相邻写成单一因果。

本条首先回到《元史》卷010〈本纪第10卷本年条〉；《元史》为明初仓促官修，本纪保存大量实录条目但有编次与纪年问题；行省执行、数字、族称及对外关系须以条格、多语文书和对方史料互校。 因而元代行省、赋役、河工与多语文书研究不是装饰性的书目，而是检验这一叙述边界的路径。

\subsection{1281年}

\noindent\eventanchor{YUAN-1281-EVENT-179}{元组织第二次对日远征，舰队损失严重}\textbf{元与日本。}
在1281年的材料中，元与日本面对的是元组织第二次对日远征，舰队损失严重。这里应先看行动所依赖的关隘、城镇、河道或粮道，再讨论政治后果；军队抵达某处，并不等于它已能稳定征收、安置或控制周边。

从账本的前后项看，““元组织第二次对日远征，舰队损失严重”发生在忽必烈政权、征宋和元朝建制的具体压力与机会之中，相关行动者的选择不能脱离此前事件链解释。”提供了背景压力；““元组织第二次对日远征，舰队损失严重”为后续政治、边防或资源配置留下条件；其社会影响与实际覆盖范围仍须以异类材料限定。”标出后来需要追踪的变化，二者之间仍存在未被材料填满的环节。

核心线索为《元史》相关本纪；《元朝秘史》、波斯文史籍或《高丽史》对应叙事；《元典章》《通制条格》、多语碑刻与文书（据事核）。《元史》明初速修，纪年与人名有异同；蒙古大汗政权不等同所有汗国，征服、赋役和身份分类须按地区及材料语言分列。 现代研究可从蒙古帝国、元朝行省财政、多语文书与区域社会研究继续追问。

\subsection{1287年}

\noindent\eventanchor{YUAN-1287-EVENT-180}{乃颜反叛，元廷在东北平定诸王叛乱}\textbf{元。}
在1287年的可见材料中，元留下“乃颜反叛，元廷在东北平定诸王叛乱”。这一动作把具体人群、机构或地方条件带进年序，但一项记录不能替未被记载者概括整体反应。

从账本的前后项看，“忽必烈政权、征宋和元朝建制中既有的联盟、交通与补给压力，为“乃颜反叛，元廷在东北平定诸王叛乱”提供了直接的军事或动员背景。”提供了背景压力；““乃颜反叛，元廷在东北平定诸王叛乱”改变了相关城镇、道路或政权间的军事选择；伤亡、俘获和控制范围仅可按各方材料分别确认。”标出后来需要追踪的变化，二者之间仍存在未被材料填满的环节。

核心线索为《元史》相关本纪；《元朝秘史》、波斯文史籍或《高丽史》对应叙事；《元典章》《通制条格》、多语碑刻与文书（据事核）。《元史》明初速修，纪年与人名有异同；蒙古大汗政权不等同所有汗国，征服、赋役和身份分类须按地区及材料语言分列。 现代研究可从蒙古帝国、元朝行省财政、多语文书与区域社会研究继续追问。

\subsection{1288年（元至元25年；公元换算据《元史》本纪纪年）}

\noindent\eventanchor{SYD-1288-YUAN-009}{:間者，行中書省右丞相伯顏遣使來奏，宋母后、幼主暨諸大臣百官，已於正月十八日齎璽綬奉表降附}\textbf{元。}
在1288年（元至元25年；公元换算据《元史》本纪纪年）的材料中，元面对的是:間者，行中書省右丞相伯顏遣使來奏，宋母后、幼主暨諸大臣百官，已於正月十八日齎璽綬奉表降附。这里应先看行动所依赖的关隘、城镇、河道或粮道，再讨论政治后果；军队抵达某处，并不等于它已能稳定征收、安置或控制周边。

此前的条件是“本纪年条记录政令或机构处置；实施背景须参照制度材料。”，此后的可见影响是“可在同年及后续政令中追踪；条文不单独证明地方执行。”。这里的“影响”须按地点、机构和材料种类分别衡量。

这一叙述依据《元史》卷009〈本纪第9卷本年条〉，但《元史》为明初仓促官修，本纪保存大量实录条目但有编次与纪年问题；行省执行、数字、族称及对外关系须以条格、多语文书和对方史料互校。。研究元代行省、赋役、河工与多语文书研究可以补问材料未直接说出的空间与社会差异。

\noindent\eventanchor{SYD-1288-YUAN-081}{伯顏既受降表、玉璽，復遣囊加帶以趙尹甫、賈餘慶等還臨安，召宰相出議降事}\textbf{元。}
在1288年（元至元25年；公元换算据《元史》本纪纪年）的材料中，元面对的是伯顏既受降表、玉璽，復遣囊加帶以趙尹甫、賈餘慶等還臨安，召宰相出議降事。这里应先看行动所依赖的关隘、城镇、河道或粮道，再讨论政治后果；军队抵达某处，并不等于它已能稳定征收、安置或控制周边。

此前的条件是“本纪年条记录政令或机构处置；实施背景须参照制度材料。”，此后的可见影响是“可在同年及后续政令中追踪；条文不单独证明地方执行。”。这里的“影响”须按地点、机构和材料种类分别衡量。

这一叙述依据《元史》卷009〈本纪第9卷本年条〉，但《元史》为明初仓促官修，本纪保存大量实录条目但有编次与纪年问题；行省执行、数字、族称及对外关系须以条格、多语文书和对方史料互校。。研究元代行省、赋役、河工与多语文书研究可以补问材料未直接说出的空间与社会差异。

\noindent\eventanchor{SYD-1288-YUAN-153}{伯顏言：「張惠守宋府庫，不俟命擅啟管鑰}\textbf{元。}
在1288年（元至元25年；公元换算据《元史》本纪纪年）的材料中，元面对的是伯顏言：「張惠守宋府庫，不俟命擅啟管鑰。这里应先看行动所依赖的关隘、城镇、河道或粮道，再讨论政治后果；军队抵达某处，并不等于它已能稳定征收、安置或控制周边。

账本把这一节点放在“本纪年条记录政令或机构处置；实施背景须参照制度材料。”与“可在同年及后续政令中追踪；条文不单独证明地方执行。”之间；两者提示前后压力，并不把时间相邻写成单一因果。

本条首先回到《元史》卷009〈本纪第9卷本年条〉；《元史》为明初仓促官修，本纪保存大量实录条目但有编次与纪年问题；行省执行、数字、族称及对外关系须以条格、多语文书和对方史料互校。 因而元代行省、赋役、河工与多语文书研究不是装饰性的书目，而是检验这一叙述边界的路径。

\noindent\eventanchor{SYD-1288-YUAN-225}{帝命董文忠答之曰：「借使似道實輕汝曹，特似道一人之過耳，且汝主何負焉}\textbf{元。}
元于1288年（元至元25年；公元换算据《元史》本纪纪年）处理“帝命董文忠答之曰：「借使似道實輕汝曹，特似道一人之過耳，且汝主何負焉”时，面对的是道路、翻译、礼仪和资源交换共同构成的关系网。书面称谓可见，地方执行和持续性则仍须由后续材料检验。

此前的条件是“本纪年条记录政令或机构处置；实施背景须参照制度材料。”，此后的可见影响是“可在同年及后续政令中追踪；条文不单独证明地方执行。”。这里的“影响”须按地点、机构和材料种类分别衡量。

这一叙述依据《元史》卷009〈本纪第9卷本年条〉，但《元史》为明初仓促官修，本纪保存大量实录条目但有编次与纪年问题；行省执行、数字、族称及对外关系须以条格、多语文书和对方史料互校。。研究元代行省、赋役、河工与多语文书研究可以补问材料未直接说出的空间与社会差异。

\noindent\eventanchor{SYD-1288-YUAN-297}{命以安慶、蘄、黃等郡宿兵，付宋都帶將之}\textbf{元。}
在1288年（元至元25年；公元换算据《元史》本纪纪年）的材料中，元面对的是命以安慶、蘄、黃等郡宿兵，付宋都帶將之。这里应先看行动所依赖的关隘、城镇、河道或粮道，再讨论政治后果；军队抵达某处，并不等于它已能稳定征收、安置或控制周边。

从账本的前后项看，“本纪年条记录政令或机构处置；实施背景须参照制度材料。”提供了背景压力；“可在同年及后续政令中追踪；条文不单独证明地方执行。”标出后来需要追踪的变化，二者之间仍存在未被材料填满的环节。

核心线索为《元史》卷009〈本纪第9卷本年条〉。《元史》为明初仓促官修，本纪保存大量实录条目但有编次与纪年问题；行省执行、数字、族称及对外关系须以条格、多语文书和对方史料互校。 现代研究可从元代行省、赋役、河工与多语文书研究继续追问。

\subsection{1289年（元至元26年；公元换算据《元史》本纪纪年）}

\noindent\eventanchor{SYD-1289-YUAN-010}{帝然其言，遽命止之}\textbf{元。}
“帝然其言，遽命止之”是元在1289年（元至元26年；公元换算据《元史》本纪纪年）可追查的节点。它提示命令、环境、宗教、迁徙或地方组织如何进入政治过程；影响的范围和速度仍要避免由一条材料外推。

其前提记为“本纪年条记录政令或机构处置；实施背景须参照制度材料。”；其后续是“可在同年及后续政令中追踪；条文不单独证明地方执行。”。这是一条待后续材料检验的事件链，而不是对所有行动者意图的替写。

可核的材料入口是《元史》卷009〈本纪第9卷本年条〉。《元史》为明初仓促官修，本纪保存大量实录条目但有编次与纪年问题；行省执行、数字、族称及对外关系须以条格、多语文书和对方史料互校。 就此而言，元代行省、赋役、河工与多语文书研究有助于把年序同区域、制度或物质材料并读。

\noindent\eventanchor{SYD-1289-YUAN-082}{壬子，寶應軍人施福殺其守將，降于淮東都元帥府，詔以福為千戶，佩金符}\textbf{元。}
“壬子，寶應軍人施福殺其守將，降于淮東都元帥府，詔以福為千戶，佩金符”把元置于军事行动的可核时点。战报能够记下攻守、入城或撤离，却很少完整保留沿途征发、居民去留和补给损耗；因此不能把一方的胜败叙事延伸成所有地方的经验。

此前的条件是“本纪年条记录政令或机构处置；实施背景须参照制度材料。”，此后的可见影响是“可在同年及后续政令中追踪；条文不单独证明地方执行。”。这里的“影响”须按地点、机构和材料种类分别衡量。

这一叙述依据《元史》卷009〈本纪第9卷本年条〉，但《元史》为明初仓促官修，本纪保存大量实录条目但有编次与纪年问题；行省执行、数字、族称及对外关系须以条格、多语文书和对方史料互校。。研究元代行省、赋役、河工与多语文书研究可以补问材料未直接说出的空间与社会差异。

\noindent\eventanchor{SYD-1289-YUAN-154}{詔賜鈔百八十錠贖還之}\textbf{元。}
“詔賜鈔百八十錠贖還之”是元在1289年（元至元26年；公元换算据《元史》本纪纪年）可追查的节点。它提示命令、环境、宗教、迁徙或地方组织如何进入政治过程；影响的范围和速度仍要避免由一条材料外推。

账本把这一节点放在“本纪年条记录政令或机构处置；实施背景须参照制度材料。”与“可在同年及后续政令中追踪；条文不单独证明地方执行。”之间；两者提示前后压力，并不把时间相邻写成单一因果。

本条首先回到《元史》卷009〈本纪第9卷本年条〉；《元史》为明初仓促官修，本纪保存大量实录条目但有编次与纪年问题；行省执行、数字、族称及对外关系须以条格、多语文书和对方史料互校。 因而元代行省、赋役、河工与多语文书研究不是装饰性的书目，而是检验这一叙述边界的路径。

\noindent\eventanchor{SYD-1289-YUAN-226}{詔鄭鼎所將侍衞軍萬人還京師，崔斌、阿里海牙同駐靜江，忽都鐵木兒、鄭鼎同駐鄂漢，賈居貞、脫博忽魯禿花同駐潭州}\textbf{元。}
“詔鄭鼎所將侍衞軍萬人還京師，崔斌、阿里海牙同駐靜江，忽都鐵木兒、鄭鼎同駐鄂漢，賈居貞、脫博忽魯禿花同駐潭州”是元在1289年（元至元26年；公元换算据《元史》本纪纪年）可追查的节点。它提示命令、环境、宗教、迁徙或地方组织如何进入政治过程；影响的范围和速度仍要避免由一条材料外推。

此前的条件是“本纪年条记录政令或机构处置；实施背景须参照制度材料。”，此后的可见影响是“可在同年及后续政令中追踪；条文不单独证明地方执行。”。这里的“影响”须按地点、机构和材料种类分别衡量。

这一叙述依据《元史》卷009〈本纪第9卷本年条〉，但《元史》为明初仓促官修，本纪保存大量实录条目但有编次与纪年问题；行省执行、数字、族称及对外关系须以条格、多语文书和对方史料互校。。研究元代行省、赋役、河工与多语文书研究可以补问材料未直接说出的空间与社会差异。

\noindent\eventanchor{SYD-1289-YUAN-298}{昂吉兒、忻都、唐兀帶等引兵攻司空山寨，破之，殺張德興，執其三子以歸}\textbf{元。}
“昂吉兒、忻都、唐兀帶等引兵攻司空山寨，破之，殺張德興，執其三子以歸”把元置于军事行动的可核时点。战报能够记下攻守、入城或撤离，却很少完整保留沿途征发、居民去留和补给损耗；因此不能把一方的胜败叙事延伸成所有地方的经验。

从账本的前后项看，“本纪从朝廷视角记录军政行动；出兵与战果需分辨。”提供了背景压力；“后续路线、控制与损失应与对方记录和遗址材料复核。”标出后来需要追踪的变化，二者之间仍存在未被材料填满的环节。

核心线索为《元史》卷009〈本纪第9卷本年条〉。《元史》为明初仓促官修，本纪保存大量实录条目但有编次与纪年问题；行省执行、数字、族称及对外关系须以条格、多语文书和对方史料互校。 现代研究可从元代行省、赋役、河工与多语文书研究继续追问。

\subsection{1294年（元至元31年；公元换算据《元史》本纪纪年）}

\noindent\eventanchor{SYD-1294-YUAN-011}{:尚念先朝庶政，悉有成規，惟慎奉行，罔敢失墜}\textbf{元。}
元于1294年（元至元31年；公元换算据《元史》本纪纪年）处理“:尚念先朝庶政，悉有成規，惟慎奉行，罔敢失墜”时，面对的是道路、翻译、礼仪和资源交换共同构成的关系网。书面称谓可见，地方执行和持续性则仍须由后续材料检验。

此前的条件是“本纪年条提供日期和中枢可见行动；前因不据单条补定。”，此后的可见影响是“可回查同年及后续条目，地方影响仍须另证。”。这里的“影响”须按地点、机构和材料种类分别衡量。

这一叙述依据《元史》卷018〈至元三十一年〉，但《元史》为明初仓促官修，本纪保存大量实录条目但有编次与纪年问题；行省执行、数字、族称及对外关系须以条格、多语文书和对方史料互校。。研究元代行省、赋役、河工与多语文书研究可以补问材料未直接说出的空间与社会差异。

\noindent\eventanchor{SYD-1294-YUAN-083}{:朕惟太祖聖武皇帝受天明命，肇造區夏，聖聖相承，光熙前緒}\textbf{元。}
元的1294年（元至元31年；公元换算据《元史》本纪纪年）记录写到“:朕惟太祖聖武皇帝受天明命，肇造區夏，聖聖相承，光熙前緒”。我们可以据此辨认一个可见的选择或压力，而不能把零散的官方记载、碑刻或后出叙事拼成无缝的社会全景。

从账本的前后项看，“本纪年条记录政令或机构处置；实施背景须参照制度材料。”提供了背景压力；“可在同年及后续政令中追踪；条文不单独证明地方执行。”标出后来需要追踪的变化，二者之间仍存在未被材料填满的环节。

核心线索为《元史》卷018〈至元三十一年〉。《元史》为明初仓促官修，本纪保存大量实录条目但有编次与纪年问题；行省执行、数字、族称及对外关系须以条格、多语文书和对方史料互校。 现代研究可从元代行省、赋役、河工与多语文书研究继续追问。

\noindent\eventanchor{SYD-1294-YUAN-155}{帝命以軍職尊崇者授之}\textbf{元。}
元的1294年（元至元31年；公元换算据《元史》本纪纪年）记录写到“帝命以軍職尊崇者授之”。我们可以据此辨认一个可见的选择或压力，而不能把零散的官方记载、碑刻或后出叙事拼成无缝的社会全景。

其前提记为“本纪年条记录政令或机构处置；实施背景须参照制度材料。”；其后续是“可在同年及后续政令中追踪；条文不单独证明地方执行。”。这是一条待后续材料检验的事件链，而不是对所有行动者意图的替写。

可核的材料入口是《元史》卷018〈至元三十一年〉。《元史》为明初仓促官修，本纪保存大量实录条目但有编次与纪年问题；行省执行、数字、族称及对外关系须以条格、多语文书和对方史料互校。 就此而言，元代行省、赋役、河工与多语文书研究有助于把年序同区域、制度或物质材料并读。

\noindent\eventanchor{SYD-1294-YUAN-227}{定西平王奧魯赤、寧遠王闊闊出、鎮南王脫歡及也先帖木而大會賞賜例，金各五百兩、銀五千兩、鈔二千錠、幣帛各二百匹}\textbf{元。}
元的1294年（元至元31年；公元换算据《元史》本纪纪年）记录写到“定西平王奧魯赤、寧遠王闊闊出、鎮南王脫歡及也先帖木而大會賞賜例，金各五百兩、銀五千兩、鈔二千錠、幣帛各二百匹”。我们可以据此辨认一个可见的选择或压力，而不能把零散的官方记载、碑刻或后出叙事拼成无缝的社会全景。

从账本的前后项看，“本纪所记财用、赈济或征发属于特定报告场景。”提供了背景压力；“可与食货志、文书和物质材料比较，不外推为全境总量。”标出后来需要追踪的变化，二者之间仍存在未被材料填满的环节。

核心线索为《元史》卷018〈至元三十一年〉。《元史》为明初仓促官修，本纪保存大量实录条目但有编次与纪年问题；行省执行、数字、族称及对外关系须以条格、多语文书和对方史料互校。 现代研究可从元代行省、赋役、河工与多语文书研究继续追问。

\noindent\eventanchor{SYD-1294-YUAN-299}{己未，復立平陽路之蒲、武鄉，保定路之博野，泰安州之新泰等縣}\textbf{元。}
在1294年（元至元31年；公元换算据《元史》本纪纪年），元所记“己未，復立平陽路之蒲、武鄉，保定路之博野，泰安州之新泰等縣”把权力安排暴露在继承压力下。后续稳定与否取决于资源和支持网络如何被重新调度，不能由即位仪式本身预先断言。

账本把这一节点放在“本纪年条记录政令或机构处置；实施背景须参照制度材料。”与“可在同年及后续政令中追踪；条文不单独证明地方执行。”之间；两者提示前后压力，并不把时间相邻写成单一因果。

本条首先回到《元史》卷018〈至元三十一年〉；《元史》为明初仓促官修，本纪保存大量实录条目但有编次与纪年问题；行省执行、数字、族称及对外关系须以条格、多语文书和对方史料互校。 因而元代行省、赋役、河工与多语文书研究不是装饰性的书目，而是检验这一叙述边界的路径。

\subsection{1294年}

\noindent\eventanchor{YUAN-1294-EVENT-181}{忽必烈去世，铁穆耳即位}\textbf{元。}
在1294年，元所记“忽必烈去世，铁穆耳即位”把权力安排暴露在继承压力下。后续稳定与否取决于资源和支持网络如何被重新调度，不能由即位仪式本身预先断言。

此前的条件是“忽必烈政权、征宋和元朝建制中前任统治的结束与宫廷、部族或官僚的权力安排相互交织，促成“忽必烈去世，铁穆耳即位”这一继承节点。”，此后的可见影响是“继承并不自动消除既有矛盾；“忽必烈去世，铁穆耳即位”后的任官、军权和对外策略需由后续条目追踪。”。这里的“影响”须按地点、机构和材料种类分别衡量。

这一叙述依据《元史》相关本纪；《元朝秘史》、波斯文史籍或《高丽史》对应叙事；《元典章》《通制条格》、多语碑刻与文书（据事核），但《元史》明初速修，纪年与人名有异同；蒙古大汗政权不等同所有汗国，征服、赋役和身份分类须按地区及材料语言分列。。研究蒙古帝国、元朝行省财政、多语文书与区域社会研究可以补问材料未直接说出的空间与社会差异。

\subsection{1295年（元元贞1年；公元换算据《元史》本纪纪年）}

\noindent\eventanchor{SYD-1295-YUAN-012}{給江南行御史臺守護軍百人}\textbf{元。}
“給江南行御史臺守護軍百人”把元置于军事行动的可核时点。战报能够记下攻守、入城或撤离，却很少完整保留沿途征发、居民去留和补给损耗；因此不能把一方的胜败叙事延伸成所有地方的经验。

此前的条件是“本纪所记财用、赈济或征发属于特定报告场景。”，此后的可见影响是“可与食货志、文书和物质材料比较，不外推为全境总量。”。这里的“影响”须按地点、机构和材料种类分别衡量。

这一叙述依据《元史》卷018〈元貞元年〉，但《元史》为明初仓促官修，本纪保存大量实录条目但有编次与纪年问题；行省执行、数字、族称及对外关系须以条格、多语文书和对方史料互校。。研究元代行省、赋役、河工与多语文书研究可以补问材料未直接说出的空间与社会差异。

\noindent\eventanchor{SYD-1295-YUAN-084}{詔自今已後，專令中書擬奏}\textbf{元。}
“詔自今已後，專令中書擬奏”是元在1295年（元元贞1年；公元换算据《元史》本纪纪年）可追查的节点。它提示命令、环境、宗教、迁徙或地方组织如何进入政治过程；影响的范围和速度仍要避免由一条材料外推。

从账本的前后项看，“本纪年条记录政令或机构处置；实施背景须参照制度材料。”提供了背景压力；“可在同年及后续政令中追踪；条文不单独证明地方执行。”标出后来需要追踪的变化，二者之间仍存在未被材料填满的环节。

核心线索为《元史》卷018〈元貞元年〉。《元史》为明初仓促官修，本纪保存大量实录条目但有编次与纪年问题；行省执行、数字、族称及对外关系须以条格、多语文书和对方史料互校。 现代研究可从元代行省、赋役、河工与多语文书研究继续追问。

\noindent\eventanchor{SYD-1295-YUAN-156}{〔己〕卯，罷四川淘金戶四千，還其元籍，罪初獻言者}\textbf{元。}
“〔己〕卯，罷四川淘金戶四千，還其元籍，罪初獻言者”是元在1295年（元元贞1年；公元换算据《元史》本纪纪年）可追查的节点。它提示命令、环境、宗教、迁徙或地方组织如何进入政治过程；影响的范围和速度仍要避免由一条材料外推。

其前提记为“本纪年条记录政令或机构处置；实施背景须参照制度材料。”；其后续是“可在同年及后续政令中追踪；条文不单独证明地方执行。”。这是一条待后续材料检验的事件链，而不是对所有行动者意图的替写。

可核的材料入口是《元史》卷018〈元貞元年〉。《元史》为明初仓促官修，本纪保存大量实录条目但有编次与纪年问题；行省执行、数字、族称及对外关系须以条格、多语文书和对方史料互校。 就此而言，元代行省、赋役、河工与多语文书研究有助于把年序同区域、制度或物质材料并读。

\noindent\eventanchor{SYD-1295-YUAN-228}{〔己〕卯，詔申飭中外：「有儒吏兼通者，各路舉之，廉訪司每道歲貢二人，省臺委官立法考試，中程者用之，所貢不公，罪其舉者}\textbf{元。}
元的“〔己〕卯，詔申飭中外：「有儒吏兼通者，各路舉之，廉訪司每道歲貢二人，省臺委官立法考試，中程者用之，所貢不公，罪其舉者”并非只改换一个统治者。任官、赏赐、军权与信息通路要在随后重新接合；材料较能确认先后，较难替当事人写出唯一动机。

从账本的前后项看，“本纪年条记录政令或机构处置；实施背景须参照制度材料。”提供了背景压力；“可在同年及后续政令中追踪；条文不单独证明地方执行。”标出后来需要追踪的变化，二者之间仍存在未被材料填满的环节。

核心线索为《元史》卷018〈元貞元年〉。《元史》为明初仓促官修，本纪保存大量实录条目但有编次与纪年问题；行省执行、数字、族称及对外关系须以条格、多语文书和对方史料互校。 现代研究可从元代行省、赋役、河工与多语文书研究继续追问。

\noindent\eventanchor{SYD-1295-YUAN-300}{罷河西軍，聽各還其所屬}\textbf{元。}
“罷河西軍，聽各還其所屬”是元在1295年（元元贞1年；公元换算据《元史》本纪纪年）可追查的节点。它提示命令、环境、宗教、迁徙或地方组织如何进入政治过程；影响的范围和速度仍要避免由一条材料外推。

账本把这一节点放在“本纪年条记录政令或机构处置；实施背景须参照制度材料。”与“可在同年及后续政令中追踪；条文不单独证明地方执行。”之间；两者提示前后压力，并不把时间相邻写成单一因果。

本条首先回到《元史》卷018〈元貞元年〉；《元史》为明初仓促官修，本纪保存大量实录条目但有编次与纪年问题；行省执行、数字、族称及对外关系须以条格、多语文书和对方史料互校。 因而元代行省、赋役、河工与多语文书研究不是装饰性的书目，而是检验这一叙述边界的路径。

\subsection{1296年（元元贞2年；公元换算据《元史》本纪纪年）}

\noindent\eventanchor{SYD-1296-YUAN-013}{甲戌，詔民間馬牛羊，百取其一，羊不滿百者亦取之，惟色目人及數乃取}\textbf{元。}
在1296年（元元贞2年；公元换算据《元史》本纪纪年）的可见材料中，元留下“甲戌，詔民間馬牛羊，百取其一，羊不滿百者亦取之，惟色目人及數乃取”。这一动作把具体人群、机构或地方条件带进年序，但一项记录不能替未被记载者概括整体反应。

此前的条件是“本纪年条记录政令或机构处置；实施背景须参照制度材料。”，此后的可见影响是“可在同年及后续政令中追踪；条文不单独证明地方执行。”。这里的“影响”须按地点、机构和材料种类分别衡量。

这一叙述依据《元史》卷019〈元貞二年〉，但《元史》为明初仓促官修，本纪保存大量实录条目但有编次与纪年问题；行省执行、数字、族称及对外关系须以条格、多语文书和对方史料互校。。研究元代行省、赋役、河工与多语文书研究可以补问材料未直接说出的空间与社会差异。

\noindent\eventanchor{SYD-1296-YUAN-085}{辛未，徙江浙行省拔都軍萬人戍潭州，潭州以南軍移戍郴州}\textbf{元。}
在1296年（元元贞2年；公元换算据《元史》本纪纪年）的可见材料中，元留下“辛未，徙江浙行省拔都軍萬人戍潭州，潭州以南軍移戍郴州”。这一动作把具体人群、机构或地方条件带进年序，但一项记录不能替未被记载者概括整体反应。

从账本的前后项看，“本纪年条记录政令或机构处置；实施背景须参照制度材料。”提供了背景压力；“可在同年及后续政令中追踪；条文不单独证明地方执行。”标出后来需要追踪的变化，二者之间仍存在未被材料填满的环节。

核心线索为《元史》卷019〈元貞二年〉。《元史》为明初仓促官修，本纪保存大量实录条目但有编次与纪年问题；行省执行、数字、族称及对外关系须以条格、多语文书和对方史料互校。 现代研究可从元代行省、赋役、河工与多语文书研究继续追问。

\noindent\eventanchor{SYD-1296-YUAN-157}{詔復立蒙樣剛等甸軍民官}\textbf{元。}
“詔復立蒙樣剛等甸軍民官”在1296年（元元贞2年；公元换算据《元史》本纪纪年）构成元的继承或权力重组节点。名位的变化可以由纪年串连，但它没有自动消除宫廷、军政和地方网络原有的拉扯。

其前提记为“本纪年条记录政令或机构处置；实施背景须参照制度材料。”；其后续是“可在同年及后续政令中追踪；条文不单独证明地方执行。”。这是一条待后续材料检验的事件链，而不是对所有行动者意图的替写。

可核的材料入口是《元史》卷019〈元貞二年〉。《元史》为明初仓促官修，本纪保存大量实录条目但有编次与纪年问题；行省执行、数字、族称及对外关系须以条格、多语文书和对方史料互校。 就此而言，元代行省、赋役、河工与多语文书研究有助于把年序同区域、制度或物质材料并读。

\noindent\eventanchor{SYD-1296-YUAN-229}{詔行省圖地形、覈軍實以聞}\textbf{元。}
在1296年（元元贞2年；公元换算据《元史》本纪纪年）的可见材料中，元留下“詔行省圖地形、覈軍實以聞”。这一动作把具体人群、机构或地方条件带进年序，但一项记录不能替未被记载者概括整体反应。

从账本的前后项看，“本纪年条记录政令或机构处置；实施背景须参照制度材料。”提供了背景压力；“可在同年及后续政令中追踪；条文不单独证明地方执行。”标出后来需要追踪的变化，二者之间仍存在未被材料填满的环节。

核心线索为《元史》卷019〈元貞二年〉。《元史》为明初仓促官修，本纪保存大量实录条目但有编次与纪年问题；行省执行、数字、族称及对外关系须以条格、多语文书和对方史料互校。 现代研究可从元代行省、赋役、河工与多语文书研究继续追问。

\noindent\eventanchor{SYD-1296-YUAN-301}{安南國遣人招誘叛賊黃勝許}\textbf{元与安南。}
在1296年（元元贞2年；公元换算据《元史》本纪纪年）的可见材料中，元与安南留下“安南國遣人招誘叛賊黃勝許”。这一动作把具体人群、机构或地方条件带进年序，但一项记录不能替未被记载者概括整体反应。

账本把这一节点放在“本纪从朝廷视角记录军政行动；出兵与战果需分辨。”与“后续路线、控制与损失应与对方记录和遗址材料复核。”之间；两者提示前后压力，并不把时间相邻写成单一因果。

本条首先回到《元史》卷019〈元貞二年〉；《安南志略》及当地碑刻材料（互校）；《元史》为明初仓促官修，本纪保存大量实录条目但有编次与纪年问题；行省执行、数字、族称及对外关系须以条格、多语文书和对方史料互校。 因而元代行省、赋役、河工与多语文书研究不是装饰性的书目，而是检验这一叙述边界的路径。

\subsection{1297年（元大德1年；公元换算据《元史》本纪纪年）}

\noindent\eventanchor{SYD-1297-YUAN-014}{帝曰：「礬課遣人覈實}\textbf{元。}
元的1297年（元大德1年；公元换算据《元史》本纪纪年）记录写到“帝曰：「礬課遣人覈實”。我们可以据此辨认一个可见的选择或压力，而不能把零散的官方记载、碑刻或后出叙事拼成无缝的社会全景。

此前的条件是“本纪年条提供日期和中枢可见行动；前因不据单条补定。”，此后的可见影响是“可回查同年及后续条目，地方影响仍须另证。”。这里的“影响”须按地点、机构和材料种类分别衡量。

这一叙述依据《元史》卷019〈大德元年〉，但《元史》为明初仓促官修，本纪保存大量实录条目但有编次与纪年问题；行省执行、数字、族称及对外关系须以条格、多语文书和对方史料互校。。研究元代行省、赋役、河工与多语文书研究可以补问材料未直接说出的空间与社会差异。

\noindent\eventanchor{SYD-1297-YUAN-086}{禁諸王、駙馬并權豪，毋奪民田，其獻田者有刑}\textbf{元。}
元在1297年（元大德1年；公元换算据《元史》本纪纪年）留下的这条记录是“禁諸王、駙馬并權豪，毋奪民田，其獻田者有刑”。它照亮资源如何被命名和调度，也暴露材料的盲点：数字、额度或禁令若无地方文书和物质材料互证，不能被写成全境同速的现实。

从账本的前后项看，“本纪年条记录政令或机构处置；实施背景须参照制度材料。”提供了背景压力；“可在同年及后续政令中追踪；条文不单独证明地方执行。”标出后来需要追踪的变化，二者之间仍存在未被材料填满的环节。

核心线索为《元史》卷019〈大德元年〉。《元史》为明初仓促官修，本纪保存大量实录条目但有编次与纪年问题；行省执行、数字、族称及对外关系须以条格、多语文书和对方史料互校。 现代研究可从元代行省、赋役、河工与多语文书研究继续追问。

\noindent\eventanchor{SYD-1297-YUAN-158}{又賜緬王弟撒邦巴一珠虎符，酋領阿散三珠虎符，從者金符及金幣，遣之}\textbf{元。}
元的1297年（元大德1年；公元换算据《元史》本纪纪年）记录写到“又賜緬王弟撒邦巴一珠虎符，酋領阿散三珠虎符，從者金符及金幣，遣之”。我们可以据此辨认一个可见的选择或压力，而不能把零散的官方记载、碑刻或后出叙事拼成无缝的社会全景。

其前提记为“本纪年条提供日期和中枢可见行动；前因不据单条补定。”；其后续是“可回查同年及后续条目，地方影响仍须另证。”。这是一条待后续材料检验的事件链，而不是对所有行动者意图的替写。

可核的材料入口是《元史》卷019〈大德元年〉。《元史》为明初仓促官修，本纪保存大量实录条目但有编次与纪年问题；行省执行、数字、族称及对外关系须以条格、多语文书和对方史料互校。 就此而言，元代行省、赋役、河工与多语文书研究有助于把年序同区域、制度或物质材料并读。

\noindent\eventanchor{SYD-1297-YUAN-230}{諸王也只里部忽剌帶於濟南商河縣侵擾居民，蹂踐禾稼，帝命詰之，走歸其部}\textbf{元。}
元的1297年（元大德1年；公元换算据《元史》本纪纪年）记录写到“諸王也只里部忽剌帶於濟南商河縣侵擾居民，蹂踐禾稼，帝命詰之，走歸其部”。我们可以据此辨认一个可见的选择或压力，而不能把零散的官方记载、碑刻或后出叙事拼成无缝的社会全景。

此前的条件是“本纪年条记录政令或机构处置；实施背景须参照制度材料。”，此后的可见影响是“可在同年及后续政令中追踪；条文不单独证明地方执行。”。这里的“影响”须按地点、机构和材料种类分别衡量。

这一叙述依据《元史》卷019〈大德元年〉，但《元史》为明初仓促官修，本纪保存大量实录条目但有编次与纪年问题；行省执行、数字、族称及对外关系须以条格、多语文书和对方史料互校。。研究元代行省、赋役、河工与多语文书研究可以补问材料未直接说出的空间与社会差异。

\noindent\eventanchor{SYD-1297-YUAN-302}{暗都剌火者所部四萬餘錠}\textbf{元。}
元的1297年（元大德1年；公元换算据《元史》本纪纪年）记录写到“暗都剌火者所部四萬餘錠”。我们可以据此辨认一个可见的选择或压力，而不能把零散的官方记载、碑刻或后出叙事拼成无缝的社会全景。

账本把这一节点放在“本纪年条提供日期和中枢可见行动；前因不据单条补定。”与“可回查同年及后续条目，地方影响仍须另证。”之间；两者提示前后压力，并不把时间相邻写成单一因果。

本条首先回到《元史》卷019〈大德元年〉；《元史》为明初仓促官修，本纪保存大量实录条目但有编次与纪年问题；行省执行、数字、族称及对外关系须以条格、多语文书和对方史料互校。 因而元代行省、赋役、河工与多语文书研究不是装饰性的书目，而是检验这一叙述边界的路径。

\subsection{1298年（元大德2年；公元换算据《元史》本纪纪年）}

\noindent\eventanchor{SYD-1298-YUAN-015}{右丞相完澤言：「歲入之數，金一萬九千兩，銀六萬兩，鈔三百六十萬錠，然猶不足於用，又於至元鈔本中借二十萬錠}\textbf{元。}
“右丞相完澤言：「歲入之數，金一萬九千兩，銀六萬兩，鈔三百六十萬錠，然猶不足於用，又於至元鈔本中借二十萬錠”是元在1298年（元大德2年；公元换算据《元史》本纪纪年）可追查的节点。它提示命令、环境、宗教、迁徙或地方组织如何进入政治过程；影响的范围和速度仍要避免由一条材料外推。

此前的条件是“本纪所记财用、赈济或征发属于特定报告场景。”，此后的可见影响是“可与食货志、文书和物质材料比较，不外推为全境总量。”。这里的“影响”须按地点、机构和材料种类分别衡量。

这一叙述依据《元史》卷019〈大德二年〉，但《元史》为明初仓促官修，本纪保存大量实录条目但有编次与纪年问题；行省执行、数字、族称及对外关系须以条格、多语文书和对方史料互校。。研究元代行省、赋役、河工与多语文书研究可以补问材料未直接说出的空间与社会差异。

\noindent\eventanchor{SYD-1298-YUAN-087}{〔衞〕輝、順德旱，大風損麥，免其田租一年}\textbf{元。}
“〔衞〕輝、順德旱，大風損麥，免其田租一年”发生时，元必须把账面规则转成仓储、市场、航道、文书和地方中介的工作。正因这些环节并不齐整，政策颁行与实际效果应分别追踪。

从账本的前后项看，“本纪保留灾害或工程的报告地点与朝廷处置；灾情范围需另证。”提供了背景压力；“可与地方文书、河工和环境材料比较，不将单点记录外推为全域因果。”标出后来需要追踪的变化，二者之间仍存在未被材料填满的环节。

核心线索为《元史》卷019〈大德二年〉。《元史》为明初仓促官修，本纪保存大量实录条目但有编次与纪年问题；行省执行、数字、族称及对外关系须以条格、多语文书和对方史料互校。 现代研究可从元代行省、赋役、河工与多语文书研究继续追问。

\noindent\eventanchor{SYD-1298-YUAN-159}{罷建康金銀銅冶轉運司，還淘金戶於元籍，歲辦金悉責有司}\textbf{元。}
“罷建康金銀銅冶轉運司，還淘金戶於元籍，歲辦金悉責有司”是元在1298年（元大德2年；公元换算据《元史》本纪纪年）可追查的节点。它提示命令、环境、宗教、迁徙或地方组织如何进入政治过程；影响的范围和速度仍要避免由一条材料外推。

其前提记为“本纪年条记录政令或机构处置；实施背景须参照制度材料。”；其后续是“可在同年及后续政令中追踪；条文不单独证明地方执行。”。这是一条待后续材料检验的事件链，而不是对所有行动者意图的替写。

可核的材料入口是《元史》卷019〈大德二年〉。《元史》为明初仓促官修，本纪保存大量实录条目但有编次与纪年问题；行省执行、数字、族称及对外关系须以条格、多语文书和对方史料互校。 就此而言，元代行省、赋役、河工与多语文书研究有助于把年序同区域、制度或物质材料并读。

\noindent\eventanchor{SYD-1298-YUAN-231}{罷徐、邳爐冶所進息錢}\textbf{元。}
“罷徐、邳爐冶所進息錢”是元在1298年（元大德2年；公元换算据《元史》本纪纪年）可追查的节点。它提示命令、环境、宗教、迁徙或地方组织如何进入政治过程；影响的范围和速度仍要避免由一条材料外推。

此前的条件是“本纪年条记录政令或机构处置；实施背景须参照制度材料。”，此后的可见影响是“可在同年及后续政令中追踪；条文不单独证明地方执行。”。这里的“影响”须按地点、机构和材料种类分别衡量。

这一叙述依据《元史》卷019〈大德二年〉，但《元史》为明初仓促官修，本纪保存大量实录条目但有编次与纪年问题；行省执行、数字、族称及对外关系须以条格、多语文书和对方史料互校。。研究元代行省、赋役、河工与多语文书研究可以补问材料未直接说出的空间与社会差异。

\noindent\eventanchor{SYD-1298-YUAN-303}{罷雲南行御史臺，置肅政廉訪司}\textbf{元。}
“罷雲南行御史臺，置肅政廉訪司”是元在1298年（元大德2年；公元换算据《元史》本纪纪年）可追查的节点。它提示命令、环境、宗教、迁徙或地方组织如何进入政治过程；影响的范围和速度仍要避免由一条材料外推。

账本把这一节点放在“本纪年条记录政令或机构处置；实施背景须参照制度材料。”与“可在同年及后续政令中追踪；条文不单独证明地方执行。”之间；两者提示前后压力，并不把时间相邻写成单一因果。

本条首先回到《元史》卷019〈大德二年〉；《元史》为明初仓促官修，本纪保存大量实录条目但有编次与纪年问题；行省执行、数字、族称及对外关系须以条格、多语文书和对方史料互校。 因而元代行省、赋役、河工与多语文书研究不是装饰性的书目，而是检验这一叙述边界的路径。

\subsection{1299年（元大德3年；公元换算据《元史》本纪纪年）}

\noindent\eventanchor{SYD-1299-YUAN-016}{罷四川、福建等處行中書省，陝西行御史臺，江東、荊南、淮西三道宣慰司}\textbf{元。}
在1299年（元大德3年；公元换算据《元史》本纪纪年）的可见材料中，元留下“罷四川、福建等處行中書省，陝西行御史臺，江東、荊南、淮西三道宣慰司”。这一动作把具体人群、机构或地方条件带进年序，但一项记录不能替未被记载者概括整体反应。

此前的条件是“本纪年条记录政令或机构处置；实施背景须参照制度材料。”，此后的可见影响是“可在同年及后续政令中追踪；条文不单独证明地方执行。”。这里的“影响”须按地点、机构和材料种类分别衡量。

这一叙述依据《元史》卷020〈大德三年〉，但《元史》为明初仓促官修，本纪保存大量实录条目但有编次与纪年问题；行省执行、数字、族称及对外关系须以条格、多语文书和对方史料互校。。研究元代行省、赋役、河工与多语文书研究可以补问材料未直接说出的空间与社会差异。

\noindent\eventanchor{SYD-1299-YUAN-088}{庚寅，詔遣使問民疾苦}\textbf{元。}
1299年（元大德3年；公元换算据《元史》本纪纪年），“庚寅，詔遣使問民疾苦”使元与相邻行动者的关系进入可核的外交时点。名号、贡赐和盟约是各方政治语言的一部分；它们能够确认交涉，却不能单凭措辞判定实际隶属或边地生活。

从账本的前后项看，“本纪年条记录政令或机构处置；实施背景须参照制度材料。”提供了背景压力；“可在同年及后续政令中追踪；条文不单独证明地方执行。”标出后来需要追踪的变化，二者之间仍存在未被材料填满的环节。

核心线索为《元史》卷020〈大德三年〉。《元史》为明初仓促官修，本纪保存大量实录条目但有编次与纪年问题；行省执行、数字、族称及对外关系须以条格、多语文书和对方史料互校。 现代研究可从元代行省、赋役、河工与多语文书研究继续追问。

\noindent\eventanchor{SYD-1299-YUAN-160}{甲午，命何榮祖等更定律令}\textbf{元。}
在1299年（元大德3年；公元换算据《元史》本纪纪年）的可见材料中，元留下“甲午，命何榮祖等更定律令”。这一动作把具体人群、机构或地方条件带进年序，但一项记录不能替未被记载者概括整体反应。

账本把这一节点放在“本纪年条记录政令或机构处置；实施背景须参照制度材料。”与“可在同年及后续政令中追踪；条文不单独证明地方执行。”之间；两者提示前后压力，并不把时间相邻写成单一因果。

本条首先回到《元史》卷020〈大德三年〉；《元史》为明初仓促官修，本纪保存大量实录条目但有编次与纪年问题；行省执行、数字、族称及对外关系须以条格、多语文书和对方史料互校。 因而元代行省、赋役、河工与多语文书研究不是装饰性的书目，而是检验这一叙述边界的路径。

\noindent\eventanchor{SYD-1299-YUAN-232}{丙申，揚州、淮安屬縣蝗，在地者為鶖啄食，飛者以翅擊死，詔禁捕鶖}\textbf{元。}
在1299年（元大德3年；公元换算据《元史》本纪纪年）的可见材料中，元留下“丙申，揚州、淮安屬縣蝗，在地者為鶖啄食，飛者以翅擊死，詔禁捕鶖”。这一动作把具体人群、机构或地方条件带进年序，但一项记录不能替未被记载者概括整体反应。

此前的条件是“本纪年条记录政令或机构处置；实施背景须参照制度材料。”，此后的可见影响是“可在同年及后续政令中追踪；条文不单独证明地方执行。”。这里的“影响”须按地点、机构和材料种类分别衡量。

这一叙述依据《元史》卷020〈大德三年〉，但《元史》为明初仓促官修，本纪保存大量实录条目但有编次与纪年问题；行省执行、数字、族称及对外关系须以条格、多语文书和对方史料互校。。研究元代行省、赋役、河工与多语文书研究可以补问材料未直接说出的空间与社会差异。

\noindent\eventanchor{SYD-1299-YUAN-304}{丙寅，詔各省戍軍輪次放還二年供役}\textbf{元。}
在1299年（元大德3年；公元换算据《元史》本纪纪年）的可见材料中，元留下“丙寅，詔各省戍軍輪次放還二年供役”。这一动作把具体人群、机构或地方条件带进年序，但一项记录不能替未被记载者概括整体反应。

账本把这一节点放在“本纪年条记录政令或机构处置；实施背景须参照制度材料。”与“可在同年及后续政令中追踪；条文不单独证明地方执行。”之间；两者提示前后压力，并不把时间相邻写成单一因果。

本条首先回到《元史》卷020〈大德三年〉；《元史》为明初仓促官修，本纪保存大量实录条目但有编次与纪年问题；行省执行、数字、族称及对外关系须以条格、多语文书和对方史料互校。 因而元代行省、赋役、河工与多语文书研究不是装饰性的书目，而是检验这一叙述边界的路径。

\subsection{1300年（元大德4年；公元换算据《元史》本纪纪年）}

\noindent\eventanchor{SYD-1300-YUAN-017}{從之，仍賜交趾使人職名及弓矢鞍勒}\textbf{元。}
1300年（元大德4年；公元换算据《元史》本纪纪年），“從之，仍賜交趾使人職名及弓矢鞍勒”使元与相邻行动者的关系进入可核的外交时点。名号、贡赐和盟约是各方政治语言的一部分；它们能够确认交涉，却不能单凭措辞判定实际隶属或边地生活。

此前的条件是“本纪年条提供日期和中枢可见行动；前因不据单条补定。”，此后的可见影响是“可回查同年及后续条目，地方影响仍须另证。”。这里的“影响”须按地点、机构和材料种类分别衡量。

这一叙述依据《元史》卷011〈本纪第11卷本年条〉，但《元史》为明初仓促官修，本纪保存大量实录条目但有编次与纪年问题；行省执行、数字、族称及对外关系须以条格、多语文书和对方史料互校。。研究元代行省、赋役、河工与多语文书研究可以补问材料未直接说出的空间与社会差异。

\noindent\eventanchor{SYD-1300-YUAN-089}{癸巳，詔諭和州諸城招集流移之民}\textbf{元。}
在1300年（元大德4年；公元换算据《元史》本纪纪年）的可见材料中，元留下“癸巳，詔諭和州諸城招集流移之民”。这一动作把具体人群、机构或地方条件带进年序，但一项记录不能替未被记载者概括整体反应。

从账本的前后项看，“本纪年条记录政令或机构处置；实施背景须参照制度材料。”提供了背景压力；“可在同年及后续政令中追踪；条文不单独证明地方执行。”标出后来需要追踪的变化，二者之间仍存在未被材料填满的环节。

核心线索为《元史》卷011〈本纪第11卷本年条〉。《元史》为明初仓促官修，本纪保存大量实录条目但有编次与纪年问题；行省执行、数字、族称及对外关系须以条格、多语文书和对方史料互校。 现代研究可从元代行省、赋役、河工与多语文书研究继续追问。

\noindent\eventanchor{SYD-1300-YUAN-161}{命和禮霍孫將兵與高和尚同赴北邊}\textbf{元。}
1300年（元大德4年；公元换算据《元史》本纪纪年），元的行动落在“命和禮霍孫將兵與高和尚同赴北邊”所指的战事与通道上。可见的不是抽象的推进：攻守双方必须让人员、消息与补给穿过具体的城防、道路或水路；账本所列的结果只说明这一节点改变了下一步选择。

账本把这一节点放在“本纪年条记录政令或机构处置；实施背景须参照制度材料。”与“可在同年及后续政令中追踪；条文不单独证明地方执行。”之间；两者提示前后压力，并不把时间相邻写成单一因果。

本条首先回到《元史》卷011〈本纪第11卷本年条〉；《元史》为明初仓促官修，本纪保存大量实录条目但有编次与纪年问题；行省执行、数字、族称及对外关系须以条格、多语文书和对方史料互校。 因而元代行省、赋役、河工与多语文书研究不是装饰性的书目，而是检验这一叙述边界的路径。

\noindent\eventanchor{SYD-1300-YUAN-233}{戊申，中書省臣議：「流通鈔法，凡賞賜宜多給幣帛，課程宜多收鈔}\textbf{元。}
在1300年（元大德4年；公元换算据《元史》本纪纪年）的可见材料中，元留下“戊申，中書省臣議：「流通鈔法，凡賞賜宜多給幣帛，課程宜多收鈔”。这一动作把具体人群、机构或地方条件带进年序，但一项记录不能替未被记载者概括整体反应。

此前的条件是“本纪年条记录政令或机构处置；实施背景须参照制度材料。”，此后的可见影响是“可在同年及后续政令中追踪；条文不单独证明地方执行。”。这里的“影响”须按地点、机构和材料种类分别衡量。

这一叙述依据《元史》卷011〈本纪第11卷本年条〉，但《元史》为明初仓促官修，本纪保存大量实录条目但有编次与纪年问题；行省执行、数字、族称及对外关系须以条格、多语文书和对方史料互校。。研究元代行省、赋役、河工与多语文书研究可以补问材料未直接说出的空间与社会差异。

\noindent\eventanchor{SYD-1300-YUAN-305}{詔以忙古帶仍行省福州}\textbf{元。}
在1300年（元大德4年；公元换算据《元史》本纪纪年）的可见材料中，元留下“詔以忙古帶仍行省福州”。这一动作把具体人群、机构或地方条件带进年序，但一项记录不能替未被记载者概括整体反应。

账本把这一节点放在“本纪年条记录政令或机构处置；实施背景须参照制度材料。”与“可在同年及后续政令中追踪；条文不单独证明地方执行。”之间；两者提示前后压力，并不把时间相邻写成单一因果。

本条首先回到《元史》卷011〈本纪第11卷本年条〉；《元史》为明初仓促官修，本纪保存大量实录条目但有编次与纪年问题；行省执行、数字、族称及对外关系须以条格、多语文书和对方史料互校。 因而元代行省、赋役、河工与多语文书研究不是装饰性的书目，而是检验这一叙述边界的路径。

\subsection{1301年（元大德5年；公元换算据《元史》本纪纪年）}

\noindent\eventanchor{SYD-1301-YUAN-018}{癸未，命中書省會計姚演所領漣、海屯田官給之資與歲入之數，便則行之，否則罷去}\textbf{元。}
元在1301年（元大德5年；公元换算据《元史》本纪纪年）留下的这条记录是“癸未，命中書省會計姚演所領漣、海屯田官給之資與歲入之數，便則行之，否則罷去”。它照亮资源如何被命名和调度，也暴露材料的盲点：数字、额度或禁令若无地方文书和物质材料互证，不能被写成全境同速的现实。

此前的条件是“本纪年条记录政令或机构处置；实施背景须参照制度材料。”，此后的可见影响是“可在同年及后续政令中追踪；条文不单独证明地方执行。”。这里的“影响”须按地点、机构和材料种类分别衡量。

这一叙述依据《元史》卷011〈本纪第11卷本年条〉，但《元史》为明初仓促官修，本纪保存大量实录条目但有编次与纪年问题；行省执行、数字、族称及对外关系须以条格、多语文书和对方史料互校。。研究元代行省、赋役、河工与多语文书研究可以补问材料未直接说出的空间与社会差异。

\noindent\eventanchor{SYD-1301-YUAN-090}{詔曰：「軍事卿等當自權衡之}\textbf{元。}
元的1301年（元大德5年；公元换算据《元史》本纪纪年）记录写到“詔曰：「軍事卿等當自權衡之”。我们可以据此辨认一个可见的选择或压力，而不能把零散的官方记载、碑刻或后出叙事拼成无缝的社会全景。

此前的条件是“本纪年条记录政令或机构处置；实施背景须参照制度材料。”，此后的可见影响是“可在同年及后续政令中追踪；条文不单独证明地方执行。”。这里的“影响”须按地点、机构和材料种类分别衡量。

这一叙述依据《元史》卷011〈本纪第11卷本年条〉，但《元史》为明初仓促官修，本纪保存大量实录条目但有编次与纪年问题；行省执行、数字、族称及对外关系须以条格、多语文书和对方史料互校。。研究元代行省、赋役、河工与多语文书研究可以补问材料未直接说出的空间与社会差异。

\noindent\eventanchor{SYD-1301-YUAN-162}{===校勘記===Category:元日戰爭‎}\textbf{元。}
元的1301年（元大德5年；公元换算据《元史》本纪纪年）记录写到“===校勘記===Category:元日戰爭‎”。我们可以据此辨认一个可见的选择或压力，而不能把零散的官方记载、碑刻或后出叙事拼成无缝的社会全景。

账本把这一节点放在“本纪从朝廷视角记录军政行动；出兵与战果需分辨。”与“后续路线、控制与损失应与对方记录和遗址材料复核。”之间；两者提示前后压力，并不把时间相邻写成单一因果。

本条首先回到《元史》卷011〈本纪第11卷本年条〉；《元史》为明初仓促官修，本纪保存大量实录条目但有编次与纪年问题；行省执行、数字、族称及对外关系须以条格、多语文书和对方史料互校。 因而元代行省、赋役、河工与多语文书研究不是装饰性的书目，而是检验这一叙述边界的路径。

\noindent\eventanchor{SYD-1301-YUAN-234}{阿塔海乞以戍三海口軍擊福建賊陳吊眼，詔以重勞不從}\textbf{元。}
元的1301年（元大德5年；公元换算据《元史》本纪纪年）记录写到“阿塔海乞以戍三海口軍擊福建賊陳吊眼，詔以重勞不從”。我们可以据此辨认一个可见的选择或压力，而不能把零散的官方记载、碑刻或后出叙事拼成无缝的社会全景。

此前的条件是“本纪年条记录政令或机构处置；实施背景须参照制度材料。”，此后的可见影响是“可在同年及后续政令中追踪；条文不单独证明地方执行。”。这里的“影响”须按地点、机构和材料种类分别衡量。

这一叙述依据《元史》卷011〈本纪第11卷本年条〉，但《元史》为明初仓促官修，本纪保存大量实录条目但有编次与纪年问题；行省执行、数字、族称及对外关系须以条格、多语文书和对方史料互校。。研究元代行省、赋役、河工与多语文书研究可以补问材料未直接说出的空间与社会差异。

\noindent\eventanchor{SYD-1301-YUAN-306}{八月甲子朔，招討使方文言擇守令、崇祀典、戢姦吏、禁盜賊、治軍旅、奬忠義六事，詔廷臣及諸老議舉行之}\textbf{元。}
在1301年（元大德5年；公元换算据《元史》本纪纪年）的材料中，元面对的是八月甲子朔，招討使方文言擇守令、崇祀典、戢姦吏、禁盜賊、治軍旅、奬忠義六事，詔廷臣及諸老議舉行之。这里应先看行动所依赖的关隘、城镇、河道或粮道，再讨论政治后果；军队抵达某处，并不等于它已能稳定征收、安置或控制周边。

账本把这一节点放在“本纪年条记录政令或机构处置；实施背景须参照制度材料。”与“可在同年及后续政令中追踪；条文不单独证明地方执行。”之间；两者提示前后压力，并不把时间相邻写成单一因果。

本条首先回到《元史》卷011〈本纪第11卷本年条〉；《元史》为明初仓促官修，本纪保存大量实录条目但有编次与纪年问题；行省执行、数字、族称及对外关系须以条格、多语文书和对方史料互校。 因而元代行省、赋役、河工与多语文书研究不是装饰性的书目，而是检验这一叙述边界的路径。

\subsection{1302年（元大德6年；公元换算据《元史》本纪纪年）}

\noindent\eventanchor{SYD-1302-YUAN-019}{復請禁諸路釀酒，減免差稅，賑濟饑民}\textbf{元。}
“復請禁諸路釀酒，減免差稅，賑濟饑民”是元在1302年（元大德6年；公元换算据《元史》本纪纪年）可追查的节点。它提示命令、环境、宗教、迁徙或地方组织如何进入政治过程；影响的范围和速度仍要避免由一条材料外推。

此前的条件是“本纪年条记录政令或机构处置；实施背景须参照制度材料。”，此后的可见影响是“可在同年及后续政令中追踪；条文不单独证明地方执行。”。这里的“影响”须按地点、机构和材料种类分别衡量。

这一叙述依据《元史》卷020〈大德六年〉，但《元史》为明初仓促官修，本纪保存大量实录条目但有编次与纪年问题；行省执行、数字、族称及对外关系须以条格、多语文书和对方史料互校。。研究元代行省、赋役、河工与多语文书研究可以补问材料未直接说出的空间与社会差异。

\noindent\eventanchor{SYD-1302-YUAN-091}{命即行之，毋越三日}\textbf{元。}
“命即行之，毋越三日”是元在1302年（元大德6年；公元换算据《元史》本纪纪年）可追查的节点。它提示命令、环境、宗教、迁徙或地方组织如何进入政治过程；影响的范围和速度仍要避免由一条材料外推。

此前的条件是“本纪年条记录政令或机构处置；实施背景须参照制度材料。”，此后的可见影响是“可在同年及后续政令中追踪；条文不单独证明地方执行。”。这里的“影响”须按地点、机构和材料种类分别衡量。

这一叙述依据《元史》卷020〈大德六年〉，但《元史》为明初仓促官修，本纪保存大量实录条目但有编次与纪年问题；行省执行、数字、族称及对外关系须以条格、多语文书和对方史料互校。。研究元代行省、赋役、河工与多语文书研究可以补问材料未直接说出的空间与社会差异。

\noindent\eventanchor{SYD-1302-YUAN-163}{〔己〕未，以諸王真童誣告濟南王，謫置劉國傑軍中自效}\textbf{元。}
“〔己〕未，以諸王真童誣告濟南王，謫置劉國傑軍中自效”是元在1302年（元大德6年；公元换算据《元史》本纪纪年）可追查的节点。它提示命令、环境、宗教、迁徙或地方组织如何进入政治过程；影响的范围和速度仍要避免由一条材料外推。

账本把这一节点放在“本纪年条记录政令或机构处置；实施背景须参照制度材料。”与“可在同年及后续政令中追踪；条文不单独证明地方执行。”之间；两者提示前后压力，并不把时间相邻写成单一因果。

本条首先回到《元史》卷020〈大德六年〉；《元史》为明初仓促官修，本纪保存大量实录条目但有编次与纪年问题；行省执行、数字、族称及对外关系须以条格、多语文书和对方史料互校。 因而元代行省、赋役、河工与多语文书研究不是装饰性的书目，而是检验这一叙述边界的路径。

\noindent\eventanchor{SYD-1302-YUAN-235}{安南國以馴象二及朱砂來獻}\textbf{元与安南。}
“安南國以馴象二及朱砂來獻”是元与安南在1302年（元大德6年；公元换算据《元史》本纪纪年）可追查的节点。它提示命令、环境、宗教、迁徙或地方组织如何进入政治过程；影响的范围和速度仍要避免由一条材料外推。

此前的条件是“本纪年条提供日期和中枢可见行动；前因不据单条补定。”，此后的可见影响是“可回查同年及后续条目，地方影响仍须另证。”。这里的“影响”须按地点、机构和材料种类分别衡量。

这一叙述依据《元史》卷020〈大德六年〉；《安南志略》及当地碑刻材料（互校），但《元史》为明初仓促官修，本纪保存大量实录条目但有编次与纪年问题；行省执行、数字、族称及对外关系须以条格、多语文书和对方史料互校。。研究元代行省、赋役、河工与多语文书研究可以补问材料未直接说出的空间与社会差异。

\noindent\eventanchor{SYD-1302-YUAN-307}{八月甲子，詔御史臺凡有司婚姻、土田文案，遇赦依例檢覆}\textbf{元。}
“八月甲子，詔御史臺凡有司婚姻、土田文案，遇赦依例檢覆”发生时，元必须把账面规则转成仓储、市场、航道、文书和地方中介的工作。正因这些环节并不齐整，政策颁行与实际效果应分别追踪。

账本把这一节点放在“本纪年条记录政令或机构处置；实施背景须参照制度材料。”与“可在同年及后续政令中追踪；条文不单独证明地方执行。”之间；两者提示前后压力，并不把时间相邻写成单一因果。

本条首先回到《元史》卷020〈大德六年〉；《元史》为明初仓促官修，本纪保存大量实录条目但有编次与纪年问题；行省执行、数字、族称及对外关系须以条格、多语文书和对方史料互校。 因而元代行省、赋役、河工与多语文书研究不是装饰性的书目，而是检验这一叙述边界的路径。

\subsection{1303年（元大德7年；公元换算据《元史》本纪纪年）}

\noindent\eventanchor{SYD-1303-YUAN-020}{罷江南財賦總管府及提舉司}\textbf{元。}
在1303年（元大德7年；公元换算据《元史》本纪纪年）的可见材料中，元留下“罷江南財賦總管府及提舉司”。这一动作把具体人群、机构或地方条件带进年序，但一项记录不能替未被记载者概括整体反应。

其前提记为“本纪年条记录政令或机构处置；实施背景须参照制度材料。”；其后续是“可在同年及后续政令中追踪；条文不单独证明地方执行。”。这是一条待后续材料检验的事件链，而不是对所有行动者意图的替写。

可核的材料入口是《元史》卷021〈大德七年〉。《元史》为明初仓促官修，本纪保存大量实录条目但有编次与纪年问题；行省执行、数字、族称及对外关系须以条格、多语文书和对方史料互校。 就此而言，元代行省、赋役、河工与多语文书研究有助于把年序同区域、制度或物质材料并读。

\noindent\eventanchor{SYD-1303-YUAN-092}{敕賜鈔五百錠，其稅課依例輸官}\textbf{元。}
在1303年（元大德7年；公元换算据《元史》本纪纪年）的可见材料中，元留下“敕賜鈔五百錠，其稅課依例輸官”。这一动作把具体人群、机构或地方条件带进年序，但一项记录不能替未被记载者概括整体反应。

此前的条件是“本纪所记财用、赈济或征发属于特定报告场景。”，此后的可见影响是“可与食货志、文书和物质材料比较，不外推为全境总量。”。这里的“影响”须按地点、机构和材料种类分别衡量。

这一叙述依据《元史》卷021〈大德七年〉，但《元史》为明初仓促官修，本纪保存大量实录条目但有编次与纪年问题；行省执行、数字、族称及对外关系须以条格、多语文书和对方史料互校。。研究元代行省、赋役、河工与多语文书研究可以补问材料未直接说出的空间与社会差异。

\noindent\eventanchor{SYD-1303-YUAN-164}{戊子，弛太原、平陽酒禁}\textbf{元。}
在1303年（元大德7年；公元换算据《元史》本纪纪年）的可见材料中，元留下“戊子，弛太原、平陽酒禁”。这一动作把具体人群、机构或地方条件带进年序，但一项记录不能替未被记载者概括整体反应。

账本把这一节点放在“本纪年条记录政令或机构处置；实施背景须参照制度材料。”与“可在同年及后续政令中追踪；条文不单独证明地方执行。”之间；两者提示前后压力，并不把时间相邻写成单一因果。

本条首先回到《元史》卷021〈大德七年〉；《元史》为明初仓促官修，本纪保存大量实录条目但有编次与纪年问题；行省执行、数字、族称及对外关系须以条格、多语文书和对方史料互校。 因而元代行省、赋役、河工与多语文书研究不是装饰性的书目，而是检验这一叙述边界的路径。

\noindent\eventanchor{SYD-1303-YUAN-236}{中書省臣言：「宜捕送其所，令省、臺、宣政院遣官雜治}\textbf{元。}
在1303年（元大德7年；公元换算据《元史》本纪纪年）的可见材料中，元留下“中書省臣言：「宜捕送其所，令省、臺、宣政院遣官雜治”。这一动作把具体人群、机构或地方条件带进年序，但一项记录不能替未被记载者概括整体反应。

此前的条件是“本纪年条记录政令或机构处置；实施背景须参照制度材料。”，此后的可见影响是“可在同年及后续政令中追踪；条文不单独证明地方执行。”。这里的“影响”须按地点、机构和材料种类分别衡量。

这一叙述依据《元史》卷021〈大德七年〉，但《元史》为明初仓促官修，本纪保存大量实录条目但有编次与纪年问题；行省执行、数字、族称及对外关系须以条格、多语文书和对方史料互校。。研究元代行省、赋役、河工与多语文书研究可以补问材料未直接说出的空间与社会差异。

\noindent\eventanchor{SYD-1303-YUAN-308}{安西轉運司於常課外增算五萬七千四百錠，人賜衣一襲，以勸其功}\textbf{元。}
在1303年（元大德7年；公元换算据《元史》本纪纪年）的可见材料中，元留下“安西轉運司於常課外增算五萬七千四百錠，人賜衣一襲，以勸其功”。这一动作把具体人群、机构或地方条件带进年序，但一项记录不能替未被记载者概括整体反应。

账本把这一节点放在“本纪年条提供日期和中枢可见行动；前因不据单条补定。”与“可回查同年及后续条目，地方影响仍须另证。”之间；两者提示前后压力，并不把时间相邻写成单一因果。

本条首先回到《元史》卷021〈大德七年〉；《元史》为明初仓促官修，本纪保存大量实录条目但有编次与纪年问题；行省执行、数字、族称及对外关系须以条格、多语文书和对方史料互校。 因而元代行省、赋役、河工与多语文书研究不是装饰性的书目，而是检验这一叙述边界的路径。

\subsection{1304年（元大德8年；公元换算据《元史》本纪纪年）}

\noindent\eventanchor{SYD-1304-YUAN-021}{敕軍民逃奴有獲者即付其主，主在他所者，赴所在官司給之，仍追逃奴鈔充獲者賞}\textbf{元。}
元的1304年（元大德8年；公元换算据《元史》本纪纪年）记录写到“敕軍民逃奴有獲者即付其主，主在他所者，赴所在官司給之，仍追逃奴鈔充獲者賞”。我们可以据此辨认一个可见的选择或压力，而不能把零散的官方记载、碑刻或后出叙事拼成无缝的社会全景。

其前提记为“本纪所记财用、赈济或征发属于特定报告场景。”；其后续是“可与食货志、文书和物质材料比较，不外推为全境总量。”。这是一条待后续材料检验的事件链，而不是对所有行动者意图的替写。

可核的材料入口是《元史》卷021〈大德八年〉。《元史》为明初仓促官修，本纪保存大量实录条目但有编次与纪年问题；行省执行、数字、族称及对外关系须以条格、多语文书和对方史料互校。 就此而言，元代行省、赋役、河工与多语文书研究有助于把年序同区域、制度或物质材料并读。

\noindent\eventanchor{SYD-1304-YUAN-093}{庚子，以永平、清、滄、柳林屯田被水，其逋租及民貸食者皆勿徵}\textbf{元。}
元在1304年（元大德8年；公元换算据《元史》本纪纪年）留下的这条记录是“庚子，以永平、清、滄、柳林屯田被水，其逋租及民貸食者皆勿徵”。它照亮资源如何被命名和调度，也暴露材料的盲点：数字、额度或禁令若无地方文书和物质材料互证，不能被写成全境同速的现实。

此前的条件是“本纪所记财用、赈济或征发属于特定报告场景。”，此后的可见影响是“可与食货志、文书和物质材料比较，不外推为全境总量。”。这里的“影响”须按地点、机构和材料种类分别衡量。

这一叙述依据《元史》卷021〈大德八年〉，但《元史》为明初仓促官修，本纪保存大量实录条目但有编次与纪年问题；行省执行、数字、族称及对外关系须以条格、多语文书和对方史料互校。。研究元代行省、赋役、河工与多语文书研究可以补问材料未直接说出的空间与社会差异。

\noindent\eventanchor{SYD-1304-YUAN-165}{以平陽、太原去歲地大震，免其稅課一年}\textbf{元。}
元的1304年（元大德8年；公元换算据《元史》本纪纪年）记录写到“以平陽、太原去歲地大震，免其稅課一年”。我们可以据此辨认一个可见的选择或压力，而不能把零散的官方记载、碑刻或后出叙事拼成无缝的社会全景。

账本把这一节点放在“本纪所记财用、赈济或征发属于特定报告场景。”与“可与食货志、文书和物质材料比较，不外推为全境总量。”之间；两者提示前后压力，并不把时间相邻写成单一因果。

本条首先回到《元史》卷021〈大德八年〉；《元史》为明初仓促官修，本纪保存大量实录条目但有编次与纪年问题；行省执行、数字、族称及对外关系须以条格、多语文书和对方史料互校。 因而元代行省、赋役、河工与多语文书研究不是装饰性的书目，而是检验这一叙述边界的路径。

\noindent\eventanchor{SYD-1304-YUAN-237}{自滎澤至睢州，築河防十有八所，給其夫鈔人十貫}\textbf{元。}
元的1304年（元大德8年；公元换算据《元史》本纪纪年）记录写到“自滎澤至睢州，築河防十有八所，給其夫鈔人十貫”。我们可以据此辨认一个可见的选择或压力，而不能把零散的官方记载、碑刻或后出叙事拼成无缝的社会全景。

此前的条件是“本纪所记财用、赈济或征发属于特定报告场景。”，此后的可见影响是“可与食货志、文书和物质材料比较，不外推为全境总量。”。这里的“影响”须按地点、机构和材料种类分别衡量。

这一叙述依据《元史》卷021〈大德八年〉，但《元史》为明初仓促官修，本纪保存大量实录条目但有编次与纪年问题；行省执行、数字、族称及对外关系须以条格、多语文书和对方史料互校。。研究元代行省、赋役、河工与多语文书研究可以补问材料未直接说出的空间与社会差异。

\noindent\eventanchor{SYD-1304-YUAN-309}{八年春正月己未，以災異故，詔天下恤民隱，省刑罰}\textbf{元。}
元的1304年（元大德8年；公元换算据《元史》本纪纪年）记录写到“八年春正月己未，以災異故，詔天下恤民隱，省刑罰”。我们可以据此辨认一个可见的选择或压力，而不能把零散的官方记载、碑刻或后出叙事拼成无缝的社会全景。

账本把这一节点放在“本纪年条记录政令或机构处置；实施背景须参照制度材料。”与“可在同年及后续政令中追踪；条文不单独证明地方执行。”之间；两者提示前后压力，并不把时间相邻写成单一因果。

本条首先回到《元史》卷021〈大德八年〉；《元史》为明初仓促官修，本纪保存大量实录条目但有编次与纪年问题；行省执行、数字、族称及对外关系须以条格、多语文书和对方史料互校。 因而元代行省、赋役、河工与多语文书研究不是装饰性的书目，而是检验这一叙述边界的路径。

\subsection{1305年（元大德9年；公元换算据《元史》本纪纪年）}

\noindent\eventanchor{SYD-1305-YUAN-022}{戊戌，詔芍陂、洪澤等屯田為豪右占據者，悉令輸租}\textbf{元。}
“戊戌，詔芍陂、洪澤等屯田為豪右占據者，悉令輸租”发生时，元必须把账面规则转成仓储、市场、航道、文书和地方中介的工作。正因这些环节并不齐整，政策颁行与实际效果应分别追踪。

其前提记为“本纪年条记录政令或机构处置；实施背景须参照制度材料。”；其后续是“可在同年及后续政令中追踪；条文不单独证明地方执行。”。这是一条待后续材料检验的事件链，而不是对所有行动者意图的替写。

可核的材料入口是《元史》卷021〈大德九年〉。《元史》为明初仓促官修，本纪保存大量实录条目但有编次与纪年问题；行省执行、数字、族称及对外关系须以条格、多语文书和对方史料互校。 就此而言，元代行省、赋役、河工与多语文书研究有助于把年序同区域、制度或物质材料并读。

\noindent\eventanchor{SYD-1305-YUAN-094}{八月乙亥朔，省孛可孫冗員}\textbf{元。}
“八月乙亥朔，省孛可孫冗員”是元在1305年（元大德9年；公元换算据《元史》本纪纪年）可追查的节点。它提示命令、环境、宗教、迁徙或地方组织如何进入政治过程；影响的范围和速度仍要避免由一条材料外推。

此前的条件是“本纪年条记录政令或机构处置；实施背景须参照制度材料。”，此后的可见影响是“可在同年及后续政令中追踪；条文不单独证明地方执行。”。这里的“影响”须按地点、机构和材料种类分别衡量。

这一叙述依据《元史》卷021〈大德九年〉，但《元史》为明初仓促官修，本纪保存大量实录条目但有编次与纪年问题；行省执行、数字、族称及对外关系须以条格、多语文书和对方史料互校。。研究元代行省、赋役、河工与多语文书研究可以补问材料未直接说出的空间与社会差异。

\noindent\eventanchor{SYD-1305-YUAN-166}{罷福建蒙古字提舉司及醫學提舉司}\textbf{元。}
“罷福建蒙古字提舉司及醫學提舉司”是元在1305年（元大德9年；公元换算据《元史》本纪纪年）可追查的节点。它提示命令、环境、宗教、迁徙或地方组织如何进入政治过程；影响的范围和速度仍要避免由一条材料外推。

账本把这一节点放在“本纪年条记录政令或机构处置；实施背景须参照制度材料。”与“可在同年及后续政令中追踪；条文不单独证明地方执行。”之间；两者提示前后压力，并不把时间相邻写成单一因果。

本条首先回到《元史》卷021〈大德九年〉；《元史》为明初仓促官修，本纪保存大量实录条目但有编次与纪年问题；行省执行、数字、族称及对外关系须以条格、多语文书和对方史料互校。 因而元代行省、赋役、河工与多语文书研究不是装饰性的书目，而是检验这一叙述边界的路径。

\noindent\eventanchor{SYD-1305-YUAN-238}{寶慶路饑，發粟五千石賑之}\textbf{元。}
“寶慶路饑，發粟五千石賑之”是元在1305年（元大德9年；公元换算据《元史》本纪纪年）可追查的节点。它提示命令、环境、宗教、迁徙或地方组织如何进入政治过程；影响的范围和速度仍要避免由一条材料外推。

此前的条件是“本纪所记财用、赈济或征发属于特定报告场景。”，此后的可见影响是“可与食货志、文书和物质材料比较，不外推为全境总量。”。这里的“影响”须按地点、机构和材料种类分别衡量。

这一叙述依据《元史》卷021〈大德九年〉，但《元史》为明初仓促官修，本纪保存大量实录条目但有编次与纪年问题；行省执行、数字、族称及对外关系须以条格、多语文书和对方史料互校。。研究元代行省、赋役、河工与多语文书研究可以补问材料未直接说出的空间与社会差异。

\noindent\eventanchor{SYD-1305-YUAN-310}{孛可孫專治芻粟，初惟數人，後以各位增入，遂至繁宂}\textbf{元。}
“孛可孫專治芻粟，初惟數人，後以各位增入，遂至繁宂”是元在1305年（元大德9年；公元换算据《元史》本纪纪年）可追查的节点。它提示命令、环境、宗教、迁徙或地方组织如何进入政治过程；影响的范围和速度仍要避免由一条材料外推。

从账本的前后项看，“本纪年条提供日期和中枢可见行动；前因不据单条补定。”提供了背景压力；“可回查同年及后续条目，地方影响仍须另证。”标出后来需要追踪的变化，二者之间仍存在未被材料填满的环节。

核心线索为《元史》卷021〈大德九年〉。《元史》为明初仓促官修，本纪保存大量实录条目但有编次与纪年问题；行省执行、数字、族称及对外关系须以条格、多语文书和对方史料互校。 现代研究可从元代行省、赋役、河工与多语文书研究继续追问。

\subsection{1306年（元大德10年；公元换算据《元史》本纪纪年）}

\noindent\eventanchor{SYD-1306-YUAN-023}{繼有旨，給瀛國公衣糧發遣之，唯與〔𤍟〕勿行}\textbf{元。}
“繼有旨，給瀛國公衣糧發遣之，唯與〔𤍟〕勿行”是元在1306年（元大德10年；公元换算据《元史》本纪纪年）可追查的节点。它提示命令、环境、宗教、迁徙或地方组织如何进入政治过程；影响的范围和速度仍要避免由一条材料外推。

其前提记为“本纪所记财用、赈济或征发属于特定报告场景。”；其后续是“可与食货志、文书和物质材料比较，不外推为全境总量。”。这是一条待后续材料检验的事件链，而不是对所有行动者意图的替写。

可核的材料入口是《元史》卷012〈本纪第12卷本年条〉。《元史》为明初仓促官修，本纪保存大量实录条目但有编次与纪年问题；行省执行、数字、族称及对外关系须以条格、多语文书和对方史料互校。 就此而言，元代行省、赋役、河工与多语文书研究有助于把年序同区域、制度或物质材料并读。

\noindent\eventanchor{SYD-1306-YUAN-095}{有旨降民還之有司，征討所得，籍其數量，賜臣下有功者}\textbf{元。}
“有旨降民還之有司，征討所得，籍其數量，賜臣下有功者”把元置于军事行动的可核时点。战报能够记下攻守、入城或撤离，却很少完整保留沿途征发、居民去留和补给损耗；因此不能把一方的胜败叙事延伸成所有地方的经验。

此前的条件是“本纪从朝廷视角记录军政行动；出兵与战果需分辨。”，此后的可见影响是“后续路线、控制与损失应与对方记录和遗址材料复核。”。这里的“影响”须按地点、机构和材料种类分别衡量。

这一叙述依据《元史》卷012〈本纪第12卷本年条〉，但《元史》为明初仓促官修，本纪保存大量实录条目但有编次与纪年问题；行省执行、数字、族称及对外关系须以条格、多语文书和对方史料互校。。研究元代行省、赋役、河工与多语文书研究可以补问材料未直接说出的空间与社会差异。

\noindent\eventanchor{SYD-1306-YUAN-167}{阿里海牙復鎮遠軍，發軍千人戍守，以其地與西川行省接，就以隸焉}\textbf{元。}
“阿里海牙復鎮遠軍，發軍千人戍守，以其地與西川行省接，就以隸焉”把元置于军事行动的可核时点。战报能够记下攻守、入城或撤离，却很少完整保留沿途征发、居民去留和补给损耗；因此不能把一方的胜败叙事延伸成所有地方的经验。

账本把这一节点放在“本纪年条记录政令或机构处置；实施背景须参照制度材料。”与“可在同年及后续政令中追踪；条文不单独证明地方执行。”之间；两者提示前后压力，并不把时间相邻写成单一因果。

本条首先回到《元史》卷012〈本纪第12卷本年条〉；《元史》为明初仓促官修，本纪保存大量实录条目但有编次与纪年问题；行省执行、数字、族称及对外关系须以条格、多语文书和对方史料互校。 因而元代行省、赋役、河工与多语文书研究不是装饰性的书目，而是检验这一叙述边界的路径。

\noindent\eventanchor{SYD-1306-YUAN-239}{八月丁亥朔，給乾山造船軍匠冬衣，及新附軍鈔}\textbf{元。}
“八月丁亥朔，給乾山造船軍匠冬衣，及新附軍鈔”是元在1306年（元大德10年；公元换算据《元史》本纪纪年）可追查的节点。它提示命令、环境、宗教、迁徙或地方组织如何进入政治过程；影响的范围和速度仍要避免由一条材料外推。

此前的条件是“本纪所记财用、赈济或征发属于特定报告场景。”，此后的可见影响是“可与食货志、文书和物质材料比较，不外推为全境总量。”。这里的“影响”须按地点、机构和材料种类分别衡量。

这一叙述依据《元史》卷012〈本纪第12卷本年条〉，但《元史》为明初仓促官修，本纪保存大量实录条目但有编次与纪年问题；行省执行、数字、族称及对外关系须以条格、多语文书和对方史料互校。。研究元代行省、赋役、河工与多语文书研究可以补问材料未直接说出的空间与社会差异。

\noindent\eventanchor{SYD-1306-YUAN-311}{罷都功德使脫烈，其修設佛事妄費官物，皆徵還之}\textbf{元。}
在“罷都功德使脫烈，其修設佛事妄費官物，皆徵還之”中，元的选择经由使节、礼物、边市或协商被记录下来。跨境文本的观察位置不同，因而应把可确认的往来与对动机、服从程度的推断分开。

从账本的前后项看，“本纪年条记录政令或机构处置；实施背景须参照制度材料。”提供了背景压力；“可在同年及后续政令中追踪；条文不单独证明地方执行。”标出后来需要追踪的变化，二者之间仍存在未被材料填满的环节。

核心线索为《元史》卷012〈本纪第12卷本年条〉。《元史》为明初仓促官修，本纪保存大量实录条目但有编次与纪年问题；行省执行、数字、族称及对外关系须以条格、多语文书和对方史料互校。 现代研究可从元代行省、赋役、河工与多语文书研究继续追问。

\subsection{1307年（元大德11年；公元换算据《元史》本纪纪年）}

\noindent\eventanchor{SYD-1307-YUAN-024}{併省江淮、雲南州郡}\textbf{元。}
在1307年（元大德11年；公元换算据《元史》本纪纪年）的可见材料中，元留下“併省江淮、雲南州郡”。这一动作把具体人群、机构或地方条件带进年序，但一项记录不能替未被记载者概括整体反应。

其前提记为“本纪年条记录政令或机构处置；实施背景须参照制度材料。”；其后续是“可在同年及后续政令中追踪；条文不单独证明地方执行。”。这是一条待后续材料检验的事件链，而不是对所有行动者意图的替写。

可核的材料入口是《元史》卷012〈本纪第12卷本年条〉。《元史》为明初仓促官修，本纪保存大量实录条目但有编次与纪年问题；行省执行、数字、族称及对外关系须以条格、多语文书和对方史料互校。 就此而言，元代行省、赋役、河工与多语文书研究有助于把年序同区域、制度或物质材料并读。

\noindent\eventanchor{SYD-1307-YUAN-096}{己未，吏部尚書劉好禮以吉利吉思風俗事宜來上}\textbf{元。}
在1307年（元大德11年；公元换算据《元史》本纪纪年）的可见材料中，元留下“己未，吏部尚書劉好禮以吉利吉思風俗事宜來上”。这一动作把具体人群、机构或地方条件带进年序，但一项记录不能替未被记载者概括整体反应。

此前的条件是“本纪年条提供日期和中枢可见行动；前因不据单条补定。”，此后的可见影响是“可回查同年及后续条目，地方影响仍须另证。”。这里的“影响”须按地点、机构和材料种类分别衡量。

这一叙述依据《元史》卷012〈本纪第12卷本年条〉，但《元史》为明初仓促官修，本纪保存大量实录条目但有编次与纪年问题；行省执行、数字、族称及对外关系须以条格、多语文书和对方史料互校。。研究元代行省、赋役、河工与多语文书研究可以补问材料未直接说出的空间与社会差异。

\noindent\eventanchor{SYD-1307-YUAN-168}{甲寅，降太醫院為尚醫監，改給銅印}\textbf{元。}
1307年（元大德11年；公元换算据《元史》本纪纪年），元的行动落在“甲寅，降太醫院為尚醫監，改給銅印”所指的战事与通道上。可见的不是抽象的推进：攻守双方必须让人员、消息与补给穿过具体的城防、道路或水路；账本所列的结果只说明这一节点改变了下一步选择。

账本把这一节点放在“本纪年条记录政令或机构处置；实施背景须参照制度材料。”与“可在同年及后续政令中追踪；条文不单独证明地方执行。”之间；两者提示前后压力，并不把时间相邻写成单一因果。

本条首先回到《元史》卷012〈本纪第12卷本年条〉；《元史》为明初仓促官修，本纪保存大量实录条目但有编次与纪年问题；行省执行、数字、族称及对外关系须以条格、多语文书和对方史料互校。 因而元代行省、赋役、河工与多语文书研究不是装饰性的书目，而是检验这一叙述边界的路径。

\noindent\eventanchor{SYD-1307-YUAN-240}{前後衞軍自願征日本者，命選留五衞漢軍千餘，其新附軍令悉行}\textbf{元。}
在1307年（元大德11年；公元换算据《元史》本纪纪年）的可见材料中，元留下“前後衞軍自願征日本者，命選留五衞漢軍千餘，其新附軍令悉行”。这一动作把具体人群、机构或地方条件带进年序，但一项记录不能替未被记载者概括整体反应。

其前提记为“本纪年条记录政令或机构处置；实施背景须参照制度材料。”；其后续是“可在同年及后续政令中追踪；条文不单独证明地方执行。”。这是一条待后续材料检验的事件链，而不是对所有行动者意图的替写。

可核的材料入口是《元史》卷012〈本纪第12卷本年条〉。《元史》为明初仓促官修，本纪保存大量实录条目但有编次与纪年问题；行省执行、数字、族称及对外关系须以条格、多语文书和对方史料互校。 就此而言，元代行省、赋役、河工与多语文书研究有助于把年序同区域、制度或物质材料并读。

\noindent\eventanchor{SYD-1307-YUAN-312}{時各衞議先遣七人，而以三人自代，從之}\textbf{元。}
在1307年（元大德11年；公元换算据《元史》本纪纪年）的可见材料中，元留下“時各衞議先遣七人，而以三人自代，從之”。这一动作把具体人群、机构或地方条件带进年序，但一项记录不能替未被记载者概括整体反应。

从账本的前后项看，“本纪年条提供日期和中枢可见行动；前因不据单条补定。”提供了背景压力；“可回查同年及后续条目，地方影响仍须另证。”标出后来需要追踪的变化，二者之间仍存在未被材料填满的环节。

核心线索为《元史》卷012〈本纪第12卷本年条〉。《元史》为明初仓促官修，本纪保存大量实录条目但有编次与纪年问题；行省执行、数字、族称及对外关系须以条格、多语文书和对方史料互校。 现代研究可从元代行省、赋役、河工与多语文书研究继续追问。

\subsection{1307年}

\noindent\eventanchor{YUAN-1307-EVENT-182}{海山即位}\textbf{元。}
元的“海山即位”并非只改换一个统治者。任官、赏赐、军权与信息通路要在随后重新接合；材料较能确认先后，较难替当事人写出唯一动机。

账本把这一节点放在“元中后期继承、财政与地方危机中前任统治的结束与宫廷、部族或官僚的权力安排相互交织，促成“海山即位”这一继承节点。”与“继承并不自动消除既有矛盾；“海山即位”后的任官、军权和对外策略需由后续条目追踪。”之间；两者提示前后压力，并不把时间相邻写成单一因果。

本条首先回到《元史》相关本纪；《元朝秘史》、波斯文史籍或《高丽史》对应叙事；《元典章》《通制条格》、多语碑刻与文书（据事核）；《元史》明初速修，纪年与人名有异同；蒙古大汗政权不等同所有汗国，征服、赋役和身份分类须按地区及材料语言分列。 因而蒙古帝国、元朝行省财政、多语文书与区域社会研究不是装饰性的书目，而是检验这一叙述边界的路径。

\subsection{1308年（元至大1年；公元换算据《元史》本纪纪年）}

\noindent\eventanchor{SYD-1308-YUAN-025}{帝曰：「礬課遣人覈實}\textbf{元。}
元的1308年（元至大1年；公元换算据《元史》本纪纪年）记录写到“帝曰：「礬課遣人覈實”。我们可以据此辨认一个可见的选择或压力，而不能把零散的官方记载、碑刻或后出叙事拼成无缝的社会全景。

其前提记为“本纪年条提供日期和中枢可见行动；前因不据单条补定。”；其后续是“可回查同年及后续条目，地方影响仍须另证。”。这是一条待后续材料检验的事件链，而不是对所有行动者意图的替写。

可核的材料入口是《元史》卷019〈本纪第19卷本年条〉。《元史》为明初仓促官修，本纪保存大量实录条目但有编次与纪年问题；行省执行、数字、族称及对外关系须以条格、多语文书和对方史料互校。 就此而言，元代行省、赋役、河工与多语文书研究有助于把年序同区域、制度或物质材料并读。

\noindent\eventanchor{SYD-1308-YUAN-097}{甲戌，詔民間馬牛羊，百取其一，羊不滿百者亦取之，惟色目人及數乃取}\textbf{元。}
元的1308年（元至大1年；公元换算据《元史》本纪纪年）记录写到“甲戌，詔民間馬牛羊，百取其一，羊不滿百者亦取之，惟色目人及數乃取”。我们可以据此辨认一个可见的选择或压力，而不能把零散的官方记载、碑刻或后出叙事拼成无缝的社会全景。

此前的条件是“本纪年条记录政令或机构处置；实施背景须参照制度材料。”，此后的可见影响是“可在同年及后续政令中追踪；条文不单独证明地方执行。”。这里的“影响”须按地点、机构和材料种类分别衡量。

这一叙述依据《元史》卷019〈本纪第19卷本年条〉，但《元史》为明初仓促官修，本纪保存大量实录条目但有编次与纪年问题；行省执行、数字、族称及对外关系须以条格、多语文书和对方史料互校。。研究元代行省、赋役、河工与多语文书研究可以补问材料未直接说出的空间与社会差异。

\noindent\eventanchor{SYD-1308-YUAN-169}{禁諸王、駙馬并權豪，毋奪民田，其獻田者有刑}\textbf{元。}
元在1308年（元至大1年；公元换算据《元史》本纪纪年）留下的这条记录是“禁諸王、駙馬并權豪，毋奪民田，其獻田者有刑”。它照亮资源如何被命名和调度，也暴露材料的盲点：数字、额度或禁令若无地方文书和物质材料互证，不能被写成全境同速的现实。

账本把这一节点放在“本纪年条记录政令或机构处置；实施背景须参照制度材料。”与“可在同年及后续政令中追踪；条文不单独证明地方执行。”之间；两者提示前后压力，并不把时间相邻写成单一因果。

本条首先回到《元史》卷019〈本纪第19卷本年条〉；《元史》为明初仓促官修，本纪保存大量实录条目但有编次与纪年问题；行省执行、数字、族称及对外关系须以条格、多语文书和对方史料互校。 因而元代行省、赋役、河工与多语文书研究不是装饰性的书目，而是检验这一叙述边界的路径。

\noindent\eventanchor{SYD-1308-YUAN-241}{辛未，徙江浙行省拔都軍萬人戍潭州，潭州以南軍移戍郴州}\textbf{元。}
元的1308年（元至大1年；公元换算据《元史》本纪纪年）记录写到“辛未，徙江浙行省拔都軍萬人戍潭州，潭州以南軍移戍郴州”。我们可以据此辨认一个可见的选择或压力，而不能把零散的官方记载、碑刻或后出叙事拼成无缝的社会全景。

其前提记为“本纪年条记录政令或机构处置；实施背景须参照制度材料。”；其后续是“可在同年及后续政令中追踪；条文不单独证明地方执行。”。这是一条待后续材料检验的事件链，而不是对所有行动者意图的替写。

可核的材料入口是《元史》卷019〈本纪第19卷本年条〉。《元史》为明初仓促官修，本纪保存大量实录条目但有编次与纪年问题；行省执行、数字、族称及对外关系须以条格、多语文书和对方史料互校。 就此而言，元代行省、赋役、河工与多语文书研究有助于把年序同区域、制度或物质材料并读。

\noindent\eventanchor{SYD-1308-YUAN-313}{又賜緬王弟撒邦巴一珠虎符，酋領阿散三珠虎符，從者金符及金幣，遣之}\textbf{元。}
元的1308年（元至大1年；公元换算据《元史》本纪纪年）记录写到“又賜緬王弟撒邦巴一珠虎符，酋領阿散三珠虎符，從者金符及金幣，遣之”。我们可以据此辨认一个可见的选择或压力，而不能把零散的官方记载、碑刻或后出叙事拼成无缝的社会全景。

从账本的前后项看，“本纪年条提供日期和中枢可见行动；前因不据单条补定。”提供了背景压力；“可回查同年及后续条目，地方影响仍须另证。”标出后来需要追踪的变化，二者之间仍存在未被材料填满的环节。

核心线索为《元史》卷019〈本纪第19卷本年条〉。《元史》为明初仓促官修，本纪保存大量实录条目但有编次与纪年问题；行省执行、数字、族称及对外关系须以条格、多语文书和对方史料互校。 现代研究可从元代行省、赋役、河工与多语文书研究继续追问。

\subsection{1309年（元至大2年；公元换算据《元史》本纪纪年）}

\noindent\eventanchor{SYD-1309-YUAN-026}{頒行至大銀鈔，詔曰：「昔我世祖皇帝既登大寶，始造中統交鈔，以便民用，歲久法隳，亦既更張，印造至元寶鈔}\textbf{元。}
“頒行至大銀鈔，詔曰：「昔我世祖皇帝既登大寶，始造中統交鈔，以便民用，歲久法隳，亦既更張，印造至元寶鈔”是元在1309年（元至大2年；公元换算据《元史》本纪纪年）可追查的节点。它提示命令、环境、宗教、迁徙或地方组织如何进入政治过程；影响的范围和速度仍要避免由一条材料外推。

其前提记为“本纪年条记录政令或机构处置；实施背景须参照制度材料。”；其后续是“可在同年及后续政令中追踪；条文不单独证明地方执行。”。这是一条待后续材料检验的事件链，而不是对所有行动者意图的替写。

可核的材料入口是《元史》卷023〈至大二年〉。《元史》为明初仓促官修，本纪保存大量实录条目但有编次与纪年问题；行省执行、数字、族称及对外关系须以条格、多语文书和对方史料互校。 就此而言，元代行省、赋役、河工与多语文书研究有助于把年序同区域、制度或物质材料并读。

\noindent\eventanchor{SYD-1309-YUAN-098}{大都立資國院，秩正二品}\textbf{元。}
元的“大都立資國院，秩正二品”并非只改换一个统治者。任官、赏赐、军权与信息通路要在随后重新接合；材料较能确认先后，较难替当事人写出唯一动机。

此前的条件是“本纪年条提供日期和中枢可见行动；前因不据单条补定。”，此后的可见影响是“可回查同年及后续条目，地方影响仍须另证。”。这里的“影响”须按地点、机构和材料种类分别衡量。

这一叙述依据《元史》卷023〈至大二年〉，但《元史》为明初仓促官修，本纪保存大量实录条目但有编次与纪年问题；行省执行、数字、族称及对外关系须以条格、多语文书和对方史料互校。。研究元代行省、赋役、河工与多语文书研究可以补问材料未直接说出的空间与社会差异。

\noindent\eventanchor{SYD-1309-YUAN-170}{帝曰：「大事也，其速擇使復齎璽書往招諭}\textbf{元。}
在“帝曰：「大事也，其速擇使復齎璽書往招諭”中，元的选择经由使节、礼物、边市或协商被记录下来。跨境文本的观察位置不同，因而应把可确认的往来与对动机、服从程度的推断分开。

从账本的前后项看，“本纪年条记录政令或机构处置；实施背景须参照制度材料。”提供了背景压力；“可在同年及后续政令中追踪；条文不单独证明地方执行。”标出后来需要追踪的变化，二者之间仍存在未被材料填满的环节。

核心线索为《元史》卷023〈至大二年〉。《元史》为明初仓促官修，本纪保存大量实录条目但有编次与纪年问题；行省执行、数字、族称及对外关系须以条格、多语文书和对方史料互校。 现代研究可从元代行省、赋役、河工与多语文书研究继续追问。

\noindent\eventanchor{SYD-1309-YUAN-242}{甲申，賜寧肅王脫脫金印}\textbf{元。}
“甲申，賜寧肅王脫脫金印”是元在1309年（元至大2年；公元换算据《元史》本纪纪年）可追查的节点。它提示命令、环境、宗教、迁徙或地方组织如何进入政治过程；影响的范围和速度仍要避免由一条材料外推。

其前提记为“本纪年条提供日期和中枢可见行动；前因不据单条补定。”；其后续是“可回查同年及后续条目，地方影响仍须另证。”。这是一条待后续材料检验的事件链，而不是对所有行动者意图的替写。

可核的材料入口是《元史》卷023〈至大二年〉。《元史》为明初仓促官修，本纪保存大量实录条目但有编次与纪年问题；行省执行、数字、族称及对外关系须以条格、多语文书和对方史料互校。 就此而言，元代行省、赋役、河工与多语文书研究有助于把年序同区域、制度或物质材料并读。

\noindent\eventanchor{SYD-1309-YUAN-314}{壬午，江南行臺劾：「平章政事教化，詐言家貧，冒受賜貨物，折鈔二萬錠}\textbf{元。}
“壬午，江南行臺劾：「平章政事教化，詐言家貧，冒受賜貨物，折鈔二萬錠”是元在1309年（元至大2年；公元换算据《元史》本纪纪年）可追查的节点。它提示命令、环境、宗教、迁徙或地方组织如何进入政治过程；影响的范围和速度仍要避免由一条材料外推。

从账本的前后项看，“本纪所记财用、赈济或征发属于特定报告场景。”提供了背景压力；“可与食货志、文书和物质材料比较，不外推为全境总量。”标出后来需要追踪的变化，二者之间仍存在未被材料填满的环节。

核心线索为《元史》卷023〈至大二年〉。《元史》为明初仓促官修，本纪保存大量实录条目但有编次与纪年问题；行省执行、数字、族称及对外关系须以条格、多语文书和对方史料互校。 现代研究可从元代行省、赋役、河工与多语文书研究继续追问。

\subsection{1310年（元至大3年；公元换算据《元史》本纪纪年）}

\noindent\eventanchor{SYD-1310-YUAN-027}{汴梁、懷孟、衞輝、彰德、歸德、汝寧、南陽、河南等路蝗}\textbf{元。}
“汴梁、懷孟、衞輝、彰德、歸德、汝寧、南陽、河南等路蝗”是元在1310年（元至大3年；公元换算据《元史》本纪纪年）可追查的节点。它提示命令、环境、宗教、迁徙或地方组织如何进入政治过程；影响的范围和速度仍要避免由一条材料外推。

此前的条件是“本纪保留灾害或工程的报告地点与朝廷处置；灾情范围需另证。”，此后的可见影响是“可与地方文书、河工和环境材料比较，不将单点记录外推为全域因果。”。这里的“影响”须按地点、机构和材料种类分别衡量。

这一叙述依据《元史》卷023〈至大三年〉，但《元史》为明初仓促官修，本纪保存大量实录条目但有编次与纪年问题；行省执行、数字、族称及对外关系须以条格、多语文书和对方史料互校。。研究元代行省、赋役、河工与多语文书研究可以补问材料未直接说出的空间与社会差异。

\noindent\eventanchor{SYD-1310-YUAN-099}{敕臺臣遣官往鞫，毋徇私情}\textbf{元。}
“敕臺臣遣官往鞫，毋徇私情”是元在1310年（元至大3年；公元换算据《元史》本纪纪年）可追查的节点。它提示命令、环境、宗教、迁徙或地方组织如何进入政治过程；影响的范围和速度仍要避免由一条材料外推。

从账本的前后项看，“本纪年条提供日期和中枢可见行动；前因不据单条补定。”提供了背景压力；“可回查同年及后续条目，地方影响仍须另证。”标出后来需要追踪的变化，二者之间仍存在未被材料填满的环节。

核心线索为《元史》卷023〈至大三年〉。《元史》为明初仓促官修，本纪保存大量实录条目但有编次与纪年问题；行省执行、数字、族称及对外关系须以条格、多语文书和对方史料互校。 现代研究可从元代行省、赋役、河工与多语文书研究继续追问。

\noindent\eventanchor{SYD-1310-YUAN-171}{從之，仍賜鈔三千錠賑其部貧民}\textbf{元。}
“從之，仍賜鈔三千錠賑其部貧民”是元在1310年（元至大3年；公元换算据《元史》本纪纪年）可追查的节点。它提示命令、环境、宗教、迁徙或地方组织如何进入政治过程；影响的范围和速度仍要避免由一条材料外推。

账本把这一节点放在“本纪所记财用、赈济或征发属于特定报告场景。”与“可与食货志、文书和物质材料比较，不外推为全境总量。”之间；两者提示前后压力，并不把时间相邻写成单一因果。

本条首先回到《元史》卷023〈至大三年〉；《元史》为明初仓促官修，本纪保存大量实录条目但有编次与纪年问题；行省执行、数字、族称及对外关系须以条格、多语文书和对方史料互校。 因而元代行省、赋役、河工与多语文书研究不是装饰性的书目，而是检验这一叙述边界的路径。

\noindent\eventanchor{SYD-1310-YUAN-243}{帝曰：「世祖謀慮深遠若是，待諸王朝會，頒賞既畢，卿等備述其故，然後與之，使彼知愧}\textbf{元。}
在“帝曰：「世祖謀慮深遠若是，待諸王朝會，頒賞既畢，卿等備述其故，然後與之，使彼知愧”中，元的选择经由使节、礼物、边市或协商被记录下来。跨境文本的观察位置不同，因而应把可确认的往来与对动机、服从程度的推断分开。

此前的条件是“本纪年条提供日期和中枢可见行动；前因不据单条补定。”，此后的可见影响是“可回查同年及后续条目，地方影响仍须另证。”。这里的“影响”须按地点、机构和材料种类分别衡量。

这一叙述依据《元史》卷023〈至大三年〉，但《元史》为明初仓促官修，本纪保存大量实录条目但有编次与纪年问题；行省执行、数字、族称及对外关系须以条格、多语文书和对方史料互校。。研究元代行省、赋役、河工与多语文书研究可以补问材料未直接说出的空间与社会差异。

\noindent\eventanchor{SYD-1310-YUAN-315}{封僧亦憐真乞烈思為文國公，賜金印}\textbf{元。}
“封僧亦憐真乞烈思為文國公，賜金印”是元在1310年（元至大3年；公元换算据《元史》本纪纪年）可追查的节点。它提示命令、环境、宗教、迁徙或地方组织如何进入政治过程；影响的范围和速度仍要避免由一条材料外推。

账本把这一节点放在“本纪年条提供日期和中枢可见行动；前因不据单条补定。”与“可回查同年及后续条目，地方影响仍须另证。”之间；两者提示前后压力，并不把时间相邻写成单一因果。

本条首先回到《元史》卷023〈至大三年〉；《元史》为明初仓促官修，本纪保存大量实录条目但有编次与纪年问题；行省执行、数字、族称及对外关系须以条格、多语文书和对方史料互校。 因而元代行省、赋役、河工与多语文书研究不是装饰性的书目，而是检验这一叙述边界的路径。

\subsection{1311年（元至大4年；公元换算据《元史》本纪纪年）}

\noindent\eventanchor{SYD-1311-YUAN-028}{壬午，靈駕發引，葬起輦谷，從諸帝陵}\textbf{元。}
在1311年（元至大4年；公元换算据《元史》本纪纪年）的可见材料中，元留下“壬午，靈駕發引，葬起輦谷，從諸帝陵”。这一动作把具体人群、机构或地方条件带进年序，但一项记录不能替未被记载者概括整体反应。

此前的条件是“本纪年条提供日期和中枢可见行动；前因不据单条补定。”，此后的可见影响是“可回查同年及后续条目，地方影响仍须另证。”。这里的“影响”须按地点、机构和材料种类分别衡量。

这一叙述依据《元史》卷023〈至大四年〉，但《元史》为明初仓促官修，本纪保存大量实录条目但有编次与纪年问题；行省执行、数字、族称及对外关系须以条格、多语文书和对方史料互校。。研究元代行省、赋役、河工与多语文书研究可以补问材料未直接说出的空间与社会差异。

\noindent\eventanchor{SYD-1311-YUAN-100}{四年春正月癸酉朔，帝不豫，免朝賀，大赦天下}\textbf{元。}
1311年（元至大4年；公元换算据《元史》本纪纪年），“四年春正月癸酉朔，帝不豫，免朝賀，大赦天下”使元与相邻行动者的关系进入可核的外交时点。名号、贡赐和盟约是各方政治语言的一部分；它们能够确认交涉，却不能单凭措辞判定实际隶属或边地生活。

从账本的前后项看，“本纪年条提供日期和中枢可见行动；前因不据单条补定。”提供了背景压力；“可回查同年及后续条目，地方影响仍须另证。”标出后来需要追踪的变化，二者之间仍存在未被材料填满的环节。

核心线索为《元史》卷023〈至大四年〉。《元史》为明初仓促官修，本纪保存大量实录条目但有编次与纪年问题；行省执行、数字、族称及对外关系须以条格、多语文书和对方史料互校。 现代研究可从元代行省、赋役、河工与多语文书研究继续追问。

\noindent\eventanchor{SYD-1311-YUAN-172}{武宗當富有之大業，慨然欲創治改法而有為，故其封爵太盛，而遙授之官眾，錫賚太隆，而泛賞之恩溥，至元、大德之政，於是稍有變更云}\textbf{元。}
在1311年（元至大4年；公元换算据《元史》本纪纪年）的可见材料中，元留下“武宗當富有之大業，慨然欲創治改法而有為，故其封爵太盛，而遙授之官眾，錫賚太隆，而泛賞之恩溥，至元、大德之政，於是稍有變更云”。这一动作把具体人群、机构或地方条件带进年序，但一项记录不能替未被记载者概括整体反应。

账本把这一节点放在“本纪年条记录政令或机构处置；实施背景须参照制度材料。”与“可在同年及后续政令中追踪；条文不单独证明地方执行。”之间；两者提示前后压力，并不把时间相邻写成单一因果。

本条首先回到《元史》卷023〈至大四年〉；《元史》为明初仓促官修，本纪保存大量实录条目但有编次与纪年问题；行省执行、数字、族称及对外关系须以条格、多语文书和对方史料互校。 因而元代行省、赋役、河工与多语文书研究不是装饰性的书目，而是检验这一叙述边界的路径。

\noindent\eventanchor{SYD-1311-YUAN-244}{敕：「諸王、駙馬戶在縉山、懷來、永興縣者，與民均服徭役}\textbf{元。}
在1311年（元至大4年；公元换算据《元史》本纪纪年）的可见材料中，元留下“敕：「諸王、駙馬戶在縉山、懷來、永興縣者，與民均服徭役”。这一动作把具体人群、机构或地方条件带进年序，但一项记录不能替未被记载者概括整体反应。

此前的条件是“本纪年条提供日期和中枢可见行动；前因不据单条补定。”，此后的可见影响是“可回查同年及后续条目，地方影响仍须另证。”。这里的“影响”须按地点、机构和材料种类分别衡量。

这一叙述依据《元史》卷024〈至大四年〉，但《元史》为明初仓促官修，本纪保存大量实录条目但有编次与纪年问题；行省执行、数字、族称及对外关系须以条格、多语文书和对方史料互校。。研究元代行省、赋役、河工与多语文书研究可以补问材料未直接说出的空间与社会差异。

\noindent\eventanchor{SYD-1311-YUAN-316}{敕直省舍人以其半給事殿庭，半聽中書差遣}\textbf{元。}
在1311年（元至大4年；公元换算据《元史》本纪纪年）的可见材料中，元留下“敕直省舍人以其半給事殿庭，半聽中書差遣”。这一动作把具体人群、机构或地方条件带进年序，但一项记录不能替未被记载者概括整体反应。

账本把这一节点放在“本纪年条记录政令或机构处置；实施背景须参照制度材料。”与“可在同年及后续政令中追踪；条文不单独证明地方执行。”之间；两者提示前后压力，并不把时间相邻写成单一因果。

本条首先回到《元史》卷024〈至大四年〉；《元史》为明初仓促官修，本纪保存大量实录条目但有编次与纪年问题；行省执行、数字、族称及对外关系须以条格、多语文书和对方史料互校。 因而元代行省、赋役、河工与多语文书研究不是装饰性的书目，而是检验这一叙述边界的路径。

\subsection{1311年}

\noindent\eventanchor{YUAN-1311-EVENT-183}{爱育黎拔力八达即位}\textbf{元。}
在1311年，元所记“爱育黎拔力八达即位”把权力安排暴露在继承压力下。后续稳定与否取决于资源和支持网络如何被重新调度，不能由即位仪式本身预先断言。

此前的条件是“元中后期继承、财政与地方危机中前任统治的结束与宫廷、部族或官僚的权力安排相互交织，促成“爱育黎拔力八达即位”这一继承节点。”，此后的可见影响是“继承并不自动消除既有矛盾；“爱育黎拔力八达即位”后的任官、军权和对外策略需由后续条目追踪。”。这里的“影响”须按地点、机构和材料种类分别衡量。

这一叙述依据《元史》相关本纪；《元朝秘史》、波斯文史籍或《高丽史》对应叙事；《元典章》《通制条格》、多语碑刻与文书（据事核），但《元史》明初速修，纪年与人名有异同；蒙古大汗政权不等同所有汗国，征服、赋役和身份分类须按地区及材料语言分列。。研究蒙古帝国、元朝行省财政、多语文书与区域社会研究可以补问材料未直接说出的空间与社会差异。

\subsection{1312年（元皇庆1年；公元换算据《元史》本纪纪年）}

\noindent\eventanchor{SYD-1312-YUAN-029}{己卯，馬八兒國遣使貢珍珠、異寶、縑段}\textbf{元。}
元于1312年（元皇庆1年；公元换算据《元史》本纪纪年）处理“己卯，馬八兒國遣使貢珍珠、異寶、縑段”时，面对的是道路、翻译、礼仪和资源交换共同构成的关系网。书面称谓可见，地方执行和持续性则仍须由后续材料检验。

此前的条件是“本纪记录使节、贡赐或往来，不等于稳定统属关系。”，此后的可见影响是“可与对方编年和交通、文书材料核对其后续落实。”。这里的“影响”须按地点、机构和材料种类分别衡量。

这一叙述依据《元史》卷013〈本纪第13卷本年条〉，但《元史》为明初仓促官修，本纪保存大量实录条目但有编次与纪年问题；行省执行、数字、族称及对外关系须以条格、多语文书和对方史料互校。。研究元代行省、赋役、河工与多语文书研究可以补问材料未直接说出的空间与社会差异。

\noindent\eventanchor{SYD-1312-YUAN-101}{己未，罷雲南都元帥府，所管軍民隸行省}\textbf{元。}
元的1312年（元皇庆1年；公元换算据《元史》本纪纪年）记录写到“己未，罷雲南都元帥府，所管軍民隸行省”。我们可以据此辨认一个可见的选择或压力，而不能把零散的官方记载、碑刻或后出叙事拼成无缝的社会全景。

从账本的前后项看，“本纪年条记录政令或机构处置；实施背景须参照制度材料。”提供了背景压力；“可在同年及后续政令中追踪；条文不单独证明地方执行。”标出后来需要追踪的变化，二者之间仍存在未被材料填满的环节。

核心线索为《元史》卷013〈本纪第13卷本年条〉。《元史》为明初仓促官修，本纪保存大量实录条目但有编次与纪年问题；行省执行、数字、族称及对外关系须以条格、多语文书和对方史料互校。 现代研究可从元代行省、赋役、河工与多语文书研究继续追问。

\noindent\eventanchor{SYD-1312-YUAN-173}{己未，荊湖占城行省言：「忽都虎、忽馬兒等將兵征占城，前鋒舟師至舒眉蓮港不知所向，令萬戶劉君慶進軍次新州，獲占蠻，始知我軍已還矣}\textbf{元。}
在1312年（元皇庆1年；公元换算据《元史》本纪纪年）的材料中，元面对的是己未，荊湖占城行省言：「忽都虎、忽馬兒等將兵征占城，前鋒舟師至舒眉蓮港不知所向，令萬戶劉君慶進軍次新州，獲占蠻，始知我軍已還矣。这里应先看行动所依赖的关隘、城镇、河道或粮道，再讨论政治后果；军队抵达某处，并不等于它已能稳定征收、安置或控制周边。

账本把这一节点放在“本纪年条记录政令或机构处置；实施背景须参照制度材料。”与“可在同年及后续政令中追踪；条文不单独证明地方执行。”之间；两者提示前后压力，并不把时间相邻写成单一因果。

本条首先回到《元史》卷013〈本纪第13卷本年条〉；《元史》为明初仓促官修，本纪保存大量实录条目但有编次与纪年问题；行省执行、数字、族称及对外关系须以条格、多语文书和对方史料互校。 因而元代行省、赋役、河工与多语文书研究不是装饰性的书目，而是检验这一叙述边界的路径。

\noindent\eventanchor{SYD-1312-YUAN-245}{詔依所言汰選，毋徇私情}\textbf{元。}
元的1312年（元皇庆1年；公元换算据《元史》本纪纪年）记录写到“詔依所言汰選，毋徇私情”。我们可以据此辨认一个可见的选择或压力，而不能把零散的官方记载、碑刻或后出叙事拼成无缝的社会全景。

此前的条件是“本纪年条记录政令或机构处置；实施背景须参照制度材料。”，此后的可见影响是“可在同年及后续政令中追踪；条文不单独证明地方执行。”。这里的“影响”须按地点、机构和材料种类分别衡量。

这一叙述依据《元史》卷013〈本纪第13卷本年条〉，但《元史》为明初仓促官修，本纪保存大量实录条目但有编次与纪年问题；行省执行、数字、族称及对外关系须以条格、多语文书和对方史料互校。。研究元代行省、赋役、河工与多语文书研究可以补问材料未直接说出的空间与社会差异。

\noindent\eventanchor{SYD-1312-YUAN-317}{詔以中書右丞、行省事忙兀台為江淮等處行中書省平章政事，其行省左丞忽剌出、蒲壽庚，參政管如德分省泉州}\textbf{元。}
元的1312年（元皇庆1年；公元换算据《元史》本纪纪年）记录写到“詔以中書右丞、行省事忙兀台為江淮等處行中書省平章政事，其行省左丞忽剌出、蒲壽庚，參政管如德分省泉州”。我们可以据此辨认一个可见的选择或压力，而不能把零散的官方记载、碑刻或后出叙事拼成无缝的社会全景。

账本把这一节点放在“本纪年条记录政令或机构处置；实施背景须参照制度材料。”与“可在同年及后续政令中追踪；条文不单独证明地方执行。”之间；两者提示前后压力，并不把时间相邻写成单一因果。

本条首先回到《元史》卷013〈本纪第13卷本年条〉；《元史》为明初仓促官修，本纪保存大量实录条目但有编次与纪年问题；行省执行、数字、族称及对外关系须以条格、多语文书和对方史料互校。 因而元代行省、赋役、河工与多语文书研究不是装饰性的书目，而是检验这一叙述边界的路径。

\subsection{1313年（元皇庆2年；公元换算据《元史》本纪纪年）}

\noindent\eventanchor{SYD-1313-YUAN-030}{安童與諸老臣議世榮所行，當罷者罷之，更者更之，其所用人實無罪者，朕自裁決}\textbf{元。}
“安童與諸老臣議世榮所行，當罷者罷之，更者更之，其所用人實無罪者，朕自裁決”是元在1313年（元皇庆2年；公元换算据《元史》本纪纪年）可追查的节点。它提示命令、环境、宗教、迁徙或地方组织如何进入政治过程；影响的范围和速度仍要避免由一条材料外推。

其前提记为“本纪年条记录政令或机构处置；实施背景须参照制度材料。”；其后续是“可在同年及后续政令中追踪；条文不单独证明地方执行。”。这是一条待后续材料检验的事件链，而不是对所有行动者意图的替写。

可核的材料入口是《元史》卷013〈本纪第13卷本年条〉。《元史》为明初仓促官修，本纪保存大量实录条目但有编次与纪年问题；行省执行、数字、族称及对外关系须以条格、多语文书和对方史料互校。 就此而言，元代行省、赋役、河工与多语文书研究有助于把年序同区域、制度或物质材料并读。

\noindent\eventanchor{SYD-1313-YUAN-102}{丙子，真蠟、占城貢樂工十人及藥材、鰐魚皮諸物}\textbf{元。}
“丙子，真蠟、占城貢樂工十人及藥材、鰐魚皮諸物”是元在1313年（元皇庆2年；公元换算据《元史》本纪纪年）可追查的节点。它提示命令、环境、宗教、迁徙或地方组织如何进入政治过程；影响的范围和速度仍要避免由一条材料外推。

从账本的前后项看，“本纪年条提供日期和中枢可见行动；前因不据单条补定。”提供了背景压力；“可回查同年及后续条目，地方影响仍须另证。”标出后来需要追踪的变化，二者之间仍存在未被材料填满的环节。

核心线索为《元史》卷013〈本纪第13卷本年条〉。《元史》为明初仓促官修，本纪保存大量实录条目但有编次与纪年问题；行省执行、数字、族称及对外关系须以条格、多语文书和对方史料互校。 现代研究可从元代行省、赋役、河工与多语文书研究继续追问。

\noindent\eventanchor{SYD-1313-YUAN-174}{癸亥，敕以麥朮丁所行清潔，與安童治省事}\textbf{元。}
“癸亥，敕以麥朮丁所行清潔，與安童治省事”是元在1313年（元皇庆2年；公元换算据《元史》本纪纪年）可追查的节点。它提示命令、环境、宗教、迁徙或地方组织如何进入政治过程；影响的范围和速度仍要避免由一条材料外推。

账本把这一节点放在“本纪年条记录政令或机构处置；实施背景须参照制度材料。”与“可在同年及后续政令中追踪；条文不单独证明地方执行。”之间；两者提示前后压力，并不把时间相邻写成单一因果。

本条首先回到《元史》卷013〈本纪第13卷本年条〉；《元史》为明初仓促官修，本纪保存大量实录条目但有编次与纪年问题；行省执行、数字、族称及对外关系须以条格、多语文书和对方史料互校。 因而元代行省、赋役、河工与多语文书研究不是装饰性的书目，而是检验这一叙述边界的路径。

\noindent\eventanchor{SYD-1313-YUAN-246}{江淮行省以戰船千艘習水戰江中}\textbf{元。}
“江淮行省以戰船千艘習水戰江中”是元在1313年（元皇庆2年；公元换算据《元史》本纪纪年）可追查的节点。它提示命令、环境、宗教、迁徙或地方组织如何进入政治过程；影响的范围和速度仍要避免由一条材料外推。

此前的条件是“本纪年条记录政令或机构处置；实施背景须参照制度材料。”，此后的可见影响是“可在同年及后续政令中追踪；条文不单独证明地方执行。”。这里的“影响”须按地点、机构和材料种类分别衡量。

这一叙述依据《元史》卷013〈本纪第13卷本年条〉，但《元史》为明初仓促官修，本纪保存大量实录条目但有编次与纪年问题；行省执行、数字、族称及对外关系须以条格、多语文书和对方史料互校。。研究元代行省、赋役、河工与多语文书研究可以补问材料未直接说出的空间与社会差异。

\noindent\eventanchor{SYD-1313-YUAN-318}{盧世榮請罷福建行中書省，立宣慰司，隸江西行中書省}\textbf{元。}
元的“盧世榮請罷福建行中書省，立宣慰司，隸江西行中書省”并非只改换一个统治者。任官、赏赐、军权与信息通路要在随后重新接合；材料较能确认先后，较难替当事人写出唯一动机。

账本把这一节点放在“本纪年条记录政令或机构处置；实施背景须参照制度材料。”与“可在同年及后续政令中追踪；条文不单独证明地方执行。”之间；两者提示前后压力，并不把时间相邻写成单一因果。

本条首先回到《元史》卷013〈本纪第13卷本年条〉；《元史》为明初仓促官修，本纪保存大量实录条目但有编次与纪年问题；行省执行、数字、族称及对外关系须以条格、多语文书和对方史料互校。 因而元代行省、赋役、河工与多语文书研究不是装饰性的书目，而是检验这一叙述边界的路径。

\subsection{1313年}

\noindent\eventanchor{YUAN-1313-EVENT-184}{元廷恢复科举制度}\textbf{元。}
元的1313年记录写到“元廷恢复科举制度”。我们可以据此辨认一个可见的选择或压力，而不能把零散的官方记载、碑刻或后出叙事拼成无缝的社会全景。

从账本的前后项看，“元中后期继承、财政与地方危机对中枢资源、交通或行政协同的要求，推动了“元廷恢复科举制度”这一制度或空间安排。”提供了背景压力；““元廷恢复科举制度”提供新的治理框架或资源通道，但颁行、复申与不同地区的实际执行须分开核验。”标出后来需要追踪的变化，二者之间仍存在未被材料填满的环节。

核心线索为《元史》相关本纪；《元朝秘史》、波斯文史籍或《高丽史》对应叙事；《元典章》《通制条格》、多语碑刻与文书（据事核）。《元史》明初速修，纪年与人名有异同；蒙古大汗政权不等同所有汗国，征服、赋役和身份分类须按地区及材料语言分列。 现代研究可从蒙古帝国、元朝行省财政、多语文书与区域社会研究继续追问。

\subsection{1314年（元延祐1年；公元换算据《元史》本纪纪年）}

\noindent\eventanchor{SYD-1314-YUAN-031}{大寧路地震，有聲如雷}\textbf{元。}
在1314年（元延祐1年；公元换算据《元史》本纪纪年）的可见材料中，元留下“大寧路地震，有聲如雷”。这一动作把具体人群、机构或地方条件带进年序，但一项记录不能替未被记载者概括整体反应。

其前提记为“本纪保留灾害或工程的报告地点与朝廷处置；灾情范围需另证。”；其后续是“可与地方文书、河工和环境材料比较，不将单点记录外推为全域因果。”。这是一条待后续材料检验的事件链，而不是对所有行动者意图的替写。

可核的材料入口是《元史》卷025〈延祐元年〉。《元史》为明初仓促官修，本纪保存大量实录条目但有编次与纪年问题；行省执行、数字、族称及对外关系须以条格、多语文书和对方史料互校。 就此而言，元代行省、赋役、河工与多语文书研究有助于把年序同区域、制度或物质材料并读。

\noindent\eventanchor{SYD-1314-YUAN-103}{帝曰：「如更不悛，則罷不敍}\textbf{元。}
在1314年（元延祐1年；公元换算据《元史》本纪纪年）的可见材料中，元留下“帝曰：「如更不悛，則罷不敍”。这一动作把具体人群、机构或地方条件带进年序，但一项记录不能替未被记载者概括整体反应。

从账本的前后项看，“本纪年条记录政令或机构处置；实施背景须参照制度材料。”提供了背景压力；“可在同年及后续政令中追踪；条文不单独证明地方执行。”标出后来需要追踪的变化，二者之间仍存在未被材料填满的环节。

核心线索为《元史》卷025〈延祐元年〉。《元史》为明初仓促官修，本纪保存大量实录条目但有编次与纪年问题；行省执行、数字、族称及对外关系须以条格、多语文书和对方史料互校。 现代研究可从元代行省、赋役、河工与多语文书研究继续追问。

\noindent\eventanchor{SYD-1314-YUAN-175}{衡州、郴州、興國、永州路、耒陽州饑，發廪減價賑糶}\textbf{元。}
在1314年（元延祐1年；公元换算据《元史》本纪纪年）的可见材料中，元留下“衡州、郴州、興國、永州路、耒陽州饑，發廪減價賑糶”。这一动作把具体人群、机构或地方条件带进年序，但一项记录不能替未被记载者概括整体反应。

账本把这一节点放在“本纪所记财用、赈济或征发属于特定报告场景。”与“可与食货志、文书和物质材料比较，不外推为全境总量。”之间；两者提示前后压力，并不把时间相邻写成单一因果。

本条首先回到《元史》卷025〈延祐元年〉；《元史》为明初仓促官修，本纪保存大量实录条目但有编次与纪年问题；行省执行、数字、族称及对外关系须以条格、多语文书和对方史料互校。 因而元代行省、赋役、河工与多语文书研究不是装饰性的书目，而是检验这一叙述边界的路径。

\noindent\eventanchor{SYD-1314-YUAN-247}{遣官括淮民所佃閑田不輸稅者}\textbf{元。}
1314年（元延祐1年；公元换算据《元史》本纪纪年）的“遣官括淮民所佃閑田不輸稅者”将元的财政或生产安排推到可见处。制度文本能说明官府打算如何收取、转运或调节，却不能直接等同于不同地区的实收、流通或家庭负担。

此前的条件是“本纪所记财用、赈济或征发属于特定报告场景。”，此后的可见影响是“可与食货志、文书和物质材料比较，不外推为全境总量。”。这里的“影响”须按地点、机构和材料种类分别衡量。

这一叙述依据《元史》卷025〈延祐元年〉，但《元史》为明初仓促官修，本纪保存大量实录条目但有编次与纪年问题；行省执行、数字、族称及对外关系须以条格、多语文书和对方史料互校。。研究元代行省、赋役、河工与多语文书研究可以补问材料未直接说出的空间与社会差异。

\noindent\eventanchor{SYD-1314-YUAN-319}{申飭內侍及諸司隔越中書奏請之禁}\textbf{元。}
在1314年（元延祐1年；公元换算据《元史》本纪纪年）的可见材料中，元留下“申飭內侍及諸司隔越中書奏請之禁”。这一动作把具体人群、机构或地方条件带进年序，但一项记录不能替未被记载者概括整体反应。

账本把这一节点放在“本纪年条记录政令或机构处置；实施背景须参照制度材料。”与“可在同年及后续政令中追踪；条文不单独证明地方执行。”之间；两者提示前后压力，并不把时间相邻写成单一因果。

本条首先回到《元史》卷025〈延祐元年〉；《元史》为明初仓促官修，本纪保存大量实录条目但有编次与纪年问题；行省执行、数字、族称及对外关系须以条格、多语文书和对方史料互校。 因而元代行省、赋役、河工与多语文书研究不是装饰性的书目，而是检验这一叙述边界的路径。

\subsection{1315年（元延祐2年；公元换算据《元史》本纪纪年）}

\noindent\eventanchor{SYD-1315-YUAN-032}{帝曰：「此朕之愆，豈卿等所致，其復乃職}\textbf{元。}
元的1315年（元延祐2年；公元换算据《元史》本纪纪年）记录写到“帝曰：「此朕之愆，豈卿等所致，其復乃職”。我们可以据此辨认一个可见的选择或压力，而不能把零散的官方记载、碑刻或后出叙事拼成无缝的社会全景。

其前提记为“本纪年条记录政令或机构处置；实施背景须参照制度材料。”；其后续是“可在同年及后续政令中追踪；条文不单独证明地方执行。”。这是一条待后续材料检验的事件链，而不是对所有行动者意图的替写。

可核的材料入口是《元史》卷025〈延祐二年〉。《元史》为明初仓促官修，本纪保存大量实录条目但有编次与纪年问题；行省执行、数字、族称及对外关系须以条格、多语文书和对方史料互校。 就此而言，元代行省、赋役、河工与多语文书研究有助于把年序同区域、制度或物质材料并读。

\noindent\eventanchor{SYD-1315-YUAN-104}{詔明言其事當行者以聞}\textbf{元。}
元的1315年（元延祐2年；公元换算据《元史》本纪纪年）记录写到“詔明言其事當行者以聞”。我们可以据此辨认一个可见的选择或压力，而不能把零散的官方记载、碑刻或后出叙事拼成无缝的社会全景。

从账本的前后项看，“本纪年条记录政令或机构处置；实施背景须参照制度材料。”提供了背景压力；“可在同年及后续政令中追踪；条文不单独证明地方执行。”标出后来需要追踪的变化，二者之间仍存在未被材料填满的环节。

核心线索为《元史》卷025〈延祐二年〉。《元史》为明初仓促官修，本纪保存大量实录条目但有编次与纪年问题；行省执行、数字、族称及对外关系须以条格、多语文书和对方史料互校。 现代研究可从元代行省、赋役、河工与多语文书研究继续追问。

\noindent\eventanchor{SYD-1315-YUAN-176}{八月丙戌，贛州賊蔡五九陷汀州寧花縣，僭稱王號，詔遣江浙行省平章張驢等率兵討之}\textbf{元。}
在1315年（元延祐2年；公元换算据《元史》本纪纪年）的材料中，元面对的是八月丙戌，贛州賊蔡五九陷汀州寧花縣，僭稱王號，詔遣江浙行省平章張驢等率兵討之。这里应先看行动所依赖的关隘、城镇、河道或粮道，再讨论政治后果；军队抵达某处，并不等于它已能稳定征收、安置或控制周边。

账本把这一节点放在“本纪年条记录政令或机构处置；实施背景须参照制度材料。”与“可在同年及后续政令中追踪；条文不单独证明地方执行。”之间；两者提示前后压力，并不把时间相邻写成单一因果。

本条首先回到《元史》卷025〈延祐二年〉；《元史》为明初仓促官修，本纪保存大量实录条目但有编次与纪年问题；行省执行、数字、族称及对外关系须以条格、多语文书和对方史料互校。 因而元代行省、赋役、河工与多语文书研究不是装饰性的书目，而是检验这一叙述边界的路径。

\noindent\eventanchor{SYD-1315-YUAN-248}{丙申，賜諸王納忽答兒金五十兩、銀二百五十兩、鈔五百錠}\textbf{元。}
元的1315年（元延祐2年；公元换算据《元史》本纪纪年）记录写到“丙申，賜諸王納忽答兒金五十兩、銀二百五十兩、鈔五百錠”。我们可以据此辨认一个可见的选择或压力，而不能把零散的官方记载、碑刻或后出叙事拼成无缝的社会全景。

此前的条件是“本纪所记财用、赈济或征发属于特定报告场景。”，此后的可见影响是“可与食货志、文书和物质材料比较，不外推为全境总量。”。这里的“影响”须按地点、机构和材料种类分别衡量。

这一叙述依据《元史》卷025〈延祐二年〉，但《元史》为明初仓促官修，本纪保存大量实录条目但有编次与纪年问题；行省执行、数字、族称及对外关系须以条格、多语文书和对方史料互校。。研究元代行省、赋役、河工与多语文书研究可以补问材料未直接说出的空间与社会差异。

\noindent\eventanchor{SYD-1315-YUAN-320}{丙午，封諸王察八兒為汝寧王}\textbf{元。}
元的1315年（元延祐2年；公元换算据《元史》本纪纪年）记录写到“丙午，封諸王察八兒為汝寧王”。我们可以据此辨认一个可见的选择或压力，而不能把零散的官方记载、碑刻或后出叙事拼成无缝的社会全景。

从账本的前后项看，“本纪年条提供日期和中枢可见行动；前因不据单条补定。”提供了背景压力；“可回查同年及后续条目，地方影响仍须另证。”标出后来需要追踪的变化，二者之间仍存在未被材料填满的环节。

核心线索为《元史》卷025〈延祐二年〉。《元史》为明初仓促官修，本纪保存大量实录条目但有编次与纪年问题；行省执行、数字、族称及对外关系须以条格、多语文书和对方史料互校。 现代研究可从元代行省、赋役、河工与多语文书研究继续追问。

\subsection{1316年（元延祐3年；公元换算据《元史》本纪纪年）}

\noindent\eventanchor{SYD-1316-YUAN-033}{罷江南財賦總管府及提舉司}\textbf{元。}
“罷江南財賦總管府及提舉司”是元在1316年（元延祐3年；公元换算据《元史》本纪纪年）可追查的节点。它提示命令、环境、宗教、迁徙或地方组织如何进入政治过程；影响的范围和速度仍要避免由一条材料外推。

其前提记为“本纪年条记录政令或机构处置；实施背景须参照制度材料。”；其后续是“可在同年及后续政令中追踪；条文不单独证明地方执行。”。这是一条待后续材料检验的事件链，而不是对所有行动者意图的替写。

可核的材料入口是《元史》卷021〈本纪第21卷本年条〉。《元史》为明初仓促官修，本纪保存大量实录条目但有编次与纪年问题；行省执行、数字、族称及对外关系须以条格、多语文书和对方史料互校。 就此而言，元代行省、赋役、河工与多语文书研究有助于把年序同区域、制度或物质材料并读。

\noindent\eventanchor{SYD-1316-YUAN-105}{敕賜鈔五百錠，其稅課依例輸官}\textbf{元。}
“敕賜鈔五百錠，其稅課依例輸官”是元在1316年（元延祐3年；公元换算据《元史》本纪纪年）可追查的节点。它提示命令、环境、宗教、迁徙或地方组织如何进入政治过程；影响的范围和速度仍要避免由一条材料外推。

从账本的前后项看，“本纪所记财用、赈济或征发属于特定报告场景。”提供了背景压力；“可与食货志、文书和物质材料比较，不外推为全境总量。”标出后来需要追踪的变化，二者之间仍存在未被材料填满的环节。

核心线索为《元史》卷021〈本纪第21卷本年条〉。《元史》为明初仓促官修，本纪保存大量实录条目但有编次与纪年问题；行省执行、数字、族称及对外关系须以条格、多语文书和对方史料互校。 现代研究可从元代行省、赋役、河工与多语文书研究继续追问。

\noindent\eventanchor{SYD-1316-YUAN-177}{敕軍民逃奴有獲者即付其主，主在他所者，赴所在官司給之，仍追逃奴鈔充獲者賞}\textbf{元。}
“敕軍民逃奴有獲者即付其主，主在他所者，赴所在官司給之，仍追逃奴鈔充獲者賞”是元在1316年（元延祐3年；公元换算据《元史》本纪纪年）可追查的节点。它提示命令、环境、宗教、迁徙或地方组织如何进入政治过程；影响的范围和速度仍要避免由一条材料外推。

账本把这一节点放在“本纪所记财用、赈济或征发属于特定报告场景。”与“可与食货志、文书和物质材料比较，不外推为全境总量。”之间；两者提示前后压力，并不把时间相邻写成单一因果。

本条首先回到《元史》卷021〈本纪第21卷本年条〉；《元史》为明初仓促官修，本纪保存大量实录条目但有编次与纪年问题；行省执行、数字、族称及对外关系须以条格、多语文书和对方史料互校。 因而元代行省、赋役、河工与多语文书研究不是装饰性的书目，而是检验这一叙述边界的路径。

\noindent\eventanchor{SYD-1316-YUAN-249}{庚子，以永平、清、滄、柳林屯田被水，其逋租及民貸食者皆勿徵}\textbf{元。}
“庚子，以永平、清、滄、柳林屯田被水，其逋租及民貸食者皆勿徵”发生时，元必须把账面规则转成仓储、市场、航道、文书和地方中介的工作。正因这些环节并不齐整，政策颁行与实际效果应分别追踪。

此前的条件是“本纪所记财用、赈济或征发属于特定报告场景。”，此后的可见影响是“可与食货志、文书和物质材料比较，不外推为全境总量。”。这里的“影响”须按地点、机构和材料种类分别衡量。

这一叙述依据《元史》卷021〈本纪第21卷本年条〉，但《元史》为明初仓促官修，本纪保存大量实录条目但有编次与纪年问题；行省执行、数字、族称及对外关系须以条格、多语文书和对方史料互校。。研究元代行省、赋役、河工与多语文书研究可以补问材料未直接说出的空间与社会差异。

\noindent\eventanchor{SYD-1316-YUAN-321}{戊戌，詔芍陂、洪澤等屯田為豪右占據者，悉令輸租}\textbf{元。}
“戊戌，詔芍陂、洪澤等屯田為豪右占據者，悉令輸租”发生时，元必须把账面规则转成仓储、市场、航道、文书和地方中介的工作。正因这些环节并不齐整，政策颁行与实际效果应分别追踪。

从账本的前后项看，“本纪年条记录政令或机构处置；实施背景须参照制度材料。”提供了背景压力；“可在同年及后续政令中追踪；条文不单独证明地方执行。”标出后来需要追踪的变化，二者之间仍存在未被材料填满的环节。

核心线索为《元史》卷021〈本纪第21卷本年条〉。《元史》为明初仓促官修，本纪保存大量实录条目但有编次与纪年问题；行省执行、数字、族称及对外关系须以条格、多语文书和对方史料互校。 现代研究可从元代行省、赋役、河工与多语文书研究继续追问。

\subsection{1317年（元延祐4年；公元换算据《元史》本纪纪年）}

\noindent\eventanchor{SYD-1317-YUAN-034}{敕免雲南從征交趾蒙古軍屯田租}\textbf{元。}
1317年（元延祐4年；公元换算据《元史》本纪纪年）的“敕免雲南從征交趾蒙古軍屯田租”将元的财政或生产安排推到可见处。制度文本能说明官府打算如何收取、转运或调节，却不能直接等同于不同地区的实收、流通或家庭负担。

其前提记为“本纪年条提供日期和中枢可见行动；前因不据单条补定。”；其后续是“可回查同年及后续条目，地方影响仍须另证。”。这是一条待后续材料检验的事件链，而不是对所有行动者意图的替写。

可核的材料入口是《元史》卷014〈本纪第14卷本年条〉。《元史》为明初仓促官修，本纪保存大量实录条目但有编次与纪年问题；行省执行、数字、族称及对外关系须以条格、多语文书和对方史料互校。 就此而言，元代行省、赋役、河工与多语文书研究有助于把年序同区域、制度或物质材料并读。

\noindent\eventanchor{SYD-1317-YUAN-106}{帝曰：「比遣人往，事已緩矣}\textbf{元。}
在1317年（元延祐4年；公元换算据《元史》本纪纪年）的可见材料中，元留下“帝曰：「比遣人往，事已緩矣”。这一动作把具体人群、机构或地方条件带进年序，但一项记录不能替未被记载者概括整体反应。

从账本的前后项看，“本纪年条提供日期和中枢可见行动；前因不据单条补定。”提供了背景压力；“可回查同年及后续条目，地方影响仍须另证。”标出后来需要追踪的变化，二者之间仍存在未被材料填满的环节。

核心线索为《元史》卷014〈本纪第14卷本年条〉。《元史》为明初仓促官修，本纪保存大量实录条目但有编次与纪年问题；行省执行、数字、族称及对外关系须以条格、多语文书和对方史料互校。 现代研究可从元代行省、赋役、河工与多语文书研究继续追问。

\noindent\eventanchor{SYD-1317-YUAN-178}{帝曰：「苟事力未損，卽遣之}\textbf{元。}
在1317年（元延祐4年；公元换算据《元史》本纪纪年）的可见材料中，元留下“帝曰：「苟事力未損，卽遣之”。这一动作把具体人群、机构或地方条件带进年序，但一项记录不能替未被记载者概括整体反应。

账本把这一节点放在“本纪年条提供日期和中枢可见行动；前因不据单条补定。”与“可回查同年及后续条目，地方影响仍须另证。”之间；两者提示前后压力，并不把时间相邻写成单一因果。

本条首先回到《元史》卷014〈本纪第14卷本年条〉；《元史》为明初仓促官修，本纪保存大量实录条目但有编次与纪年问题；行省执行、数字、族称及对外关系须以条格、多语文书和对方史料互校。 因而元代行省、赋役、河工与多语文书研究不是装饰性的书目，而是检验这一叙述边界的路径。

\noindent\eventanchor{SYD-1317-YUAN-250}{帝曰：「汝漢人用事者，豈皆賢邪？」江南諸路學田昔皆隸官，詔復給本學，以便教養}\textbf{元。}
1317年（元延祐4年；公元换算据《元史》本纪纪年）的“帝曰：「汝漢人用事者，豈皆賢邪？」江南諸路學田昔皆隸官，詔復給本學，以便教養”将元的财政或生产安排推到可见处。制度文本能说明官府打算如何收取、转运或调节，却不能直接等同于不同地区的实收、流通或家庭负担。

其前提记为“本纪年条记录政令或机构处置；实施背景须参照制度材料。”；其后续是“可在同年及后续政令中追踪；条文不单独证明地方执行。”。这是一条待后续材料检验的事件链，而不是对所有行动者意图的替写。

可核的材料入口是《元史》卷014〈本纪第14卷本年条〉。《元史》为明初仓促官修，本纪保存大量实录条目但有编次与纪年问题；行省执行、数字、族称及对外关系须以条格、多语文书和对方史料互校。 就此而言，元代行省、赋役、河工与多语文书研究有助于把年序同区域、制度或物质材料并读。

\noindent\eventanchor{SYD-1317-YUAN-322}{帝曰：「朕以漢人徇私，用泰和律處事，致盜賊滋眾，故有是言}\textbf{元。}
在1317年（元延祐4年；公元换算据《元史》本纪纪年）的可见材料中，元留下“帝曰：「朕以漢人徇私，用泰和律處事，致盜賊滋眾，故有是言”。这一动作把具体人群、机构或地方条件带进年序，但一项记录不能替未被记载者概括整体反应。

从账本的前后项看，“本纪从朝廷视角记录军政行动；出兵与战果需分辨。”提供了背景压力；“后续路线、控制与损失应与对方记录和遗址材料复核。”标出后来需要追踪的变化，二者之间仍存在未被材料填满的环节。

核心线索为《元史》卷014〈本纪第14卷本年条〉。《元史》为明初仓促官修，本纪保存大量实录条目但有编次与纪年问题；行省执行、数字、族称及对外关系须以条格、多语文书和对方史料互校。 现代研究可从元代行省、赋役、河工与多语文书研究继续追问。

\subsection{1318年（元延祐5年；公元换算据《元史》本纪纪年）}

\noindent\eventanchor{SYD-1318-YUAN-035}{帝皆從其議，詔行之}\textbf{元。}
元的1318年（元延祐5年；公元换算据《元史》本纪纪年）记录写到“帝皆從其議，詔行之”。我们可以据此辨认一个可见的选择或压力，而不能把零散的官方记载、碑刻或后出叙事拼成无缝的社会全景。

其前提记为“本纪年条记录政令或机构处置；实施背景须参照制度材料。”；其后续是“可在同年及后续政令中追踪；条文不单独证明地方执行。”。这是一条待后续材料检验的事件链，而不是对所有行动者意图的替写。

可核的材料入口是《元史》卷014〈本纪第14卷本年条〉。《元史》为明初仓促官修，本纪保存大量实录条目但有编次与纪年问题；行省执行、数字、族称及对外关系须以条格、多语文书和对方史料互校。 就此而言，元代行省、赋役、河工与多语文书研究有助于把年序同区域、制度或物质材料并读。

\noindent\eventanchor{SYD-1318-YUAN-107}{帝曰：「除陳巖、呂師夔、管如德、范文虎四人，餘從卿議}\textbf{元。}
元的1318年（元延祐5年；公元换算据《元史》本纪纪年）记录写到“帝曰：「除陳巖、呂師夔、管如德、范文虎四人，餘從卿議”。我们可以据此辨认一个可见的选择或压力，而不能把零散的官方记载、碑刻或后出叙事拼成无缝的社会全景。

从账本的前后项看，“本纪年条记录政令或机构处置；实施背景须参照制度材料。”提供了背景压力；“可在同年及后续政令中追踪；条文不单独证明地方执行。”标出后来需要追踪的变化，二者之间仍存在未被材料填满的环节。

核心线索为《元史》卷014〈本纪第14卷本年条〉。《元史》为明初仓促官修，本纪保存大量实录条目但有编次与纪年问题；行省执行、数字、族称及对外关系须以条格、多语文书和对方史料互校。 现代研究可从元代行省、赋役、河工与多语文书研究继续追问。

\noindent\eventanchor{SYD-1318-YUAN-179}{帝曰：「囚非羣羊，豈可遽殺耶！宜悉配隸淘金}\textbf{元。}
元的1318年（元延祐5年；公元换算据《元史》本纪纪年）记录写到“帝曰：「囚非羣羊，豈可遽殺耶！宜悉配隸淘金”。我们可以据此辨认一个可见的选择或压力，而不能把零散的官方记载、碑刻或后出叙事拼成无缝的社会全景。

账本把这一节点放在“本纪从朝廷视角记录军政行动；出兵与战果需分辨。”与“后续路线、控制与损失应与对方记录和遗址材料复核。”之间；两者提示前后压力，并不把时间相邻写成单一因果。

本条首先回到《元史》卷014〈本纪第14卷本年条〉；《元史》为明初仓促官修，本纪保存大量实录条目但有编次与纪年问题；行省执行、数字、族称及对外关系须以条格、多语文书和对方史料互校。 因而元代行省、赋役、河工与多语文书研究不是装饰性的书目，而是检验这一叙述边界的路径。

\noindent\eventanchor{SYD-1318-YUAN-251}{帝自將征乃顏，發上都}\textbf{元。}
元的1318年（元延祐5年；公元换算据《元史》本纪纪年）记录写到“帝自將征乃顏，發上都”。我们可以据此辨认一个可见的选择或压力，而不能把零散的官方记载、碑刻或后出叙事拼成无缝的社会全景。

其前提记为“本纪年条提供日期和中枢可见行动；前因不据单条补定。”；其后续是“可回查同年及后续条目，地方影响仍须另证。”。这是一条待后续材料检验的事件链，而不是对所有行动者意图的替写。

可核的材料入口是《元史》卷014〈本纪第14卷本年条〉。《元史》为明初仓促官修，本纪保存大量实录条目但有编次与纪年问题；行省执行、数字、族称及对外关系须以条格、多语文书和对方史料互校。 就此而言，元代行省、赋役、河工与多语文书研究有助于把年序同区域、制度或物质材料并读。

\noindent\eventanchor{SYD-1318-YUAN-323}{詔從之，各設官六員}\textbf{元。}
元的1318年（元延祐5年；公元换算据《元史》本纪纪年）记录写到“詔從之，各設官六員”。我们可以据此辨认一个可见的选择或压力，而不能把零散的官方记载、碑刻或后出叙事拼成无缝的社会全景。

从账本的前后项看，“本纪年条记录政令或机构处置；实施背景须参照制度材料。”提供了背景压力；“可在同年及后续政令中追踪；条文不单独证明地方执行。”标出后来需要追踪的变化，二者之间仍存在未被材料填满的环节。

核心线索为《元史》卷014〈本纪第14卷本年条〉。《元史》为明初仓促官修，本纪保存大量实录条目但有编次与纪年问题；行省执行、数字、族称及对外关系须以条格、多语文书和对方史料互校。 现代研究可从元代行省、赋役、河工与多语文书研究继续追问。

\subsection{1319年（元延祐6年；公元换算据《元史》本纪纪年）}

\noindent\eventanchor{SYD-1319-YUAN-036}{丙寅，改懷孟路為懷慶路}\textbf{元。}
“丙寅，改懷孟路為懷慶路”是元在1319年（元延祐6年；公元换算据《元史》本纪纪年）可追查的节点。它提示命令、环境、宗教、迁徙或地方组织如何进入政治过程；影响的范围和速度仍要避免由一条材料外推。

其前提记为“本纪年条记录政令或机构处置；实施背景须参照制度材料。”；其后续是“可在同年及后续政令中追踪；条文不单独证明地方执行。”。这是一条待后续材料检验的事件链，而不是对所有行动者意图的替写。

可核的材料入口是《元史》卷026〈延祐六年〉。《元史》为明初仓促官修，本纪保存大量实录条目但有编次与纪年问题；行省执行、数字、族称及对外关系须以条格、多语文书和对方史料互校。 就此而言，元代行省、赋役、河工与多语文书研究有助于把年序同区域、制度或物质材料并读。

\noindent\eventanchor{SYD-1319-YUAN-108}{敕：「諸司有受命不之官及避繁劇託故去職者，奪其宣敕}\textbf{元。}
“敕：「諸司有受命不之官及避繁劇託故去職者，奪其宣敕”是元在1319年（元延祐6年；公元换算据《元史》本纪纪年）可追查的节点。它提示命令、环境、宗教、迁徙或地方组织如何进入政治过程；影响的范围和速度仍要避免由一条材料外推。

从账本的前后项看，“本纪年条记录政令或机构处置；实施背景须参照制度材料。”提供了背景压力；“可在同年及后续政令中追踪；条文不单独证明地方执行。”标出后来需要追踪的变化，二者之间仍存在未被材料填满的环节。

核心线索为《元史》卷026〈延祐六年〉。《元史》为明初仓促官修，本纪保存大量实录条目但有编次与纪年问题；行省执行、数字、族称及对外关系须以条格、多语文书和对方史料互校。 现代研究可从元代行省、赋役、河工与多语文书研究继续追问。

\noindent\eventanchor{SYD-1319-YUAN-180}{敕上都、大都冬夏設食于路，以食饑者}\textbf{元。}
“敕上都、大都冬夏設食于路，以食饑者”是元在1319年（元延祐6年；公元换算据《元史》本纪纪年）可追查的节点。它提示命令、环境、宗教、迁徙或地方组织如何进入政治过程；影响的范围和速度仍要避免由一条材料外推。

从账本的前后项看，“本纪所记财用、赈济或征发属于特定报告场景。”提供了背景压力；“可与食货志、文书和物质材料比较，不外推为全境总量。”标出后来需要追踪的变化，二者之间仍存在未被材料填满的环节。

核心线索为《元史》卷026〈延祐六年〉。《元史》为明初仓促官修，本纪保存大量实录条目但有编次与纪年问题；行省执行、数字、族称及对外关系须以条格、多语文书和对方史料互校。 现代研究可从元代行省、赋役、河工与多语文书研究继续追问。

\noindent\eventanchor{SYD-1319-YUAN-252}{封諸王月魯鐵木兒為恩王，給印，置王傅官}\textbf{元。}
“封諸王月魯鐵木兒為恩王，給印，置王傅官”是元在1319年（元延祐6年；公元换算据《元史》本纪纪年）可追查的节点。它提示命令、环境、宗教、迁徙或地方组织如何进入政治过程；影响的范围和速度仍要避免由一条材料外推。

其前提记为“本纪年条记录政令或机构处置；实施背景须参照制度材料。”；其后续是“可在同年及后续政令中追踪；条文不单独证明地方执行。”。这是一条待后续材料检验的事件链，而不是对所有行动者意图的替写。

可核的材料入口是《元史》卷026〈延祐六年〉。《元史》为明初仓促官修，本纪保存大量实录条目但有编次与纪年问题；行省执行、数字、族称及对外关系须以条格、多语文书和对方史料互校。 就此而言，元代行省、赋役、河工与多语文书研究有助于把年序同区域、制度或物质材料并读。

\noindent\eventanchor{SYD-1319-YUAN-324}{河間民饑，發粟賑之}\textbf{元。}
“河間民饑，發粟賑之”是元在1319年（元延祐6年；公元换算据《元史》本纪纪年）可追查的节点。它提示命令、环境、宗教、迁徙或地方组织如何进入政治过程；影响的范围和速度仍要避免由一条材料外推。

从账本的前后项看，“本纪所记财用、赈济或征发属于特定报告场景。”提供了背景压力；“可与食货志、文书和物质材料比较，不外推为全境总量。”标出后来需要追踪的变化，二者之间仍存在未被材料填满的环节。

核心线索为《元史》卷026〈延祐六年〉。《元史》为明初仓促官修，本纪保存大量实录条目但有编次与纪年问题；行省执行、数字、族称及对外关系须以条格、多语文书和对方史料互校。 现代研究可从元代行省、赋役、河工与多语文书研究继续追问。

\subsection{1320年（元延祐7年；公元换算据《元史》本纪纪年）}

\noindent\eventanchor{SYD-1320-YUAN-037}{癸未，帝御大明殿，受諸王、百官朝賀}\textbf{元。}
在“癸未，帝御大明殿，受諸王、百官朝賀”中，元的选择经由使节、礼物、边市或协商被记录下来。跨境文本的观察位置不同，因而应把可确认的往来与对动机、服从程度的推断分开。

此前的条件是“本纪年条提供日期和中枢可见行动；前因不据单条补定。”，此后的可见影响是“可回查同年及后续条目，地方影响仍须另证。”。这里的“影响”须按地点、机构和材料种类分别衡量。

这一叙述依据《元史》卷026〈延祐七年〉，但《元史》为明初仓促官修，本纪保存大量实录条目但有编次与纪年问题；行省执行、数字、族称及对外关系须以条格、多语文书和对方史料互校。。研究元代行省、赋役、河工与多语文书研究可以补问材料未直接说出的空间与社会差异。

\noindent\eventanchor{SYD-1320-YUAN-109}{中書省臣進曰：「臺臣所言良是，若非振理朝綱，法度愈壞}\textbf{元。}
在“中書省臣進曰：「臺臣所言良是，若非振理朝綱，法度愈壞”中，元的选择经由使节、礼物、边市或协商被记录下来。跨境文本的观察位置不同，因而应把可确认的往来与对动机、服从程度的推断分开。

账本把这一节点放在“本纪年条记录政令或机构处置；实施背景须参照制度材料。”与“可在同年及后续政令中追踪；条文不单独证明地方执行。”之间；两者提示前后压力，并不把时间相邻写成单一因果。

本条首先回到《元史》卷026〈延祐七年〉；《元史》为明初仓促官修，本纪保存大量实录条目但有编次与纪年问题；行省执行、数字、族称及对外关系须以条格、多语文书和对方史料互校。 因而元代行省、赋役、河工与多语文书研究不是装饰性的书目，而是检验这一叙述边界的路径。

\noindent\eventanchor{SYD-1320-YUAN-181}{臣等乞賜罷黜，選任賢者}\textbf{元。}
“臣等乞賜罷黜，選任賢者”是元在1320年（元延祐7年；公元换算据《元史》本纪纪年）可追查的节点。它提示命令、环境、宗教、迁徙或地方组织如何进入政治过程；影响的范围和速度仍要避免由一条材料外推。

账本把这一节点放在“本纪年条记录政令或机构处置；实施背景须参照制度材料。”与“可在同年及后续政令中追踪；条文不单独证明地方执行。”之间；两者提示前后压力，并不把时间相邻写成单一因果。

本条首先回到《元史》卷026〈延祐七年〉；《元史》为明初仓促官修，本纪保存大量实录条目但有编次与纪年问题；行省执行、数字、族称及对外关系须以条格、多语文书和对方史料互校。 因而元代行省、赋役、河工与多语文书研究不是装饰性的书目，而是检验这一叙述边界的路径。

\noindent\eventanchor{SYD-1320-YUAN-253}{壬午，御史臺臣言：「比賜不兒罕丁山場、完者不花海舶稅，會計其鈔，皆數十萬錠，諸王軍民貧乏者，所賜未嘗若是，苟不撙節，漸致帑藏虛竭，民益困矣}\textbf{元。}
“壬午，御史臺臣言：「比賜不兒罕丁山場、完者不花海舶稅，會計其鈔，皆數十萬錠，諸王軍民貧乏者，所賜未嘗若是，苟不撙節，漸致帑藏虛竭，民益困矣”是元在1320年（元延祐7年；公元换算据《元史》本纪纪年）可追查的节点。它提示命令、环境、宗教、迁徙或地方组织如何进入政治过程；影响的范围和速度仍要避免由一条材料外推。

此前的条件是“本纪所记财用、赈济或征发属于特定报告场景。”，此后的可见影响是“可与食货志、文书和物质材料比较，不外推为全境总量。”。这里的“影响”须按地点、机构和材料种类分别衡量。

这一叙述依据《元史》卷026〈延祐七年〉，但《元史》为明初仓促官修，本纪保存大量实录条目但有编次与纪年问题；行省执行、数字、族称及对外关系须以条格、多语文书和对方史料互校。。研究元代行省、赋役、河工与多语文书研究可以补问材料未直接说出的空间与社会差异。

\noindent\eventanchor{SYD-1320-YUAN-325}{太官進膳，必分賜貴近}\textbf{元。}
“太官進膳，必分賜貴近”是元在1320年（元延祐7年；公元换算据《元史》本纪纪年）可追查的节点。它提示命令、环境、宗教、迁徙或地方组织如何进入政治过程；影响的范围和速度仍要避免由一条材料外推。

账本把这一节点放在“本纪年条提供日期和中枢可见行动；前因不据单条补定。”与“可回查同年及后续条目，地方影响仍须另证。”之间；两者提示前后压力，并不把时间相邻写成单一因果。

本条首先回到《元史》卷026〈延祐七年〉；《元史》为明初仓促官修，本纪保存大量实录条目但有编次与纪年问题；行省执行、数字、族称及对外关系须以条格、多语文书和对方史料互校。 因而元代行省、赋役、河工与多语文书研究不是装饰性的书目，而是检验这一叙述边界的路径。

\subsection{1320年}

\noindent\eventanchor{YUAN-1320-EVENT-185}{英宗即位，试图调整宫廷与官僚权力}\textbf{元。}
在1320年，元所记“英宗即位，试图调整宫廷与官僚权力”把权力安排暴露在继承压力下。后续稳定与否取决于资源和支持网络如何被重新调度，不能由即位仪式本身预先断言。

此前的条件是“元中后期继承、财政与地方危机中前任统治的结束与宫廷、部族或官僚的权力安排相互交织，促成“英宗即位，试图调整宫廷与官僚权力”这一继承节点。”，此后的可见影响是“继承并不自动消除既有矛盾；“英宗即位，试图调整宫廷与官僚权力”后的任官、军权和对外策略需由后续条目追踪。”。这里的“影响”须按地点、机构和材料种类分别衡量。

这一叙述依据《元史》相关本纪；《元朝秘史》、波斯文史籍或《高丽史》对应叙事；《元典章》《通制条格》、多语碑刻与文书（据事核），但《元史》明初速修，纪年与人名有异同；蒙古大汗政权不等同所有汗国，征服、赋役和身份分类须按地区及材料语言分列。。研究蒙古帝国、元朝行省财政、多语文书与区域社会研究可以补问材料未直接说出的空间与社会差异。

\subsection{1321年（元至治1年；公元换算据《元史》本纪纪年）}

\noindent\eventanchor{SYD-1321-YUAN-038}{成宗大德三年，以寧遠王闊闊出總兵北邊，怠於備禦，命帝即軍中代之}\textbf{元。}
1321年（元至治1年；公元换算据《元史》本纪纪年），元的行动落在“成宗大德三年，以寧遠王闊闊出總兵北邊，怠於備禦，命帝即軍中代之”所指的战事与通道上。可见的不是抽象的推进：攻守双方必须让人员、消息与补给穿过具体的城防、道路或水路；账本所列的结果只说明这一节点改变了下一步选择。

此前的条件是“本纪年条记录政令或机构处置；实施背景须参照制度材料。”，此后的可见影响是“可在同年及后续政令中追踪；条文不单独证明地方执行。”。这里的“影响”须按地点、机构和材料种类分别衡量。

这一叙述依据《元史》卷022〈本纪第22卷本年条〉，但《元史》为明初仓促官修，本纪保存大量实录条目但有编次与纪年问题；行省执行、数字、族称及对外关系须以条格、多语文书和对方史料互校。。研究元代行省、赋役、河工与多语文书研究可以补问材料未直接说出的空间与社会差异。

\noindent\eventanchor{SYD-1321-YUAN-110}{海都不得志去，旋亦死}\textbf{元。}
在1321年（元至治1年；公元换算据《元史》本纪纪年）的可见材料中，元留下“海都不得志去，旋亦死”。这一动作把具体人群、机构或地方条件带进年序，但一项记录不能替未被记载者概括整体反应。

从账本的前后项看，“本纪年条提供日期和中枢可见行动；前因不据单条补定。”提供了背景压力；“可回查同年及后续条目，地方影响仍须另证。”标出后来需要追踪的变化，二者之间仍存在未被材料填满的环节。

核心线索为《元史》卷022〈本纪第22卷本年条〉。《元史》为明初仓促官修，本纪保存大量实录条目但有编次与纪年问题；行省执行、数字、族称及对外关系须以条格、多语文书和对方史料互校。 现代研究可从元代行省、赋役、河工与多语文书研究继续追问。

\noindent\eventanchor{SYD-1321-YUAN-182}{明日復戰，軍少却，海都乘之，帝揮軍力戰，突出敵陣後，全軍而還}\textbf{元。}
在1321年（元至治1年；公元换算据《元史》本纪纪年）的可见材料中，元留下“明日復戰，軍少却，海都乘之，帝揮軍力戰，突出敵陣後，全軍而還”。这一动作把具体人群、机构或地方条件带进年序，但一项记录不能替未被记载者概括整体反应。

账本把这一节点放在“本纪年条记录政令或机构处置；实施背景须参照制度材料。”与“可在同年及后续政令中追踪；条文不单独证明地方执行。”之间；两者提示前后压力，并不把时间相邻写成单一因果。

本条首先回到《元史》卷022〈本纪第22卷本年条〉；《元史》为明初仓促官修，本纪保存大量实录条目但有编次与纪年问题；行省执行、数字、族称及对外关系须以条格、多语文书和对方史料互校。 因而元代行省、赋役、河工与多语文书研究不是装饰性的书目，而是检验这一叙述边界的路径。

\noindent\eventanchor{SYD-1321-YUAN-254}{師失利，親出陣，力戰大敗之，盡獲其輜重，悉援諸王、駙馬眾軍以出}\textbf{元。}
在1321年（元至治1年；公元换算据《元史》本纪纪年）的可见材料中，元留下“師失利，親出陣，力戰大敗之，盡獲其輜重，悉援諸王、駙馬眾軍以出”。这一动作把具体人群、机构或地方条件带进年序，但一项记录不能替未被记载者概括整体反应。

此前的条件是“本纪从朝廷视角记录军政行动；出兵与战果需分辨。”，此后的可见影响是“后续路线、控制与损失应与对方记录和遗址材料复核。”。这里的“影响”须按地点、机构和材料种类分别衡量。

这一叙述依据《元史》卷022〈本纪第22卷本年条〉，但《元史》为明初仓促官修，本纪保存大量实录条目但有编次与纪年问题；行省执行、数字、族称及对外关系须以条格、多语文书和对方史料互校。。研究元代行省、赋役、河工与多语文书研究可以补问材料未直接说出的空间与社会差异。

\noindent\eventanchor{SYD-1321-YUAN-326}{十二月，軍至按台山，乃蠻帶部落降}\textbf{元。}
1321年（元至治1年；公元换算据《元史》本纪纪年），元的行动落在“十二月，軍至按台山，乃蠻帶部落降”所指的战事与通道上。可见的不是抽象的推进：攻守双方必须让人员、消息与补给穿过具体的城防、道路或水路；账本所列的结果只说明这一节点改变了下一步选择。

账本把这一节点放在“本纪从朝廷视角记录军政行动；出兵与战果需分辨。”与“后续路线、控制与损失应与对方记录和遗址材料复核。”之间；两者提示前后压力，并不把时间相邻写成单一因果。

本条首先回到《元史》卷022〈本纪第22卷本年条〉；《元史》为明初仓促官修，本纪保存大量实录条目但有编次与纪年问题；行省执行、数字、族称及对外关系须以条格、多语文书和对方史料互校。 因而元代行省、赋役、河工与多语文书研究不是装饰性的书目，而是检验这一叙述边界的路径。

\subsection{1322年（元至治2年；公元换算据《元史》本纪纪年）}

\noindent\eventanchor{SYD-1322-YUAN-039}{帝曰：「父兄雖死事，子弟不勝任者，安可用之？苟賢矣，則病故者亦不可降也}\textbf{元。}
在1322年（元至治2年；公元换算据《元史》本纪纪年）的材料中，元面对的是帝曰：「父兄雖死事，子弟不勝任者，安可用之？苟賢矣，則病故者亦不可降也。这里应先看行动所依赖的关隘、城镇、河道或粮道，再讨论政治后果；军队抵达某处，并不等于它已能稳定征收、安置或控制周边。

此前的条件是“本纪从朝廷视角记录军政行动；出兵与战果需分辨。”，此后的可见影响是“后续路线、控制与损失应与对方记录和遗址材料复核。”。这里的“影响”须按地点、机构和材料种类分别衡量。

这一叙述依据《元史》卷015〈本纪第15卷本年条〉，但《元史》为明初仓促官修，本纪保存大量实录条目但有编次与纪年问题；行省执行、数字、族称及对外关系须以条格、多语文书和对方史料互校。。研究元代行省、赋役、河工与多语文书研究可以补问材料未直接说出的空间与社会差异。

\noindent\eventanchor{SYD-1322-YUAN-111}{帝曰：「自今不當給者汝即畫之，當給者宜覆奏，朕自處之}\textbf{元。}
元的1322年（元至治2年；公元换算据《元史》本纪纪年）记录写到“帝曰：「自今不當給者汝即畫之，當給者宜覆奏，朕自處之”。我们可以据此辨认一个可见的选择或压力，而不能把零散的官方记载、碑刻或后出叙事拼成无缝的社会全景。

从账本的前后项看，“本纪所记财用、赈济或征发属于特定报告场景。”提供了背景压力；“可与食货志、文书和物质材料比较，不外推为全境总量。”标出后来需要追踪的变化，二者之间仍存在未被材料填满的环节。

核心线索为《元史》卷015〈本纪第15卷本年条〉。《元史》为明初仓促官修，本纪保存大量实录条目但有编次与纪年问题；行省执行、数字、族称及对外关系须以条格、多语文书和对方史料互校。 现代研究可从元代行省、赋役、河工与多语文书研究继续追问。

\noindent\eventanchor{SYD-1322-YUAN-183}{甲申，詔皇孫撫諸軍討叛王火魯火孫、合丹禿魯干}\textbf{元。}
元的1322年（元至治2年；公元换算据《元史》本纪纪年）记录写到“甲申，詔皇孫撫諸軍討叛王火魯火孫、合丹禿魯干”。我们可以据此辨认一个可见的选择或压力，而不能把零散的官方记载、碑刻或后出叙事拼成无缝的社会全景。

账本把这一节点放在“本纪年条记录政令或机构处置；实施背景须参照制度材料。”与“可在同年及后续政令中追踪；条文不单独证明地方执行。”之间；两者提示前后压力，并不把时间相邻写成单一因果。

本条首先回到《元史》卷015〈本纪第15卷本年条〉；《元史》为明初仓促官修，本纪保存大量实录条目但有编次与纪年问题；行省执行、数字、族称及对外关系须以条格、多语文书和对方史料互校。 因而元代行省、赋役、河工与多语文书研究不是装饰性的书目，而是检验这一叙述边界的路径。

\noindent\eventanchor{SYD-1322-YUAN-255}{募民能耕江南曠土及公田者，免其差役三年，其輸租免三分之一}\textbf{元。}
元在1322年（元至治2年；公元换算据《元史》本纪纪年）留下的这条记录是“募民能耕江南曠土及公田者，免其差役三年，其輸租免三分之一”。它照亮资源如何被命名和调度，也暴露材料的盲点：数字、额度或禁令若无地方文书和物质材料互证，不能被写成全境同速的现实。

此前的条件是“本纪所记财用、赈济或征发属于特定报告场景。”，此后的可见影响是“可与食货志、文书和物质材料比较，不外推为全境总量。”。这里的“影响”须按地点、机构和材料种类分别衡量。

这一叙述依据《元史》卷015〈本纪第15卷本年条〉，但《元史》为明初仓促官修，本纪保存大量实录条目但有编次与纪年问题；行省执行、数字、族称及对外关系须以条格、多语文书和对方史料互校。。研究元代行省、赋役、河工与多语文书研究可以补问材料未直接说出的空间与社会差异。

\noindent\eventanchor{SYD-1322-YUAN-327}{辛卯，以六衞漢兵千二百、新附軍四百、屯田兵四百造尚書省}\textbf{元。}
在1322年（元至治2年；公元换算据《元史》本纪纪年）的材料中，元面对的是辛卯，以六衞漢兵千二百、新附軍四百、屯田兵四百造尚書省。这里应先看行动所依赖的关隘、城镇、河道或粮道，再讨论政治后果；军队抵达某处，并不等于它已能稳定征收、安置或控制周边。

账本把这一节点放在“本纪年条记录政令或机构处置；实施背景须参照制度材料。”与“可在同年及后续政令中追踪；条文不单独证明地方执行。”之间；两者提示前后压力，并不把时间相邻写成单一因果。

本条首先回到《元史》卷015〈本纪第15卷本年条〉；《元史》为明初仓促官修，本纪保存大量实录条目但有编次与纪年问题；行省执行、数字、族称及对外关系须以条格、多语文书和对方史料互校。 因而元代行省、赋役、河工与多语文书研究不是装饰性的书目，而是检验这一叙述边界的路径。

\subsection{1323年（元至治3年；公元换算据《元史》本纪纪年）}

\noindent\eventanchor{SYD-1323-YUAN-040}{帝曰：「汝家不與它漢人比，弓矢不汝禁也，任汝執之}\textbf{元。}
“帝曰：「汝家不與它漢人比，弓矢不汝禁也，任汝執之”是元在1323年（元至治3年；公元换算据《元史》本纪纪年）可追查的节点。它提示命令、环境、宗教、迁徙或地方组织如何进入政治过程；影响的范围和速度仍要避免由一条材料外推。

其前提记为“本纪年条记录政令或机构处置；实施背景须参照制度材料。”；其后续是“可在同年及后续政令中追踪；条文不单独证明地方执行。”。这是一条待后续材料检验的事件链，而不是对所有行动者意图的替写。

可核的材料入口是《元史》卷015〈本纪第15卷本年条〉。《元史》为明初仓促官修，本纪保存大量实录条目但有编次与纪年问题；行省执行、数字、族称及对外关系须以条格、多语文书和对方史料互校。 就此而言，元代行省、赋役、河工与多语文书研究有助于把年序同区域、制度或物质材料并读。

\noindent\eventanchor{SYD-1323-YUAN-112}{丁亥，安南國王陳日烜遣使來貢方物}\textbf{元与安南。}
在“丁亥，安南國王陳日烜遣使來貢方物”中，元与安南的选择经由使节、礼物、边市或协商被记录下来。跨境文本的观察位置不同，因而应把可确认的往来与对动机、服从程度的推断分开。

从账本的前后项看，“本纪记录使节、贡赐或往来，不等于稳定统属关系。”提供了背景压力；“可与对方编年和交通、文书材料核对其后续落实。”标出后来需要追踪的变化，二者之间仍存在未被材料填满的环节。

核心线索为《元史》卷015〈本纪第15卷本年条〉；《安南志略》及当地碑刻材料（互校）。《元史》为明初仓促官修，本纪保存大量实录条目但有编次与纪年问题；行省执行、数字、族称及对外关系须以条格、多语文书和对方史料互校。 现代研究可从元代行省、赋役、河工与多语文书研究继续追问。

\noindent\eventanchor{SYD-1323-YUAN-184}{皇孫甘不剌所部軍乏食，發大同路榷場糧賑之}\textbf{元。}
“皇孫甘不剌所部軍乏食，發大同路榷場糧賑之”是元在1323年（元至治3年；公元换算据《元史》本纪纪年）可追查的节点。它提示命令、环境、宗教、迁徙或地方组织如何进入政治过程；影响的范围和速度仍要避免由一条材料外推。

账本把这一节点放在“本纪所记财用、赈济或征发属于特定报告场景。”与“可与食货志、文书和物质材料比较，不外推为全境总量。”之间；两者提示前后压力，并不把时间相邻写成单一因果。

本条首先回到《元史》卷015〈本纪第15卷本年条〉；《元史》为明初仓促官修，本纪保存大量实录条目但有编次与纪年问题；行省执行、数字、族称及对外关系须以条格、多语文书和对方史料互校。 因而元代行省、赋役、河工与多语文书研究不是装饰性的书目，而是检验这一叙述边界的路径。

\noindent\eventanchor{SYD-1323-YUAN-256}{命贛州路發米千八百九十石賑之}\textbf{元。}
“命贛州路發米千八百九十石賑之”是元在1323年（元至治3年；公元换算据《元史》本纪纪年）可追查的节点。它提示命令、环境、宗教、迁徙或地方组织如何进入政治过程；影响的范围和速度仍要避免由一条材料外推。

此前的条件是“本纪年条记录政令或机构处置；实施背景须参照制度材料。”，此后的可见影响是“可在同年及后续政令中追踪；条文不单独证明地方执行。”。这里的“影响”须按地点、机构和材料种类分别衡量。

这一叙述依据《元史》卷015〈本纪第15卷本年条〉，但《元史》为明初仓促官修，本纪保存大量实录条目但有编次与纪年问题；行省执行、数字、族称及对外关系须以条格、多语文书和对方史料互校。。研究元代行省、赋役、河工与多语文书研究可以补问材料未直接说出的空间与社会差异。

\noindent\eventanchor{SYD-1323-YUAN-328}{桑哥以聞，請擢絜矩為尚書省舍人，從之}\textbf{元。}
“桑哥以聞，請擢絜矩為尚書省舍人，從之”是元在1323年（元至治3年；公元换算据《元史》本纪纪年）可追查的节点。它提示命令、环境、宗教、迁徙或地方组织如何进入政治过程；影响的范围和速度仍要避免由一条材料外推。

账本把这一节点放在“本纪年条记录政令或机构处置；实施背景须参照制度材料。”与“可在同年及后续政令中追踪；条文不单独证明地方执行。”之间；两者提示前后压力，并不把时间相邻写成单一因果。

本条首先回到《元史》卷015〈本纪第15卷本年条〉；《元史》为明初仓促官修，本纪保存大量实录条目但有编次与纪年问题；行省执行、数字、族称及对外关系须以条格、多语文书和对方史料互校。 因而元代行省、赋役、河工与多语文书研究不是装饰性的书目，而是检验这一叙述边界的路径。

\subsection{1323年}

\noindent\eventanchor{YUAN-1323-EVENT-186}{英宗在南坡遇刺，泰定帝即位}\textbf{元。}
元的“英宗在南坡遇刺，泰定帝即位”并非只改换一个统治者。任官、赏赐、军权与信息通路要在随后重新接合；材料较能确认先后，较难替当事人写出唯一动机。

此前的条件是“元中后期继承、财政与地方危机中前任统治的结束与宫廷、部族或官僚的权力安排相互交织，促成“英宗在南坡遇刺，泰定帝即位”这一继承节点。”，此后的可见影响是“继承并不自动消除既有矛盾；“英宗在南坡遇刺，泰定帝即位”后的任官、军权和对外策略需由后续条目追踪。”。这里的“影响”须按地点、机构和材料种类分别衡量。

这一叙述依据《元史》相关本纪；《元朝秘史》、波斯文史籍或《高丽史》对应叙事；《元典章》《通制条格》、多语碑刻与文书（据事核），但《元史》明初速修，纪年与人名有异同；蒙古大汗政权不等同所有汗国，征服、赋役和身份分类须按地区及材料语言分列。。研究蒙古帝国、元朝行省财政、多语文书与区域社会研究可以补问材料未直接说出的空间与社会差异。

\subsection{1324年（元泰定1年；公元换算据《元史》本纪纪年）}

\noindent\eventanchor{SYD-1324-YUAN-041}{頒行至大銀鈔，詔曰：「昔我世祖皇帝既登大寶，始造中統交鈔，以便民用，歲久法隳，亦既更張，印造至元寶鈔}\textbf{元。}
在1324年（元泰定1年；公元换算据《元史》本纪纪年）的可见材料中，元留下“頒行至大銀鈔，詔曰：「昔我世祖皇帝既登大寶，始造中統交鈔，以便民用，歲久法隳，亦既更張，印造至元寶鈔”。这一动作把具体人群、机构或地方条件带进年序，但一项记录不能替未被记载者概括整体反应。

其前提记为“本纪年条记录政令或机构处置；实施背景须参照制度材料。”；其后续是“可在同年及后续政令中追踪；条文不单独证明地方执行。”。这是一条待后续材料检验的事件链，而不是对所有行动者意图的替写。

可核的材料入口是《元史》卷023〈本纪第23卷本年条〉。《元史》为明初仓促官修，本纪保存大量实录条目但有编次与纪年问题；行省执行、数字、族称及对外关系须以条格、多语文书和对方史料互校。 就此而言，元代行省、赋役、河工与多语文书研究有助于把年序同区域、制度或物质材料并读。

\noindent\eventanchor{SYD-1324-YUAN-113}{汴梁、懷孟、衞輝、彰德、歸德、汝寧、南陽、河南等路蝗}\textbf{元。}
在1324年（元泰定1年；公元换算据《元史》本纪纪年）的可见材料中，元留下“汴梁、懷孟、衞輝、彰德、歸德、汝寧、南陽、河南等路蝗”。这一动作把具体人群、机构或地方条件带进年序，但一项记录不能替未被记载者概括整体反应。

从账本的前后项看，“本纪保留灾害或工程的报告地点与朝廷处置；灾情范围需另证。”提供了背景压力；“可与地方文书、河工和环境材料比较，不将单点记录外推为全域因果。”标出后来需要追踪的变化，二者之间仍存在未被材料填满的环节。

核心线索为《元史》卷023〈本纪第23卷本年条〉。《元史》为明初仓促官修，本纪保存大量实录条目但有编次与纪年问题；行省执行、数字、族称及对外关系须以条格、多语文书和对方史料互校。 现代研究可从元代行省、赋役、河工与多语文书研究继续追问。

\noindent\eventanchor{SYD-1324-YUAN-185}{敕臺臣遣官往鞫，毋徇私情}\textbf{元。}
在1324年（元泰定1年；公元换算据《元史》本纪纪年）的可见材料中，元留下“敕臺臣遣官往鞫，毋徇私情”。这一动作把具体人群、机构或地方条件带进年序，但一项记录不能替未被记载者概括整体反应。

账本把这一节点放在“本纪年条提供日期和中枢可见行动；前因不据单条补定。”与“可回查同年及后续条目，地方影响仍须另证。”之间；两者提示前后压力，并不把时间相邻写成单一因果。

本条首先回到《元史》卷023〈本纪第23卷本年条〉；《元史》为明初仓促官修，本纪保存大量实录条目但有编次与纪年问题；行省执行、数字、族称及对外关系须以条格、多语文书和对方史料互校。 因而元代行省、赋役、河工与多语文书研究不是装饰性的书目，而是检验这一叙述边界的路径。

\noindent\eventanchor{SYD-1324-YUAN-257}{從之，仍賜鈔三千錠賑其部貧民}\textbf{元。}
在1324年（元泰定1年；公元换算据《元史》本纪纪年）的可见材料中，元留下“從之，仍賜鈔三千錠賑其部貧民”。这一动作把具体人群、机构或地方条件带进年序，但一项记录不能替未被记载者概括整体反应。

此前的条件是“本纪所记财用、赈济或征发属于特定报告场景。”，此后的可见影响是“可与食货志、文书和物质材料比较，不外推为全境总量。”。这里的“影响”须按地点、机构和材料种类分别衡量。

这一叙述依据《元史》卷023〈本纪第23卷本年条〉，但《元史》为明初仓促官修，本纪保存大量实录条目但有编次与纪年问题；行省执行、数字、族称及对外关系须以条格、多语文书和对方史料互校。。研究元代行省、赋役、河工与多语文书研究可以补问材料未直接说出的空间与社会差异。

\noindent\eventanchor{SYD-1324-YUAN-329}{大都立資國院，秩正二品}\textbf{元。}
“大都立資國院，秩正二品”在1324年（元泰定1年；公元换算据《元史》本纪纪年）构成元的继承或权力重组节点。名位的变化可以由纪年串连，但它没有自动消除宫廷、军政和地方网络原有的拉扯。

账本把这一节点放在“本纪年条提供日期和中枢可见行动；前因不据单条补定。”与“可回查同年及后续条目，地方影响仍须另证。”之间；两者提示前后压力，并不把时间相邻写成单一因果。

本条首先回到《元史》卷023〈本纪第23卷本年条〉；《元史》为明初仓促官修，本纪保存大量实录条目但有编次与纪年问题；行省执行、数字、族称及对外关系须以条格、多语文书和对方史料互校。 因而元代行省、赋役、河工与多语文书研究不是装饰性的书目，而是检验这一叙述边界的路径。

\subsection{1327年（元泰定4年；公元换算据《元史》本纪纪年）}

\noindent\eventanchor{SYD-1327-YUAN-042}{敕遼陽行省驗實給之}\textbf{元。}
“敕遼陽行省驗實給之”是元在1327年（元泰定4年；公元换算据《元史》本纪纪年）可追查的节点。它提示命令、环境、宗教、迁徙或地方组织如何进入政治过程；影响的范围和速度仍要避免由一条材料外推。

从账本的前后项看，“本纪年条记录政令或机构处置；实施背景须参照制度材料。”提供了背景压力；“可在同年及后续政令中追踪；条文不单独证明地方执行。”标出后来需要追踪的变化，二者之间仍存在未被材料填满的环节。

核心线索为《元史》卷016〈本纪第16卷本年条〉。《元史》为明初仓促官修，本纪保存大量实录条目但有编次与纪年问题；行省执行、数字、族称及对外关系须以条格、多语文书和对方史料互校。 现代研究可从元代行省、赋役、河工与多语文书研究继续追问。

\noindent\eventanchor{SYD-1327-YUAN-114}{帝曰：「此亦何待上聞，當速賑之！」凡出粟五十八萬二千八百八十九石}\textbf{元。}
“帝曰：「此亦何待上聞，當速賑之！」凡出粟五十八萬二千八百八十九石”是元在1327年（元泰定4年；公元换算据《元史》本纪纪年）可追查的节点。它提示命令、环境、宗教、迁徙或地方组织如何进入政治过程；影响的范围和速度仍要避免由一条材料外推。

其前提记为“本纪所记财用、赈济或征发属于特定报告场景。”；其后续是“可与食货志、文书和物质材料比较，不外推为全境总量。”。这是一条待后续材料检验的事件链，而不是对所有行动者意图的替写。

可核的材料入口是《元史》卷016〈本纪第16卷本年条〉。《元史》为明初仓促官修，本纪保存大量实录条目但有编次与纪年问题；行省执行、数字、族称及对外关系须以条格、多语文书和对方史料互校。 就此而言，元代行省、赋役、河工与多语文书研究有助于把年序同区域、制度或物质材料并读。

\noindent\eventanchor{SYD-1327-YUAN-186}{平灤民萬五千四百六十五戶饑，賑粟五千石}\textbf{元。}
“平灤民萬五千四百六十五戶饑，賑粟五千石”是元在1327年（元泰定4年；公元换算据《元史》本纪纪年）可追查的节点。它提示命令、环境、宗教、迁徙或地方组织如何进入政治过程；影响的范围和速度仍要避免由一条材料外推。

此前的条件是“本纪所记财用、赈济或征发属于特定报告场景。”，此后的可见影响是“可与食货志、文书和物质材料比较，不外推为全境总量。”。这里的“影响”须按地点、机构和材料种类分别衡量。

这一叙述依据《元史》卷016〈本纪第16卷本年条〉，但《元史》为明初仓促官修，本纪保存大量实录条目但有编次与纪年问题；行省执行、数字、族称及对外关系须以条格、多语文书和对方史料互校。。研究元代行省、赋役、河工与多语文书研究可以补问材料未直接说出的空间与社会差异。

\noindent\eventanchor{SYD-1327-YUAN-258}{阿速敦等二百九十五人乏食，命驗其實，給糧賑之}\textbf{元。}
“阿速敦等二百九十五人乏食，命驗其實，給糧賑之”是元在1327年（元泰定4年；公元换算据《元史》本纪纪年）可追查的节点。它提示命令、环境、宗教、迁徙或地方组织如何进入政治过程；影响的范围和速度仍要避免由一条材料外推。

账本把这一节点放在“本纪年条记录政令或机构处置；实施背景须参照制度材料。”与“可在同年及后续政令中追踪；条文不单独证明地方执行。”之间；两者提示前后压力，并不把时间相邻写成单一因果。

本条首先回到《元史》卷016〈本纪第16卷本年条〉；《元史》为明初仓促官修，本纪保存大量实录条目但有编次与纪年问题；行省执行、数字、族称及对外关系须以条格、多语文书和对方史料互校。 因而元代行省、赋役、河工与多语文书研究不是装饰性的书目，而是检验这一叙述边界的路径。

\noindent\eventanchor{SYD-1327-YUAN-330}{安南國王陳日烜遣其中大夫陳克用來貢方物}\textbf{元与安南。}
“安南國王陳日烜遣其中大夫陳克用來貢方物”是元与安南在1327年（元泰定4年；公元换算据《元史》本纪纪年）可追查的节点。它提示命令、环境、宗教、迁徙或地方组织如何进入政治过程；影响的范围和速度仍要避免由一条材料外推。

其前提记为“本纪年条提供日期和中枢可见行动；前因不据单条补定。”；其后续是“可回查同年及后续条目，地方影响仍须另证。”。这是一条待后续材料检验的事件链，而不是对所有行动者意图的替写。

可核的材料入口是《元史》卷016〈本纪第16卷本年条〉；《安南志略》及当地碑刻材料（互校）。《元史》为明初仓促官修，本纪保存大量实录条目但有编次与纪年问题；行省执行、数字、族称及对外关系须以条格、多语文书和对方史料互校。 就此而言，元代行省、赋役、河工与多语文书研究有助于把年序同区域、制度或物质材料并读。

\subsection{1328年（元天历1年；公元换算据《元史》本纪纪年）}

\noindent\eventanchor{SYD-1328-YUAN-043}{罷江南諸省買銀提舉司}\textbf{元。}
在1328年（元天历1年；公元换算据《元史》本纪纪年）的可见材料中，元留下“罷江南諸省買銀提舉司”。这一动作把具体人群、机构或地方条件带进年序，但一项记录不能替未被记载者概括整体反应。

从账本的前后项看，“本纪年条记录政令或机构处置；实施背景须参照制度材料。”提供了背景压力；“可在同年及后续政令中追踪；条文不单独证明地方执行。”标出后来需要追踪的变化，二者之间仍存在未被材料填满的环节。

核心线索为《元史》卷016〈本纪第16卷本年条〉。《元史》为明初仓促官修，本纪保存大量实录条目但有编次与纪年问题；行省执行、数字、族称及对外关系须以条格、多语文书和对方史料互校。 现代研究可从元代行省、赋役、河工与多语文书研究继续追问。

\noindent\eventanchor{SYD-1328-YUAN-115}{帝曰：「桑哥已誅，納〔速〕剌丁滅里在獄，唯沙不丁朕姑釋之耳}\textbf{元。}
在1328年（元天历1年；公元换算据《元史》本纪纪年）的可见材料中，元留下“帝曰：「桑哥已誅，納〔速〕剌丁滅里在獄，唯沙不丁朕姑釋之耳”。这一动作把具体人群、机构或地方条件带进年序，但一项记录不能替未被记载者概括整体反应。

其前提记为“本纪年条提供日期和中枢可见行动；前因不据单条补定。”；其后续是“可回查同年及后续条目，地方影响仍须另证。”。这是一条待后续材料检验的事件链，而不是对所有行动者意图的替写。

可核的材料入口是《元史》卷016〈本纪第16卷本年条〉。《元史》为明初仓促官修，本纪保存大量实录条目但有编次与纪年问题；行省执行、数字、族称及对外关系须以条格、多语文书和对方史料互校。 就此而言，元代行省、赋役、河工与多语文书研究有助于把年序同区域、制度或物质材料并读。

\noindent\eventanchor{SYD-1328-YUAN-187}{何榮祖言：「宜召各省官任錢穀者詣京師，集議科取之法以聞}\textbf{元。}
在1328年（元天历1年；公元换算据《元史》本纪纪年）的可见材料中，元留下“何榮祖言：「宜召各省官任錢穀者詣京師，集議科取之法以聞”。这一动作把具体人群、机构或地方条件带进年序，但一项记录不能替未被记载者概括整体反应。

此前的条件是“本纪年条记录政令或机构处置；实施背景须参照制度材料。”，此后的可见影响是“可在同年及后续政令中追踪；条文不单独证明地方执行。”。这里的“影响”须按地点、机构和材料种类分别衡量。

这一叙述依据《元史》卷016〈本纪第16卷本年条〉，但《元史》为明初仓促官修，本纪保存大量实录条目但有编次与纪年问题；行省执行、数字、族称及对外关系须以条格、多语文书和对方史料互校。。研究元代行省、赋役、河工与多语文书研究可以补问材料未直接说出的空间与社会差异。

\noindent\eventanchor{SYD-1328-YUAN-259}{壬申，立河南江北行中書省，治汴梁}\textbf{元。}
“壬申，立河南江北行中書省，治汴梁”在1328年（元天历1年；公元换算据《元史》本纪纪年）构成元的继承或权力重组节点。名位的变化可以由纪年串连，但它没有自动消除宫廷、军政和地方网络原有的拉扯。

账本把这一节点放在“本纪年条记录政令或机构处置；实施背景须参照制度材料。”与“可在同年及后续政令中追踪；条文不单独证明地方执行。”之间；两者提示前后压力，并不把时间相邻写成单一因果。

本条首先回到《元史》卷016〈本纪第16卷本年条〉；《元史》为明初仓促官修，本纪保存大量实录条目但有编次与纪年问题；行省执行、数字、族称及对外关系须以条格、多语文书和对方史料互校。 因而元代行省、赋役、河工与多语文书研究不是装饰性的书目，而是检验这一叙述边界的路径。

\noindent\eventanchor{SYD-1328-YUAN-331}{塔叉兒、塔帶民饑，發米賑之}\textbf{元。}
在1328年（元天历1年；公元换算据《元史》本纪纪年）的可见材料中，元留下“塔叉兒、塔帶民饑，發米賑之”。这一动作把具体人群、机构或地方条件带进年序，但一项记录不能替未被记载者概括整体反应。

其前提记为“本纪所记财用、赈济或征发属于特定报告场景。”；其后续是“可与食货志、文书和物质材料比较，不外推为全境总量。”。这是一条待后续材料检验的事件链，而不是对所有行动者意图的替写。

可核的材料入口是《元史》卷016〈本纪第16卷本年条〉。《元史》为明初仓促官修，本纪保存大量实录条目但有编次与纪年问题；行省执行、数字、族称及对外关系须以条格、多语文书和对方史料互校。 就此而言，元代行省、赋役、河工与多语文书研究有助于把年序同区域、制度或物质材料并读。

\subsection{1328年}

\noindent\eventanchor{YUAN-1328-EVENT-187}{元廷爆发两都之战}\textbf{元。}
在1328年的材料中，元面对的是元廷爆发两都之战。这里应先看行动所依赖的关隘、城镇、河道或粮道，再讨论政治后果；军队抵达某处，并不等于它已能稳定征收、安置或控制周边。

账本把这一节点放在“元中后期继承、财政与地方危机中既有的联盟、交通与补给压力，为“元廷爆发两都之战”提供了直接的军事或动员背景。”与““元廷爆发两都之战”改变了相关城镇、道路或政权间的军事选择；伤亡、俘获和控制范围仅可按各方材料分别确认。”之间；两者提示前后压力，并不把时间相邻写成单一因果。

本条首先回到《元史》相关本纪；《元朝秘史》、波斯文史籍或《高丽史》对应叙事；《元典章》《通制条格》、多语碑刻与文书（据事核）；《元史》明初速修，纪年与人名有异同；蒙古大汗政权不等同所有汗国，征服、赋役和身份分类须按地区及材料语言分列。 因而蒙古帝国、元朝行省财政、多语文书与区域社会研究不是装饰性的书目，而是检验这一叙述边界的路径。

\subsection{1329年（元天历2年；公元换算据《元史》本纪纪年）}

\noindent\eventanchor{SYD-1329-YUAN-044}{丙寅，改懷孟路為懷慶路}\textbf{元。}
在1329年（元天历2年；公元换算据《元史》本纪纪年）的可见材料中，元留下“丙寅，改懷孟路為懷慶路”。这一动作把具体人群、机构或地方条件带进年序，但一项记录不能替未被记载者概括整体反应。

从账本的前后项看，“本纪年条记录政令或机构处置；实施背景须参照制度材料。”提供了背景压力；“可在同年及后续政令中追踪；条文不单独证明地方执行。”标出后来需要追踪的变化，二者之间仍存在未被材料填满的环节。

核心线索为《元史》卷026〈本纪第26卷本年条〉。《元史》为明初仓促官修，本纪保存大量实录条目但有编次与纪年问题；行省执行、数字、族称及对外关系须以条格、多语文书和对方史料互校。 现代研究可从元代行省、赋役、河工与多语文书研究继续追问。

\noindent\eventanchor{SYD-1329-YUAN-116}{播州南寧長官洛麼作亂，思州守臣換住哥招諭之，洛麼遣人以方物來覲}\textbf{元。}
1329年（元天历2年；公元换算据《元史》本纪纪年），元的行动落在“播州南寧長官洛麼作亂，思州守臣換住哥招諭之，洛麼遣人以方物來覲”所指的战事与通道上。可见的不是抽象的推进：攻守双方必须让人员、消息与补给穿过具体的城防、道路或水路；账本所列的结果只说明这一节点改变了下一步选择。

其前提记为“本纪从朝廷视角记录军政行动；出兵与战果需分辨。”；其后续是“后续路线、控制与损失应与对方记录和遗址材料复核。”。这是一条待后续材料检验的事件链，而不是对所有行动者意图的替写。

可核的材料入口是《元史》卷026〈本纪第26卷本年条〉。《元史》为明初仓促官修，本纪保存大量实录条目但有编次与纪年问题；行省执行、数字、族称及对外关系须以条格、多语文书和对方史料互校。 就此而言，元代行省、赋役、河工与多语文书研究有助于把年序同区域、制度或物质材料并读。

\noindent\eventanchor{SYD-1329-YUAN-188}{敕：「諸司有受命不之官及避繁劇託故去職者，奪其宣敕}\textbf{元。}
在1329年（元天历2年；公元换算据《元史》本纪纪年）的可见材料中，元留下“敕：「諸司有受命不之官及避繁劇託故去職者，奪其宣敕”。这一动作把具体人群、机构或地方条件带进年序，但一项记录不能替未被记载者概括整体反应。

此前的条件是“本纪年条记录政令或机构处置；实施背景须参照制度材料。”，此后的可见影响是“可在同年及后续政令中追踪；条文不单独证明地方执行。”。这里的“影响”须按地点、机构和材料种类分别衡量。

这一叙述依据《元史》卷026〈本纪第26卷本年条〉，但《元史》为明初仓促官修，本纪保存大量实录条目但有编次与纪年问题；行省执行、数字、族称及对外关系须以条格、多语文书和对方史料互校。。研究元代行省、赋役、河工与多语文书研究可以补问材料未直接说出的空间与社会差异。

\noindent\eventanchor{SYD-1329-YUAN-260}{敕上都、大都冬夏設食于路，以食饑者}\textbf{元。}
在1329年（元天历2年；公元换算据《元史》本纪纪年）的可见材料中，元留下“敕上都、大都冬夏設食于路，以食饑者”。这一动作把具体人群、机构或地方条件带进年序，但一项记录不能替未被记载者概括整体反应。

从账本的前后项看，“本纪所记财用、赈济或征发属于特定报告场景。”提供了背景压力；“可与食货志、文书和物质材料比较，不外推为全境总量。”标出后来需要追踪的变化，二者之间仍存在未被材料填满的环节。

核心线索为《元史》卷026〈本纪第26卷本年条〉。《元史》为明初仓促官修，本纪保存大量实录条目但有编次与纪年问题；行省执行、数字、族称及对外关系须以条格、多语文书和对方史料互校。 现代研究可从元代行省、赋役、河工与多语文书研究继续追问。

\subsection{1330年（元至顺1年；公元换算据《元史》本纪纪年）}

\noindent\eventanchor{SYD-1330-YUAN-045}{帝謂中書省臣曰：「至尊委我以天下事，日夜寅畏，惟恐弗堪}\textbf{元。}
在1330年（元至顺1年；公元换算据《元史》本纪纪年）的可见材料中，元留下“帝謂中書省臣曰：「至尊委我以天下事，日夜寅畏，惟恐弗堪”。这一动作把具体人群、机构或地方条件带进年序，但一项记录不能替未被记载者概括整体反应。

账本把这一节点放在“本纪年条记录政令或机构处置；实施背景须参照制度材料。”与“可在同年及后续政令中追踪；条文不单独证明地方执行。”之间；两者提示前后压力，并不把时间相邻写成单一因果。

本条首先回到《元史》卷027〈本纪第27卷本年条〉；《元史》为明初仓促官修，本纪保存大量实录条目但有编次与纪年问题；行省执行、数字、族称及对外关系须以条格、多语文书和对方史料互校。 因而元代行省、赋役、河工与多语文书研究不是装饰性的书目，而是检验这一叙述边界的路径。

\noindent\eventanchor{SYD-1330-YUAN-117}{監察御史段輔、太子詹事郭貫等，首請近賢人，擇師傅，帝嘉納之}\textbf{元。}
在1330年（元至顺1年；公元换算据《元史》本纪纪年）的可见材料中，元留下“監察御史段輔、太子詹事郭貫等，首請近賢人，擇師傅，帝嘉納之”。这一动作把具体人群、机构或地方条件带进年序，但一项记录不能替未被记载者概括整体反应。

此前的条件是“本纪年条提供日期和中枢可见行动；前因不据单条补定。”，此后的可见影响是“可回查同年及后续条目，地方影响仍须另证。”。这里的“影响”须按地点、机构和材料种类分别衡量。

这一叙述依据《元史》卷027〈本纪第27卷本年条〉，但《元史》为明初仓促官修，本纪保存大量实录条目但有编次与纪年问题；行省执行、数字、族称及对外关系须以条格、多语文书和对方史料互校。。研究元代行省、赋役、河工与多语文书研究可以补问材料未直接说出的空间与社会差异。

\noindent\eventanchor{SYD-1330-YUAN-189}{六年十月戊午，受玉冊，詔命百司庶務必先啟太子，然後奏聞}\textbf{元。}
在1330年（元至顺1年；公元换算据《元史》本纪纪年）的可见材料中，元留下“六年十月戊午，受玉冊，詔命百司庶務必先啟太子，然後奏聞”。这一动作把具体人群、机构或地方条件带进年序，但一项记录不能替未被记载者概括整体反应。

从账本的前后项看，“本纪年条记录政令或机构处置；实施背景须参照制度材料。”提供了背景压力；“可在同年及后续政令中追踪；条文不单独证明地方执行。”标出后来需要追踪的变化，二者之间仍存在未被材料填满的环节。

核心线索为《元史》卷027〈本纪第27卷本年条〉。《元史》为明初仓促官修，本纪保存大量实录条目但有编次与纪年问题；行省执行、数字、族称及对外关系须以条格、多语文书和对方史料互校。 现代研究可从元代行省、赋役、河工与多语文书研究继续追问。

\noindent\eventanchor{SYD-1330-YUAN-261}{仁宗欲立為太子，帝入謁太后固辭，曰：「臣幼無能，且有兄在，宜立兄，以臣輔之}\textbf{元。}
“仁宗欲立為太子，帝入謁太后固辭，曰：「臣幼無能，且有兄在，宜立兄，以臣輔之”在1330年（元至顺1年；公元换算据《元史》本纪纪年）构成元的继承或权力重组节点。名位的变化可以由纪年串连，但它没有自动消除宫廷、军政和地方网络原有的拉扯。

账本把这一节点放在“本纪年条提供日期和中枢可见行动；前因不据单条补定。”与“可回查同年及后续条目，地方影响仍须另证。”之间；两者提示前后压力，并不把时间相邻写成单一因果。

本条首先回到《元史》卷027〈本纪第27卷本年条〉；《元史》为明初仓促官修，本纪保存大量实录条目但有编次与纪年问题；行省执行、数字、族称及对外关系须以条格、多语文书和对方史料互校。 因而元代行省、赋役、河工与多语文书研究不是装饰性的书目，而是检验这一叙述边界的路径。

\subsection{1332年（元至顺3年；公元换算据《元史》本纪纪年）}

\noindent\eventanchor{SYD-1332-YUAN-046}{丙午，中書省臣言：「京畿荐饑，宜免今歲田租}\textbf{元。}
1332年（元至顺3年；公元换算据《元史》本纪纪年）的“丙午，中書省臣言：「京畿荐饑，宜免今歲田租”将元的财政或生产安排推到可见处。制度文本能说明官府打算如何收取、转运或调节，却不能直接等同于不同地区的实收、流通或家庭负担。

其前提记为“本纪年条记录政令或机构处置；实施背景须参照制度材料。”；其后续是“可在同年及后续政令中追踪；条文不单独证明地方执行。”。这是一条待后续材料检验的事件链，而不是对所有行动者意图的替写。

可核的材料入口是《元史》卷017〈本纪第17卷本年条〉。《元史》为明初仓促官修，本纪保存大量实录条目但有编次与纪年问题；行省执行、数字、族称及对外关系须以条格、多语文书和对方史料互校。 就此而言，元代行省、赋役、河工与多语文书研究有助于把年序同区域、制度或物质材料并读。

\noindent\eventanchor{SYD-1332-YUAN-118}{敕依東平例，發附近官廪，計口以給}\textbf{元。}
在1332年（元至顺3年；公元换算据《元史》本纪纪年）的可见材料中，元留下“敕依東平例，發附近官廪，計口以給”。这一动作把具体人群、机构或地方条件带进年序，但一项记录不能替未被记载者概括整体反应。

从账本的前后项看，“本纪所记财用、赈济或征发属于特定报告场景。”提供了背景压力；“可与食货志、文书和物质材料比较，不外推为全境总量。”标出后来需要追踪的变化，二者之间仍存在未被材料填满的环节。

核心线索为《元史》卷017〈本纪第17卷本年条〉。《元史》为明初仓促官修，本纪保存大量实录条目但有编次与纪年问题；行省执行、数字、族称及对外关系须以条格、多语文书和对方史料互校。 现代研究可从元代行省、赋役、河工与多语文书研究继续追问。

\noindent\eventanchor{SYD-1332-YUAN-190}{敕應昌府給乞答帶糧五百石，以賑饑民}\textbf{元。}
在1332年（元至顺3年；公元换算据《元史》本纪纪年）的可见材料中，元留下“敕應昌府給乞答帶糧五百石，以賑饑民”。这一动作把具体人群、机构或地方条件带进年序，但一项记录不能替未被记载者概括整体反应。

从账本的前后项看，“本纪所记财用、赈济或征发属于特定报告场景。”提供了背景压力；“可与食货志、文书和物质材料比较，不外推为全境总量。”标出后来需要追踪的变化，二者之间仍存在未被材料填满的环节。

核心线索为《元史》卷017〈本纪第17卷本年条〉。《元史》为明初仓促官修，本纪保存大量实录条目但有编次与纪年问题；行省执行、数字、族称及对外关系须以条格、多语文书和对方史料互校。 现代研究可从元代行省、赋役、河工与多语文书研究继续追问。

\noindent\eventanchor{SYD-1332-YUAN-262}{帝曰：「死者勿論，其存者罰不可恕也}\textbf{元。}
在1332年（元至顺3年；公元换算据《元史》本纪纪年）的可见材料中，元留下“帝曰：「死者勿論，其存者罰不可恕也”。这一动作把具体人群、机构或地方条件带进年序，但一项记录不能替未被记载者概括整体反应。

其前提记为“本纪年条提供日期和中枢可见行动；前因不据单条补定。”；其后续是“可回查同年及后续条目，地方影响仍须另证。”。这是一条待后续材料检验的事件链，而不是对所有行动者意图的替写。

可核的材料入口是《元史》卷017〈本纪第17卷本年条〉。《元史》为明初仓促官修，本纪保存大量实录条目但有编次与纪年问题；行省执行、数字、族称及对外关系须以条格、多语文书和对方史料互校。 就此而言，元代行省、赋役、河工与多语文书研究有助于把年序同区域、制度或物质材料并读。

\subsection{1333年（元元统1年；公元换算据《元史》本纪纪年）}

\noindent\eventanchor{SYD-1333-YUAN-047}{敕海運米十萬石給遼陽戍兵，仍諭其省官薛闍干，令伯鐵木部欽察等耕漁自養，糧不須給}\textbf{元。}
在1333年（元元统1年；公元换算据《元史》本纪纪年）的材料中，元面对的是敕海運米十萬石給遼陽戍兵，仍諭其省官薛闍干，令伯鐵木部欽察等耕漁自養，糧不須給。这里应先看行动所依赖的关隘、城镇、河道或粮道，再讨论政治后果；军队抵达某处，并不等于它已能稳定征收、安置或控制周边。

其前提记为“本纪年条记录政令或机构处置；实施背景须参照制度材料。”；其后续是“可在同年及后续政令中追踪；条文不单独证明地方执行。”。这是一条待后续材料检验的事件链，而不是对所有行动者意图的替写。

可核的材料入口是《元史》卷017〈本纪第17卷本年条〉。《元史》为明初仓促官修，本纪保存大量实录条目但有编次与纪年问题；行省执行、数字、族称及对外关系须以条格、多语文书和对方史料互校。 就此而言，元代行省、赋役、河工与多语文书研究有助于把年序同区域、制度或物质材料并读。

\noindent\eventanchor{SYD-1333-YUAN-119}{帝曰：「非其人不行，乃朕中止之耳，勿徵}\textbf{元。}
元的1333年（元元统1年；公元换算据《元史》本纪纪年）记录写到“帝曰：「非其人不行，乃朕中止之耳，勿徵”。我们可以据此辨认一个可见的选择或压力，而不能把零散的官方记载、碑刻或后出叙事拼成无缝的社会全景。

从账本的前后项看，“本纪所记财用、赈济或征发属于特定报告场景。”提供了背景压力；“可与食货志、文书和物质材料比较，不外推为全境总量。”标出后来需要追踪的变化，二者之间仍存在未被材料填满的环节。

核心线索为《元史》卷017〈本纪第17卷本年条〉。《元史》为明初仓促官修，本纪保存大量实录条目但有编次与纪年问题；行省执行、数字、族称及对外关系须以条格、多语文书和对方史料互校。 现代研究可从元代行省、赋役、河工与多语文书研究继续追问。

\noindent\eventanchor{SYD-1333-YUAN-191}{高麗國王王〔賰〕請易名曰昛，其僉議府請陞僉議司，降二品印，從之}\textbf{元与高丽。}
在1333年（元元统1年；公元换算据《元史》本纪纪年）的材料中，元与高丽面对的是高麗國王王〔賰〕請易名曰昛，其僉議府請陞僉議司，降二品印，從之。这里应先看行动所依赖的关隘、城镇、河道或粮道，再讨论政治后果；军队抵达某处，并不等于它已能稳定征收、安置或控制周边。

从账本的前后项看，“本纪从朝廷视角记录军政行动；出兵与战果需分辨。”提供了背景压力；“后续路线、控制与损失应与对方记录和遗址材料复核。”标出后来需要追踪的变化，二者之间仍存在未被材料填满的环节。

核心线索为《元史》卷017〈本纪第17卷本年条〉；《高丽史》同年条（互校）。《元史》为明初仓促官修，本纪保存大量实录条目但有编次与纪年问题；行省执行、数字、族称及对外关系须以条格、多语文书和对方史料互校。 现代研究可从元代行省、赋役、河工与多语文书研究继续追问。

\noindent\eventanchor{SYD-1333-YUAN-263}{甲戌，河南江北行省平章伯顏言：「揚州忙兀臺所立屯田，為田四萬餘頃，官種外，宜聽民耕墾}\textbf{元。}
元在1333年（元元统1年；公元换算据《元史》本纪纪年）留下的这条记录是“甲戌，河南江北行省平章伯顏言：「揚州忙兀臺所立屯田，為田四萬餘頃，官種外，宜聽民耕墾”。它照亮资源如何被命名和调度，也暴露材料的盲点：数字、额度或禁令若无地方文书和物质材料互证，不能被写成全境同速的现实。

其前提记为“本纪年条记录政令或机构处置；实施背景须参照制度材料。”；其后续是“可在同年及后续政令中追踪；条文不单独证明地方执行。”。这是一条待后续材料检验的事件链，而不是对所有行动者意图的替写。

可核的材料入口是《元史》卷017〈本纪第17卷本年条〉。《元史》为明初仓促官修，本纪保存大量实录条目但有编次与纪年问题；行省执行、数字、族称及对外关系须以条格、多语文书和对方史料互校。 就此而言，元代行省、赋役、河工与多语文书研究有助于把年序同区域、制度或物质材料并读。

\subsection{1333年}

\noindent\eventanchor{YUAN-1333-EVENT-188}{妥懽帖睦尔即位}\textbf{元。}
“妥懽帖睦尔即位”在1333年构成元的继承或权力重组节点。名位的变化可以由纪年串连，但它没有自动消除宫廷、军政和地方网络原有的拉扯。

其前提记为“元中后期继承、财政与地方危机中前任统治的结束与宫廷、部族或官僚的权力安排相互交织，促成“妥懽帖睦尔即位”这一继承节点。”；其后续是“继承并不自动消除既有矛盾；“妥懽帖睦尔即位”后的任官、军权和对外策略需由后续条目追踪。”。这是一条待后续材料检验的事件链，而不是对所有行动者意图的替写。

可核的材料入口是《元史》相关本纪；《元朝秘史》、波斯文史籍或《高丽史》对应叙事；《元典章》《通制条格》、多语碑刻与文书（据事核）。《元史》明初速修，纪年与人名有异同；蒙古大汗政权不等同所有汗国，征服、赋役和身份分类须按地区及材料语言分列。 就此而言，蒙古帝国、元朝行省财政、多语文书与区域社会研究有助于把年序同区域、制度或物质材料并读。

\subsection{1334年（元元统2年；公元换算据《元史》本纪纪年）}

\noindent\eventanchor{SYD-1334-YUAN-048}{甲辰，遣司徒兀都帶、平章政事不忽木、左丞張九思，率百官請諡于南郊}\textbf{元。}
“甲辰，遣司徒兀都帶、平章政事不忽木、左丞張九思，率百官請諡于南郊”是元在1334年（元元统2年；公元换算据《元史》本纪纪年）可追查的节点。它提示命令、环境、宗教、迁徙或地方组织如何进入政治过程；影响的范围和速度仍要避免由一条材料外推。

其前提记为“本纪年条提供日期和中枢可见行动；前因不据单条补定。”；其后续是“可回查同年及后续条目，地方影响仍须另证。”。这是一条待后续材料检验的事件链，而不是对所有行动者意图的替写。

可核的材料入口是《元史》卷017〈本纪第17卷本年条〉。《元史》为明初仓促官修，本纪保存大量实录条目但有编次与纪年问题；行省执行、数字、族称及对外关系须以条格、多语文书和对方史料互校。 就此而言，元代行省、赋役、河工与多语文书研究有助于把年序同区域、制度或物质材料并读。

\noindent\eventanchor{SYD-1334-YUAN-120}{親王、諸大臣發使告哀于皇孫}\textbf{元。}
在“親王、諸大臣發使告哀于皇孫”中，元的选择经由使节、礼物、边市或协商被记录下来。跨境文本的观察位置不同，因而应把可确认的往来与对动机、服从程度的推断分开。

此前的条件是“本纪年条提供日期和中枢可见行动；前因不据单条补定。”，此后的可见影响是“可回查同年及后续条目，地方影响仍须另证。”。这里的“影响”须按地点、机构和材料种类分别衡量。

这一叙述依据《元史》卷017〈本纪第17卷本年条〉，但《元史》为明初仓促官修，本纪保存大量实录条目但有编次与纪年问题；行省执行、数字、族称及对外关系须以条格、多语文书和对方史料互校。。研究元代行省、赋役、河工与多语文书研究可以补问材料未直接说出的空间与社会差异。

\noindent\eventanchor{SYD-1334-YUAN-192}{三十一年春正月壬子朔，帝不豫，免朝賀}\textbf{元。}
在“三十一年春正月壬子朔，帝不豫，免朝賀”中，元的选择经由使节、礼物、边市或协商被记录下来。跨境文本的观察位置不同，因而应把可确认的往来与对动机、服从程度的推断分开。

从账本的前后项看，“本纪年条提供日期和中枢可见行动；前因不据单条补定。”提供了背景压力；“可回查同年及后续条目，地方影响仍须另证。”标出后来需要追踪的变化，二者之间仍存在未被材料填满的环节。

核心线索为《元史》卷017〈本纪第17卷本年条〉。《元史》为明初仓促官修，本纪保存大量实录条目但有编次与纪年问题；行省执行、数字、族称及对外关系须以条格、多语文书和对方史料互校。 现代研究可从元代行省、赋役、河工与多语文书研究继续追问。

\noindent\eventanchor{SYD-1334-YUAN-264}{世祖度量弘廣，知人善任使，信用儒術，用能以夏變夷，立經陳紀，所以為一代之制者，規模宏遠矣}\textbf{元。}
在“世祖度量弘廣，知人善任使，信用儒術，用能以夏變夷，立經陳紀，所以為一代之制者，規模宏遠矣”中，元的选择经由使节、礼物、边市或协商被记录下来。跨境文本的观察位置不同，因而应把可确认的往来与对动机、服从程度的推断分开。

其前提记为“本纪年条提供日期和中枢可见行动；前因不据单条补定。”；其后续是“可回查同年及后续条目，地方影响仍须另证。”。这是一条待后续材料检验的事件链，而不是对所有行动者意图的替写。

可核的材料入口是《元史》卷017〈本纪第17卷本年条〉。《元史》为明初仓促官修，本纪保存大量实录条目但有编次与纪年问题；行省执行、数字、族称及对外关系须以条格、多语文书和对方史料互校。 就此而言，元代行省、赋役、河工与多语文书研究有助于把年序同区域、制度或物质材料并读。

\subsection{1335年（元至元1年；公元换算据《元史》本纪纪年）}

\noindent\eventanchor{SYD-1335-YUAN-049}{丁丑，頒農桑舊制十四條于天下，仍詔勵有司以察勤惰}\textbf{元。}
在1335年（元至元1年；公元换算据《元史》本纪纪年）的可见材料中，元留下“丁丑，頒農桑舊制十四條于天下，仍詔勵有司以察勤惰”。这一动作把具体人群、机构或地方条件带进年序，但一项记录不能替未被记载者概括整体反应。

其前提记为“本纪年条记录政令或机构处置；实施背景须参照制度材料。”；其后续是“可在同年及后续政令中追踪；条文不单独证明地方执行。”。这是一条待后续材料检验的事件链，而不是对所有行动者意图的替写。

可核的材料入口是《元史》卷030〈本纪第30卷本年条〉。《元史》为明初仓促官修，本纪保存大量实录条目但有编次与纪年问题；行省执行、数字、族称及对外关系须以条格、多语文书和对方史料互校。 就此而言，元代行省、赋役、河工与多语文书研究有助于把年序同区域、制度或物质材料并读。

\noindent\eventanchor{SYD-1335-YUAN-121}{甲申，廣西花角蠻為寇，命所部討之}\textbf{元。}
在1335年（元至元1年；公元换算据《元史》本纪纪年）的可见材料中，元留下“甲申，廣西花角蠻為寇，命所部討之”。这一动作把具体人群、机构或地方条件带进年序，但一项记录不能替未被记载者概括整体反应。

此前的条件是“本纪年条记录政令或机构处置；实施背景须参照制度材料。”，此后的可见影响是“可在同年及后续政令中追踪；条文不单独证明地方执行。”。这里的“影响”须按地点、机构和材料种类分别衡量。

这一叙述依据《元史》卷030〈本纪第30卷本年条〉，但《元史》为明初仓促官修，本纪保存大量实录条目但有编次与纪年问题；行省执行、数字、族称及对外关系须以条格、多语文书和对方史料互校。。研究元代行省、赋役、河工与多语文书研究可以补问材料未直接说出的空间与社会差异。

\noindent\eventanchor{SYD-1335-YUAN-193}{禁僧道買民田，違者坐罪，沒其直}\textbf{元。}
1335年（元至元1年；公元换算据《元史》本纪纪年）的“禁僧道買民田，違者坐罪，沒其直”将元的财政或生产安排推到可见处。制度文本能说明官府打算如何收取、转运或调节，却不能直接等同于不同地区的实收、流通或家庭负担。

从账本的前后项看，“本纪年条记录政令或机构处置；实施背景须参照制度材料。”提供了背景压力；“可在同年及后续政令中追踪；条文不单独证明地方执行。”标出后来需要追踪的变化，二者之间仍存在未被材料填满的环节。

核心线索为《元史》卷030〈本纪第30卷本年条〉。《元史》为明初仓促官修，本纪保存大量实录条目但有编次与纪年问题；行省执行、数字、族称及对外关系须以条格、多语文书和对方史料互校。 现代研究可从元代行省、赋役、河工与多语文书研究继续追问。

\noindent\eventanchor{SYD-1335-YUAN-265}{乙亥，賜公主不答昔你媵戶鈔四千錠}\textbf{元。}
在1335年（元至元1年；公元换算据《元史》本纪纪年）的可见材料中，元留下“乙亥，賜公主不答昔你媵戶鈔四千錠”。这一动作把具体人群、机构或地方条件带进年序，但一项记录不能替未被记载者概括整体反应。

其前提记为“本纪所记财用、赈济或征发属于特定报告场景。”；其后续是“可与食货志、文书和物质材料比较，不外推为全境总量。”。这是一条待后续材料检验的事件链，而不是对所有行动者意图的替写。

可核的材料入口是《元史》卷030〈本纪第30卷本年条〉。《元史》为明初仓促官修，本纪保存大量实录条目但有编次与纪年问题；行省执行、数字、族称及对外关系须以条格、多语文书和对方史料互校。 就此而言，元代行省、赋役、河工与多语文书研究有助于把年序同区域、制度或物质材料并读。

\subsection{1336年（元至元2年；公元换算据《元史》本纪纪年）}

\noindent\eventanchor{SYD-1336-YUAN-050}{:洪惟太祖皇帝膺期撫運，肇開帝業}\textbf{元。}
元的1336年（元至元2年；公元换算据《元史》本纪纪年）记录写到“:洪惟太祖皇帝膺期撫運，肇開帝業”。我们可以据此辨认一个可见的选择或压力，而不能把零散的官方记载、碑刻或后出叙事拼成无缝的社会全景。

账本把这一节点放在“本纪年条提供日期和中枢可见行动；前因不据单条补定。”与“可回查同年及后续条目，地方影响仍须另证。”之间；两者提示前后压力，并不把时间相邻写成单一因果。

本条首先回到《元史》卷027〈本纪第27卷本年条〉；《元史》为明初仓促官修，本纪保存大量实录条目但有编次与纪年问题；行省执行、数字、族称及对外关系须以条格、多语文书和对方史料互校。 因而元代行省、赋役、河工与多语文书研究不是装饰性的书目，而是检验这一叙述边界的路径。

\noindent\eventanchor{SYD-1336-YUAN-122}{安南國遣其臣鄧恭儉來貢方物}\textbf{元与安南。}
元与安南的1336年（元至元2年；公元换算据《元史》本纪纪年）记录写到“安南國遣其臣鄧恭儉來貢方物”。我们可以据此辨认一个可见的选择或压力，而不能把零散的官方记载、碑刻或后出叙事拼成无缝的社会全景。

此前的条件是“本纪年条提供日期和中枢可见行动；前因不据单条补定。”，此后的可见影响是“可回查同年及后续条目，地方影响仍须另证。”。这里的“影响”须按地点、机构和材料种类分别衡量。

这一叙述依据《元史》卷027〈本纪第27卷本年条〉；《安南志略》及当地碑刻材料（互校），但《元史》为明初仓促官修，本纪保存大量实录条目但有编次与纪年问题；行省执行、数字、族称及对外关系须以条格、多语文书和对方史料互校。。研究元代行省、赋役、河工与多语文书研究可以补问材料未直接说出的空间与社会差异。

\noindent\eventanchor{SYD-1336-YUAN-194}{帝問曰：「所賜為誰？」對曰：「左丞相阿散所得為多}\textbf{元。}
元的1336年（元至元2年；公元换算据《元史》本纪纪年）记录写到“帝問曰：「所賜為誰？」對曰：「左丞相阿散所得為多”。我们可以据此辨认一个可见的选择或压力，而不能把零散的官方记载、碑刻或后出叙事拼成无缝的社会全景。

从账本的前后项看，“本纪年条提供日期和中枢可见行动；前因不据单条补定。”提供了背景压力；“可回查同年及后续条目，地方影响仍须另证。”标出后来需要追踪的变化，二者之间仍存在未被材料填满的环节。

核心线索为《元史》卷027〈本纪第27卷本年条〉。《元史》为明初仓促官修，本纪保存大量实录条目但有编次与纪年问题；行省执行、数字、族称及对外关系须以条格、多语文书和对方史料互校。 现代研究可从元代行省、赋役、河工与多语文书研究继续追问。

\noindent\eventanchor{SYD-1336-YUAN-266}{帝曰：「言事者直至朕前可也，如細民輒訴訟者則禁之}\textbf{元。}
元的1336年（元至元2年；公元换算据《元史》本纪纪年）记录写到“帝曰：「言事者直至朕前可也，如細民輒訴訟者則禁之”。我们可以据此辨认一个可见的选择或压力，而不能把零散的官方记载、碑刻或后出叙事拼成无缝的社会全景。

其前提记为“本纪年条记录政令或机构处置；实施背景须参照制度材料。”；其后续是“可在同年及后续政令中追踪；条文不单独证明地方执行。”。这是一条待后续材料检验的事件链，而不是对所有行动者意图的替写。

可核的材料入口是《元史》卷027〈本纪第27卷本年条〉。《元史》为明初仓促官修，本纪保存大量实录条目但有编次与纪年问题；行省执行、数字、族称及对外关系须以条格、多语文书和对方史料互校。 就此而言，元代行省、赋役、河工与多语文书研究有助于把年序同区域、制度或物质材料并读。

\subsection{1338年（元至元4年；公元换算据《元史》本纪纪年）}

\noindent\eventanchor{SYD-1338-YUAN-051}{敕：「諸王、駙馬戶在縉山、懷來、永興縣者，與民均服徭役}\textbf{元。}
元的1338年（元至元4年；公元换算据《元史》本纪纪年）记录写到“敕：「諸王、駙馬戶在縉山、懷來、永興縣者，與民均服徭役”。我们可以据此辨认一个可见的选择或压力，而不能把零散的官方记载、碑刻或后出叙事拼成无缝的社会全景。

其前提记为“本纪年条提供日期和中枢可见行动；前因不据单条补定。”；其后续是“可回查同年及后续条目，地方影响仍须另证。”。这是一条待后续材料检验的事件链，而不是对所有行动者意图的替写。

可核的材料入口是《元史》卷024〈本纪第24卷本年条〉。《元史》为明初仓促官修，本纪保存大量实录条目但有编次与纪年问题；行省执行、数字、族称及对外关系须以条格、多语文书和对方史料互校。 就此而言，元代行省、赋役、河工与多语文书研究有助于把年序同区域、制度或物质材料并读。

\noindent\eventanchor{SYD-1338-YUAN-123}{敕直省舍人以其半給事殿庭，半聽中書差遣}\textbf{元。}
元的1338年（元至元4年；公元换算据《元史》本纪纪年）记录写到“敕直省舍人以其半給事殿庭，半聽中書差遣”。我们可以据此辨认一个可见的选择或压力，而不能把零散的官方记载、碑刻或后出叙事拼成无缝的社会全景。

从账本的前后项看，“本纪年条记录政令或机构处置；实施背景须参照制度材料。”提供了背景压力；“可在同年及后续政令中追踪；条文不单独证明地方执行。”标出后来需要追踪的变化，二者之间仍存在未被材料填满的环节。

核心线索为《元史》卷024〈本纪第24卷本年条〉。《元史》为明初仓促官修，本纪保存大量实录条目但有编次与纪年问题；行省执行、数字、族称及对外关系须以条格、多语文书和对方史料互校。 现代研究可从元代行省、赋役、河工与多语文书研究继续追问。

\noindent\eventanchor{SYD-1338-YUAN-195}{帝納其言，凡營繕悉罷之}\textbf{元。}
元的1338年（元至元4年；公元换算据《元史》本纪纪年）记录写到“帝納其言，凡營繕悉罷之”。我们可以据此辨认一个可见的选择或压力，而不能把零散的官方记载、碑刻或后出叙事拼成无缝的社会全景。

账本把这一节点放在“本纪年条记录政令或机构处置；实施背景须参照制度材料。”与“可在同年及后续政令中追踪；条文不单独证明地方执行。”之间；两者提示前后压力，并不把时间相邻写成单一因果。

本条首先回到《元史》卷024〈本纪第24卷本年条〉；《元史》为明初仓促官修，本纪保存大量实录条目但有编次与纪年问题；行省执行、数字、族称及对外关系须以条格、多语文书和对方史料互校。 因而元代行省、赋役、河工与多语文书研究不是装饰性的书目，而是检验这一叙述边界的路径。

\noindent\eventanchor{SYD-1338-YUAN-267}{帝謂省臣曰：「安南國王慕義來歸，宜厚其賜，以懷遠人，其進勳爵、受田如故}\textbf{元与安南。}
元与安南在1338年（元至元4年；公元换算据《元史》本纪纪年）留下的这条记录是“帝謂省臣曰：「安南國王慕義來歸，宜厚其賜，以懷遠人，其進勳爵、受田如故”。它照亮资源如何被命名和调度，也暴露材料的盲点：数字、额度或禁令若无地方文书和物质材料互证，不能被写成全境同速的现实。

此前的条件是“本纪年条记录政令或机构处置；实施背景须参照制度材料。”，此后的可见影响是“可在同年及后续政令中追踪；条文不单独证明地方执行。”。这里的“影响”须按地点、机构和材料种类分别衡量。

这一叙述依据《元史》卷024〈本纪第24卷本年条〉；《安南志略》及当地碑刻材料（互校），但《元史》为明初仓促官修，本纪保存大量实录条目但有编次与纪年问题；行省执行、数字、族称及对外关系须以条格、多语文书和对方史料互校。。研究元代行省、赋役、河工与多语文书研究可以补问材料未直接说出的空间与社会差异。

\subsection{1339年（元至元5年；公元换算据《元史》本纪纪年）}

\noindent\eventanchor{SYD-1339-YUAN-052}{:大都、上都、興和三路，差稅免三年}\textbf{元。}
“:大都、上都、興和三路，差稅免三年”是元在1339年（元至元5年；公元换算据《元史》本纪纪年）可追查的节点。它提示命令、环境、宗教、迁徙或地方组织如何进入政治过程；影响的范围和速度仍要避免由一条材料外推。

其前提记为“本纪所记财用、赈济或征发属于特定报告场景。”；其后续是“可与食货志、文书和物质材料比较，不外推为全境总量。”。这是一条待后续材料检验的事件链，而不是对所有行动者意图的替写。

可核的材料入口是《元史》卷037〈本纪第37卷本年条〉。《元史》为明初仓促官修，本纪保存大量实录条目但有编次与纪年问题；行省执行、数字、族称及对外关系须以条格、多语文书和对方史料互校。 就此而言，元代行省、赋役、河工与多语文书研究有助于把年序同区域、制度或物质材料并读。

\noindent\eventanchor{SYD-1339-YUAN-124}{丙寅，楚丘縣河堤壞，發民丁二千三百五十人修之}\textbf{元。}
“丙寅，楚丘縣河堤壞，發民丁二千三百五十人修之”是元在1339年（元至元5年；公元换算据《元史》本纪纪年）可追查的节点。它提示命令、环境、宗教、迁徙或地方组织如何进入政治过程；影响的范围和速度仍要避免由一条材料外推。

从账本的前后项看，“本纪保留灾害或工程的报告地点与朝廷处置；灾情范围需另证。”提供了背景压力；“可与地方文书、河工和环境材料比较，不将单点记录外推为全域因果。”标出后来需要追踪的变化，二者之间仍存在未被材料填满的环节。

核心线索为《元史》卷037〈本纪第37卷本年条〉。《元史》为明初仓促官修，本纪保存大量实录条目但有编次与纪年问题；行省执行、数字、族称及对外关系须以条格、多语文书和对方史料互校。 现代研究可从元代行省、赋役、河工与多语文书研究继续追问。

\noindent\eventanchor{SYD-1339-YUAN-196}{賓天之日，皇后傳顧命於太師、太平王、右丞相、答剌罕燕帖木兒，太保、浚寧王、知樞密院事伯顏等，謂聖體彌留，益推固讓之初志，以宗社之重，屬諸大兄忽都篤皇帝之世嫡}\textbf{元。}
“賓天之日，皇后傳顧命於太師、太平王、右丞相、答剌罕燕帖木兒，太保、浚寧王、知樞密院事伯顏等，謂聖體彌留，益推固讓之初志，以宗社之重，屬諸大兄忽都篤皇帝之世嫡”是元在1339年（元至元5年；公元换算据《元史》本纪纪年）可追查的节点。它提示命令、环境、宗教、迁徙或地方组织如何进入政治过程；影响的范围和速度仍要避免由一条材料外推。

账本把这一节点放在“本纪年条记录政令或机构处置；实施背景须参照制度材料。”与“可在同年及后续政令中追踪；条文不单独证明地方执行。”之间；两者提示前后压力，并不把时间相邻写成单一因果。

本条首先回到《元史》卷037〈本纪第37卷本年条〉；《元史》为明初仓促官修，本纪保存大量实录条目但有编次与纪年问题；行省执行、数字、族称及对外关系须以条格、多语文书和对方史料互校。 因而元代行省、赋役、河工与多语文书研究不是装饰性的书目，而是检验这一叙述边界的路径。

\noindent\eventanchor{SYD-1339-YUAN-268}{丙辰，給宿衞士、蒙古、漢軍三萬人禦寒衣}\textbf{元。}
“丙辰，給宿衞士、蒙古、漢軍三萬人禦寒衣”是元在1339年（元至元5年；公元换算据《元史》本纪纪年）可追查的节点。它提示命令、环境、宗教、迁徙或地方组织如何进入政治过程；影响的范围和速度仍要避免由一条材料外推。

此前的条件是“本纪所记财用、赈济或征发属于特定报告场景。”，此后的可见影响是“可与食货志、文书和物质材料比较，不外推为全境总量。”。这里的“影响”须按地点、机构和材料种类分别衡量。

这一叙述依据《元史》卷037〈本纪第37卷本年条〉，但《元史》为明初仓促官修，本纪保存大量实录条目但有编次与纪年问题；行省执行、数字、族称及对外关系须以条格、多语文书和对方史料互校。。研究元代行省、赋役、河工与多语文书研究可以补问材料未直接说出的空间与社会差异。

\subsection{1340年（元至元6年；公元换算据《元史》本纪纪年）}

\noindent\eventanchor{SYD-1340-YUAN-053}{敕更賜璽書三十二、圓符四，仍究詰所失者}\textbf{元。}
“敕更賜璽書三十二、圓符四，仍究詰所失者”是元在1340年（元至元6年；公元换算据《元史》本纪纪年）可追查的节点。它提示命令、环境、宗教、迁徙或地方组织如何进入政治过程；影响的范围和速度仍要避免由一条材料外推。

此前的条件是“本纪年条提供日期和中枢可见行动；前因不据单条补定。”，此后的可见影响是“可回查同年及后续条目，地方影响仍须另证。”。这里的“影响”须按地点、机构和材料种类分别衡量。

这一叙述依据《元史》卷036〈本纪第36卷本年条〉，但《元史》为明初仓促官修，本纪保存大量实录条目但有编次与纪年问题；行省执行、数字、族称及对外关系须以条格、多语文书和对方史料互校。。研究元代行省、赋役、河工与多语文书研究可以补问材料未直接说出的空间与社会差异。

\noindent\eventanchor{SYD-1340-YUAN-125}{庚辰，以安陸府賜并王晃火〔帖木兒〕為食邑}\textbf{元。}
“庚辰，以安陸府賜并王晃火〔帖木兒〕為食邑”是元在1340年（元至元6年；公元换算据《元史》本纪纪年）可追查的节点。它提示命令、环境、宗教、迁徙或地方组织如何进入政治过程；影响的范围和速度仍要避免由一条材料外推。

账本把这一节点放在“本纪年条提供日期和中枢可见行动；前因不据单条补定。”与“可回查同年及后续条目，地方影响仍须另证。”之间；两者提示前后压力，并不把时间相邻写成单一因果。

本条首先回到《元史》卷036〈本纪第36卷本年条〉；《元史》为明初仓促官修，本纪保存大量实录条目但有编次与纪年问题；行省执行、数字、族称及对外关系须以条格、多语文书和对方史料互校。 因而元代行省、赋役、河工与多语文书研究不是装饰性的书目，而是检验这一叙述边界的路径。

\noindent\eventanchor{SYD-1340-YUAN-197}{壬午，江西行省造螺鈿几榻遺燕鐵木兒，詔賜匠者幣帛各一}\textbf{元。}
“壬午，江西行省造螺鈿几榻遺燕鐵木兒，詔賜匠者幣帛各一”是元在1340年（元至元6年；公元换算据《元史》本纪纪年）可追查的节点。它提示命令、环境、宗教、迁徙或地方组织如何进入政治过程；影响的范围和速度仍要避免由一条材料外推。

其前提记为“本纪年条记录政令或机构处置；实施背景须参照制度材料。”；其后续是“可在同年及后续政令中追踪；条文不单独证明地方执行。”。这是一条待后续材料检验的事件链，而不是对所有行动者意图的替写。

可核的材料入口是《元史》卷036〈本纪第36卷本年条〉。《元史》为明初仓促官修，本纪保存大量实录条目但有编次与纪年问题；行省执行、数字、族称及对外关系须以条格、多语文书和对方史料互校。 就此而言，元代行省、赋役、河工与多语文书研究有助于把年序同区域、制度或物质材料并读。

\noindent\eventanchor{SYD-1340-YUAN-269}{順元宣撫司亦言：「賊列行營為十六所，乞調兵分道備禦}\textbf{元。}
“順元宣撫司亦言：「賊列行營為十六所，乞調兵分道備禦”把元置于军事行动的可核时点。战报能够记下攻守、入城或撤离，却很少完整保留沿途征发、居民去留和补给损耗；因此不能把一方的胜败叙事延伸成所有地方的经验。

从账本的前后项看，“本纪从朝廷视角记录军政行动；出兵与战果需分辨。”提供了背景压力；“后续路线、控制与损失应与对方记录和遗址材料复核。”标出后来需要追踪的变化，二者之间仍存在未被材料填满的环节。

核心线索为《元史》卷036〈本纪第36卷本年条〉。《元史》为明初仓促官修，本纪保存大量实录条目但有编次与纪年问题；行省执行、数字、族称及对外关系须以条格、多语文书和对方史料互校。 现代研究可从元代行省、赋役、河工与多语文书研究继续追问。

\subsection{1340年}

\noindent\eventanchor{YUAN-1340-EVENT-189}{脱脱掌政后推动修史、钞法和行政整顿}\textbf{元。}
1340年的“脱脱掌政后推动修史、钞法和行政整顿”将元的财政或生产安排推到可见处。制度文本能说明官府打算如何收取、转运或调节，却不能直接等同于不同地区的实收、流通或家庭负担。

账本把这一节点放在“元中后期继承、财政与地方危机对中枢资源、交通或行政协同的要求，推动了“脱脱掌政后推动修史、钞法和行政整顿”这一制度或空间安排。”与““脱脱掌政后推动修史、钞法和行政整顿”提供新的治理框架或资源通道，但颁行、复申与不同地区的实际执行须分开核验。”之间；两者提示前后压力，并不把时间相邻写成单一因果。

本条首先回到《元史》相关本纪；《元朝秘史》、波斯文史籍或《高丽史》对应叙事；《元典章》《通制条格》、多语碑刻与文书（据事核）；《元史》明初速修，纪年与人名有异同；蒙古大汗政权不等同所有汗国，征服、赋役和身份分类须按地区及材料语言分列。 因而蒙古帝国、元朝行省财政、多语文书与区域社会研究不是装饰性的书目，而是检验这一叙述边界的路径。

\subsection{1342年（元至正2年；公元换算据《元史》本纪纪年）}

\noindent\eventanchor{SYD-1342-YUAN-054}{命威順王寬徹不花還鎮湖廣}\textbf{元。}
“命威順王寬徹不花還鎮湖廣”是元在1342年（元至正2年；公元换算据《元史》本纪纪年）可追查的节点。它提示命令、环境、宗教、迁徙或地方组织如何进入政治过程；影响的范围和速度仍要避免由一条材料外推。

从账本的前后项看，“本纪年条记录政令或机构处置；实施背景须参照制度材料。”提供了背景压力；“可在同年及后续政令中追踪；条文不单独证明地方执行。”标出后来需要追踪的变化，二者之间仍存在未被材料填满的环节。

核心线索为《元史》卷039〈本纪第39卷本年条〉。《元史》为明初仓促官修，本纪保存大量实录条目但有编次与纪年问题；行省执行、数字、族称及对外关系须以条格、多语文书和对方史料互校。 现代研究可从元代行省、赋役、河工与多语文书研究继续追问。

\noindent\eventanchor{SYD-1342-YUAN-126}{安豐路饑，賑糶麥四萬二千四百石}\textbf{元。}
“安豐路饑，賑糶麥四萬二千四百石”是元在1342年（元至正2年；公元换算据《元史》本纪纪年）可追查的节点。它提示命令、环境、宗教、迁徙或地方组织如何进入政治过程；影响的范围和速度仍要避免由一条材料外推。

其前提记为“本纪所记财用、赈济或征发属于特定报告场景。”；其后续是“可与食货志、文书和物质材料比较，不外推为全境总量。”。这是一条待后续材料检验的事件链，而不是对所有行动者意图的替写。

可核的材料入口是《元史》卷039〈本纪第39卷本年条〉。《元史》为明初仓促官修，本纪保存大量实录条目但有编次与纪年问题；行省执行、数字、族称及对外关系须以条格、多语文书和对方史料互校。 就此而言，元代行省、赋役、河工与多语文书研究有助于把年序同区域、制度或物质材料并读。

\noindent\eventanchor{SYD-1342-YUAN-198}{丙申，命參知政事納麟監繪明宗皇帝御容}\textbf{元。}
“丙申，命參知政事納麟監繪明宗皇帝御容”是元在1342年（元至正2年；公元换算据《元史》本纪纪年）可追查的节点。它提示命令、环境、宗教、迁徙或地方组织如何进入政治过程；影响的范围和速度仍要避免由一条材料外推。

此前的条件是“本纪年条记录政令或机构处置；实施背景须参照制度材料。”，此后的可见影响是“可在同年及后续政令中追踪；条文不单独证明地方执行。”。这里的“影响”须按地点、机构和材料种类分别衡量。

这一叙述依据《元史》卷039〈本纪第39卷本年条〉，但《元史》为明初仓促官修，本纪保存大量实录条目但有编次与纪年问题；行省执行、数字、族称及对外关系须以条格、多语文书和对方史料互校。。研究元代行省、赋役、河工与多语文书研究可以补问材料未直接说出的空间与社会差异。

\noindent\eventanchor{SYD-1342-YUAN-270}{丙子，命興元府鳳州留壩鎮及晉寧路遼山縣十八盤各立巡檢司}\textbf{元。}
元的“丙子，命興元府鳳州留壩鎮及晉寧路遼山縣十八盤各立巡檢司”并非只改换一个统治者。任官、赏赐、军权与信息通路要在随后重新接合；材料较能确认先后，较难替当事人写出唯一动机。

从账本的前后项看，“本纪年条记录政令或机构处置；实施背景须参照制度材料。”提供了背景压力；“可在同年及后续政令中追踪；条文不单独证明地方执行。”标出后来需要追踪的变化，二者之间仍存在未被材料填满的环节。

核心线索为《元史》卷039〈本纪第39卷本年条〉。《元史》为明初仓促官修，本纪保存大量实录条目但有编次与纪年问题；行省执行、数字、族称及对外关系须以条格、多语文书和对方史料互校。 现代研究可从元代行省、赋役、河工与多语文书研究继续追问。

\subsection{1343年（元至正3年；公元换算据《元史》本纪纪年）}

\noindent\eventanchor{SYD-1343-YUAN-055}{是月，臨江路新淦州、新喻州，瑞州民饑，賑糶米二萬石}\textbf{元。}
在1343年（元至正3年；公元换算据《元史》本纪纪年）的可见材料中，元留下“是月，臨江路新淦州、新喻州，瑞州民饑，賑糶米二萬石”。这一动作把具体人群、机构或地方条件带进年序，但一项记录不能替未被记载者概括整体反应。

从账本的前后项看，“本纪所记财用、赈济或征发属于特定报告场景。”提供了背景压力；“可与食货志、文书和物质材料比较，不外推为全境总量。”标出后来需要追踪的变化，二者之间仍存在未被材料填满的环节。

核心线索为《元史》卷039〈本纪第39卷本年条〉。《元史》为明初仓促官修，本纪保存大量实录条目但有编次与纪年问题；行省执行、数字、族称及对外关系须以条格、多语文书和对方史料互校。 现代研究可从元代行省、赋役、河工与多语文书研究继续追问。

\noindent\eventanchor{SYD-1343-YUAN-127}{以米八千石、鈔二千八百錠，賑哈剌奴兒饑民}\textbf{元。}
在1343年（元至正3年；公元换算据《元史》本纪纪年）的可见材料中，元留下“以米八千石、鈔二千八百錠，賑哈剌奴兒饑民”。这一动作把具体人群、机构或地方条件带进年序，但一项记录不能替未被记载者概括整体反应。

其前提记为“本纪所记财用、赈济或征发属于特定报告场景。”；其后续是“可与食货志、文书和物质材料比较，不外推为全境总量。”。这是一条待后续材料检验的事件链，而不是对所有行动者意图的替写。

可核的材料入口是《元史》卷039〈本纪第39卷本年条〉。《元史》为明初仓促官修，本纪保存大量实录条目但有编次与纪年问题；行省执行、数字、族称及对外关系须以条格、多语文书和对方史料互校。 就此而言，元代行省、赋役、河工与多语文书研究有助于把年序同区域、制度或物质材料并读。

\noindent\eventanchor{SYD-1343-YUAN-199}{〔乙〕丑，命宗王燕帖木兒為大宗正府札魯忽赤}\textbf{元。}
在1343年（元至正3年；公元换算据《元史》本纪纪年）的可见材料中，元留下“〔乙〕丑，命宗王燕帖木兒為大宗正府札魯忽赤”。这一动作把具体人群、机构或地方条件带进年序，但一项记录不能替未被记载者概括整体反应。

此前的条件是“本纪年条记录政令或机构处置；实施背景须参照制度材料。”，此后的可见影响是“可在同年及后续政令中追踪；条文不单独证明地方执行。”。这里的“影响”须按地点、机构和材料种类分别衡量。

这一叙述依据《元史》卷039〈本纪第39卷本年条〉，但《元史》为明初仓促官修，本纪保存大量实录条目但有编次与纪年问题；行省执行、数字、族称及对外关系须以条格、多语文书和对方史料互校。。研究元代行省、赋役、河工与多语文书研究可以补问材料未直接说出的空间与社会差异。

\noindent\eventanchor{SYD-1343-YUAN-271}{棒胡本陳州人，名閏兒，以燒香惑眾，妄造妖言作亂，破歸德府鹿邑，焚陳州，屯營於杏岡，命河南行省左丞慶童領兵討之}\textbf{元。}
1343年（元至正3年；公元换算据《元史》本纪纪年），元的行动落在“棒胡本陳州人，名閏兒，以燒香惑眾，妄造妖言作亂，破歸德府鹿邑，焚陳州，屯營於杏岡，命河南行省左丞慶童領兵討之”所指的战事与通道上。可见的不是抽象的推进：攻守双方必须让人员、消息与补给穿过具体的城防、道路或水路；账本所列的结果只说明这一节点改变了下一步选择。

从账本的前后项看，“本纪年条记录政令或机构处置；实施背景须参照制度材料。”提供了背景压力；“可在同年及后续政令中追踪；条文不单独证明地方执行。”标出后来需要追踪的变化，二者之间仍存在未被材料填满的环节。

核心线索为《元史》卷039〈本纪第39卷本年条〉。《元史》为明初仓促官修，本纪保存大量实录条目但有编次与纪年问题；行省执行、数字、族称及对外关系须以条格、多语文书和对方史料互校。 现代研究可从元代行省、赋役、河工与多语文书研究继续追问。

\subsection{1344年（元至正4年；公元换算据《元史》本纪纪年）}

\noindent\eventanchor{SYD-1344-YUAN-056}{:洪惟我太祖皇帝混一海宇，爰立定制，以一統緒，宗親各受分地，勿敢妄生覬覦，此不易之成規，萬世所共守者也}\textbf{元。}
在1344年（元至正4年；公元换算据《元史》本纪纪年）的材料中，元面对的是:洪惟我太祖皇帝混一海宇，爰立定制，以一統緒，宗親各受分地，勿敢妄生覬覦，此不易之成規，萬世所共守者也。这里应先看行动所依赖的关隘、城镇、河道或粮道，再讨论政治后果；军队抵达某处，并不等于它已能稳定征收、安置或控制周边。

从账本的前后项看，“本纪年条提供日期和中枢可见行动；前因不据单条补定。”提供了背景压力；“可回查同年及后续条目，地方影响仍须另证。”标出后来需要追踪的变化，二者之间仍存在未被材料填满的环节。

核心线索为《元史》卷032〈本纪第32卷本年条〉。《元史》为明初仓促官修，本纪保存大量实录条目但有编次与纪年问题；行省执行、数字、族称及对外关系须以条格、多语文书和对方史料互校。 现代研究可从元代行省、赋役、河工与多语文书研究继续追问。

\noindent\eventanchor{SYD-1344-YUAN-128}{:於戲，朕豈有意於天下哉！重念祖宗開創之艱，恐隳大業，是以勉徇輿情}\textbf{元。}
元的1344年（元至正4年；公元换算据《元史》本纪纪年）记录写到“:於戲，朕豈有意於天下哉！重念祖宗開創之艱，恐隳大業，是以勉徇輿情”。我们可以据此辨认一个可见的选择或压力，而不能把零散的官方记载、碑刻或后出叙事拼成无缝的社会全景。

其前提记为“本纪年条提供日期和中枢可见行动；前因不据单条补定。”；其后续是“可回查同年及后续条目，地方影响仍须另证。”。这是一条待后续材料检验的事件链，而不是对所有行动者意图的替写。

可核的材料入口是《元史》卷032〈本纪第32卷本年条〉。《元史》为明初仓促官修，本纪保存大量实录条目但有编次与纪年问题；行省执行、数字、族称及对外关系须以条格、多语文书和对方史料互校。 就此而言，元代行省、赋役、河工与多语文书研究有助于把年序同区域、制度或物质材料并读。

\noindent\eventanchor{SYD-1344-YUAN-200}{賜西安王阿剌忒納失里、鎮南王帖木兒不花、威順王寬徹不花、宣靖王買奴等，金各五十兩、銀各五百兩、幣各三十匹}\textbf{元。}
元的1344年（元至正4年；公元换算据《元史》本纪纪年）记录写到“賜西安王阿剌忒納失里、鎮南王帖木兒不花、威順王寬徹不花、宣靖王買奴等，金各五十兩、銀各五百兩、幣各三十匹”。我们可以据此辨认一个可见的选择或压力，而不能把零散的官方记载、碑刻或后出叙事拼成无缝的社会全景。

此前的条件是“本纪年条提供日期和中枢可见行动；前因不据单条补定。”，此后的可见影响是“可回查同年及后续条目，地方影响仍须另证。”。这里的“影响”须按地点、机构和材料种类分别衡量。

这一叙述依据《元史》卷032〈本纪第32卷本年条〉，但《元史》为明初仓促官修，本纪保存大量实录条目但有编次与纪年问题；行省执行、数字、族称及对外关系须以条格、多语文书和对方史料互校。。研究元代行省、赋役、河工与多语文书研究可以补问材料未直接说出的空间与社会差异。

\noindent\eventanchor{SYD-1344-YUAN-272}{帝謂中書省臣曰：「朕在瓊州、建康時，撒迪皆從，備極艱苦，其賜鹽引六萬，俾規利以贍其家}\textbf{元。}
元的1344年（元至正4年；公元换算据《元史》本纪纪年）记录写到“帝謂中書省臣曰：「朕在瓊州、建康時，撒迪皆從，備極艱苦，其賜鹽引六萬，俾規利以贍其家”。我们可以据此辨认一个可见的选择或压力，而不能把零散的官方记载、碑刻或后出叙事拼成无缝的社会全景。

从账本的前后项看，“本纪年条记录政令或机构处置；实施背景须参照制度材料。”提供了背景压力；“可在同年及后续政令中追踪；条文不单独证明地方执行。”标出后来需要追踪的变化，二者之间仍存在未被材料填满的环节。

核心线索为《元史》卷032〈本纪第32卷本年条〉。《元史》为明初仓促官修，本纪保存大量实录条目但有编次与纪年问题；行省执行、数字、族称及对外关系须以条格、多语文书和对方史料互校。 现代研究可从元代行省、赋役、河工与多语文书研究继续追问。

\subsection{1344年}

\noindent\eventanchor{YUAN-1344-EVENT-190}{黄河决溢及相关灾害加深北方社会与财政压力}\textbf{元。}
元的1344年记录写到“黄河决溢及相关灾害加深北方社会与财政压力”。我们可以据此辨认一个可见的选择或压力，而不能把零散的官方记载、碑刻或后出叙事拼成无缝的社会全景。

此前的条件是““黄河决溢及相关灾害加深北方社会与财政压力”发生在元中后期继承、财政与地方危机的具体压力与机会之中，相关行动者的选择不能脱离此前事件链解释。”，此后的可见影响是““黄河决溢及相关灾害加深北方社会与财政压力”为后续政治、边防或资源配置留下条件；其社会影响与实际覆盖范围仍须以异类材料限定。”。这里的“影响”须按地点、机构和材料种类分别衡量。

这一叙述依据《元史》相关本纪；《元朝秘史》、波斯文史籍或《高丽史》对应叙事；《元典章》《通制条格》、多语碑刻与文书（据事核），但《元史》明初速修，纪年与人名有异同；蒙古大汗政权不等同所有汗国，征服、赋役和身份分类须按地区及材料语言分列。。研究蒙古帝国、元朝行省财政、多语文书与区域社会研究可以补问材料未直接说出的空间与社会差异。

\subsection{1346年（元至正6年；公元换算据《元史》本纪纪年）}

\noindent\eventanchor{SYD-1346-YUAN-057}{丙辰，賜雲南行省參知政事不老三珠虎符，以兵討死可伐}\textbf{元。}
在1346年（元至正6年；公元换算据《元史》本纪纪年）的材料中，元面对的是丙辰，賜雲南行省參知政事不老三珠虎符，以兵討死可伐。这里应先看行动所依赖的关隘、城镇、河道或粮道，再讨论政治后果；军队抵达某处，并不等于它已能稳定征收、安置或控制周边。

账本把这一节点放在“本纪年条记录政令或机构处置；实施背景须参照制度材料。”与“可在同年及后续政令中追踪；条文不单独证明地方执行。”之间；两者提示前后压力，并不把时间相邻写成单一因果。

本条首先回到《元史》卷040〈本纪第40卷本年条〉；《元史》为明初仓促官修，本纪保存大量实录条目但有编次与纪年问题；行省执行、数字、族称及对外关系须以条格、多语文书和对方史料互校。 因而元代行省、赋役、河工与多语文书研究不是装饰性的书目，而是检验这一叙述边界的路径。

\noindent\eventanchor{SYD-1346-YUAN-129}{丙午，命中書右丞太平、樞密副使姚庸、御史中丞張起巖知經筵事}\textbf{元。}
元于1346年（元至正6年；公元换算据《元史》本纪纪年）处理“丙午，命中書右丞太平、樞密副使姚庸、御史中丞張起巖知經筵事”时，面对的是道路、翻译、礼仪和资源交换共同构成的关系网。书面称谓可见，地方执行和持续性则仍须由后续材料检验。

此前的条件是“本纪年条记录政令或机构处置；实施背景须参照制度材料。”，此后的可见影响是“可在同年及后续政令中追踪；条文不单独证明地方执行。”。这里的“影响”须按地点、机构和材料种类分别衡量。

这一叙述依据《元史》卷040〈本纪第40卷本年条〉，但《元史》为明初仓促官修，本纪保存大量实录条目但有编次与纪年问题；行省执行、数字、族称及对外关系须以条格、多语文书和对方史料互校。。研究元代行省、赋役、河工与多语文书研究可以补问材料未直接说出的空间与社会差异。

\noindent\eventanchor{SYD-1346-YUAN-201}{丙戌，開京師金口河，深五十尺，廣一百五十尺，役夫一十萬}\textbf{元。}
元的1346年（元至正6年；公元换算据《元史》本纪纪年）记录写到“丙戌，開京師金口河，深五十尺，廣一百五十尺，役夫一十萬”。我们可以据此辨认一个可见的选择或压力，而不能把零散的官方记载、碑刻或后出叙事拼成无缝的社会全景。

从账本的前后项看，“本纪保留灾害或工程的报告地点与朝廷处置；灾情范围需另证。”提供了背景压力；“可与地方文书、河工和环境材料比较，不将单点记录外推为全域因果。”标出后来需要追踪的变化，二者之间仍存在未被材料填满的环节。

核心线索为《元史》卷040〈本纪第40卷本年条〉。《元史》为明初仓促官修，本纪保存大量实录条目但有编次与纪年问题；行省执行、数字、族称及对外关系须以条格、多语文书和对方史料互校。 现代研究可从元代行省、赋役、河工与多语文书研究继续追问。

\noindent\eventanchor{SYD-1346-YUAN-273}{大名路饑，以鈔萬二千錠賑之}\textbf{元。}
元的1346年（元至正6年；公元换算据《元史》本纪纪年）记录写到“大名路饑，以鈔萬二千錠賑之”。我们可以据此辨认一个可见的选择或压力，而不能把零散的官方记载、碑刻或后出叙事拼成无缝的社会全景。

账本把这一节点放在“本纪所记财用、赈济或征发属于特定报告场景。”与“可与食货志、文书和物质材料比较，不外推为全境总量。”之间；两者提示前后压力，并不把时间相邻写成单一因果。

本条首先回到《元史》卷040〈本纪第40卷本年条〉；《元史》为明初仓促官修，本纪保存大量实录条目但有编次与纪年问题；行省执行、数字、族称及对外关系须以条格、多语文书和对方史料互校。 因而元代行省、赋役、河工与多语文书研究不是装饰性的书目，而是检验这一叙述边界的路径。

\subsection{1348年}

\noindent\eventanchor{YUAN-1348-EVENT-191}{方国珍在浙东海域起事}\textbf{元。}
1348年，元的行动落在“方国珍在浙东海域起事”所指的战事与通道上。可见的不是抽象的推进：攻守双方必须让人员、消息与补给穿过具体的城防、道路或水路；账本所列的结果只说明这一节点改变了下一步选择。

其前提记为“元中后期继承、财政与地方危机中既有的联盟、交通与补给压力，为“方国珍在浙东海域起事”提供了直接的军事或动员背景。”；其后续是““方国珍在浙东海域起事”改变了相关城镇、道路或政权间的军事选择；伤亡、俘获和控制范围仅可按各方材料分别确认。”。这是一条待后续材料检验的事件链，而不是对所有行动者意图的替写。

可核的材料入口是《元史》相关本纪；《元朝秘史》、波斯文史籍或《高丽史》对应叙事；《元典章》《通制条格》、多语碑刻与文书（据事核）。《元史》明初速修，纪年与人名有异同；蒙古大汗政权不等同所有汗国，征服、赋役和身份分类须按地区及材料语言分列。 就此而言，蒙古帝国、元朝行省财政、多语文书与区域社会研究有助于把年序同区域、制度或物质材料并读。

\subsection{1349年（元至正9年；公元换算据《元史》本纪纪年）}

\noindent\eventanchor{SYD-1349-YUAN-058}{甲申，常州宜興山水出，勢高一丈，壞民廬}\textbf{元。}
在1349年（元至正9年；公元换算据《元史》本纪纪年）的可见材料中，元留下“甲申，常州宜興山水出，勢高一丈，壞民廬”。这一动作把具体人群、机构或地方条件带进年序，但一项记录不能替未被记载者概括整体反应。

此前的条件是“本纪保留灾害或工程的报告地点与朝廷处置；灾情范围需另证。”，此后的可见影响是“可与地方文书、河工和环境材料比较，不将单点记录外推为全域因果。”。这里的“影响”须按地点、机构和材料种类分别衡量。

这一叙述依据《元史》卷040〈本纪第40卷本年条〉，但《元史》为明初仓促官修，本纪保存大量实录条目但有编次与纪年问题；行省执行、数字、族称及对外关系须以条格、多语文书和对方史料互校。。研究元代行省、赋役、河工与多语文书研究可以补问材料未直接说出的空间与社会差异。

\noindent\eventanchor{SYD-1349-YUAN-130}{丙戌，加封瀏陽州道吾山龍神崇惠昭應靈顯廣濟侯}\textbf{元。}
在1349年（元至正9年；公元换算据《元史》本纪纪年）的可见材料中，元留下“丙戌，加封瀏陽州道吾山龍神崇惠昭應靈顯廣濟侯”。这一动作把具体人群、机构或地方条件带进年序，但一项记录不能替未被记载者概括整体反应。

从账本的前后项看，“本纪年条提供日期和中枢可见行动；前因不据单条补定。”提供了背景压力；“可回查同年及后续条目，地方影响仍须另证。”标出后来需要追踪的变化，二者之间仍存在未被材料填满的环节。

核心线索为《元史》卷040〈本纪第40卷本年条〉。《元史》为明初仓促官修，本纪保存大量实录条目但有编次与纪年问题；行省执行、数字、族称及对外关系须以条格、多语文书和对方史料互校。 现代研究可从元代行省、赋役、河工与多语文书研究继续追问。

\noindent\eventanchor{SYD-1349-YUAN-202}{丙子，開上都、興和等處酒禁}\textbf{元。}
在1349年（元至正9年；公元换算据《元史》本纪纪年）的可见材料中，元留下“丙子，開上都、興和等處酒禁”。这一动作把具体人群、机构或地方条件带进年序，但一项记录不能替未被记载者概括整体反应。

其前提记为“本纪年条记录政令或机构处置；实施背景须参照制度材料。”；其后续是“可在同年及后续政令中追踪；条文不单独证明地方执行。”。这是一条待后续材料检验的事件链，而不是对所有行动者意图的替写。

可核的材料入口是《元史》卷040〈本纪第40卷本年条〉。《元史》为明初仓促官修，本纪保存大量实录条目但有编次与纪年问题；行省执行、数字、族称及对外关系须以条格、多语文书和对方史料互校。 就此而言，元代行省、赋役、河工与多语文书研究有助于把年序同区域、制度或物质材料并读。

\noindent\eventanchor{SYD-1349-YUAN-274}{丁丑，封皇姊月魯公主為昌國大長公主}\textbf{元。}
在1349年（元至正9年；公元换算据《元史》本纪纪年）的可见材料中，元留下“丁丑，封皇姊月魯公主為昌國大長公主”。这一动作把具体人群、机构或地方条件带进年序，但一项记录不能替未被记载者概括整体反应。

此前的条件是“本纪年条提供日期和中枢可见行动；前因不据单条补定。”，此后的可见影响是“可回查同年及后续条目，地方影响仍须另证。”。这里的“影响”须按地点、机构和材料种类分别衡量。

这一叙述依据《元史》卷040〈本纪第40卷本年条〉，但《元史》为明初仓促官修，本纪保存大量实录条目但有编次与纪年问题；行省执行、数字、族称及对外关系须以条格、多语文书和对方史料互校。。研究元代行省、赋役、河工与多语文书研究可以补问材料未直接说出的空间与社会差异。

\subsection{1350年（元至正10年；公元换算据《元史》本纪纪年）}

\noindent\eventanchor{SYD-1350-YUAN-059}{:每念治必本於盡孝，事莫先於正名，賴天之靈，權奸屏黜，盡孝正名，不容復緩，永惟鞠育罔極之恩，忍忘不共戴天之義}\textbf{元。}
在1350年（元至正10年；公元换算据《元史》本纪纪年）的可见材料中，元留下“:每念治必本於盡孝，事莫先於正名，賴天之靈，權奸屏黜，盡孝正名，不容復緩，永惟鞠育罔極之恩，忍忘不共戴天之義”。这一动作把具体人群、机构或地方条件带进年序，但一项记录不能替未被记载者概括整体反应。

从账本的前后项看，“本纪年条记录政令或机构处置；实施背景须参照制度材料。”提供了背景压力；“可在同年及后续政令中追踪；条文不单独证明地方执行。”标出后来需要追踪的变化，二者之间仍存在未被材料填满的环节。

核心线索为《元史》卷040〈本纪第40卷本年条〉。《元史》为明初仓促官修，本纪保存大量实录条目但有编次与纪年问题；行省执行、数字、族称及对外关系须以条格、多语文书和对方史料互校。 现代研究可从元代行省、赋役、河工与多语文书研究继续追问。

\noindent\eventanchor{SYD-1350-YUAN-131}{:昔我皇祖武宗皇帝昇遐之後，祖母太皇太后惑於憸慝，俾皇考明宗皇帝出封雲南}\textbf{元。}
在1350年（元至正10年；公元换算据《元史》本纪纪年）的可见材料中，元留下“:昔我皇祖武宗皇帝昇遐之後，祖母太皇太后惑於憸慝，俾皇考明宗皇帝出封雲南”。这一动作把具体人群、机构或地方条件带进年序，但一项记录不能替未被记载者概括整体反应。

账本把这一节点放在“本纪年条提供日期和中枢可见行动；前因不据单条补定。”与“可回查同年及后续条目，地方影响仍须另证。”之间；两者提示前后压力，并不把时间相邻写成单一因果。

本条首先回到《元史》卷040〈本纪第40卷本年条〉；《元史》为明初仓促官修，本纪保存大量实录条目但有编次与纪年问题；行省执行、数字、族称及对外关系须以条格、多语文书和对方史料互校。 因而元代行省、赋役、河工与多语文书研究不是装饰性的书目，而是检验这一叙述边界的路径。

\noindent\eventanchor{SYD-1350-YUAN-203}{以太保馬札兒台為太師、中書右丞相，太尉塔失海牙為太傅，知樞密院事塔馬赤為太保，御史大夫脫脫為知樞密院事，汪家奴為中書平章政事，嶺北行省平章政事也先帖木兒為御史大夫}\textbf{元。}
在1350年（元至正10年；公元换算据《元史》本纪纪年）的可见材料中，元留下“以太保馬札兒台為太師、中書右丞相，太尉塔失海牙為太傅，知樞密院事塔馬赤為太保，御史大夫脫脫為知樞密院事，汪家奴為中書平章政事，嶺北行省平章政事也先帖木兒為御史大夫”。这一动作把具体人群、机构或地方条件带进年序，但一项记录不能替未被记载者概括整体反应。

此前的条件是“本纪年条记录政令或机构处置；实施背景须参照制度材料。”，此后的可见影响是“可在同年及后续政令中追踪；条文不单独证明地方执行。”。这里的“影响”须按地点、机构和材料种类分别衡量。

这一叙述依据《元史》卷040〈本纪第40卷本年条〉，但《元史》为明初仓促官修，本纪保存大量实录条目但有编次与纪年问题；行省执行、数字、族称及对外关系须以条格、多语文书和对方史料互校。。研究元代行省、赋役、河工与多语文书研究可以补问材料未直接说出的空间与社会差异。

\noindent\eventanchor{SYD-1350-YUAN-275}{罷通州、河西務等處抽分按利房，大都東裏山查提領所}\textbf{元。}
在1350年（元至正10年；公元换算据《元史》本纪纪年）的可见材料中，元留下“罷通州、河西務等處抽分按利房，大都東裏山查提領所”。这一动作把具体人群、机构或地方条件带进年序，但一项记录不能替未被记载者概括整体反应。

从账本的前后项看，“本纪年条记录政令或机构处置；实施背景须参照制度材料。”提供了背景压力；“可在同年及后续政令中追踪；条文不单独证明地方执行。”标出后来需要追踪的变化，二者之间仍存在未被材料填满的环节。

核心线索为《元史》卷040〈本纪第40卷本年条〉。《元史》为明初仓促官修，本纪保存大量实录条目但有编次与纪年问题；行省执行、数字、族称及对外关系须以条格、多语文书和对方史料互校。 现代研究可从元代行省、赋役、河工与多语文书研究继续追问。

\subsection{1351年（元至正11年；公元换算据《元史》本纪纪年）}

\noindent\eventanchor{SYD-1351-YUAN-060}{戊寅，詔：「作新風憲}\textbf{元。}
元的1351年（元至正11年；公元换算据《元史》本纪纪年）记录写到“戊寅，詔：「作新風憲”。我们可以据此辨认一个可见的选择或压力，而不能把零散的官方记载、碑刻或后出叙事拼成无缝的社会全景。

此前的条件是“本纪年条记录政令或机构处置；实施背景须参照制度材料。”，此后的可见影响是“可在同年及后续政令中追踪；条文不单独证明地方执行。”。这里的“影响”须按地点、机构和材料种类分别衡量。

这一叙述依据《元史》卷041〈本纪第41卷本年条〉，但《元史》为明初仓促官修，本纪保存大量实录条目但有编次与纪年问题；行省执行、数字、族称及对外关系须以条格、多语文书和对方史料互校。。研究元代行省、赋役、河工与多语文书研究可以补问材料未直接说出的空间与社会差异。

\noindent\eventanchor{SYD-1351-YUAN-132}{寶慶路饑，判官文殊奴以所受敕牒貸官糧萬石賑之}\textbf{元。}
元的1351年（元至正11年；公元换算据《元史》本纪纪年）记录写到“寶慶路饑，判官文殊奴以所受敕牒貸官糧萬石賑之”。我们可以据此辨认一个可见的选择或压力，而不能把零散的官方记载、碑刻或后出叙事拼成无缝的社会全景。

账本把这一节点放在“本纪所记财用、赈济或征发属于特定报告场景。”与“可与食货志、文书和物质材料比较，不外推为全境总量。”之间；两者提示前后压力，并不把时间相邻写成单一因果。

本条首先回到《元史》卷041〈本纪第41卷本年条〉；《元史》为明初仓促官修，本纪保存大量实录条目但有编次与纪年问题；行省执行、数字、族称及对外关系须以条格、多语文书和对方史料互校。 因而元代行省、赋役、河工与多语文书研究不是装饰性的书目，而是检验这一叙述边界的路径。

\noindent\eventanchor{SYD-1351-YUAN-204}{丁未，立四川省檢校官}\textbf{元。}
在1351年（元至正11年；公元换算据《元史》本纪纪年），元所记“丁未，立四川省檢校官”把权力安排暴露在继承压力下。后续稳定与否取决于资源和支持网络如何被重新调度，不能由即位仪式本身预先断言。

此前的条件是“本纪年条记录政令或机构处置；实施背景须参照制度材料。”，此后的可见影响是“可在同年及后续政令中追踪；条文不单独证明地方执行。”。这里的“影响”须按地点、机构和材料种类分别衡量。

这一叙述依据《元史》卷041〈本纪第41卷本年条〉，但《元史》为明初仓促官修，本纪保存大量实录条目但有编次与纪年问题；行省执行、数字、族称及对外关系须以条格、多语文书和对方史料互校。。研究元代行省、赋役、河工与多语文书研究可以补问材料未直接说出的空间与社会差异。

\noindent\eventanchor{SYD-1351-YUAN-276}{冬十月乙未，增立巡防捕盜所於永昌}\textbf{元。}
在1351年（元至正11年；公元换算据《元史》本纪纪年），元所记“冬十月乙未，增立巡防捕盜所於永昌”把权力安排暴露在继承压力下。后续稳定与否取决于资源和支持网络如何被重新调度，不能由即位仪式本身预先断言。

从账本的前后项看，“本纪年条提供日期和中枢可见行动；前因不据单条补定。”提供了背景压力；“可回查同年及后续条目，地方影响仍须另证。”标出后来需要追踪的变化，二者之间仍存在未被材料填满的环节。

核心线索为《元史》卷041〈本纪第41卷本年条〉。《元史》为明初仓促官修，本纪保存大量实录条目但有编次与纪年问题；行省执行、数字、族称及对外关系须以条格、多语文书和对方史料互校。 现代研究可从元代行省、赋役、河工与多语文书研究继续追问。

\subsection{1351年}

\noindent\eventanchor{YUAN-1351-EVENT-192}{红巾军起事，元廷同时发动大规模河工}\textbf{元。}
在1351年的材料中，元面对的是红巾军起事，元廷同时发动大规模河工。这里应先看行动所依赖的关隘、城镇、河道或粮道，再讨论政治后果；军队抵达某处，并不等于它已能稳定征收、安置或控制周边。

账本把这一节点放在“元中后期继承、财政与地方危机中既有的联盟、交通与补给压力，为“红巾军起事，元廷同时发动大规模河工”提供了直接的军事或动员背景。”与““红巾军起事，元廷同时发动大规模河工”改变了相关城镇、道路或政权间的军事选择；伤亡、俘获和控制范围仅可按各方材料分别确认。”之间；两者提示前后压力，并不把时间相邻写成单一因果。

本条首先回到《元史》相关本纪；《元朝秘史》、波斯文史籍或《高丽史》对应叙事；《元典章》《通制条格》、多语碑刻与文书（据事核）；《元史》明初速修，纪年与人名有异同；蒙古大汗政权不等同所有汗国，征服、赋役和身份分类须按地区及材料语言分列。 因而蒙古帝国、元朝行省财政、多语文书与区域社会研究不是装饰性的书目，而是检验这一叙述边界的路径。

\subsection{1352年（元至正12年；公元换算据《元史》本纪纪年）}

\noindent\eventanchor{SYD-1352-YUAN-061}{庚寅，河決曹州，雇夫萬五千八百修築之}\textbf{元。}
“庚寅，河決曹州，雇夫萬五千八百修築之”是元在1352年（元至正12年；公元换算据《元史》本纪纪年）可追查的节点。它提示命令、环境、宗教、迁徙或地方组织如何进入政治过程；影响的范围和速度仍要避免由一条材料外推。

此前的条件是“本纪保留灾害或工程的报告地点与朝廷处置；灾情范围需另证。”，此后的可见影响是“可与地方文书、河工和环境材料比较，不将单点记录外推为全域因果。”。这里的“影响”须按地点、机构和材料种类分别衡量。

这一叙述依据《元史》卷041〈本纪第41卷本年条〉，但《元史》为明初仓促官修，本纪保存大量实录条目但有编次与纪年问题；行省执行、数字、族称及对外关系须以条格、多语文书和对方史料互校。。研究元代行省、赋役、河工与多语文书研究可以补问材料未直接说出的空间与社会差异。

\noindent\eventanchor{SYD-1352-YUAN-133}{丙午，命太平提調都水監}\textbf{元。}
“丙午，命太平提調都水監”是元在1352年（元至正12年；公元换算据《元史》本纪纪年）可追查的节点。它提示命令、环境、宗教、迁徙或地方组织如何进入政治过程；影响的范围和速度仍要避免由一条材料外推。

账本把这一节点放在“本纪年条记录政令或机构处置；实施背景须参照制度材料。”与“可在同年及后续政令中追踪；条文不单独证明地方执行。”之间；两者提示前后压力，并不把时间相邻写成单一因果。

本条首先回到《元史》卷041〈本纪第41卷本年条〉；《元史》为明初仓促官修，本纪保存大量实录条目但有编次与纪年问题；行省执行、数字、族称及对外关系须以条格、多语文书和对方史料互校。 因而元代行省、赋役、河工与多语文书研究不是装饰性的书目，而是检验这一叙述边界的路径。

\noindent\eventanchor{SYD-1352-YUAN-205}{丙戌，賜脫脫金十錠、銀五十錠、鈔萬錠、幣帛二百匹，辭不受}\textbf{元。}
“丙戌，賜脫脫金十錠、銀五十錠、鈔萬錠、幣帛二百匹，辭不受”是元在1352年（元至正12年；公元换算据《元史》本纪纪年）可追查的节点。它提示命令、环境、宗教、迁徙或地方组织如何进入政治过程；影响的范围和速度仍要避免由一条材料外推。

此前的条件是“本纪所记财用、赈济或征发属于特定报告场景。”，此后的可见影响是“可与食货志、文书和物质材料比较，不外推为全境总量。”。这里的“影响”须按地点、机构和材料种类分别衡量。

这一叙述依据《元史》卷041〈本纪第41卷本年条〉，但《元史》为明初仓促官修，本纪保存大量实录条目但有编次与纪年问题；行省执行、数字、族称及对外关系须以条格、多语文书和对方史料互校。。研究元代行省、赋役、河工与多语文书研究可以补问材料未直接说出的空间与社会差异。

\noindent\eventanchor{SYD-1352-YUAN-277}{丁卯，山東霖雨，民饑相食，賑之}\textbf{元。}
“丁卯，山東霖雨，民饑相食，賑之”是元在1352年（元至正12年；公元换算据《元史》本纪纪年）可追查的节点。它提示命令、环境、宗教、迁徙或地方组织如何进入政治过程；影响的范围和速度仍要避免由一条材料外推。

从账本的前后项看，“本纪所记财用、赈济或征发属于特定报告场景。”提供了背景压力；“可与食货志、文书和物质材料比较，不外推为全境总量。”标出后来需要追踪的变化，二者之间仍存在未被材料填满的环节。

核心线索为《元史》卷041〈本纪第41卷本年条〉。《元史》为明初仓促官修，本纪保存大量实录条目但有编次与纪年问题；行省执行、数字、族称及对外关系须以条格、多语文书和对方史料互校。 现代研究可从元代行省、赋役、河工与多语文书研究继续追问。

\subsection{1353年（元至正13年；公元换算据《元史》本纪纪年）}

\noindent\eventanchor{SYD-1353-YUAN-062}{:朕自踐祚以來，至今十有餘年，託身億兆之上，端居九重之中，耳目所及，豈能周知}\textbf{元。}
在1353年（元至正13年；公元换算据《元史》本纪纪年）的可见材料中，元留下“:朕自踐祚以來，至今十有餘年，託身億兆之上，端居九重之中，耳目所及，豈能周知”。这一动作把具体人群、机构或地方条件带进年序，但一项记录不能替未被记载者概括整体反应。

此前的条件是“本纪年条提供日期和中枢可见行动；前因不据单条补定。”，此后的可见影响是“可回查同年及后续条目，地方影响仍须另证。”。这里的“影响”须按地点、机构和材料种类分别衡量。

这一叙述依据《元史》卷041〈本纪第41卷本年条〉，但《元史》为明初仓促官修，本纪保存大量实录条目但有编次与纪年问题；行省执行、数字、族称及对外关系须以条格、多语文书和对方史料互校。。研究元代行省、赋役、河工与多语文书研究可以补问材料未直接说出的空间与社会差异。

\noindent\eventanchor{SYD-1353-YUAN-134}{己卯，監察御史不答失里請罷造作不急之務}\textbf{元。}
在1353年（元至正13年；公元换算据《元史》本纪纪年）的可见材料中，元留下“己卯，監察御史不答失里請罷造作不急之務”。这一动作把具体人群、机构或地方条件带进年序，但一项记录不能替未被记载者概括整体反应。

账本把这一节点放在“本纪年条记录政令或机构处置；实施背景须参照制度材料。”与“可在同年及后续政令中追踪；条文不单独证明地方执行。”之间；两者提示前后压力，并不把时间相邻写成单一因果。

本条首先回到《元史》卷041〈本纪第41卷本年条〉；《元史》为明初仓促官修，本纪保存大量实录条目但有编次与纪年问题；行省执行、数字、族称及对外关系须以条格、多语文书和对方史料互校。 因而元代行省、赋役、河工与多语文书研究不是装饰性的书目，而是检验这一叙述边界的路径。

\noindent\eventanchor{SYD-1353-YUAN-206}{丙午，命也先帖木兒、鐵木兒塔識並為御史大夫}\textbf{元。}
在1353年（元至正13年；公元换算据《元史》本纪纪年）的可见材料中，元留下“丙午，命也先帖木兒、鐵木兒塔識並為御史大夫”。这一动作把具体人群、机构或地方条件带进年序，但一项记录不能替未被记载者概括整体反应。

此前的条件是“本纪年条记录政令或机构处置；实施背景须参照制度材料。”，此后的可见影响是“可在同年及后续政令中追踪；条文不单独证明地方执行。”。这里的“影响”须按地点、机构和材料种类分别衡量。

这一叙述依据《元史》卷041〈本纪第41卷本年条〉，但《元史》为明初仓促官修，本纪保存大量实录条目但有编次与纪年问题；行省执行、数字、族称及对外关系须以条格、多语文书和对方史料互校。。研究元代行省、赋役、河工与多语文书研究可以补问材料未直接说出的空间与社会差异。

\noindent\eventanchor{SYD-1353-YUAN-278}{大都、永平、鞏昌、興國、安陸等處并桃溫萬戶府各翼人民饑，賑之}\textbf{元。}
在1353年（元至正13年；公元换算据《元史》本纪纪年）的可见材料中，元留下“大都、永平、鞏昌、興國、安陸等處并桃溫萬戶府各翼人民饑，賑之”。这一动作把具体人群、机构或地方条件带进年序，但一项记录不能替未被记载者概括整体反应。

从账本的前后项看，“本纪所记财用、赈济或征发属于特定报告场景。”提供了背景压力；“可与食货志、文书和物质材料比较，不外推为全境总量。”标出后来需要追踪的变化，二者之间仍存在未被材料填满的环节。

核心线索为《元史》卷041〈本纪第41卷本年条〉。《元史》为明初仓促官修，本纪保存大量实录条目但有编次与纪年问题；行省执行、数字、族称及对外关系须以条格、多语文书和对方史料互校。 现代研究可从元代行省、赋役、河工与多语文书研究继续追问。

\subsection{1354年（元至正14年；公元换算据《元史》本纪纪年）}

\noindent\eventanchor{SYD-1354-YUAN-063}{復立八百宣慰司，以土官韓部襲其父爵}\textbf{元。}
“復立八百宣慰司，以土官韓部襲其父爵”在1354年（元至正14年；公元换算据《元史》本纪纪年）构成元的继承或权力重组节点。名位的变化可以由纪年串连，但它没有自动消除宫廷、军政和地方网络原有的拉扯。

此前的条件是“本纪年条记录政令或机构处置；实施背景须参照制度材料。”，此后的可见影响是“可在同年及后续政令中追踪；条文不单独证明地方执行。”。这里的“影响”须按地点、机构和材料种类分别衡量。

这一叙述依据《元史》卷041〈本纪第41卷本年条〉，但《元史》为明初仓促官修，本纪保存大量实录条目但有编次与纪年问题；行省执行、数字、族称及对外关系须以条格、多语文书和对方史料互校。。研究元代行省、赋役、河工与多语文书研究可以补问材料未直接说出的空间与社会差异。

\noindent\eventanchor{SYD-1354-YUAN-135}{八月丙午〔朔〕，命江浙行省右丞忽都不花、江西行省右丞禿魯統軍合討羅天麟}\textbf{元。}
在1354年（元至正14年；公元换算据《元史》本纪纪年）的可见材料中，元留下“八月丙午〔朔〕，命江浙行省右丞忽都不花、江西行省右丞禿魯統軍合討羅天麟”。这一动作把具体人群、机构或地方条件带进年序，但一项记录不能替未被记载者概括整体反应。

账本把这一节点放在“本纪年条记录政令或机构处置；实施背景须参照制度材料。”与“可在同年及后续政令中追踪；条文不单独证明地方执行。”之间；两者提示前后压力，并不把时间相邻写成单一因果。

本条首先回到《元史》卷041〈本纪第41卷本年条〉；《元史》为明初仓促官修，本纪保存大量实录条目但有编次与纪年问题；行省执行、数字、族称及对外关系须以条格、多语文书和对方史料互校。 因而元代行省、赋役、河工与多语文书研究不是装饰性的书目，而是检验这一叙述边界的路径。

\noindent\eventanchor{SYD-1354-YUAN-207}{丙申，以朵兒直班為中書右丞，答兒麻為參知政事}\textbf{元。}
在1354年（元至正14年；公元换算据《元史》本纪纪年）的可见材料中，元留下“丙申，以朵兒直班為中書右丞，答兒麻為參知政事”。这一动作把具体人群、机构或地方条件带进年序，但一项记录不能替未被记载者概括整体反应。

此前的条件是“本纪年条记录政令或机构处置；实施背景须参照制度材料。”，此后的可见影响是“可在同年及后续政令中追踪；条文不单独证明地方执行。”。这里的“影响”须按地点、机构和材料种类分别衡量。

这一叙述依据《元史》卷041〈本纪第41卷本年条〉，但《元史》为明初仓促官修，本纪保存大量实录条目但有编次与纪年问题；行省执行、数字、族称及对外关系须以条格、多语文书和对方史料互校。。研究元代行省、赋役、河工与多语文书研究可以补问材料未直接说出的空间与社会差异。

\noindent\eventanchor{SYD-1354-YUAN-279}{丙戌，以遼陽吾者野人等未靖，命太保伯撒里為遼陽行省左丞相鎮之}\textbf{元。}
在1354年（元至正14年；公元换算据《元史》本纪纪年）的可见材料中，元留下“丙戌，以遼陽吾者野人等未靖，命太保伯撒里為遼陽行省左丞相鎮之”。这一动作把具体人群、机构或地方条件带进年序，但一项记录不能替未被记载者概括整体反应。

从账本的前后项看，“本纪年条记录政令或机构处置；实施背景须参照制度材料。”提供了背景压力；“可在同年及后续政令中追踪；条文不单独证明地方执行。”标出后来需要追踪的变化，二者之间仍存在未被材料填满的环节。

核心线索为《元史》卷041〈本纪第41卷本年条〉。《元史》为明初仓促官修，本纪保存大量实录条目但有编次与纪年问题；行省执行、数字、族称及对外关系须以条格、多语文书和对方史料互校。 现代研究可从元代行省、赋役、河工与多语文书研究继续追问。

\subsection{1354年}

\noindent\eventanchor{YUAN-1354-EVENT-193}{脱脱在镇压起事期间失势被罢}\textbf{元。}
在1354年的材料中，元面对的是脱脱在镇压起事期间失势被罢。这里应先看行动所依赖的关隘、城镇、河道或粮道，再讨论政治后果；军队抵达某处，并不等于它已能稳定征收、安置或控制周边。

账本把这一节点放在“元中后期继承、财政与地方危机中既有的联盟、交通与补给压力，为“脱脱在镇压起事期间失势被罢”提供了直接的军事或动员背景。”与““脱脱在镇压起事期间失势被罢”改变了相关城镇、道路或政权间的军事选择；伤亡、俘获和控制范围仅可按各方材料分别确认。”之间；两者提示前后压力，并不把时间相邻写成单一因果。

本条首先回到《元史》相关本纪；《元朝秘史》、波斯文史籍或《高丽史》对应叙事；《元典章》《通制条格》、多语碑刻与文书（据事核）；《元史》明初速修，纪年与人名有异同；蒙古大汗政权不等同所有汗国，征服、赋役和身份分类须按地区及材料语言分列。 因而蒙古帝国、元朝行省财政、多语文书与区域社会研究不是装饰性的书目，而是检验这一叙述边界的路径。

\subsection{1355年（元至正15年；公元换算据《元史》本纪纪年）}

\noindent\eventanchor{SYD-1355-YUAN-064}{帝牽於眾請，令三年後減之}\textbf{元。}
元的1355年（元至正15年；公元换算据《元史》本纪纪年）记录写到“帝牽於眾請，令三年後減之”。我们可以据此辨认一个可见的选择或压力，而不能把零散的官方记载、碑刻或后出叙事拼成无缝的社会全景。

此前的条件是“本纪年条提供日期和中枢可见行动；前因不据单条补定。”，此后的可见影响是“可回查同年及后续条目，地方影响仍须另证。”。这里的“影响”须按地点、机构和材料种类分别衡量。

这一叙述依据《元史》卷041〈本纪第41卷本年条〉，但《元史》为明初仓促官修，本纪保存大量实录条目但有编次与纪年问题；行省执行、数字、族称及对外关系须以条格、多语文书和对方史料互校。。研究元代行省、赋役、河工与多语文书研究可以补问材料未直接说出的空间与社会差异。

\noindent\eventanchor{SYD-1355-YUAN-136}{庚辰，詔建木華黎、伯顏祠堂於東平}\textbf{元。}
元的1355年（元至正15年；公元换算据《元史》本纪纪年）记录写到“庚辰，詔建木華黎、伯顏祠堂於東平”。我们可以据此辨认一个可见的选择或压力，而不能把零散的官方记载、碑刻或后出叙事拼成无缝的社会全景。

账本把这一节点放在“本纪年条记录政令或机构处置；实施背景须参照制度材料。”与“可在同年及后续政令中追踪；条文不单独证明地方执行。”之间；两者提示前后压力，并不把时间相邻写成单一因果。

本条首先回到《元史》卷041〈本纪第41卷本年条〉；《元史》为明初仓促官修，本纪保存大量实录条目但有编次与纪年问题；行省执行、数字、族称及对外关系须以条格、多语文书和对方史料互校。 因而元代行省、赋役、河工与多语文书研究不是装饰性的书目，而是检验这一叙述边界的路径。

\noindent\eventanchor{SYD-1355-YUAN-208}{丙戌，亦憐只答兒反，遣兵討之}\textbf{元。}
在1355年（元至正15年；公元换算据《元史》本纪纪年）的材料中，元面对的是丙戌，亦憐只答兒反，遣兵討之。这里应先看行动所依赖的关隘、城镇、河道或粮道，再讨论政治后果；军队抵达某处，并不等于它已能稳定征收、安置或控制周边。

此前的条件是“本纪年条提供日期和中枢可见行动；前因不据单条补定。”，此后的可见影响是“可回查同年及后续条目，地方影响仍须另证。”。这里的“影响”须按地点、机构和材料种类分别衡量。

这一叙述依据《元史》卷041〈本纪第41卷本年条〉，但《元史》为明初仓促官修，本纪保存大量实录条目但有编次与纪年问题；行省执行、数字、族称及对外关系须以条格、多语文书和对方史料互校。。研究元代行省、赋役、河工与多语文书研究可以补问材料未直接说出的空间与社会差异。

\noindent\eventanchor{SYD-1355-YUAN-280}{丙戌，中書省臣建議，以河南盜賊出入無常，宜分撥達達軍與揚州舊軍於河南水陸關隘戍守，東至徐、邳，北至夾馬營，遇賊掩捕，從之}\textbf{元。}
在1355年（元至正15年；公元换算据《元史》本纪纪年）的材料中，元面对的是丙戌，中書省臣建議，以河南盜賊出入無常，宜分撥達達軍與揚州舊軍於河南水陸關隘戍守，東至徐、邳，北至夾馬營，遇賊掩捕，從之。这里应先看行动所依赖的关隘、城镇、河道或粮道，再讨论政治后果；军队抵达某处，并不等于它已能稳定征收、安置或控制周边。

此前的条件是“本纪年条记录政令或机构处置；实施背景须参照制度材料。”，此后的可见影响是“可在同年及后续政令中追踪；条文不单独证明地方执行。”。这里的“影响”须按地点、机构和材料种类分别衡量。

这一叙述依据《元史》卷041〈本纪第41卷本年条〉，但《元史》为明初仓促官修，本纪保存大量实录条目但有编次与纪年问题；行省执行、数字、族称及对外关系须以条格、多语文书和对方史料互校。。研究元代行省、赋役、河工与多语文书研究可以补问材料未直接说出的空间与社会差异。

\subsection{1356年（元至正16年；公元换算据《元史》本纪纪年）}

\noindent\eventanchor{SYD-1356-YUAN-065}{平江、松江水災，給海運糧十萬石賑之}\textbf{元。}
“平江、松江水災，給海運糧十萬石賑之”是元在1356年（元至正16年；公元换算据《元史》本纪纪年）可追查的节点。它提示命令、环境、宗教、迁徙或地方组织如何进入政治过程；影响的范围和速度仍要避免由一条材料外推。

此前的条件是“本纪所记财用、赈济或征发属于特定报告场景。”，此后的可见影响是“可与食货志、文书和物质材料比较，不外推为全境总量。”。这里的“影响”须按地点、机构和材料种类分别衡量。

这一叙述依据《元史》卷041〈本纪第41卷本年条〉，但《元史》为明初仓促官修，本纪保存大量实录条目但有编次与纪年问题；行省执行、数字、族称及对外关系须以条格、多语文书和对方史料互校。。研究元代行省、赋役、河工与多语文书研究可以补问材料未直接说出的空间与社会差异。

\noindent\eventanchor{SYD-1356-YUAN-137}{詔：「京官三品以上，歲舉守令一人，守令到任三月，亦舉一人自代}\textbf{元。}
“詔：「京官三品以上，歲舉守令一人，守令到任三月，亦舉一人自代”把元置于军事行动的可核时点。战报能够记下攻守、入城或撤离，却很少完整保留沿途征发、居民去留和补给损耗；因此不能把一方的胜败叙事延伸成所有地方的经验。

账本把这一节点放在“本纪年条记录政令或机构处置；实施背景须参照制度材料。”与“可在同年及后续政令中追踪；条文不单独证明地方执行。”之间；两者提示前后压力，并不把时间相邻写成单一因果。

本条首先回到《元史》卷041〈本纪第41卷本年条〉；《元史》为明初仓促官修，本纪保存大量实录条目但有编次与纪年问题；行省执行、数字、族称及对外关系须以条格、多语文书和对方史料互校。 因而元代行省、赋役、河工与多语文书研究不是装饰性的书目，而是检验这一叙述边界的路径。

\noindent\eventanchor{SYD-1356-YUAN-209}{詔翰林國史院纂修后妃、功臣列傳，學士承旨張起巖、學士楊宗瑞、侍講學士黃溍為總裁官，左丞相太平、左丞呂思誠領其事}\textbf{元。}
“詔翰林國史院纂修后妃、功臣列傳，學士承旨張起巖、學士楊宗瑞、侍講學士黃溍為總裁官，左丞相太平、左丞呂思誠領其事”是元在1356年（元至正16年；公元换算据《元史》本纪纪年）可追查的节点。它提示命令、环境、宗教、迁徙或地方组织如何进入政治过程；影响的范围和速度仍要避免由一条材料外推。

此前的条件是“本纪年条记录政令或机构处置；实施背景须参照制度材料。”，此后的可见影响是“可在同年及后续政令中追踪；条文不单独证明地方执行。”。这里的“影响”须按地点、机构和材料种类分别衡量。

这一叙述依据《元史》卷041〈本纪第41卷本年条〉，但《元史》为明初仓促官修，本纪保存大量实录条目但有编次与纪年问题；行省执行、数字、族称及对外关系须以条格、多语文书和对方史料互校。。研究元代行省、赋役、河工与多语文书研究可以补问材料未直接说出的空间与社会差异。

\noindent\eventanchor{SYD-1356-YUAN-281}{八年春正月戊戌朔，命也先帖木兒知樞密院事}\textbf{元。}
“八年春正月戊戌朔，命也先帖木兒知樞密院事”是元在1356年（元至正16年；公元换算据《元史》本纪纪年）可追查的节点。它提示命令、环境、宗教、迁徙或地方组织如何进入政治过程；影响的范围和速度仍要避免由一条材料外推。

此前的条件是“本纪年条记录政令或机构处置；实施背景须参照制度材料。”，此后的可见影响是“可在同年及后续政令中追踪；条文不单独证明地方执行。”。这里的“影响”须按地点、机构和材料种类分别衡量。

这一叙述依据《元史》卷041〈本纪第41卷本年条〉，但《元史》为明初仓促官修，本纪保存大量实录条目但有编次与纪年问题；行省执行、数字、族称及对外关系须以条格、多语文书和对方史料互校。。研究元代行省、赋役、河工与多语文书研究可以补问材料未直接说出的空间与社会差异。

\subsection{1356年}

\noindent\eventanchor{YUAN-1356-EVENT-194}{朱元璋攻取集庆并改名应天}\textbf{元与江南起事者。}
“朱元璋攻取集庆并改名应天”把元与江南起事者置于军事行动的可核时点。战报能够记下攻守、入城或撤离，却很少完整保留沿途征发、居民去留和补给损耗；因此不能把一方的胜败叙事延伸成所有地方的经验。

其前提记为“元中后期继承、财政与地方危机中既有的联盟、交通与补给压力，为“朱元璋攻取集庆并改名应天”提供了直接的军事或动员背景。”；其后续是““朱元璋攻取集庆并改名应天”改变了相关城镇、道路或政权间的军事选择；伤亡、俘获和控制范围仅可按各方材料分别确认。”。这是一条待后续材料检验的事件链，而不是对所有行动者意图的替写。

可核的材料入口是《元史》相关本纪；《元朝秘史》、波斯文史籍或《高丽史》对应叙事；《元典章》《通制条格》、多语碑刻与文书（据事核）。《元史》明初速修，纪年与人名有异同；蒙古大汗政权不等同所有汗国，征服、赋役和身份分类须按地区及材料语言分列。 就此而言，蒙古帝国、元朝行省财政、多语文书与区域社会研究有助于把年序同区域、制度或物质材料并读。

\subsection{1361年（元至正21年；公元换算据《元史》本纪纪年）}

\noindent\eventanchor{SYD-1361-YUAN-066}{〔徭〕賊吳天保復寇沅州}\textbf{元。}
元的1361年（元至正21年；公元换算据《元史》本纪纪年）记录写到“〔徭〕賊吳天保復寇沅州”。我们可以据此辨认一个可见的选择或压力，而不能把零散的官方记载、碑刻或后出叙事拼成无缝的社会全景。

从账本的前后项看，“本纪年条记录政令或机构处置；实施背景须参照制度材料。”提供了背景压力；“可在同年及后续政令中追踪；条文不单独证明地方执行。”标出后来需要追踪的变化，二者之间仍存在未被材料填满的环节。

核心线索为《元史》卷042〈本纪第42卷本年条〉。《元史》为明初仓促官修，本纪保存大量实录条目但有编次与纪年问题；行省执行、数字、族称及对外关系须以条格、多语文书和对方史料互校。 现代研究可从元代行省、赋役、河工与多语文书研究继续追问。

\noindent\eventanchor{SYD-1361-YUAN-138}{八月甲辰，以集賢大學士柏顏為中書平章政事，河南行省平章政事月魯不花為宣政院使}\textbf{元。}
元于1361年（元至正21年；公元换算据《元史》本纪纪年）处理“八月甲辰，以集賢大學士柏顏為中書平章政事，河南行省平章政事月魯不花為宣政院使”时，面对的是道路、翻译、礼仪和资源交换共同构成的关系网。书面称谓可见，地方执行和持续性则仍须由后续材料检验。

其前提记为“本纪年条记录政令或机构处置；实施背景须参照制度材料。”；其后续是“可在同年及后续政令中追踪；条文不单独证明地方执行。”。这是一条待后续材料检验的事件链，而不是对所有行动者意图的替写。

可核的材料入口是《元史》卷042〈本纪第42卷本年条〉。《元史》为明初仓促官修，本纪保存大量实录条目但有编次与纪年问题；行省执行、数字、族称及对外关系须以条格、多语文书和对方史料互校。 就此而言，元代行省、赋役、河工与多语文书研究有助于把年序同区域、制度或物质材料并读。

\noindent\eventanchor{SYD-1361-YUAN-210}{丙午，命中書平章政事太不花提調會同館}\textbf{元。}
元的1361年（元至正21年；公元换算据《元史》本纪纪年）记录写到“丙午，命中書平章政事太不花提調會同館”。我们可以据此辨认一个可见的选择或压力，而不能把零散的官方记载、碑刻或后出叙事拼成无缝的社会全景。

此前的条件是“本纪年条记录政令或机构处置；实施背景须参照制度材料。”，此后的可见影响是“可在同年及后续政令中追踪；条文不单独证明地方执行。”。这里的“影响”须按地点、机构和材料种类分别衡量。

这一叙述依据《元史》卷042〈本纪第42卷本年条〉，但《元史》为明初仓促官修，本纪保存大量实录条目但有编次与纪年问题；行省执行、数字、族称及对外关系须以条格、多语文书和对方史料互校。。研究元代行省、赋役、河工与多语文书研究可以补问材料未直接说出的空间与社会差异。

\noindent\eventanchor{SYD-1361-YUAN-282}{丙寅，命平章政事柏顏提調留守司}\textbf{元。}
在1361年（元至正21年；公元换算据《元史》本纪纪年）的材料中，元面对的是丙寅，命平章政事柏顏提調留守司。这里应先看行动所依赖的关隘、城镇、河道或粮道，再讨论政治后果；军队抵达某处，并不等于它已能稳定征收、安置或控制周边。

从账本的前后项看，“本纪年条记录政令或机构处置；实施背景须参照制度材料。”提供了背景压力；“可在同年及后续政令中追踪；条文不单独证明地方执行。”标出后来需要追踪的变化，二者之间仍存在未被材料填满的环节。

核心线索为《元史》卷042〈本纪第42卷本年条〉。《元史》为明初仓促官修，本纪保存大量实录条目但有编次与纪年问题；行省执行、数字、族称及对外关系须以条格、多语文书和对方史料互校。 现代研究可从元代行省、赋役、河工与多语文书研究继续追问。

\subsection{1362年（元至正22年；公元换算据《元史》本纪纪年）}

\noindent\eventanchor{SYD-1362-YUAN-067}{二月丙戌〔朔〕，詔加封天妃父種德積慶侯，母育聖顯慶夫人}\textbf{元。}
“二月丙戌〔朔〕，詔加封天妃父種德積慶侯，母育聖顯慶夫人”是元在1362年（元至正22年；公元换算据《元史》本纪纪年）可追查的节点。它提示命令、环境、宗教、迁徙或地方组织如何进入政治过程；影响的范围和速度仍要避免由一条材料外推。

从账本的前后项看，“本纪年条记录政令或机构处置；实施背景须参照制度材料。”提供了背景压力；“可在同年及后续政令中追踪；条文不单独证明地方执行。”标出后来需要追踪的变化，二者之间仍存在未被材料填满的环节。

核心线索为《元史》卷042〈本纪第42卷本年条〉。《元史》为明初仓促官修，本纪保存大量实录条目但有编次与纪年问题；行省执行、数字、族称及对外关系须以条格、多语文书和对方史料互校。 现代研究可从元代行省、赋役、河工与多语文书研究继续追问。

\noindent\eventanchor{SYD-1362-YUAN-139}{庚午，命樞密院以軍士五百修築白河堤}\textbf{元。}
“庚午，命樞密院以軍士五百修築白河堤”是元在1362年（元至正22年；公元换算据《元史》本纪纪年）可追查的节点。它提示命令、环境、宗教、迁徙或地方组织如何进入政治过程；影响的范围和速度仍要避免由一条材料外推。

其前提记为“本纪年条记录政令或机构处置；实施背景须参照制度材料。”；其后续是“可在同年及后续政令中追踪；条文不单独证明地方执行。”。这是一条待后续材料检验的事件链，而不是对所有行动者意图的替写。

可核的材料入口是《元史》卷042〈本纪第42卷本年条〉。《元史》为明初仓促官修，本纪保存大量实录条目但有编次与纪年问题；行省执行、数字、族称及对外关系须以条格、多语文书和对方史料互校。 就此而言，元代行省、赋役、河工与多语文书研究有助于把年序同区域、制度或物质材料并读。

\noindent\eventanchor{SYD-1362-YUAN-211}{己巳，詔天下以中統交鈔壹貫文權銅錢壹千文，準至元寶鈔貳貫，仍鑄至正通寶錢並用，以實鈔法，至元寶鈔通行如故}\textbf{元。}
“己巳，詔天下以中統交鈔壹貫文權銅錢壹千文，準至元寶鈔貳貫，仍鑄至正通寶錢並用，以實鈔法，至元寶鈔通行如故”是元在1362年（元至正22年；公元换算据《元史》本纪纪年）可追查的节点。它提示命令、环境、宗教、迁徙或地方组织如何进入政治过程；影响的范围和速度仍要避免由一条材料外推。

此前的条件是“本纪年条记录政令或机构处置；实施背景须参照制度材料。”，此后的可见影响是“可在同年及后续政令中追踪；条文不单独证明地方执行。”。这里的“影响”须按地点、机构和材料种类分别衡量。

这一叙述依据《元史》卷042〈本纪第42卷本年条〉，但《元史》为明初仓促官修，本纪保存大量实录条目但有编次与纪年问题；行省执行、数字、族称及对外关系须以条格、多语文书和对方史料互校。。研究元代行省、赋役、河工与多语文书研究可以补问材料未直接说出的空间与社会差异。

\noindent\eventanchor{SYD-1362-YUAN-283}{六月壬子，星大如月，入北斗，震聲若雷，三日復還}\textbf{元。}
“六月壬子，星大如月，入北斗，震聲若雷，三日復還”是元在1362年（元至正22年；公元换算据《元史》本纪纪年）可追查的节点。它提示命令、环境、宗教、迁徙或地方组织如何进入政治过程；影响的范围和速度仍要避免由一条材料外推。

从账本的前后项看，“本纪年条记录政令或机构处置；实施背景须参照制度材料。”提供了背景压力；“可在同年及后续政令中追踪；条文不单独证明地方执行。”标出后来需要追踪的变化，二者之间仍存在未被材料填满的环节。

核心线索为《元史》卷042〈本纪第42卷本年条〉。《元史》为明初仓促官修，本纪保存大量实录条目但有编次与纪年问题；行省执行、数字、族称及对外关系须以条格、多语文书和对方史料互校。 现代研究可从元代行省、赋役、河工与多语文书研究继续追问。

\subsection{1363年（元至正23年；公元换算据《元史》本纪纪年）}

\noindent\eventanchor{SYD-1363-YUAN-068}{八月丁丑朔，中興地震}\textbf{元。}
在1363年（元至正23年；公元换算据《元史》本纪纪年）的可见材料中，元留下“八月丁丑朔，中興地震”。这一动作把具体人群、机构或地方条件带进年序，但一项记录不能替未被记载者概括整体反应。

从账本的前后项看，“本纪保留灾害或工程的报告地点与朝廷处置；灾情范围需另证。”提供了背景压力；“可与地方文书、河工和环境材料比较，不将单点记录外推为全域因果。”标出后来需要追踪的变化，二者之间仍存在未被材料填满的环节。

核心线索为《元史》卷042〈本纪第42卷本年条〉。《元史》为明初仓促官修，本纪保存大量实录条目但有编次与纪年问题；行省执行、数字、族称及对外关系须以条格、多语文书和对方史料互校。 现代研究可从元代行省、赋役、河工与多语文书研究继续追问。

\noindent\eventanchor{SYD-1363-YUAN-140}{丙辰，親策進士八十三人，賜朵烈圖、文允中進士及第，其餘賜出身有差}\textbf{元。}
在1363年（元至正23年；公元换算据《元史》本纪纪年）的可见材料中，元留下“丙辰，親策進士八十三人，賜朵烈圖、文允中進士及第，其餘賜出身有差”。这一动作把具体人群、机构或地方条件带进年序，但一项记录不能替未被记载者概括整体反应。

账本把这一节点放在“本纪年条提供日期和中枢可见行动；前因不据单条补定。”与“可回查同年及后续条目，地方影响仍须另证。”之间；两者提示前后压力，并不把时间相邻写成单一因果。

本条首先回到《元史》卷042〈本纪第42卷本年条〉；《元史》为明初仓促官修，本纪保存大量实录条目但有编次与纪年问题；行省执行、数字、族称及对外关系须以条格、多语文书和对方史料互校。 因而元代行省、赋役、河工与多语文书研究不是装饰性的书目，而是检验这一叙述边界的路径。

\noindent\eventanchor{SYD-1363-YUAN-212}{丙戌，蕭縣李二及老彭、趙君用攻陷徐州}\textbf{元。}
1363年（元至正23年；公元换算据《元史》本纪纪年），元的行动落在“丙戌，蕭縣李二及老彭、趙君用攻陷徐州”所指的战事与通道上。可见的不是抽象的推进：攻守双方必须让人员、消息与补给穿过具体的城防、道路或水路；账本所列的结果只说明这一节点改变了下一步选择。

此前的条件是“本纪从朝廷视角记录军政行动；出兵与战果需分辨。”，此后的可见影响是“后续路线、控制与损失应与对方记录和遗址材料复核。”。这里的“影响”须按地点、机构和材料种类分别衡量。

这一叙述依据《元史》卷042〈本纪第42卷本年条〉，但《元史》为明初仓促官修，本纪保存大量实录条目但有编次与纪年问题；行省执行、数字、族称及对外关系须以条格、多语文书和对方史料互校。。研究元代行省、赋役、河工与多语文书研究可以补问材料未直接说出的空间与社会差异。

\noindent\eventanchor{SYD-1363-YUAN-284}{丙子，命大都至汴梁二十四驛，凡馬一匹助給鈔五錠}\textbf{元。}
在1363年（元至正23年；公元换算据《元史》本纪纪年）的可见材料中，元留下“丙子，命大都至汴梁二十四驛，凡馬一匹助給鈔五錠”。这一动作把具体人群、机构或地方条件带进年序，但一项记录不能替未被记载者概括整体反应。

从账本的前后项看，“本纪年条记录政令或机构处置；实施背景须参照制度材料。”提供了背景压力；“可在同年及后续政令中追踪；条文不单独证明地方执行。”标出后来需要追踪的变化，二者之间仍存在未被材料填满的环节。

核心线索为《元史》卷042〈本纪第42卷本年条〉。《元史》为明初仓促官修，本纪保存大量实录条目但有编次与纪年问题；行省执行、数字、族称及对外关系须以条格、多语文书和对方史料互校。 现代研究可从元代行省、赋役、河工与多语文书研究继续追问。

\subsection{1363年}

\noindent\eventanchor{YUAN-1363-EVENT-195}{鄱阳湖之战改变江南群雄力量对比}\textbf{元末群雄。}
“鄱阳湖之战改变江南群雄力量对比”把元末群雄置于军事行动的可核时点。战报能够记下攻守、入城或撤离，却很少完整保留沿途征发、居民去留和补给损耗；因此不能把一方的胜败叙事延伸成所有地方的经验。

账本把这一节点放在“元中后期继承、财政与地方危机中既有的联盟、交通与补给压力，为“鄱阳湖之战改变江南群雄力量对比”提供了直接的军事或动员背景。”与““鄱阳湖之战改变江南群雄力量对比”改变了相关城镇、道路或政权间的军事选择；伤亡、俘获和控制范围仅可按各方材料分别确认。”之间；两者提示前后压力，并不把时间相邻写成单一因果。

本条首先回到《元史》相关本纪；《元朝秘史》、波斯文史籍或《高丽史》对应叙事；《元典章》《通制条格》、多语碑刻与文书（据事核）；《元史》明初速修，纪年与人名有异同；蒙古大汗政权不等同所有汗国，征服、赋役和身份分类须按地区及材料语言分列。 因而蒙古帝国、元朝行省财政、多语文书与区域社会研究不是装饰性的书目，而是检验这一叙述边界的路径。

\subsection{1364年（元至正24年；公元换算据《元史》本纪纪年）}

\noindent\eventanchor{SYD-1364-YUAN-069}{是月，方國珍復劫其黨下海，入黃巖港，台州路達魯花赤泰不花率官軍與戰，死之}\textbf{元。}
元的1364年（元至正24年；公元换算据《元史》本纪纪年）记录写到“是月，方國珍復劫其黨下海，入黃巖港，台州路達魯花赤泰不花率官軍與戰，死之”。我们可以据此辨认一个可见的选择或压力，而不能把零散的官方记载、碑刻或后出叙事拼成无缝的社会全景。

从账本的前后项看，“本纪年条记录政令或机构处置；实施背景须参照制度材料。”提供了背景压力；“可在同年及后续政令中追踪；条文不单独证明地方执行。”标出后来需要追踪的变化，二者之间仍存在未被材料填满的环节。

核心线索为《元史》卷042〈本纪第42卷本年条〉。《元史》为明初仓促官修，本纪保存大量实录条目但有编次与纪年问题；行省执行、数字、族称及对外关系须以条格、多语文书和对方史料互校。 现代研究可从元代行省、赋役、河工与多语文书研究继续追问。

\noindent\eventanchor{SYD-1364-YUAN-141}{置安東、安豐分元帥府}\textbf{元。}
元的1364年（元至正24年；公元换算据《元史》本纪纪年）记录写到“置安東、安豐分元帥府”。我们可以据此辨认一个可见的选择或压力，而不能把零散的官方记载、碑刻或后出叙事拼成无缝的社会全景。

账本把这一节点放在“本纪年条记录政令或机构处置；实施背景须参照制度材料。”与“可在同年及后续政令中追踪；条文不单独证明地方执行。”之间；两者提示前后压力，并不把时间相邻写成单一因果。

本条首先回到《元史》卷042〈本纪第42卷本年条〉；《元史》为明初仓促官修，本纪保存大量实录条目但有编次与纪年问题；行省执行、数字、族称及对外关系须以条格、多语文书和对方史料互校。 因而元代行省、赋役、河工与多语文书研究不是装饰性的书目，而是检验这一叙述边界的路径。

\noindent\eventanchor{SYD-1364-YUAN-213}{中書省臣言：「張理獻言，饒州德興三處，膽水浸鐵，可以成銅，宜即其地各立銅冶場，直隸寶泉提舉司，宜以張理就為銅冶場官}\textbf{元。}
在1364年（元至正24年；公元换算据《元史》本纪纪年），元所记“中書省臣言：「張理獻言，饒州德興三處，膽水浸鐵，可以成銅，宜即其地各立銅冶場，直隸寶泉提舉司，宜以張理就為銅冶場官”把权力安排暴露在继承压力下。后续稳定与否取决于资源和支持网络如何被重新调度，不能由即位仪式本身预先断言。

此前的条件是“本纪年条记录政令或机构处置；实施背景须参照制度材料。”，此后的可见影响是“可在同年及后续政令中追踪；条文不单独证明地方执行。”。这里的“影响”须按地点、机构和材料种类分别衡量。

这一叙述依据《元史》卷042〈本纪第42卷本年条〉，但《元史》为明初仓促官修，本纪保存大量实录条目但有编次与纪年问题；行省执行、数字、族称及对外关系须以条格、多语文书和对方史料互校。。研究元代行省、赋役、河工与多语文书研究可以补问材料未直接说出的空间与社会差异。

\noindent\eventanchor{SYD-1364-YUAN-285}{安陸賊將俞君正復陷荊門州，知州聶炳死之}\textbf{元。}
在1364年（元至正24年；公元换算据《元史》本纪纪年）的材料中，元面对的是安陸賊將俞君正復陷荊門州，知州聶炳死之。这里应先看行动所依赖的关隘、城镇、河道或粮道，再讨论政治后果；军队抵达某处，并不等于它已能稳定征收、安置或控制周边。

从账本的前后项看，“本纪年条记录政令或机构处置；实施背景须参照制度材料。”提供了背景压力；“可在同年及后续政令中追踪；条文不单独证明地方执行。”标出后来需要追踪的变化，二者之间仍存在未被材料填满的环节。

核心线索为《元史》卷042〈本纪第42卷本年条〉。《元史》为明初仓促官修，本纪保存大量实录条目但有编次与纪年问题；行省执行、数字、族称及对外关系须以条格、多语文书和对方史料互校。 现代研究可从元代行省、赋役、河工与多语文书研究继续追问。

\subsection{1366年（元至正26年；公元换算据《元史》本纪纪年）}

\noindent\eventanchor{SYD-1366-YUAN-070}{皇后肅良合氏生日，百官進箋，皇后諭沙藍答里等曰：「自世祖以來，正宮皇后壽日，不曾進箋，近年雖行，不合典故}\textbf{元。}
元的1366年（元至正26年；公元换算据《元史》本纪纪年）记录写到“皇后肅良合氏生日，百官進箋，皇后諭沙藍答里等曰：「自世祖以來，正宮皇后壽日，不曾進箋，近年雖行，不合典故”。我们可以据此辨认一个可见的选择或压力，而不能把零散的官方记载、碑刻或后出叙事拼成无缝的社会全景。

从账本的前后项看，“本纪年条提供日期和中枢可见行动；前因不据单条补定。”提供了背景压力；“可回查同年及后续条目，地方影响仍须另证。”标出后来需要追踪的变化，二者之间仍存在未被材料填满的环节。

核心线索为《元史》卷047〈至正二十六年〉。《元史》为明初仓促官修，本纪保存大量实录条目但有编次与纪年问题；行省执行、数字、族称及对外关系须以条格、多语文书和对方史料互校。 现代研究可从元代行省、赋役、河工与多语文书研究继续追问。

\noindent\eventanchor{SYD-1366-YUAN-142}{八月戊寅，以李國鳳為中書左丞，陳有定為福建行省平章政事}\textbf{元。}
元的1366年（元至正26年；公元换算据《元史》本纪纪年）记录写到“八月戊寅，以李國鳳為中書左丞，陳有定為福建行省平章政事”。我们可以据此辨认一个可见的选择或压力，而不能把零散的官方记载、碑刻或后出叙事拼成无缝的社会全景。

其前提记为“本纪年条记录政令或机构处置；实施背景须参照制度材料。”；其后续是“可在同年及后续政令中追踪；条文不单独证明地方执行。”。这是一条待后续材料检验的事件链，而不是对所有行动者意图的替写。

可核的材料入口是《元史》卷047〈至正二十六年〉。《元史》为明初仓促官修，本纪保存大量实录条目但有编次与纪年问题；行省执行、数字、族称及对外关系须以条格、多语文书和对方史料互校。 就此而言，元代行省、赋役、河工与多语文书研究有助于把年序同区域、制度或物质材料并读。

\noindent\eventanchor{SYD-1366-YUAN-214}{丙申，擴廓帖木兒遣朱珍、盧旺屯兵河中，遣關保、虎林赤合兵渡河，會竹貞、商暠，且約李思齊以攻張良弼}\textbf{元。}
在1366年（元至正26年；公元换算据《元史》本纪纪年）的材料中，元面对的是丙申，擴廓帖木兒遣朱珍、盧旺屯兵河中，遣關保、虎林赤合兵渡河，會竹貞、商暠，且約李思齊以攻張良弼。这里应先看行动所依赖的关隘、城镇、河道或粮道，再讨论政治后果；军队抵达某处，并不等于它已能稳定征收、安置或控制周边。

从账本的前后项看，“本纪保留灾害或工程的报告地点与朝廷处置；灾情范围需另证。”提供了背景压力；“可与地方文书、河工和环境材料比较，不将单点记录外推为全域因果。”标出后来需要追踪的变化，二者之间仍存在未被材料填满的环节。

核心线索为《元史》卷047〈至正二十六年〉。《元史》为明初仓促官修，本纪保存大量实录条目但有编次与纪年问题；行省执行、数字、族称及对外关系须以条格、多语文书和对方史料互校。 现代研究可从元代行省、赋役、河工与多语文书研究继续追问。

\noindent\eventanchor{SYD-1366-YUAN-286}{丙戌，以方國珍為江浙行省左丞相，弟國瑛、國珉，姪明善，並為江浙行省平章政事}\textbf{元。}
元的1366年（元至正26年；公元换算据《元史》本纪纪年）记录写到“丙戌，以方國珍為江浙行省左丞相，弟國瑛、國珉，姪明善，並為江浙行省平章政事”。我们可以据此辨认一个可见的选择或压力，而不能把零散的官方记载、碑刻或后出叙事拼成无缝的社会全景。

账本把这一节点放在“本纪年条记录政令或机构处置；实施背景须参照制度材料。”与“可在同年及后续政令中追踪；条文不单独证明地方执行。”之间；两者提示前后压力，并不把时间相邻写成单一因果。

本条首先回到《元史》卷047〈至正二十六年〉；《元史》为明初仓促官修，本纪保存大量实录条目但有编次与纪年问题；行省执行、数字、族称及对外关系须以条格、多语文书和对方史料互校。 因而元代行省、赋役、河工与多语文书研究不是装饰性的书目，而是检验这一叙述边界的路径。

\subsection{1367年（元至正27年；公元换算据《元史》本纪纪年）}

\noindent\eventanchor{SYD-1367-YUAN-071}{庚戌，貊高殺衞輝守禦官余仁輔、彰德守禦官范國英，引軍至清化，聞懷慶有備，遂還彰德，上疏言：「人臣以尊君為本，以盡忠為心，以愛民為務}\textbf{元。}
“庚戌，貊高殺衞輝守禦官余仁輔、彰德守禦官范國英，引軍至清化，聞懷慶有備，遂還彰德，上疏言：「人臣以尊君為本，以盡忠為心，以愛民為務”把元置于军事行动的可核时点。战报能够记下攻守、入城或撤离，却很少完整保留沿途征发、居民去留和补给损耗；因此不能把一方的胜败叙事延伸成所有地方的经验。

从账本的前后项看，“本纪从朝廷视角记录军政行动；出兵与战果需分辨。”提供了背景压力；“后续路线、控制与损失应与对方记录和遗址材料复核。”标出后来需要追踪的变化，二者之间仍存在未被材料填满的环节。

核心线索为《元史》卷047〈至正二十七年〉。《元史》为明初仓促官修，本纪保存大量实录条目但有编次与纪年问题；行省执行、数字、族称及对外关系须以条格、多语文书和对方史料互校。 现代研究可从元代行省、赋役、河工与多语文书研究继续追问。

\noindent\eventanchor{SYD-1367-YUAN-143}{詔以擴廓帖木兒不遵君命，宜黜其兵權，就命貊高討之}\textbf{元。}
“詔以擴廓帖木兒不遵君命，宜黜其兵權，就命貊高討之”把元置于军事行动的可核时点。战报能够记下攻守、入城或撤离，却很少完整保留沿途征发、居民去留和补给损耗；因此不能把一方的胜败叙事延伸成所有地方的经验。

其前提记为“本纪年条记录政令或机构处置；实施背景须参照制度材料。”；其后续是“可在同年及后续政令中追踪；条文不单独证明地方执行。”。这是一条待后续材料检验的事件链，而不是对所有行动者意图的替写。

可核的材料入口是《元史》卷047〈至正二十七年〉。《元史》为明初仓促官修，本纪保存大量实录条目但有编次与纪年问题；行省执行、数字、族称及对外关系须以条格、多语文书和对方史料互校。 就此而言，元代行省、赋役、河工与多语文书研究有助于把年序同区域、制度或物质材料并读。

\noindent\eventanchor{SYD-1367-YUAN-215}{愛猷識理達臘計安宗社，累請出師}\textbf{元。}
“愛猷識理達臘計安宗社，累請出師”是元在1367年（元至正27年；公元换算据《元史》本纪纪年）可追查的节点。它提示命令、环境、宗教、迁徙或地方组织如何进入政治过程；影响的范围和速度仍要避免由一条材料外推。

从账本的前后项看，“本纪年条提供日期和中枢可见行动；前因不据单条补定。”提供了背景压力；“可回查同年及后续条目，地方影响仍须另证。”标出后来需要追踪的变化，二者之间仍存在未被材料填满的环节。

核心线索为《元史》卷047〈至正二十七年〉。《元史》为明初仓促官修，本纪保存大量实录条目但有编次与纪年问题；行省执行、数字、族称及对外关系须以条格、多语文书和对方史料互校。 现代研究可从元代行省、赋役、河工与多语文书研究继续追问。

\noindent\eventanchor{SYD-1367-YUAN-287}{八月丙午，詔命皇太子總天下兵馬，其略曰：「元良重任，職在撫軍，稽古徵今，卓有成憲}\textbf{元。}
“八月丙午，詔命皇太子總天下兵馬，其略曰：「元良重任，職在撫軍，稽古徵今，卓有成憲”把元置于军事行动的可核时点。战报能够记下攻守、入城或撤离，却很少完整保留沿途征发、居民去留和补给损耗；因此不能把一方的胜败叙事延伸成所有地方的经验。

账本把这一节点放在“本纪年条记录政令或机构处置；实施背景须参照制度材料。”与“可在同年及后续政令中追踪；条文不单独证明地方执行。”之间；两者提示前后压力，并不把时间相邻写成单一因果。

本条首先回到《元史》卷047〈至正二十七年〉；《元史》为明初仓促官修，本纪保存大量实录条目但有编次与纪年问题；行省执行、数字、族称及对外关系须以条格、多语文书和对方史料互校。 因而元代行省、赋役、河工与多语文书研究不是装饰性的书目，而是检验这一叙述边界的路径。

\subsection{1367年}

\noindent\eventanchor{YUAN-1367-EVENT-196}{明军发动北伐，元在华北防线动摇}\textbf{元与明军。}
1367年，元与明军的行动落在“明军发动北伐，元在华北防线动摇”所指的战事与通道上。可见的不是抽象的推进：攻守双方必须让人员、消息与补给穿过具体的城防、道路或水路；账本所列的结果只说明这一节点改变了下一步选择。

从账本的前后项看，“元中后期继承、财政与地方危机中既有的联盟、交通与补给压力，为“明军发动北伐，元在华北防线动摇”提供了直接的军事或动员背景。”提供了背景压力；““明军发动北伐，元在华北防线动摇”改变了相关城镇、道路或政权间的军事选择；伤亡、俘获和控制范围仅可按各方材料分别确认。”标出后来需要追踪的变化，二者之间仍存在未被材料填满的环节。

核心线索为《元史》相关本纪；《元朝秘史》、波斯文史籍或《高丽史》对应叙事；《元典章》《通制条格》、多语碑刻与文书（据事核）。《元史》明初速修，纪年与人名有异同；蒙古大汗政权不等同所有汗国，征服、赋役和身份分类须按地区及材料语言分列。 现代研究可从蒙古帝国、元朝行省财政、多语文书与区域社会研究继续追问。

\subsection{1368年（元至正28年；公元换算据《元史》本纪纪年）}

\noindent\eventanchor{SYD-1368-YUAN-072}{丙辰，詔皇太子位下立儀衞司，設指揮二員，給二珠金牌，副指揮二員，一珠金牌}\textbf{元。}
“丙辰，詔皇太子位下立儀衞司，設指揮二員，給二珠金牌，副指揮二員，一珠金牌”在1368年（元至正28年；公元换算据《元史》本纪纪年）构成元的继承或权力重组节点。名位的变化可以由纪年串连，但它没有自动消除宫廷、军政和地方网络原有的拉扯。

从账本的前后项看，“本纪年条记录政令或机构处置；实施背景须参照制度材料。”提供了背景压力；“可在同年及后续政令中追踪；条文不单独证明地方执行。”标出后来需要追踪的变化，二者之间仍存在未被材料填满的环节。

核心线索为《元史》卷043〈本纪第43卷本年条〉。《元史》为明初仓促官修，本纪保存大量实录条目但有编次与纪年问题；行省执行、数字、族称及对外关系须以条格、多语文书和对方史料互校。 现代研究可从元代行省、赋役、河工与多语文书研究继续追问。

\noindent\eventanchor{SYD-1368-YUAN-144}{丙戌，以武衞所管鹽臺屯田八百頃，除軍見種外，荒閑之地，盡付分司農司}\textbf{元。}
1368年（元至正28年；公元换算据《元史》本纪纪年）的“丙戌，以武衞所管鹽臺屯田八百頃，除軍見種外，荒閑之地，盡付分司農司”将元的财政或生产安排推到可见处。制度文本能说明官府打算如何收取、转运或调节，却不能直接等同于不同地区的实收、流通或家庭负担。

其前提记为“本纪年条记录政令或机构处置；实施背景须参照制度材料。”；其后续是“可在同年及后续政令中追踪；条文不单独证明地方执行。”。这是一条待后续材料检验的事件链，而不是对所有行动者意图的替写。

可核的材料入口是《元史》卷043〈本纪第43卷本年条〉。《元史》为明初仓促官修，本纪保存大量实录条目但有编次与纪年问题；行省执行、数字、族称及对外关系须以条格、多语文书和对方史料互校。 就此而言，元代行省、赋役、河工与多语文书研究有助于把年序同区域、制度或物质材料并读。

\noindent\eventanchor{SYD-1368-YUAN-216}{賜吳王搠思監金二錠、銀五錠、鈔二千錠、幣帛各九匹}\textbf{元。}
在1368年（元至正28年；公元换算据《元史》本纪纪年）的可见材料中，元留下“賜吳王搠思監金二錠、銀五錠、鈔二千錠、幣帛各九匹”。这一动作把具体人群、机构或地方条件带进年序，但一项记录不能替未被记载者概括整体反应。

从账本的前后项看，“本纪所记财用、赈济或征发属于特定报告场景。”提供了背景压力；“可与食货志、文书和物质材料比较，不外推为全境总量。”标出后来需要追踪的变化，二者之间仍存在未被材料填满的环节。

核心线索为《元史》卷043〈本纪第43卷本年条〉。《元史》为明初仓促官修，本纪保存大量实录条目但有编次与纪年问题；行省执行、数字、族称及对外关系须以条格、多语文书和对方史料互校。 现代研究可从元代行省、赋役、河工与多语文书研究继续追问。

\noindent\eventanchor{SYD-1368-YUAN-288}{答失八都魯克復襄陽、樊城有功，陞四川行省右丞，賜金繫腰一}\textbf{元。}
在1368年（元至正28年；公元换算据《元史》本纪纪年）的可见材料中，元留下“答失八都魯克復襄陽、樊城有功，陞四川行省右丞，賜金繫腰一”。这一动作把具体人群、机构或地方条件带进年序，但一项记录不能替未被记载者概括整体反应。

账本把这一节点放在“本纪年条记录政令或机构处置；实施背景须参照制度材料。”与“可在同年及后续政令中追踪；条文不单独证明地方执行。”之间；两者提示前后压力，并不把时间相邻写成单一因果。

本条首先回到《元史》卷043〈本纪第43卷本年条〉；《元史》为明初仓促官修，本纪保存大量实录条目但有编次与纪年问题；行省执行、数字、族称及对外关系须以条格、多语文书和对方史料互校。 因而元代行省、赋役、河工与多语文书研究不是装饰性的书目，而是检验这一叙述边界的路径。

\subsection{1368年}

\noindent\eventanchor{YUAN-1368-EVENT-197}{明军进入大都，元顺帝北撤}\textbf{元与明。}
1368年，元与明的行动落在“明军进入大都，元顺帝北撤”所指的战事与通道上。可见的不是抽象的推进：攻守双方必须让人员、消息与补给穿过具体的城防、道路或水路；账本所列的结果只说明这一节点改变了下一步选择。

其前提记为““明军进入大都，元顺帝北撤”发生在元中后期继承、财政与地方危机的具体压力与机会之中，相关行动者的选择不能脱离此前事件链解释。”；其后续是““明军进入大都，元顺帝北撤”为后续政治、边防或资源配置留下条件；其社会影响与实际覆盖范围仍须以异类材料限定。”。这是一条待后续材料检验的事件链，而不是对所有行动者意图的替写。

可核的材料入口是《元史》相关本纪；《元朝秘史》、波斯文史籍或《高丽史》对应叙事；《元典章》《通制条格》、多语碑刻与文书（据事核）。《元史》明初速修，纪年与人名有异同；蒙古大汗政权不等同所有汗国，征服、赋役和身份分类须按地区及材料语言分列。 就此而言，蒙古帝国、元朝行省财政、多语文书与区域社会研究有助于把年序同区域、制度或物质材料并读。

\subsection{1280 年（共享年序桥接）}

\noindent\eventanchor{SY-1280-CHRONOLOGY-BRIDGE}{【共享年序日期坐标】保留1280 年的多政权并行检索入口；本条不主张所列政权在当年发生直接互动，具体事件须另以单项材料入账。}\textbf{元与高丽、安南及地方政权（年序桥接）。}
作为共享年序的接点，“【共享年序日期坐标】保留1280 年的多政权并行检索入口；本条不主张所列政权在当年发生直接互动，具体事件须另以单项材料入账。”让元与高丽、安南及地方政权（年序桥接）与同年其他地区保持可比的日期位置。可确认的是文书使用这一时间框架；年内何事发生、何处落实，仍须由各自条目另证。

从账本的前后项看，“相邻已入账事件之间缺少该年度的共享年序入口，易使跨政权材料被单一政权叙事遮蔽。”提供了背景压力；“保留该年度的检索与补账位置；后续仅在获得可核单项材料时拆分为具体政治、财政、军事、社会或文化事件。”标出后来需要追踪的变化，二者之间仍存在未被材料填满的环节。

核心线索为《元史》；《高丽史》；《安南志略》及碑刻、地方文书（据事核）。本条只确认该年位于可追踪的共享年序中；正史、地方文书和外部材料的保存密度不一，不能由桥接标签推断行动、统属、疆界或社会影响。 现代研究可从元代区域治理、海域联通与地方社会的并行年序研究继续追问。

\subsection{1282 年（共享年序桥接）}

\noindent\eventanchor{SY-1282-CHRONOLOGY-BRIDGE}{【共享年序日期坐标】保留1282 年的多政权并行检索入口；本条不主张所列政权在当年发生直接互动，具体事件须另以单项材料入账。}\textbf{元与高丽、安南及地方政权（年序桥接）。}
作为共享年序的接点，“【共享年序日期坐标】保留1282 年的多政权并行检索入口；本条不主张所列政权在当年发生直接互动，具体事件须另以单项材料入账。”让元与高丽、安南及地方政权（年序桥接）与同年其他地区保持可比的日期位置。可确认的是文书使用这一时间框架；年内何事发生、何处落实，仍须由各自条目另证。

其前提记为“相邻已入账事件之间缺少该年度的共享年序入口，易使跨政权材料被单一政权叙事遮蔽。”；其后续是“保留该年度的检索与补账位置；后续仅在获得可核单项材料时拆分为具体政治、财政、军事、社会或文化事件。”。这是一条待后续材料检验的事件链，而不是对所有行动者意图的替写。

可核的材料入口是《元史》；《高丽史》；《安南志略》及碑刻、地方文书（据事核）。本条只确认该年位于可追踪的共享年序中；正史、地方文书和外部材料的保存密度不一，不能由桥接标签推断行动、统属、疆界或社会影响。 就此而言，元代区域治理、海域联通与地方社会的并行年序研究有助于把年序同区域、制度或物质材料并读。

\subsection{1283 年（共享年序桥接）}

\noindent\eventanchor{SY-1283-CHRONOLOGY-BRIDGE}{【共享年序日期坐标】保留1283 年的多政权并行检索入口；本条不主张所列政权在当年发生直接互动，具体事件须另以单项材料入账。}\textbf{元与高丽、安南及地方政权（年序桥接）。}
元与高丽、安南及地方政权（年序桥接）在1283 年（共享年序桥接）的这条年序记录提供的是检索坐标：“【共享年序日期坐标】保留1283 年的多政权并行检索入口；本条不主张所列政权在当年发生直接互动，具体事件须另以单项材料入账。”。它保留材料稀薄处的空白，而不以一条王朝纪年替代地方、边疆或竞争政权的行动。

账本把这一节点放在“相邻已入账事件之间缺少该年度的共享年序入口，易使跨政权材料被单一政权叙事遮蔽。”与“保留该年度的检索与补账位置；后续仅在获得可核单项材料时拆分为具体政治、财政、军事、社会或文化事件。”之间；两者提示前后压力，并不把时间相邻写成单一因果。

本条首先回到《元史》；《高丽史》；《安南志略》及碑刻、地方文书（据事核）；本条只确认该年位于可追踪的共享年序中；正史、地方文书和外部材料的保存密度不一，不能由桥接标签推断行动、统属、疆界或社会影响。 因而元代区域治理、海域联通与地方社会的并行年序研究不是装饰性的书目，而是检验这一叙述边界的路径。

\subsection{1284 年（共享年序桥接）}

\noindent\eventanchor{SY-1284-CHRONOLOGY-BRIDGE}{【共享年序日期坐标】保留1284 年的多政权并行检索入口；本条不主张所列政权在当年发生直接互动，具体事件须另以单项材料入账。}\textbf{元与高丽、安南及地方政权（年序桥接）。}
1284 年（共享年序桥接）的元与高丽、安南及地方政权（年序桥接）材料只留下“【共享年序日期坐标】保留1284 年的多政权并行检索入口；本条不主张所列政权在当年发生直接互动，具体事件须另以单项材料入账。”这一纪年接点。它使同年的诏令、交涉或地方材料可以被放回并行时间轴，却不应被误写成一件已经具备完整情节的政治行动。

从账本的前后项看，“相邻已入账事件之间缺少该年度的共享年序入口，易使跨政权材料被单一政权叙事遮蔽。”提供了背景压力；“保留该年度的检索与补账位置；后续仅在获得可核单项材料时拆分为具体政治、财政、军事、社会或文化事件。”标出后来需要追踪的变化，二者之间仍存在未被材料填满的环节。

核心线索为《元史》；《高丽史》；《安南志略》及碑刻、地方文书（据事核）。本条只确认该年位于可追踪的共享年序中；正史、地方文书和外部材料的保存密度不一，不能由桥接标签推断行动、统属、疆界或社会影响。 现代研究可从元代区域治理、海域联通与地方社会的并行年序研究继续追问。

\subsection{1285 年（共享年序桥接）}

\noindent\eventanchor{SY-1285-CHRONOLOGY-BRIDGE}{【共享年序日期坐标】保留1285 年的多政权并行检索入口；本条不主张所列政权在当年发生直接互动，具体事件须另以单项材料入账。}\textbf{元与高丽、安南及地方政权（年序桥接）。}
作为共享年序的接点，“【共享年序日期坐标】保留1285 年的多政权并行检索入口；本条不主张所列政权在当年发生直接互动，具体事件须另以单项材料入账。”让元与高丽、安南及地方政权（年序桥接）与同年其他地区保持可比的日期位置。可确认的是文书使用这一时间框架；年内何事发生、何处落实，仍须由各自条目另证。

此前的条件是“相邻已入账事件之间缺少该年度的共享年序入口，易使跨政权材料被单一政权叙事遮蔽。”，此后的可见影响是“保留该年度的检索与补账位置；后续仅在获得可核单项材料时拆分为具体政治、财政、军事、社会或文化事件。”。这里的“影响”须按地点、机构和材料种类分别衡量。

这一叙述依据《元史》；《高丽史》；《安南志略》及碑刻、地方文书（据事核），但本条只确认该年位于可追踪的共享年序中；正史、地方文书和外部材料的保存密度不一，不能由桥接标签推断行动、统属、疆界或社会影响。。研究元代区域治理、海域联通与地方社会的并行年序研究可以补问材料未直接说出的空间与社会差异。

\subsection{1286 年（共享年序桥接）}

\noindent\eventanchor{SY-1286-CHRONOLOGY-BRIDGE}{【共享年序日期坐标】保留1286 年的多政权并行检索入口；本条不主张所列政权在当年发生直接互动，具体事件须另以单项材料入账。}\textbf{元与高丽、安南及地方政权（年序桥接）。}
1286 年（共享年序桥接）的元与高丽、安南及地方政权（年序桥接）材料只留下“【共享年序日期坐标】保留1286 年的多政权并行检索入口；本条不主张所列政权在当年发生直接互动，具体事件须另以单项材料入账。”这一纪年接点。它使同年的诏令、交涉或地方材料可以被放回并行时间轴，却不应被误写成一件已经具备完整情节的政治行动。

其前提记为“相邻已入账事件之间缺少该年度的共享年序入口，易使跨政权材料被单一政权叙事遮蔽。”；其后续是“保留该年度的检索与补账位置；后续仅在获得可核单项材料时拆分为具体政治、财政、军事、社会或文化事件。”。这是一条待后续材料检验的事件链，而不是对所有行动者意图的替写。

可核的材料入口是《元史》；《高丽史》；《安南志略》及碑刻、地方文书（据事核）。本条只确认该年位于可追踪的共享年序中；正史、地方文书和外部材料的保存密度不一，不能由桥接标签推断行动、统属、疆界或社会影响。 就此而言，元代区域治理、海域联通与地方社会的并行年序研究有助于把年序同区域、制度或物质材料并读。

\subsection{1288 年（共享年序桥接）}

\noindent\eventanchor{SY-1288-CHRONOLOGY-BRIDGE}{【共享年序日期坐标】保留1288 年的多政权并行检索入口；本条不主张所列政权在当年发生直接互动，具体事件须另以单项材料入账。}\textbf{元与高丽、安南及地方政权（年序桥接）。}
元与高丽、安南及地方政权（年序桥接）在1288 年（共享年序桥接）的这条年序记录提供的是检索坐标：“【共享年序日期坐标】保留1288 年的多政权并行检索入口；本条不主张所列政权在当年发生直接互动，具体事件须另以单项材料入账。”。它保留材料稀薄处的空白，而不以一条王朝纪年替代地方、边疆或竞争政权的行动。

从账本的前后项看，“相邻已入账事件之间缺少该年度的共享年序入口，易使跨政权材料被单一政权叙事遮蔽。”提供了背景压力；“保留该年度的检索与补账位置；后续仅在获得可核单项材料时拆分为具体政治、财政、军事、社会或文化事件。”标出后来需要追踪的变化，二者之间仍存在未被材料填满的环节。

核心线索为《元史》；《高丽史》；《安南志略》及碑刻、地方文书（据事核）。本条只确认该年位于可追踪的共享年序中；正史、地方文书和外部材料的保存密度不一，不能由桥接标签推断行动、统属、疆界或社会影响。 现代研究可从元代区域治理、海域联通与地方社会的并行年序研究继续追问。

\subsection{1289 年（共享年序桥接）}

\noindent\eventanchor{SY-1289-CHRONOLOGY-BRIDGE}{【共享年序日期坐标】保留1289 年的多政权并行检索入口；本条不主张所列政权在当年发生直接互动，具体事件须另以单项材料入账。}\textbf{元与高丽、安南及地方政权（年序桥接）。}
1289 年（共享年序桥接）的元与高丽、安南及地方政权（年序桥接）材料只留下“【共享年序日期坐标】保留1289 年的多政权并行检索入口；本条不主张所列政权在当年发生直接互动，具体事件须另以单项材料入账。”这一纪年接点。它使同年的诏令、交涉或地方材料可以被放回并行时间轴，却不应被误写成一件已经具备完整情节的政治行动。

此前的条件是“相邻已入账事件之间缺少该年度的共享年序入口，易使跨政权材料被单一政权叙事遮蔽。”，此后的可见影响是“保留该年度的检索与补账位置；后续仅在获得可核单项材料时拆分为具体政治、财政、军事、社会或文化事件。”。这里的“影响”须按地点、机构和材料种类分别衡量。

这一叙述依据《元史》；《高丽史》；《安南志略》及碑刻、地方文书（据事核），但本条只确认该年位于可追踪的共享年序中；正史、地方文书和外部材料的保存密度不一，不能由桥接标签推断行动、统属、疆界或社会影响。。研究元代区域治理、海域联通与地方社会的并行年序研究可以补问材料未直接说出的空间与社会差异。

\subsection{1290 年（共享年序桥接）}

\noindent\eventanchor{SY-1290-CHRONOLOGY-BRIDGE}{【共享年序日期坐标】保留1290 年的多政权并行检索入口；本条不主张所列政权在当年发生直接互动，具体事件须另以单项材料入账。}\textbf{元与高丽、安南及地方政权（年序桥接）。}
元与高丽、安南及地方政权（年序桥接）在1290 年（共享年序桥接）的这条年序记录提供的是检索坐标：“【共享年序日期坐标】保留1290 年的多政权并行检索入口；本条不主张所列政权在当年发生直接互动，具体事件须另以单项材料入账。”。它保留材料稀薄处的空白，而不以一条王朝纪年替代地方、边疆或竞争政权的行动。

账本把这一节点放在“相邻已入账事件之间缺少该年度的共享年序入口，易使跨政权材料被单一政权叙事遮蔽。”与“保留该年度的检索与补账位置；后续仅在获得可核单项材料时拆分为具体政治、财政、军事、社会或文化事件。”之间；两者提示前后压力，并不把时间相邻写成单一因果。

本条首先回到《元史》；《高丽史》；《安南志略》及碑刻、地方文书（据事核）；本条只确认该年位于可追踪的共享年序中；正史、地方文书和外部材料的保存密度不一，不能由桥接标签推断行动、统属、疆界或社会影响。 因而元代区域治理、海域联通与地方社会的并行年序研究不是装饰性的书目，而是检验这一叙述边界的路径。

\subsection{1291 年（共享年序桥接）}

\noindent\eventanchor{SY-1291-CHRONOLOGY-BRIDGE}{【共享年序日期坐标】保留1291 年的多政权并行检索入口；本条不主张所列政权在当年发生直接互动，具体事件须另以单项材料入账。}\textbf{元与高丽、安南及地方政权（年序桥接）。}
作为共享年序的接点，“【共享年序日期坐标】保留1291 年的多政权并行检索入口；本条不主张所列政权在当年发生直接互动，具体事件须另以单项材料入账。”让元与高丽、安南及地方政权（年序桥接）与同年其他地区保持可比的日期位置。可确认的是文书使用这一时间框架；年内何事发生、何处落实，仍须由各自条目另证。

从账本的前后项看，“相邻已入账事件之间缺少该年度的共享年序入口，易使跨政权材料被单一政权叙事遮蔽。”提供了背景压力；“保留该年度的检索与补账位置；后续仅在获得可核单项材料时拆分为具体政治、财政、军事、社会或文化事件。”标出后来需要追踪的变化，二者之间仍存在未被材料填满的环节。

核心线索为《元史》；《高丽史》；《安南志略》及碑刻、地方文书（据事核）。本条只确认该年位于可追踪的共享年序中；正史、地方文书和外部材料的保存密度不一，不能由桥接标签推断行动、统属、疆界或社会影响。 现代研究可从元代区域治理、海域联通与地方社会的并行年序研究继续追问。

\subsection{1292 年（共享年序桥接）}

\noindent\eventanchor{SY-1292-CHRONOLOGY-BRIDGE}{【共享年序日期坐标】保留1292 年的多政权并行检索入口；本条不主张所列政权在当年发生直接互动，具体事件须另以单项材料入账。}\textbf{元与高丽、安南及地方政权（年序桥接）。}
元与高丽、安南及地方政权（年序桥接）在1292 年（共享年序桥接）的这条年序记录提供的是检索坐标：“【共享年序日期坐标】保留1292 年的多政权并行检索入口；本条不主张所列政权在当年发生直接互动，具体事件须另以单项材料入账。”。它保留材料稀薄处的空白，而不以一条王朝纪年替代地方、边疆或竞争政权的行动。

此前的条件是“相邻已入账事件之间缺少该年度的共享年序入口，易使跨政权材料被单一政权叙事遮蔽。”，此后的可见影响是“保留该年度的检索与补账位置；后续仅在获得可核单项材料时拆分为具体政治、财政、军事、社会或文化事件。”。这里的“影响”须按地点、机构和材料种类分别衡量。

这一叙述依据《元史》；《高丽史》；《安南志略》及碑刻、地方文书（据事核），但本条只确认该年位于可追踪的共享年序中；正史、地方文书和外部材料的保存密度不一，不能由桥接标签推断行动、统属、疆界或社会影响。。研究元代区域治理、海域联通与地方社会的并行年序研究可以补问材料未直接说出的空间与社会差异。

\subsection{1293 年（共享年序桥接）}

\noindent\eventanchor{SY-1293-CHRONOLOGY-BRIDGE}{【共享年序日期坐标】保留1293 年的多政权并行检索入口；本条不主张所列政权在当年发生直接互动，具体事件须另以单项材料入账。}\textbf{元与高丽、安南及地方政权（年序桥接）。}
1293 年（共享年序桥接）的元与高丽、安南及地方政权（年序桥接）材料只留下“【共享年序日期坐标】保留1293 年的多政权并行检索入口；本条不主张所列政权在当年发生直接互动，具体事件须另以单项材料入账。”这一纪年接点。它使同年的诏令、交涉或地方材料可以被放回并行时间轴，却不应被误写成一件已经具备完整情节的政治行动。

其前提记为“相邻已入账事件之间缺少该年度的共享年序入口，易使跨政权材料被单一政权叙事遮蔽。”；其后续是“保留该年度的检索与补账位置；后续仅在获得可核单项材料时拆分为具体政治、财政、军事、社会或文化事件。”。这是一条待后续材料检验的事件链，而不是对所有行动者意图的替写。

可核的材料入口是《元史》；《高丽史》；《安南志略》及碑刻、地方文书（据事核）。本条只确认该年位于可追踪的共享年序中；正史、地方文书和外部材料的保存密度不一，不能由桥接标签推断行动、统属、疆界或社会影响。 就此而言，元代区域治理、海域联通与地方社会的并行年序研究有助于把年序同区域、制度或物质材料并读。

\subsection{1295 年（共享年序桥接）}

\noindent\eventanchor{SY-1295-CHRONOLOGY-BRIDGE}{【共享年序日期坐标】保留1295 年的多政权并行检索入口；本条不主张所列政权在当年发生直接互动，具体事件须另以单项材料入账。}\textbf{元与高丽、安南及地方政权（年序桥接）。}
元与高丽、安南及地方政权（年序桥接）在1295 年（共享年序桥接）的这条年序记录提供的是检索坐标：“【共享年序日期坐标】保留1295 年的多政权并行检索入口；本条不主张所列政权在当年发生直接互动，具体事件须另以单项材料入账。”。它保留材料稀薄处的空白，而不以一条王朝纪年替代地方、边疆或竞争政权的行动。

从账本的前后项看，“相邻已入账事件之间缺少该年度的共享年序入口，易使跨政权材料被单一政权叙事遮蔽。”提供了背景压力；“保留该年度的检索与补账位置；后续仅在获得可核单项材料时拆分为具体政治、财政、军事、社会或文化事件。”标出后来需要追踪的变化，二者之间仍存在未被材料填满的环节。

核心线索为《元史》；《高丽史》；《安南志略》及碑刻、地方文书（据事核）。本条只确认该年位于可追踪的共享年序中；正史、地方文书和外部材料的保存密度不一，不能由桥接标签推断行动、统属、疆界或社会影响。 现代研究可从元代区域治理、海域联通与地方社会的并行年序研究继续追问。

\subsection{1296 年（共享年序桥接）}

\noindent\eventanchor{SY-1296-CHRONOLOGY-BRIDGE}{【共享年序日期坐标】保留1296 年的多政权并行检索入口；本条不主张所列政权在当年发生直接互动，具体事件须另以单项材料入账。}\textbf{元与高丽、安南及地方政权（年序桥接）。}
作为共享年序的接点，“【共享年序日期坐标】保留1296 年的多政权并行检索入口；本条不主张所列政权在当年发生直接互动，具体事件须另以单项材料入账。”让元与高丽、安南及地方政权（年序桥接）与同年其他地区保持可比的日期位置。可确认的是文书使用这一时间框架；年内何事发生、何处落实，仍须由各自条目另证。

此前的条件是“相邻已入账事件之间缺少该年度的共享年序入口，易使跨政权材料被单一政权叙事遮蔽。”，此后的可见影响是“保留该年度的检索与补账位置；后续仅在获得可核单项材料时拆分为具体政治、财政、军事、社会或文化事件。”。这里的“影响”须按地点、机构和材料种类分别衡量。

这一叙述依据《元史》；《高丽史》；《安南志略》及碑刻、地方文书（据事核），但本条只确认该年位于可追踪的共享年序中；正史、地方文书和外部材料的保存密度不一，不能由桥接标签推断行动、统属、疆界或社会影响。。研究元代区域治理、海域联通与地方社会的并行年序研究可以补问材料未直接说出的空间与社会差异。

\subsection{1297 年（共享年序桥接）}

\noindent\eventanchor{SY-1297-CHRONOLOGY-BRIDGE}{【共享年序日期坐标】保留1297 年的多政权并行检索入口；本条不主张所列政权在当年发生直接互动，具体事件须另以单项材料入账。}\textbf{元与高丽、安南及地方政权（年序桥接）。}
元与高丽、安南及地方政权（年序桥接）在1297 年（共享年序桥接）的这条年序记录提供的是检索坐标：“【共享年序日期坐标】保留1297 年的多政权并行检索入口；本条不主张所列政权在当年发生直接互动，具体事件须另以单项材料入账。”。它保留材料稀薄处的空白，而不以一条王朝纪年替代地方、边疆或竞争政权的行动。

其前提记为“相邻已入账事件之间缺少该年度的共享年序入口，易使跨政权材料被单一政权叙事遮蔽。”；其后续是“保留该年度的检索与补账位置；后续仅在获得可核单项材料时拆分为具体政治、财政、军事、社会或文化事件。”。这是一条待后续材料检验的事件链，而不是对所有行动者意图的替写。

可核的材料入口是《元史》；《高丽史》；《安南志略》及碑刻、地方文书（据事核）。本条只确认该年位于可追踪的共享年序中；正史、地方文书和外部材料的保存密度不一，不能由桥接标签推断行动、统属、疆界或社会影响。 就此而言，元代区域治理、海域联通与地方社会的并行年序研究有助于把年序同区域、制度或物质材料并读。

\subsection{1298 年（共享年序桥接）}

\noindent\eventanchor{SY-1298-CHRONOLOGY-BRIDGE}{【共享年序日期坐标】保留1298 年的多政权并行检索入口；本条不主张所列政权在当年发生直接互动，具体事件须另以单项材料入账。}\textbf{元与高丽、安南及地方政权（年序桥接）。}
1298 年（共享年序桥接）的元与高丽、安南及地方政权（年序桥接）材料只留下“【共享年序日期坐标】保留1298 年的多政权并行检索入口；本条不主张所列政权在当年发生直接互动，具体事件须另以单项材料入账。”这一纪年接点。它使同年的诏令、交涉或地方材料可以被放回并行时间轴，却不应被误写成一件已经具备完整情节的政治行动。

账本把这一节点放在“相邻已入账事件之间缺少该年度的共享年序入口，易使跨政权材料被单一政权叙事遮蔽。”与“保留该年度的检索与补账位置；后续仅在获得可核单项材料时拆分为具体政治、财政、军事、社会或文化事件。”之间；两者提示前后压力，并不把时间相邻写成单一因果。

本条首先回到《元史》；《高丽史》；《安南志略》及碑刻、地方文书（据事核）；本条只确认该年位于可追踪的共享年序中；正史、地方文书和外部材料的保存密度不一，不能由桥接标签推断行动、统属、疆界或社会影响。 因而元代区域治理、海域联通与地方社会的并行年序研究不是装饰性的书目，而是检验这一叙述边界的路径。

\subsection{1299 年（共享年序桥接）}

\noindent\eventanchor{SY-1299-CHRONOLOGY-BRIDGE}{【共享年序日期坐标】保留1299 年的多政权并行检索入口；本条不主张所列政权在当年发生直接互动，具体事件须另以单项材料入账。}\textbf{元与高丽、安南及地方政权（年序桥接）。}
作为共享年序的接点，“【共享年序日期坐标】保留1299 年的多政权并行检索入口；本条不主张所列政权在当年发生直接互动，具体事件须另以单项材料入账。”让元与高丽、安南及地方政权（年序桥接）与同年其他地区保持可比的日期位置。可确认的是文书使用这一时间框架；年内何事发生、何处落实，仍须由各自条目另证。

从账本的前后项看，“相邻已入账事件之间缺少该年度的共享年序入口，易使跨政权材料被单一政权叙事遮蔽。”提供了背景压力；“保留该年度的检索与补账位置；后续仅在获得可核单项材料时拆分为具体政治、财政、军事、社会或文化事件。”标出后来需要追踪的变化，二者之间仍存在未被材料填满的环节。

核心线索为《元史》；《高丽史》；《安南志略》及碑刻、地方文书（据事核）。本条只确认该年位于可追踪的共享年序中；正史、地方文书和外部材料的保存密度不一，不能由桥接标签推断行动、统属、疆界或社会影响。 现代研究可从元代区域治理、海域联通与地方社会的并行年序研究继续追问。

\subsection{1300 年（共享年序桥接）}

\noindent\eventanchor{SY-1300-CHRONOLOGY-BRIDGE}{【共享年序日期坐标】保留1300 年的多政权并行检索入口；本条不主张所列政权在当年发生直接互动，具体事件须另以单项材料入账。}\textbf{元与高丽、安南及地方政权（年序桥接）。}
1300 年（共享年序桥接）的元与高丽、安南及地方政权（年序桥接）材料只留下“【共享年序日期坐标】保留1300 年的多政权并行检索入口；本条不主张所列政权在当年发生直接互动，具体事件须另以单项材料入账。”这一纪年接点。它使同年的诏令、交涉或地方材料可以被放回并行时间轴，却不应被误写成一件已经具备完整情节的政治行动。

此前的条件是“相邻已入账事件之间缺少该年度的共享年序入口，易使跨政权材料被单一政权叙事遮蔽。”，此后的可见影响是“保留该年度的检索与补账位置；后续仅在获得可核单项材料时拆分为具体政治、财政、军事、社会或文化事件。”。这里的“影响”须按地点、机构和材料种类分别衡量。

这一叙述依据《元史》；《高丽史》；《安南志略》及碑刻、地方文书（据事核），但本条只确认该年位于可追踪的共享年序中；正史、地方文书和外部材料的保存密度不一，不能由桥接标签推断行动、统属、疆界或社会影响。。研究元代区域治理、海域联通与地方社会的并行年序研究可以补问材料未直接说出的空间与社会差异。

\subsection{1301 年（共享年序桥接）}

\noindent\eventanchor{SY-1301-CHRONOLOGY-BRIDGE}{【共享年序日期坐标】保留1301 年的多政权并行检索入口；本条不主张所列政权在当年发生直接互动，具体事件须另以单项材料入账。}\textbf{元与高丽、安南及地方政权（年序桥接）。}
作为共享年序的接点，“【共享年序日期坐标】保留1301 年的多政权并行检索入口；本条不主张所列政权在当年发生直接互动，具体事件须另以单项材料入账。”让元与高丽、安南及地方政权（年序桥接）与同年其他地区保持可比的日期位置。可确认的是文书使用这一时间框架；年内何事发生、何处落实，仍须由各自条目另证。

其前提记为“相邻已入账事件之间缺少该年度的共享年序入口，易使跨政权材料被单一政权叙事遮蔽。”；其后续是“保留该年度的检索与补账位置；后续仅在获得可核单项材料时拆分为具体政治、财政、军事、社会或文化事件。”。这是一条待后续材料检验的事件链，而不是对所有行动者意图的替写。

可核的材料入口是《元史》；《高丽史》；《安南志略》及碑刻、地方文书（据事核）。本条只确认该年位于可追踪的共享年序中；正史、地方文书和外部材料的保存密度不一，不能由桥接标签推断行动、统属、疆界或社会影响。 就此而言，元代区域治理、海域联通与地方社会的并行年序研究有助于把年序同区域、制度或物质材料并读。

\subsection{1302 年（共享年序桥接）}

\noindent\eventanchor{SY-1302-CHRONOLOGY-BRIDGE}{【共享年序日期坐标】保留1302 年的多政权并行检索入口；本条不主张所列政权在当年发生直接互动，具体事件须另以单项材料入账。}\textbf{元与高丽、安南及地方政权（年序桥接）。}
元与高丽、安南及地方政权（年序桥接）在1302 年（共享年序桥接）的这条年序记录提供的是检索坐标：“【共享年序日期坐标】保留1302 年的多政权并行检索入口；本条不主张所列政权在当年发生直接互动，具体事件须另以单项材料入账。”。它保留材料稀薄处的空白，而不以一条王朝纪年替代地方、边疆或竞争政权的行动。

账本把这一节点放在“相邻已入账事件之间缺少该年度的共享年序入口，易使跨政权材料被单一政权叙事遮蔽。”与“保留该年度的检索与补账位置；后续仅在获得可核单项材料时拆分为具体政治、财政、军事、社会或文化事件。”之间；两者提示前后压力，并不把时间相邻写成单一因果。

本条首先回到《元史》；《高丽史》；《安南志略》及碑刻、地方文书（据事核）；本条只确认该年位于可追踪的共享年序中；正史、地方文书和外部材料的保存密度不一，不能由桥接标签推断行动、统属、疆界或社会影响。 因而元代区域治理、海域联通与地方社会的并行年序研究不是装饰性的书目，而是检验这一叙述边界的路径。

\subsection{1303 年（共享年序桥接）}

\noindent\eventanchor{SY-1303-CHRONOLOGY-BRIDGE}{【共享年序日期坐标】保留1303 年的多政权并行检索入口；本条不主张所列政权在当年发生直接互动，具体事件须另以单项材料入账。}\textbf{元与高丽、安南及地方政权（年序桥接）。}
1303 年（共享年序桥接）的元与高丽、安南及地方政权（年序桥接）材料只留下“【共享年序日期坐标】保留1303 年的多政权并行检索入口；本条不主张所列政权在当年发生直接互动，具体事件须另以单项材料入账。”这一纪年接点。它使同年的诏令、交涉或地方材料可以被放回并行时间轴，却不应被误写成一件已经具备完整情节的政治行动。

从账本的前后项看，“相邻已入账事件之间缺少该年度的共享年序入口，易使跨政权材料被单一政权叙事遮蔽。”提供了背景压力；“保留该年度的检索与补账位置；后续仅在获得可核单项材料时拆分为具体政治、财政、军事、社会或文化事件。”标出后来需要追踪的变化，二者之间仍存在未被材料填满的环节。

核心线索为《元史》；《高丽史》；《安南志略》及碑刻、地方文书（据事核）。本条只确认该年位于可追踪的共享年序中；正史、地方文书和外部材料的保存密度不一，不能由桥接标签推断行动、统属、疆界或社会影响。 现代研究可从元代区域治理、海域联通与地方社会的并行年序研究继续追问。

\subsection{1304 年（共享年序桥接）}

\noindent\eventanchor{SY-1304-CHRONOLOGY-BRIDGE}{【共享年序日期坐标】保留1304 年的多政权并行检索入口；本条不主张所列政权在当年发生直接互动，具体事件须另以单项材料入账。}\textbf{元与高丽、安南及地方政权（年序桥接）。}
作为共享年序的接点，“【共享年序日期坐标】保留1304 年的多政权并行检索入口；本条不主张所列政权在当年发生直接互动，具体事件须另以单项材料入账。”让元与高丽、安南及地方政权（年序桥接）与同年其他地区保持可比的日期位置。可确认的是文书使用这一时间框架；年内何事发生、何处落实，仍须由各自条目另证。

此前的条件是“相邻已入账事件之间缺少该年度的共享年序入口，易使跨政权材料被单一政权叙事遮蔽。”，此后的可见影响是“保留该年度的检索与补账位置；后续仅在获得可核单项材料时拆分为具体政治、财政、军事、社会或文化事件。”。这里的“影响”须按地点、机构和材料种类分别衡量。

这一叙述依据《元史》；《高丽史》；《安南志略》及碑刻、地方文书（据事核），但本条只确认该年位于可追踪的共享年序中；正史、地方文书和外部材料的保存密度不一，不能由桥接标签推断行动、统属、疆界或社会影响。。研究元代区域治理、海域联通与地方社会的并行年序研究可以补问材料未直接说出的空间与社会差异。

\subsection{1305 年（共享年序桥接）}

\noindent\eventanchor{SY-1305-CHRONOLOGY-BRIDGE}{【共享年序日期坐标】保留1305 年的多政权并行检索入口；本条不主张所列政权在当年发生直接互动，具体事件须另以单项材料入账。}\textbf{元与高丽、安南及地方政权（年序桥接）。}
1305 年（共享年序桥接）的元与高丽、安南及地方政权（年序桥接）材料只留下“【共享年序日期坐标】保留1305 年的多政权并行检索入口；本条不主张所列政权在当年发生直接互动，具体事件须另以单项材料入账。”这一纪年接点。它使同年的诏令、交涉或地方材料可以被放回并行时间轴，却不应被误写成一件已经具备完整情节的政治行动。

其前提记为“相邻已入账事件之间缺少该年度的共享年序入口，易使跨政权材料被单一政权叙事遮蔽。”；其后续是“保留该年度的检索与补账位置；后续仅在获得可核单项材料时拆分为具体政治、财政、军事、社会或文化事件。”。这是一条待后续材料检验的事件链，而不是对所有行动者意图的替写。

可核的材料入口是《元史》；《高丽史》；《安南志略》及碑刻、地方文书（据事核）。本条只确认该年位于可追踪的共享年序中；正史、地方文书和外部材料的保存密度不一，不能由桥接标签推断行动、统属、疆界或社会影响。 就此而言，元代区域治理、海域联通与地方社会的并行年序研究有助于把年序同区域、制度或物质材料并读。

\subsection{1306 年（共享年序桥接）}

\noindent\eventanchor{SY-1306-CHRONOLOGY-BRIDGE}{【共享年序日期坐标】保留1306 年的多政权并行检索入口；本条不主张所列政权在当年发生直接互动，具体事件须另以单项材料入账。}\textbf{元与高丽、安南及地方政权（年序桥接）。}
作为共享年序的接点，“【共享年序日期坐标】保留1306 年的多政权并行检索入口；本条不主张所列政权在当年发生直接互动，具体事件须另以单项材料入账。”让元与高丽、安南及地方政权（年序桥接）与同年其他地区保持可比的日期位置。可确认的是文书使用这一时间框架；年内何事发生、何处落实，仍须由各自条目另证。

账本把这一节点放在“相邻已入账事件之间缺少该年度的共享年序入口，易使跨政权材料被单一政权叙事遮蔽。”与“保留该年度的检索与补账位置；后续仅在获得可核单项材料时拆分为具体政治、财政、军事、社会或文化事件。”之间；两者提示前后压力，并不把时间相邻写成单一因果。

本条首先回到《元史》；《高丽史》；《安南志略》及碑刻、地方文书（据事核）；本条只确认该年位于可追踪的共享年序中；正史、地方文书和外部材料的保存密度不一，不能由桥接标签推断行动、统属、疆界或社会影响。 因而元代区域治理、海域联通与地方社会的并行年序研究不是装饰性的书目，而是检验这一叙述边界的路径。

\subsection{1308 年（共享年序桥接）}

\noindent\eventanchor{SY-1308-CHRONOLOGY-BRIDGE}{【共享年序日期坐标】保留1308 年的多政权并行检索入口；本条不主张所列政权在当年发生直接互动，具体事件须另以单项材料入账。}\textbf{元与高丽、安南及地方政权（年序桥接）。}
1308 年（共享年序桥接）的元与高丽、安南及地方政权（年序桥接）材料只留下“【共享年序日期坐标】保留1308 年的多政权并行检索入口；本条不主张所列政权在当年发生直接互动，具体事件须另以单项材料入账。”这一纪年接点。它使同年的诏令、交涉或地方材料可以被放回并行时间轴，却不应被误写成一件已经具备完整情节的政治行动。

此前的条件是“相邻已入账事件之间缺少该年度的共享年序入口，易使跨政权材料被单一政权叙事遮蔽。”，此后的可见影响是“保留该年度的检索与补账位置；后续仅在获得可核单项材料时拆分为具体政治、财政、军事、社会或文化事件。”。这里的“影响”须按地点、机构和材料种类分别衡量。

这一叙述依据《元史》；《高丽史》；《安南志略》及碑刻、地方文书（据事核），但本条只确认该年位于可追踪的共享年序中；正史、地方文书和外部材料的保存密度不一，不能由桥接标签推断行动、统属、疆界或社会影响。。研究元代区域治理、海域联通与地方社会的并行年序研究可以补问材料未直接说出的空间与社会差异。

\subsection{1309 年（共享年序桥接）}

\noindent\eventanchor{SY-1309-CHRONOLOGY-BRIDGE}{【共享年序日期坐标】保留1309 年的多政权并行检索入口；本条不主张所列政权在当年发生直接互动，具体事件须另以单项材料入账。}\textbf{元与高丽、安南及地方政权（年序桥接）。}
作为共享年序的接点，“【共享年序日期坐标】保留1309 年的多政权并行检索入口；本条不主张所列政权在当年发生直接互动，具体事件须另以单项材料入账。”让元与高丽、安南及地方政权（年序桥接）与同年其他地区保持可比的日期位置。可确认的是文书使用这一时间框架；年内何事发生、何处落实，仍须由各自条目另证。

其前提记为“相邻已入账事件之间缺少该年度的共享年序入口，易使跨政权材料被单一政权叙事遮蔽。”；其后续是“保留该年度的检索与补账位置；后续仅在获得可核单项材料时拆分为具体政治、财政、军事、社会或文化事件。”。这是一条待后续材料检验的事件链，而不是对所有行动者意图的替写。

可核的材料入口是《元史》；《高丽史》；《安南志略》及碑刻、地方文书（据事核）。本条只确认该年位于可追踪的共享年序中；正史、地方文书和外部材料的保存密度不一，不能由桥接标签推断行动、统属、疆界或社会影响。 就此而言，元代区域治理、海域联通与地方社会的并行年序研究有助于把年序同区域、制度或物质材料并读。

\subsection{1310 年（共享年序桥接）}

\noindent\eventanchor{SY-1310-CHRONOLOGY-BRIDGE}{【共享年序日期坐标】保留1310 年的多政权并行检索入口；本条不主张所列政权在当年发生直接互动，具体事件须另以单项材料入账。}\textbf{元与高丽、安南及地方政权（年序桥接）。}
作为共享年序的接点，“【共享年序日期坐标】保留1310 年的多政权并行检索入口；本条不主张所列政权在当年发生直接互动，具体事件须另以单项材料入账。”让元与高丽、安南及地方政权（年序桥接）与同年其他地区保持可比的日期位置。可确认的是文书使用这一时间框架；年内何事发生、何处落实，仍须由各自条目另证。

从账本的前后项看，“相邻已入账事件之间缺少该年度的共享年序入口，易使跨政权材料被单一政权叙事遮蔽。”提供了背景压力；“保留该年度的检索与补账位置；后续仅在获得可核单项材料时拆分为具体政治、财政、军事、社会或文化事件。”标出后来需要追踪的变化，二者之间仍存在未被材料填满的环节。

核心线索为《元史》；《高丽史》；《安南志略》及碑刻、地方文书（据事核）。本条只确认该年位于可追踪的共享年序中；正史、地方文书和外部材料的保存密度不一，不能由桥接标签推断行动、统属、疆界或社会影响。 现代研究可从元代区域治理、海域联通与地方社会的并行年序研究继续追问。

\subsection{1312 年（共享年序桥接）}

\noindent\eventanchor{SY-1312-CHRONOLOGY-BRIDGE}{【共享年序日期坐标】保留1312 年的多政权并行检索入口；本条不主张所列政权在当年发生直接互动，具体事件须另以单项材料入账。}\textbf{元与高丽、安南及地方政权（年序桥接）。}
1312 年（共享年序桥接）的元与高丽、安南及地方政权（年序桥接）材料只留下“【共享年序日期坐标】保留1312 年的多政权并行检索入口；本条不主张所列政权在当年发生直接互动，具体事件须另以单项材料入账。”这一纪年接点。它使同年的诏令、交涉或地方材料可以被放回并行时间轴，却不应被误写成一件已经具备完整情节的政治行动。

其前提记为“相邻已入账事件之间缺少该年度的共享年序入口，易使跨政权材料被单一政权叙事遮蔽。”；其后续是“保留该年度的检索与补账位置；后续仅在获得可核单项材料时拆分为具体政治、财政、军事、社会或文化事件。”。这是一条待后续材料检验的事件链，而不是对所有行动者意图的替写。

可核的材料入口是《元史》；《高丽史》；《安南志略》及碑刻、地方文书（据事核）。本条只确认该年位于可追踪的共享年序中；正史、地方文书和外部材料的保存密度不一，不能由桥接标签推断行动、统属、疆界或社会影响。 就此而言，元代区域治理、海域联通与地方社会的并行年序研究有助于把年序同区域、制度或物质材料并读。

\subsection{1314 年（共享年序桥接）}

\noindent\eventanchor{SY-1314-CHRONOLOGY-BRIDGE}{【共享年序日期坐标】保留1314 年的多政权并行检索入口；本条不主张所列政权在当年发生直接互动，具体事件须另以单项材料入账。}\textbf{元与高丽、安南及地方政权（年序桥接）。}
元与高丽、安南及地方政权（年序桥接）在1314 年（共享年序桥接）的这条年序记录提供的是检索坐标：“【共享年序日期坐标】保留1314 年的多政权并行检索入口；本条不主张所列政权在当年发生直接互动，具体事件须另以单项材料入账。”。它保留材料稀薄处的空白，而不以一条王朝纪年替代地方、边疆或竞争政权的行动。

从账本的前后项看，“相邻已入账事件之间缺少该年度的共享年序入口，易使跨政权材料被单一政权叙事遮蔽。”提供了背景压力；“保留该年度的检索与补账位置；后续仅在获得可核单项材料时拆分为具体政治、财政、军事、社会或文化事件。”标出后来需要追踪的变化，二者之间仍存在未被材料填满的环节。

核心线索为《元史》；《高丽史》；《安南志略》及碑刻、地方文书（据事核）。本条只确认该年位于可追踪的共享年序中；正史、地方文书和外部材料的保存密度不一，不能由桥接标签推断行动、统属、疆界或社会影响。 现代研究可从元代区域治理、海域联通与地方社会的并行年序研究继续追问。

\subsection{1315 年（共享年序桥接）}

\noindent\eventanchor{SY-1315-CHRONOLOGY-BRIDGE}{【共享年序日期坐标】保留1315 年的多政权并行检索入口；本条不主张所列政权在当年发生直接互动，具体事件须另以单项材料入账。}\textbf{元与高丽、安南及地方政权（年序桥接）。}
作为共享年序的接点，“【共享年序日期坐标】保留1315 年的多政权并行检索入口；本条不主张所列政权在当年发生直接互动，具体事件须另以单项材料入账。”让元与高丽、安南及地方政权（年序桥接）与同年其他地区保持可比的日期位置。可确认的是文书使用这一时间框架；年内何事发生、何处落实，仍须由各自条目另证。

此前的条件是“相邻已入账事件之间缺少该年度的共享年序入口，易使跨政权材料被单一政权叙事遮蔽。”，此后的可见影响是“保留该年度的检索与补账位置；后续仅在获得可核单项材料时拆分为具体政治、财政、军事、社会或文化事件。”。这里的“影响”须按地点、机构和材料种类分别衡量。

这一叙述依据《元史》；《高丽史》；《安南志略》及碑刻、地方文书（据事核），但本条只确认该年位于可追踪的共享年序中；正史、地方文书和外部材料的保存密度不一，不能由桥接标签推断行动、统属、疆界或社会影响。。研究元代区域治理、海域联通与地方社会的并行年序研究可以补问材料未直接说出的空间与社会差异。

\subsection{1316 年（共享年序桥接）}

\noindent\eventanchor{SY-1316-CHRONOLOGY-BRIDGE}{【共享年序日期坐标】保留1316 年的多政权并行检索入口；本条不主张所列政权在当年发生直接互动，具体事件须另以单项材料入账。}\textbf{元与高丽、安南及地方政权（年序桥接）。}
元与高丽、安南及地方政权（年序桥接）在1316 年（共享年序桥接）的这条年序记录提供的是检索坐标：“【共享年序日期坐标】保留1316 年的多政权并行检索入口；本条不主张所列政权在当年发生直接互动，具体事件须另以单项材料入账。”。它保留材料稀薄处的空白，而不以一条王朝纪年替代地方、边疆或竞争政权的行动。

其前提记为“相邻已入账事件之间缺少该年度的共享年序入口，易使跨政权材料被单一政权叙事遮蔽。”；其后续是“保留该年度的检索与补账位置；后续仅在获得可核单项材料时拆分为具体政治、财政、军事、社会或文化事件。”。这是一条待后续材料检验的事件链，而不是对所有行动者意图的替写。

可核的材料入口是《元史》；《高丽史》；《安南志略》及碑刻、地方文书（据事核）。本条只确认该年位于可追踪的共享年序中；正史、地方文书和外部材料的保存密度不一，不能由桥接标签推断行动、统属、疆界或社会影响。 就此而言，元代区域治理、海域联通与地方社会的并行年序研究有助于把年序同区域、制度或物质材料并读。

\subsection{1317 年（共享年序桥接）}

\noindent\eventanchor{SY-1317-CHRONOLOGY-BRIDGE}{【共享年序日期坐标】保留1317 年的多政权并行检索入口；本条不主张所列政权在当年发生直接互动，具体事件须另以单项材料入账。}\textbf{元与高丽、安南及地方政权（年序桥接）。}
1317 年（共享年序桥接）的元与高丽、安南及地方政权（年序桥接）材料只留下“【共享年序日期坐标】保留1317 年的多政权并行检索入口；本条不主张所列政权在当年发生直接互动，具体事件须另以单项材料入账。”这一纪年接点。它使同年的诏令、交涉或地方材料可以被放回并行时间轴，却不应被误写成一件已经具备完整情节的政治行动。

账本把这一节点放在“相邻已入账事件之间缺少该年度的共享年序入口，易使跨政权材料被单一政权叙事遮蔽。”与“保留该年度的检索与补账位置；后续仅在获得可核单项材料时拆分为具体政治、财政、军事、社会或文化事件。”之间；两者提示前后压力，并不把时间相邻写成单一因果。

本条首先回到《元史》；《高丽史》；《安南志略》及碑刻、地方文书（据事核）；本条只确认该年位于可追踪的共享年序中；正史、地方文书和外部材料的保存密度不一，不能由桥接标签推断行动、统属、疆界或社会影响。 因而元代区域治理、海域联通与地方社会的并行年序研究不是装饰性的书目，而是检验这一叙述边界的路径。

\subsection{1318 年（共享年序桥接）}

\noindent\eventanchor{SY-1318-CHRONOLOGY-BRIDGE}{【共享年序日期坐标】保留1318 年的多政权并行检索入口；本条不主张所列政权在当年发生直接互动，具体事件须另以单项材料入账。}\textbf{元与高丽、安南及地方政权（年序桥接）。}
作为共享年序的接点，“【共享年序日期坐标】保留1318 年的多政权并行检索入口；本条不主张所列政权在当年发生直接互动，具体事件须另以单项材料入账。”让元与高丽、安南及地方政权（年序桥接）与同年其他地区保持可比的日期位置。可确认的是文书使用这一时间框架；年内何事发生、何处落实，仍须由各自条目另证。

从账本的前后项看，“相邻已入账事件之间缺少该年度的共享年序入口，易使跨政权材料被单一政权叙事遮蔽。”提供了背景压力；“保留该年度的检索与补账位置；后续仅在获得可核单项材料时拆分为具体政治、财政、军事、社会或文化事件。”标出后来需要追踪的变化，二者之间仍存在未被材料填满的环节。

核心线索为《元史》；《高丽史》；《安南志略》及碑刻、地方文书（据事核）。本条只确认该年位于可追踪的共享年序中；正史、地方文书和外部材料的保存密度不一，不能由桥接标签推断行动、统属、疆界或社会影响。 现代研究可从元代区域治理、海域联通与地方社会的并行年序研究继续追问。

\subsection{1319 年（共享年序桥接）}

\noindent\eventanchor{SY-1319-CHRONOLOGY-BRIDGE}{【共享年序日期坐标】保留1319 年的多政权并行检索入口；本条不主张所列政权在当年发生直接互动，具体事件须另以单项材料入账。}\textbf{元与高丽、安南及地方政权（年序桥接）。}
元与高丽、安南及地方政权（年序桥接）在1319 年（共享年序桥接）的这条年序记录提供的是检索坐标：“【共享年序日期坐标】保留1319 年的多政权并行检索入口；本条不主张所列政权在当年发生直接互动，具体事件须另以单项材料入账。”。它保留材料稀薄处的空白，而不以一条王朝纪年替代地方、边疆或竞争政权的行动。

此前的条件是“相邻已入账事件之间缺少该年度的共享年序入口，易使跨政权材料被单一政权叙事遮蔽。”，此后的可见影响是“保留该年度的检索与补账位置；后续仅在获得可核单项材料时拆分为具体政治、财政、军事、社会或文化事件。”。这里的“影响”须按地点、机构和材料种类分别衡量。

这一叙述依据《元史》；《高丽史》；《安南志略》及碑刻、地方文书（据事核），但本条只确认该年位于可追踪的共享年序中；正史、地方文书和外部材料的保存密度不一，不能由桥接标签推断行动、统属、疆界或社会影响。。研究元代区域治理、海域联通与地方社会的并行年序研究可以补问材料未直接说出的空间与社会差异。

\subsection{1321 年（共享年序桥接）}

\noindent\eventanchor{SY-1321-CHRONOLOGY-BRIDGE}{【共享年序日期坐标】保留1321 年的多政权并行检索入口；本条不主张所列政权在当年发生直接互动，具体事件须另以单项材料入账。}\textbf{元与高丽、安南及地方政权（年序桥接）。}
1321 年（共享年序桥接）的元与高丽、安南及地方政权（年序桥接）材料只留下“【共享年序日期坐标】保留1321 年的多政权并行检索入口；本条不主张所列政权在当年发生直接互动，具体事件须另以单项材料入账。”这一纪年接点。它使同年的诏令、交涉或地方材料可以被放回并行时间轴，却不应被误写成一件已经具备完整情节的政治行动。

从账本的前后项看，“相邻已入账事件之间缺少该年度的共享年序入口，易使跨政权材料被单一政权叙事遮蔽。”提供了背景压力；“保留该年度的检索与补账位置；后续仅在获得可核单项材料时拆分为具体政治、财政、军事、社会或文化事件。”标出后来需要追踪的变化，二者之间仍存在未被材料填满的环节。

核心线索为《元史》；《高丽史》；《安南志略》及碑刻、地方文书（据事核）。本条只确认该年位于可追踪的共享年序中；正史、地方文书和外部材料的保存密度不一，不能由桥接标签推断行动、统属、疆界或社会影响。 现代研究可从元代区域治理、海域联通与地方社会的并行年序研究继续追问。

\subsection{1322 年（共享年序桥接）}

\noindent\eventanchor{SY-1322-CHRONOLOGY-BRIDGE}{【共享年序日期坐标】保留1322 年的多政权并行检索入口；本条不主张所列政权在当年发生直接互动，具体事件须另以单项材料入账。}\textbf{元与高丽、安南及地方政权（年序桥接）。}
作为共享年序的接点，“【共享年序日期坐标】保留1322 年的多政权并行检索入口；本条不主张所列政权在当年发生直接互动，具体事件须另以单项材料入账。”让元与高丽、安南及地方政权（年序桥接）与同年其他地区保持可比的日期位置。可确认的是文书使用这一时间框架；年内何事发生、何处落实，仍须由各自条目另证。

此前的条件是“相邻已入账事件之间缺少该年度的共享年序入口，易使跨政权材料被单一政权叙事遮蔽。”，此后的可见影响是“保留该年度的检索与补账位置；后续仅在获得可核单项材料时拆分为具体政治、财政、军事、社会或文化事件。”。这里的“影响”须按地点、机构和材料种类分别衡量。

这一叙述依据《元史》；《高丽史》；《安南志略》及碑刻、地方文书（据事核），但本条只确认该年位于可追踪的共享年序中；正史、地方文书和外部材料的保存密度不一，不能由桥接标签推断行动、统属、疆界或社会影响。。研究元代区域治理、海域联通与地方社会的并行年序研究可以补问材料未直接说出的空间与社会差异。

\subsection{1324 年（共享年序桥接）}

\noindent\eventanchor{SY-1324-CHRONOLOGY-BRIDGE}{【共享年序日期坐标】保留1324 年的多政权并行检索入口；本条不主张所列政权在当年发生直接互动，具体事件须另以单项材料入账。}\textbf{元与高丽、安南及地方政权（年序桥接）。}
1324 年（共享年序桥接）的元与高丽、安南及地方政权（年序桥接）材料只留下“【共享年序日期坐标】保留1324 年的多政权并行检索入口；本条不主张所列政权在当年发生直接互动，具体事件须另以单项材料入账。”这一纪年接点。它使同年的诏令、交涉或地方材料可以被放回并行时间轴，却不应被误写成一件已经具备完整情节的政治行动。

账本把这一节点放在“相邻已入账事件之间缺少该年度的共享年序入口，易使跨政权材料被单一政权叙事遮蔽。”与“保留该年度的检索与补账位置；后续仅在获得可核单项材料时拆分为具体政治、财政、军事、社会或文化事件。”之间；两者提示前后压力，并不把时间相邻写成单一因果。

本条首先回到《元史》；《高丽史》；《安南志略》及碑刻、地方文书（据事核）；本条只确认该年位于可追踪的共享年序中；正史、地方文书和外部材料的保存密度不一，不能由桥接标签推断行动、统属、疆界或社会影响。 因而元代区域治理、海域联通与地方社会的并行年序研究不是装饰性的书目，而是检验这一叙述边界的路径。

\subsection{1325 年（共享年序桥接）}

\noindent\eventanchor{SY-1325-CHRONOLOGY-BRIDGE}{【共享年序日期坐标】保留1325 年的多政权并行检索入口；本条不主张所列政权在当年发生直接互动，具体事件须另以单项材料入账。}\textbf{元与高丽、安南及地方政权（年序桥接）。}
元与高丽、安南及地方政权（年序桥接）在1325 年（共享年序桥接）的这条年序记录提供的是检索坐标：“【共享年序日期坐标】保留1325 年的多政权并行检索入口；本条不主张所列政权在当年发生直接互动，具体事件须另以单项材料入账。”。它保留材料稀薄处的空白，而不以一条王朝纪年替代地方、边疆或竞争政权的行动。

从账本的前后项看，“相邻已入账事件之间缺少该年度的共享年序入口，易使跨政权材料被单一政权叙事遮蔽。”提供了背景压力；“保留该年度的检索与补账位置；后续仅在获得可核单项材料时拆分为具体政治、财政、军事、社会或文化事件。”标出后来需要追踪的变化，二者之间仍存在未被材料填满的环节。

核心线索为《元史》；《高丽史》；《安南志略》及碑刻、地方文书（据事核）。本条只确认该年位于可追踪的共享年序中；正史、地方文书和外部材料的保存密度不一，不能由桥接标签推断行动、统属、疆界或社会影响。 现代研究可从元代区域治理、海域联通与地方社会的并行年序研究继续追问。

\subsection{1326 年（共享年序桥接）}

\noindent\eventanchor{SY-1326-CHRONOLOGY-BRIDGE}{【共享年序日期坐标】保留1326 年的多政权并行检索入口；本条不主张所列政权在当年发生直接互动，具体事件须另以单项材料入账。}\textbf{元与高丽、安南及地方政权（年序桥接）。}
1326 年（共享年序桥接）的元与高丽、安南及地方政权（年序桥接）材料只留下“【共享年序日期坐标】保留1326 年的多政权并行检索入口；本条不主张所列政权在当年发生直接互动，具体事件须另以单项材料入账。”这一纪年接点。它使同年的诏令、交涉或地方材料可以被放回并行时间轴，却不应被误写成一件已经具备完整情节的政治行动。

此前的条件是“相邻已入账事件之间缺少该年度的共享年序入口，易使跨政权材料被单一政权叙事遮蔽。”，此后的可见影响是“保留该年度的检索与补账位置；后续仅在获得可核单项材料时拆分为具体政治、财政、军事、社会或文化事件。”。这里的“影响”须按地点、机构和材料种类分别衡量。

这一叙述依据《元史》；《高丽史》；《安南志略》及碑刻、地方文书（据事核），但本条只确认该年位于可追踪的共享年序中；正史、地方文书和外部材料的保存密度不一，不能由桥接标签推断行动、统属、疆界或社会影响。。研究元代区域治理、海域联通与地方社会的并行年序研究可以补问材料未直接说出的空间与社会差异。

\subsection{1327 年（共享年序桥接）}

\noindent\eventanchor{SY-1327-CHRONOLOGY-BRIDGE}{【共享年序日期坐标】保留1327 年的多政权并行检索入口；本条不主张所列政权在当年发生直接互动，具体事件须另以单项材料入账。}\textbf{元与高丽、安南及地方政权（年序桥接）。}
作为共享年序的接点，“【共享年序日期坐标】保留1327 年的多政权并行检索入口；本条不主张所列政权在当年发生直接互动，具体事件须另以单项材料入账。”让元与高丽、安南及地方政权（年序桥接）与同年其他地区保持可比的日期位置。可确认的是文书使用这一时间框架；年内何事发生、何处落实，仍须由各自条目另证。

其前提记为“相邻已入账事件之间缺少该年度的共享年序入口，易使跨政权材料被单一政权叙事遮蔽。”；其后续是“保留该年度的检索与补账位置；后续仅在获得可核单项材料时拆分为具体政治、财政、军事、社会或文化事件。”。这是一条待后续材料检验的事件链，而不是对所有行动者意图的替写。

可核的材料入口是《元史》；《高丽史》；《安南志略》及碑刻、地方文书（据事核）。本条只确认该年位于可追踪的共享年序中；正史、地方文书和外部材料的保存密度不一，不能由桥接标签推断行动、统属、疆界或社会影响。 就此而言，元代区域治理、海域联通与地方社会的并行年序研究有助于把年序同区域、制度或物质材料并读。

\subsection{1329 年（共享年序桥接）}

\noindent\eventanchor{SY-1329-CHRONOLOGY-BRIDGE}{【共享年序日期坐标】保留1329 年的多政权并行检索入口；本条不主张所列政权在当年发生直接互动，具体事件须另以单项材料入账。}\textbf{元与高丽、安南及地方政权（年序桥接）。}
1329 年（共享年序桥接）的元与高丽、安南及地方政权（年序桥接）材料只留下“【共享年序日期坐标】保留1329 年的多政权并行检索入口；本条不主张所列政权在当年发生直接互动，具体事件须另以单项材料入账。”这一纪年接点。它使同年的诏令、交涉或地方材料可以被放回并行时间轴，却不应被误写成一件已经具备完整情节的政治行动。

从账本的前后项看，“相邻已入账事件之间缺少该年度的共享年序入口，易使跨政权材料被单一政权叙事遮蔽。”提供了背景压力；“保留该年度的检索与补账位置；后续仅在获得可核单项材料时拆分为具体政治、财政、军事、社会或文化事件。”标出后来需要追踪的变化，二者之间仍存在未被材料填满的环节。

核心线索为《元史》；《高丽史》；《安南志略》及碑刻、地方文书（据事核）。本条只确认该年位于可追踪的共享年序中；正史、地方文书和外部材料的保存密度不一，不能由桥接标签推断行动、统属、疆界或社会影响。 现代研究可从元代区域治理、海域联通与地方社会的并行年序研究继续追问。

\subsection{1330 年（共享年序桥接）}

\noindent\eventanchor{SY-1330-CHRONOLOGY-BRIDGE}{【共享年序日期坐标】保留1330 年的多政权并行检索入口；本条不主张所列政权在当年发生直接互动，具体事件须另以单项材料入账。}\textbf{元与高丽、安南及地方政权（年序桥接）。}
1330 年（共享年序桥接）的元与高丽、安南及地方政权（年序桥接）材料只留下“【共享年序日期坐标】保留1330 年的多政权并行检索入口；本条不主张所列政权在当年发生直接互动，具体事件须另以单项材料入账。”这一纪年接点。它使同年的诏令、交涉或地方材料可以被放回并行时间轴，却不应被误写成一件已经具备完整情节的政治行动。

其前提记为“相邻已入账事件之间缺少该年度的共享年序入口，易使跨政权材料被单一政权叙事遮蔽。”；其后续是“保留该年度的检索与补账位置；后续仅在获得可核单项材料时拆分为具体政治、财政、军事、社会或文化事件。”。这是一条待后续材料检验的事件链，而不是对所有行动者意图的替写。

可核的材料入口是《元史》；《高丽史》；《安南志略》及碑刻、地方文书（据事核）。本条只确认该年位于可追踪的共享年序中；正史、地方文书和外部材料的保存密度不一，不能由桥接标签推断行动、统属、疆界或社会影响。 就此而言，元代区域治理、海域联通与地方社会的并行年序研究有助于把年序同区域、制度或物质材料并读。

\subsection{1331 年（共享年序桥接）}

\noindent\eventanchor{SY-1331-CHRONOLOGY-BRIDGE}{【共享年序日期坐标】保留1331 年的多政权并行检索入口；本条不主张所列政权在当年发生直接互动，具体事件须另以单项材料入账。}\textbf{元与高丽、安南及地方政权（年序桥接）。}
作为共享年序的接点，“【共享年序日期坐标】保留1331 年的多政权并行检索入口；本条不主张所列政权在当年发生直接互动，具体事件须另以单项材料入账。”让元与高丽、安南及地方政权（年序桥接）与同年其他地区保持可比的日期位置。可确认的是文书使用这一时间框架；年内何事发生、何处落实，仍须由各自条目另证。

账本把这一节点放在“相邻已入账事件之间缺少该年度的共享年序入口，易使跨政权材料被单一政权叙事遮蔽。”与“保留该年度的检索与补账位置；后续仅在获得可核单项材料时拆分为具体政治、财政、军事、社会或文化事件。”之间；两者提示前后压力，并不把时间相邻写成单一因果。

本条首先回到《元史》；《高丽史》；《安南志略》及碑刻、地方文书（据事核）；本条只确认该年位于可追踪的共享年序中；正史、地方文书和外部材料的保存密度不一，不能由桥接标签推断行动、统属、疆界或社会影响。 因而元代区域治理、海域联通与地方社会的并行年序研究不是装饰性的书目，而是检验这一叙述边界的路径。

\subsection{1332 年（共享年序桥接）}

\noindent\eventanchor{SY-1332-CHRONOLOGY-BRIDGE}{【共享年序日期坐标】保留1332 年的多政权并行检索入口；本条不主张所列政权在当年发生直接互动，具体事件须另以单项材料入账。}\textbf{元与高丽、安南及地方政权（年序桥接）。}
元与高丽、安南及地方政权（年序桥接）在1332 年（共享年序桥接）的这条年序记录提供的是检索坐标：“【共享年序日期坐标】保留1332 年的多政权并行检索入口；本条不主张所列政权在当年发生直接互动，具体事件须另以单项材料入账。”。它保留材料稀薄处的空白，而不以一条王朝纪年替代地方、边疆或竞争政权的行动。

从账本的前后项看，“相邻已入账事件之间缺少该年度的共享年序入口，易使跨政权材料被单一政权叙事遮蔽。”提供了背景压力；“保留该年度的检索与补账位置；后续仅在获得可核单项材料时拆分为具体政治、财政、军事、社会或文化事件。”标出后来需要追踪的变化，二者之间仍存在未被材料填满的环节。

核心线索为《元史》；《高丽史》；《安南志略》及碑刻、地方文书（据事核）。本条只确认该年位于可追踪的共享年序中；正史、地方文书和外部材料的保存密度不一，不能由桥接标签推断行动、统属、疆界或社会影响。 现代研究可从元代区域治理、海域联通与地方社会的并行年序研究继续追问。

\subsection{1334 年（共享年序桥接）}

\noindent\eventanchor{SY-1334-CHRONOLOGY-BRIDGE}{【共享年序日期坐标】保留1334 年的多政权并行检索入口；本条不主张所列政权在当年发生直接互动，具体事件须另以单项材料入账。}\textbf{元与高丽、安南及地方政权（年序桥接）。}
元与高丽、安南及地方政权（年序桥接）在1334 年（共享年序桥接）的这条年序记录提供的是检索坐标：“【共享年序日期坐标】保留1334 年的多政权并行检索入口；本条不主张所列政权在当年发生直接互动，具体事件须另以单项材料入账。”。它保留材料稀薄处的空白，而不以一条王朝纪年替代地方、边疆或竞争政权的行动。

其前提记为“相邻已入账事件之间缺少该年度的共享年序入口，易使跨政权材料被单一政权叙事遮蔽。”；其后续是“保留该年度的检索与补账位置；后续仅在获得可核单项材料时拆分为具体政治、财政、军事、社会或文化事件。”。这是一条待后续材料检验的事件链，而不是对所有行动者意图的替写。

可核的材料入口是《元史》；《高丽史》；《安南志略》及碑刻、地方文书（据事核）。本条只确认该年位于可追踪的共享年序中；正史、地方文书和外部材料的保存密度不一，不能由桥接标签推断行动、统属、疆界或社会影响。 就此而言，元代区域治理、海域联通与地方社会的并行年序研究有助于把年序同区域、制度或物质材料并读。

\subsection{1335 年（共享年序桥接）}

\noindent\eventanchor{SY-1335-CHRONOLOGY-BRIDGE}{【共享年序日期坐标】保留1335 年的多政权并行检索入口；本条不主张所列政权在当年发生直接互动，具体事件须另以单项材料入账。}\textbf{元与高丽、安南及地方政权（年序桥接）。}
1335 年（共享年序桥接）的元与高丽、安南及地方政权（年序桥接）材料只留下“【共享年序日期坐标】保留1335 年的多政权并行检索入口；本条不主张所列政权在当年发生直接互动，具体事件须另以单项材料入账。”这一纪年接点。它使同年的诏令、交涉或地方材料可以被放回并行时间轴，却不应被误写成一件已经具备完整情节的政治行动。

账本把这一节点放在“相邻已入账事件之间缺少该年度的共享年序入口，易使跨政权材料被单一政权叙事遮蔽。”与“保留该年度的检索与补账位置；后续仅在获得可核单项材料时拆分为具体政治、财政、军事、社会或文化事件。”之间；两者提示前后压力，并不把时间相邻写成单一因果。

本条首先回到《元史》；《高丽史》；《安南志略》及碑刻、地方文书（据事核）；本条只确认该年位于可追踪的共享年序中；正史、地方文书和外部材料的保存密度不一，不能由桥接标签推断行动、统属、疆界或社会影响。 因而元代区域治理、海域联通与地方社会的并行年序研究不是装饰性的书目，而是检验这一叙述边界的路径。

\subsection{1336 年（共享年序桥接）}

\noindent\eventanchor{SY-1336-CHRONOLOGY-BRIDGE}{【共享年序日期坐标】保留1336 年的多政权并行检索入口；本条不主张所列政权在当年发生直接互动，具体事件须另以单项材料入账。}\textbf{元与高丽、安南及地方政权（年序桥接）。}
作为共享年序的接点，“【共享年序日期坐标】保留1336 年的多政权并行检索入口；本条不主张所列政权在当年发生直接互动，具体事件须另以单项材料入账。”让元与高丽、安南及地方政权（年序桥接）与同年其他地区保持可比的日期位置。可确认的是文书使用这一时间框架；年内何事发生、何处落实，仍须由各自条目另证。

从账本的前后项看，“相邻已入账事件之间缺少该年度的共享年序入口，易使跨政权材料被单一政权叙事遮蔽。”提供了背景压力；“保留该年度的检索与补账位置；后续仅在获得可核单项材料时拆分为具体政治、财政、军事、社会或文化事件。”标出后来需要追踪的变化，二者之间仍存在未被材料填满的环节。

核心线索为《元史》；《高丽史》；《安南志略》及碑刻、地方文书（据事核）。本条只确认该年位于可追踪的共享年序中；正史、地方文书和外部材料的保存密度不一，不能由桥接标签推断行动、统属、疆界或社会影响。 现代研究可从元代区域治理、海域联通与地方社会的并行年序研究继续追问。

\subsection{1337 年（共享年序桥接）}

\noindent\eventanchor{SY-1337-CHRONOLOGY-BRIDGE}{【共享年序日期坐标】保留1337 年的多政权并行检索入口；本条不主张所列政权在当年发生直接互动，具体事件须另以单项材料入账。}\textbf{元与高丽、安南及地方政权（年序桥接）。}
元与高丽、安南及地方政权（年序桥接）在1337 年（共享年序桥接）的这条年序记录提供的是检索坐标：“【共享年序日期坐标】保留1337 年的多政权并行检索入口；本条不主张所列政权在当年发生直接互动，具体事件须另以单项材料入账。”。它保留材料稀薄处的空白，而不以一条王朝纪年替代地方、边疆或竞争政权的行动。

此前的条件是“相邻已入账事件之间缺少该年度的共享年序入口，易使跨政权材料被单一政权叙事遮蔽。”，此后的可见影响是“保留该年度的检索与补账位置；后续仅在获得可核单项材料时拆分为具体政治、财政、军事、社会或文化事件。”。这里的“影响”须按地点、机构和材料种类分别衡量。

这一叙述依据《元史》；《高丽史》；《安南志略》及碑刻、地方文书（据事核），但本条只确认该年位于可追踪的共享年序中；正史、地方文书和外部材料的保存密度不一，不能由桥接标签推断行动、统属、疆界或社会影响。。研究元代区域治理、海域联通与地方社会的并行年序研究可以补问材料未直接说出的空间与社会差异。

\subsection{1338 年（共享年序桥接）}

\noindent\eventanchor{SY-1338-CHRONOLOGY-BRIDGE}{【共享年序日期坐标】保留1338 年的多政权并行检索入口；本条不主张所列政权在当年发生直接互动，具体事件须另以单项材料入账。}\textbf{元与高丽、安南及地方政权（年序桥接）。}
1338 年（共享年序桥接）的元与高丽、安南及地方政权（年序桥接）材料只留下“【共享年序日期坐标】保留1338 年的多政权并行检索入口；本条不主张所列政权在当年发生直接互动，具体事件须另以单项材料入账。”这一纪年接点。它使同年的诏令、交涉或地方材料可以被放回并行时间轴，却不应被误写成一件已经具备完整情节的政治行动。

其前提记为“相邻已入账事件之间缺少该年度的共享年序入口，易使跨政权材料被单一政权叙事遮蔽。”；其后续是“保留该年度的检索与补账位置；后续仅在获得可核单项材料时拆分为具体政治、财政、军事、社会或文化事件。”。这是一条待后续材料检验的事件链，而不是对所有行动者意图的替写。

可核的材料入口是《元史》；《高丽史》；《安南志略》及碑刻、地方文书（据事核）。本条只确认该年位于可追踪的共享年序中；正史、地方文书和外部材料的保存密度不一，不能由桥接标签推断行动、统属、疆界或社会影响。 就此而言，元代区域治理、海域联通与地方社会的并行年序研究有助于把年序同区域、制度或物质材料并读。

\subsection{1339 年（共享年序桥接）}

\noindent\eventanchor{SY-1339-CHRONOLOGY-BRIDGE}{【共享年序日期坐标】保留1339 年的多政权并行检索入口；本条不主张所列政权在当年发生直接互动，具体事件须另以单项材料入账。}\textbf{元与高丽、安南及地方政权（年序桥接）。}
元与高丽、安南及地方政权（年序桥接）在1339 年（共享年序桥接）的这条年序记录提供的是检索坐标：“【共享年序日期坐标】保留1339 年的多政权并行检索入口；本条不主张所列政权在当年发生直接互动，具体事件须另以单项材料入账。”。它保留材料稀薄处的空白，而不以一条王朝纪年替代地方、边疆或竞争政权的行动。

账本把这一节点放在“相邻已入账事件之间缺少该年度的共享年序入口，易使跨政权材料被单一政权叙事遮蔽。”与“保留该年度的检索与补账位置；后续仅在获得可核单项材料时拆分为具体政治、财政、军事、社会或文化事件。”之间；两者提示前后压力，并不把时间相邻写成单一因果。

本条首先回到《元史》；《高丽史》；《安南志略》及碑刻、地方文书（据事核）；本条只确认该年位于可追踪的共享年序中；正史、地方文书和外部材料的保存密度不一，不能由桥接标签推断行动、统属、疆界或社会影响。 因而元代区域治理、海域联通与地方社会的并行年序研究不是装饰性的书目，而是检验这一叙述边界的路径。

\subsection{1341 年（共享年序桥接）}

\noindent\eventanchor{SY-1341-CHRONOLOGY-BRIDGE}{【共享年序日期坐标】保留1341 年的多政权并行检索入口；本条不主张所列政权在当年发生直接互动，具体事件须另以单项材料入账。}\textbf{元与高丽、安南及地方政权（年序桥接）。}
元与高丽、安南及地方政权（年序桥接）在1341 年（共享年序桥接）的这条年序记录提供的是检索坐标：“【共享年序日期坐标】保留1341 年的多政权并行检索入口；本条不主张所列政权在当年发生直接互动，具体事件须另以单项材料入账。”。它保留材料稀薄处的空白，而不以一条王朝纪年替代地方、边疆或竞争政权的行动。

其前提记为“相邻已入账事件之间缺少该年度的共享年序入口，易使跨政权材料被单一政权叙事遮蔽。”；其后续是“保留该年度的检索与补账位置；后续仅在获得可核单项材料时拆分为具体政治、财政、军事、社会或文化事件。”。这是一条待后续材料检验的事件链，而不是对所有行动者意图的替写。

可核的材料入口是《元史》；《高丽史》；《安南志略》及碑刻、地方文书（据事核）。本条只确认该年位于可追踪的共享年序中；正史、地方文书和外部材料的保存密度不一，不能由桥接标签推断行动、统属、疆界或社会影响。 就此而言，元代区域治理、海域联通与地方社会的并行年序研究有助于把年序同区域、制度或物质材料并读。

\subsection{1342 年（共享年序桥接）}

\noindent\eventanchor{SY-1342-CHRONOLOGY-BRIDGE}{【共享年序日期坐标】保留1342 年的多政权并行检索入口；本条不主张所列政权在当年发生直接互动，具体事件须另以单项材料入账。}\textbf{元与高丽、安南及地方政权（年序桥接）。}
1342 年（共享年序桥接）的元与高丽、安南及地方政权（年序桥接）材料只留下“【共享年序日期坐标】保留1342 年的多政权并行检索入口；本条不主张所列政权在当年发生直接互动，具体事件须另以单项材料入账。”这一纪年接点。它使同年的诏令、交涉或地方材料可以被放回并行时间轴，却不应被误写成一件已经具备完整情节的政治行动。

账本把这一节点放在“相邻已入账事件之间缺少该年度的共享年序入口，易使跨政权材料被单一政权叙事遮蔽。”与“保留该年度的检索与补账位置；后续仅在获得可核单项材料时拆分为具体政治、财政、军事、社会或文化事件。”之间；两者提示前后压力，并不把时间相邻写成单一因果。

本条首先回到《元史》；《高丽史》；《安南志略》及碑刻、地方文书（据事核）；本条只确认该年位于可追踪的共享年序中；正史、地方文书和外部材料的保存密度不一，不能由桥接标签推断行动、统属、疆界或社会影响。 因而元代区域治理、海域联通与地方社会的并行年序研究不是装饰性的书目，而是检验这一叙述边界的路径。

\subsection{1343 年（共享年序桥接）}

\noindent\eventanchor{SY-1343-CHRONOLOGY-BRIDGE}{【共享年序日期坐标】保留1343 年的多政权并行检索入口；本条不主张所列政权在当年发生直接互动，具体事件须另以单项材料入账。}\textbf{元与高丽、安南及地方政权（年序桥接）。}
作为共享年序的接点，“【共享年序日期坐标】保留1343 年的多政权并行检索入口；本条不主张所列政权在当年发生直接互动，具体事件须另以单项材料入账。”让元与高丽、安南及地方政权（年序桥接）与同年其他地区保持可比的日期位置。可确认的是文书使用这一时间框架；年内何事发生、何处落实，仍须由各自条目另证。

从账本的前后项看，“相邻已入账事件之间缺少该年度的共享年序入口，易使跨政权材料被单一政权叙事遮蔽。”提供了背景压力；“保留该年度的检索与补账位置；后续仅在获得可核单项材料时拆分为具体政治、财政、军事、社会或文化事件。”标出后来需要追踪的变化，二者之间仍存在未被材料填满的环节。

核心线索为《元史》；《高丽史》；《安南志略》及碑刻、地方文书（据事核）。本条只确认该年位于可追踪的共享年序中；正史、地方文书和外部材料的保存密度不一，不能由桥接标签推断行动、统属、疆界或社会影响。 现代研究可从元代区域治理、海域联通与地方社会的并行年序研究继续追问。

\subsection{1345 年（共享年序桥接）}

\noindent\eventanchor{SY-1345-CHRONOLOGY-BRIDGE}{【共享年序日期坐标】保留1345 年的多政权并行检索入口；本条不主张所列政权在当年发生直接互动，具体事件须另以单项材料入账。}\textbf{元与高丽、安南及地方政权（年序桥接）。}
作为共享年序的接点，“【共享年序日期坐标】保留1345 年的多政权并行检索入口；本条不主张所列政权在当年发生直接互动，具体事件须另以单项材料入账。”让元与高丽、安南及地方政权（年序桥接）与同年其他地区保持可比的日期位置。可确认的是文书使用这一时间框架；年内何事发生、何处落实，仍须由各自条目另证。

其前提记为“相邻已入账事件之间缺少该年度的共享年序入口，易使跨政权材料被单一政权叙事遮蔽。”；其后续是“保留该年度的检索与补账位置；后续仅在获得可核单项材料时拆分为具体政治、财政、军事、社会或文化事件。”。这是一条待后续材料检验的事件链，而不是对所有行动者意图的替写。

可核的材料入口是《元史》；《高丽史》；《安南志略》及碑刻、地方文书（据事核）。本条只确认该年位于可追踪的共享年序中；正史、地方文书和外部材料的保存密度不一，不能由桥接标签推断行动、统属、疆界或社会影响。 就此而言，元代区域治理、海域联通与地方社会的并行年序研究有助于把年序同区域、制度或物质材料并读。

\subsection{1346 年（共享年序桥接）}

\noindent\eventanchor{SY-1346-CHRONOLOGY-BRIDGE}{【共享年序日期坐标】保留1346 年的多政权并行检索入口；本条不主张所列政权在当年发生直接互动，具体事件须另以单项材料入账。}\textbf{元与高丽、安南及地方政权（年序桥接）。}
元与高丽、安南及地方政权（年序桥接）在1346 年（共享年序桥接）的这条年序记录提供的是检索坐标：“【共享年序日期坐标】保留1346 年的多政权并行检索入口；本条不主张所列政权在当年发生直接互动，具体事件须另以单项材料入账。”。它保留材料稀薄处的空白，而不以一条王朝纪年替代地方、边疆或竞争政权的行动。

账本把这一节点放在“相邻已入账事件之间缺少该年度的共享年序入口，易使跨政权材料被单一政权叙事遮蔽。”与“保留该年度的检索与补账位置；后续仅在获得可核单项材料时拆分为具体政治、财政、军事、社会或文化事件。”之间；两者提示前后压力，并不把时间相邻写成单一因果。

本条首先回到《元史》；《高丽史》；《安南志略》及碑刻、地方文书（据事核）；本条只确认该年位于可追踪的共享年序中；正史、地方文书和外部材料的保存密度不一，不能由桥接标签推断行动、统属、疆界或社会影响。 因而元代区域治理、海域联通与地方社会的并行年序研究不是装饰性的书目，而是检验这一叙述边界的路径。

\subsection{1347 年（共享年序桥接）}

\noindent\eventanchor{SY-1347-CHRONOLOGY-BRIDGE}{【共享年序日期坐标】保留1347 年的多政权并行检索入口；本条不主张所列政权在当年发生直接互动，具体事件须另以单项材料入账。}\textbf{元与高丽、安南及地方政权（年序桥接）。}
1347 年（共享年序桥接）的元与高丽、安南及地方政权（年序桥接）材料只留下“【共享年序日期坐标】保留1347 年的多政权并行检索入口；本条不主张所列政权在当年发生直接互动，具体事件须另以单项材料入账。”这一纪年接点。它使同年的诏令、交涉或地方材料可以被放回并行时间轴，却不应被误写成一件已经具备完整情节的政治行动。

从账本的前后项看，“相邻已入账事件之间缺少该年度的共享年序入口，易使跨政权材料被单一政权叙事遮蔽。”提供了背景压力；“保留该年度的检索与补账位置；后续仅在获得可核单项材料时拆分为具体政治、财政、军事、社会或文化事件。”标出后来需要追踪的变化，二者之间仍存在未被材料填满的环节。

核心线索为《元史》；《高丽史》；《安南志略》及碑刻、地方文书（据事核）。本条只确认该年位于可追踪的共享年序中；正史、地方文书和外部材料的保存密度不一，不能由桥接标签推断行动、统属、疆界或社会影响。 现代研究可从元代区域治理、海域联通与地方社会的并行年序研究继续追问。

\subsection{1349 年（共享年序桥接）}

\noindent\eventanchor{SY-1349-CHRONOLOGY-BRIDGE}{【共享年序日期坐标】保留1349 年的多政权并行检索入口；本条不主张所列政权在当年发生直接互动，具体事件须另以单项材料入账。}\textbf{元与高丽、安南及地方政权（年序桥接）。}
1349 年（共享年序桥接）的元与高丽、安南及地方政权（年序桥接）材料只留下“【共享年序日期坐标】保留1349 年的多政权并行检索入口；本条不主张所列政权在当年发生直接互动，具体事件须另以单项材料入账。”这一纪年接点。它使同年的诏令、交涉或地方材料可以被放回并行时间轴，却不应被误写成一件已经具备完整情节的政治行动。

其前提记为“相邻已入账事件之间缺少该年度的共享年序入口，易使跨政权材料被单一政权叙事遮蔽。”；其后续是“保留该年度的检索与补账位置；后续仅在获得可核单项材料时拆分为具体政治、财政、军事、社会或文化事件。”。这是一条待后续材料检验的事件链，而不是对所有行动者意图的替写。

可核的材料入口是《元史》；《高丽史》；《安南志略》及碑刻、地方文书（据事核）。本条只确认该年位于可追踪的共享年序中；正史、地方文书和外部材料的保存密度不一，不能由桥接标签推断行动、统属、疆界或社会影响。 就此而言，元代区域治理、海域联通与地方社会的并行年序研究有助于把年序同区域、制度或物质材料并读。

\subsection{1350 年（共享年序桥接）}

\noindent\eventanchor{SY-1350-CHRONOLOGY-BRIDGE}{【共享年序日期坐标】保留1350 年的多政权并行检索入口；本条不主张所列政权在当年发生直接互动，具体事件须另以单项材料入账。}\textbf{元与高丽、安南及地方政权（年序桥接）。}
元与高丽、安南及地方政权（年序桥接）在1350 年（共享年序桥接）的这条年序记录提供的是检索坐标：“【共享年序日期坐标】保留1350 年的多政权并行检索入口；本条不主张所列政权在当年发生直接互动，具体事件须另以单项材料入账。”。它保留材料稀薄处的空白，而不以一条王朝纪年替代地方、边疆或竞争政权的行动。

从账本的前后项看，“相邻已入账事件之间缺少该年度的共享年序入口，易使跨政权材料被单一政权叙事遮蔽。”提供了背景压力；“保留该年度的检索与补账位置；后续仅在获得可核单项材料时拆分为具体政治、财政、军事、社会或文化事件。”标出后来需要追踪的变化，二者之间仍存在未被材料填满的环节。

核心线索为《元史》；《高丽史》；《安南志略》及碑刻、地方文书（据事核）。本条只确认该年位于可追踪的共享年序中；正史、地方文书和外部材料的保存密度不一，不能由桥接标签推断行动、统属、疆界或社会影响。 现代研究可从元代区域治理、海域联通与地方社会的并行年序研究继续追问。

\subsection{1352 年（共享年序桥接）}

\noindent\eventanchor{SY-1352-CHRONOLOGY-BRIDGE}{【共享年序日期坐标】保留1352 年的多政权并行检索入口；本条不主张所列政权在当年发生直接互动，具体事件须另以单项材料入账。}\textbf{元与高丽、安南及地方政权（年序桥接）。}
作为共享年序的接点，“【共享年序日期坐标】保留1352 年的多政权并行检索入口；本条不主张所列政权在当年发生直接互动，具体事件须另以单项材料入账。”让元与高丽、安南及地方政权（年序桥接）与同年其他地区保持可比的日期位置。可确认的是文书使用这一时间框架；年内何事发生、何处落实，仍须由各自条目另证。

其前提记为“相邻已入账事件之间缺少该年度的共享年序入口，易使跨政权材料被单一政权叙事遮蔽。”；其后续是“保留该年度的检索与补账位置；后续仅在获得可核单项材料时拆分为具体政治、财政、军事、社会或文化事件。”。这是一条待后续材料检验的事件链，而不是对所有行动者意图的替写。

可核的材料入口是《元史》；《高丽史》；《安南志略》及碑刻、地方文书（据事核）。本条只确认该年位于可追踪的共享年序中；正史、地方文书和外部材料的保存密度不一，不能由桥接标签推断行动、统属、疆界或社会影响。 就此而言，元代区域治理、海域联通与地方社会的并行年序研究有助于把年序同区域、制度或物质材料并读。

\subsection{1353 年（共享年序桥接）}

\noindent\eventanchor{SY-1353-CHRONOLOGY-BRIDGE}{【共享年序日期坐标】保留1353 年的多政权并行检索入口；本条不主张所列政权在当年发生直接互动，具体事件须另以单项材料入账。}\textbf{元与高丽、安南及地方政权（年序桥接）。}
1353 年（共享年序桥接）的元与高丽、安南及地方政权（年序桥接）材料只留下“【共享年序日期坐标】保留1353 年的多政权并行检索入口；本条不主张所列政权在当年发生直接互动，具体事件须另以单项材料入账。”这一纪年接点。它使同年的诏令、交涉或地方材料可以被放回并行时间轴，却不应被误写成一件已经具备完整情节的政治行动。

账本把这一节点放在“相邻已入账事件之间缺少该年度的共享年序入口，易使跨政权材料被单一政权叙事遮蔽。”与“保留该年度的检索与补账位置；后续仅在获得可核单项材料时拆分为具体政治、财政、军事、社会或文化事件。”之间；两者提示前后压力，并不把时间相邻写成单一因果。

本条首先回到《元史》；《高丽史》；《安南志略》及碑刻、地方文书（据事核）；本条只确认该年位于可追踪的共享年序中；正史、地方文书和外部材料的保存密度不一，不能由桥接标签推断行动、统属、疆界或社会影响。 因而元代区域治理、海域联通与地方社会的并行年序研究不是装饰性的书目，而是检验这一叙述边界的路径。

\subsection{1355 年（共享年序桥接）}

\noindent\eventanchor{SY-1355-CHRONOLOGY-BRIDGE}{【共享年序日期坐标】保留1355 年的多政权并行检索入口；本条不主张所列政权在当年发生直接互动，具体事件须另以单项材料入账。}\textbf{元与高丽、安南及地方政权（年序桥接）。}
元与高丽、安南及地方政权（年序桥接）在1355 年（共享年序桥接）的这条年序记录提供的是检索坐标：“【共享年序日期坐标】保留1355 年的多政权并行检索入口；本条不主张所列政权在当年发生直接互动，具体事件须另以单项材料入账。”。它保留材料稀薄处的空白，而不以一条王朝纪年替代地方、边疆或竞争政权的行动。

此前的条件是“相邻已入账事件之间缺少该年度的共享年序入口，易使跨政权材料被单一政权叙事遮蔽。”，此后的可见影响是“保留该年度的检索与补账位置；后续仅在获得可核单项材料时拆分为具体政治、财政、军事、社会或文化事件。”。这里的“影响”须按地点、机构和材料种类分别衡量。

这一叙述依据《元史》；《高丽史》；《安南志略》及碑刻、地方文书（据事核），但本条只确认该年位于可追踪的共享年序中；正史、地方文书和外部材料的保存密度不一，不能由桥接标签推断行动、统属、疆界或社会影响。。研究元代区域治理、海域联通与地方社会的并行年序研究可以补问材料未直接说出的空间与社会差异。

\subsection{1357 年（共享年序桥接）}

\noindent\eventanchor{SY-1357-CHRONOLOGY-BRIDGE}{【共享年序日期坐标】保留1357 年的多政权并行检索入口；本条不主张所列政权在当年发生直接互动，具体事件须另以单项材料入账。}\textbf{元与高丽、安南及地方政权（年序桥接）。}
作为共享年序的接点，“【共享年序日期坐标】保留1357 年的多政权并行检索入口；本条不主张所列政权在当年发生直接互动，具体事件须另以单项材料入账。”让元与高丽、安南及地方政权（年序桥接）与同年其他地区保持可比的日期位置。可确认的是文书使用这一时间框架；年内何事发生、何处落实，仍须由各自条目另证。

账本把这一节点放在“相邻已入账事件之间缺少该年度的共享年序入口，易使跨政权材料被单一政权叙事遮蔽。”与“保留该年度的检索与补账位置；后续仅在获得可核单项材料时拆分为具体政治、财政、军事、社会或文化事件。”之间；两者提示前后压力，并不把时间相邻写成单一因果。

本条首先回到《元史》；《高丽史》；《安南志略》及碑刻、地方文书（据事核）；本条只确认该年位于可追踪的共享年序中；正史、地方文书和外部材料的保存密度不一，不能由桥接标签推断行动、统属、疆界或社会影响。 因而元代区域治理、海域联通与地方社会的并行年序研究不是装饰性的书目，而是检验这一叙述边界的路径。

\subsection{1358 年（共享年序桥接）}

\noindent\eventanchor{SY-1358-CHRONOLOGY-BRIDGE}{【共享年序日期坐标】保留1358 年的多政权并行检索入口；本条不主张所列政权在当年发生直接互动，具体事件须另以单项材料入账。}\textbf{元与高丽、安南及地方政权（年序桥接）。}
1358 年（共享年序桥接）的元与高丽、安南及地方政权（年序桥接）材料只留下“【共享年序日期坐标】保留1358 年的多政权并行检索入口；本条不主张所列政权在当年发生直接互动，具体事件须另以单项材料入账。”这一纪年接点。它使同年的诏令、交涉或地方材料可以被放回并行时间轴，却不应被误写成一件已经具备完整情节的政治行动。

从账本的前后项看，“相邻已入账事件之间缺少该年度的共享年序入口，易使跨政权材料被单一政权叙事遮蔽。”提供了背景压力；“保留该年度的检索与补账位置；后续仅在获得可核单项材料时拆分为具体政治、财政、军事、社会或文化事件。”标出后来需要追踪的变化，二者之间仍存在未被材料填满的环节。

核心线索为《元史》；《高丽史》；《安南志略》及碑刻、地方文书（据事核）。本条只确认该年位于可追踪的共享年序中；正史、地方文书和外部材料的保存密度不一，不能由桥接标签推断行动、统属、疆界或社会影响。 现代研究可从元代区域治理、海域联通与地方社会的并行年序研究继续追问。

\subsection{1359 年（共享年序桥接）}

\noindent\eventanchor{SY-1359-CHRONOLOGY-BRIDGE}{【共享年序日期坐标】保留1359 年的多政权并行检索入口；本条不主张所列政权在当年发生直接互动，具体事件须另以单项材料入账。}\textbf{元与高丽、安南及地方政权（年序桥接）。}
作为共享年序的接点，“【共享年序日期坐标】保留1359 年的多政权并行检索入口；本条不主张所列政权在当年发生直接互动，具体事件须另以单项材料入账。”让元与高丽、安南及地方政权（年序桥接）与同年其他地区保持可比的日期位置。可确认的是文书使用这一时间框架；年内何事发生、何处落实，仍须由各自条目另证。

此前的条件是“相邻已入账事件之间缺少该年度的共享年序入口，易使跨政权材料被单一政权叙事遮蔽。”，此后的可见影响是“保留该年度的检索与补账位置；后续仅在获得可核单项材料时拆分为具体政治、财政、军事、社会或文化事件。”。这里的“影响”须按地点、机构和材料种类分别衡量。

这一叙述依据《元史》；《高丽史》；《安南志略》及碑刻、地方文书（据事核），但本条只确认该年位于可追踪的共享年序中；正史、地方文书和外部材料的保存密度不一，不能由桥接标签推断行动、统属、疆界或社会影响。。研究元代区域治理、海域联通与地方社会的并行年序研究可以补问材料未直接说出的空间与社会差异。

\subsection{1360 年（共享年序桥接）}

\noindent\eventanchor{SY-1360-CHRONOLOGY-BRIDGE}{【共享年序日期坐标】保留1360 年的多政权并行检索入口；本条不主张所列政权在当年发生直接互动，具体事件须另以单项材料入账。}\textbf{元与高丽、安南及地方政权（年序桥接）。}
1360 年（共享年序桥接）的元与高丽、安南及地方政权（年序桥接）材料只留下“【共享年序日期坐标】保留1360 年的多政权并行检索入口；本条不主张所列政权在当年发生直接互动，具体事件须另以单项材料入账。”这一纪年接点。它使同年的诏令、交涉或地方材料可以被放回并行时间轴，却不应被误写成一件已经具备完整情节的政治行动。

账本把这一节点放在“相邻已入账事件之间缺少该年度的共享年序入口，易使跨政权材料被单一政权叙事遮蔽。”与“保留该年度的检索与补账位置；后续仅在获得可核单项材料时拆分为具体政治、财政、军事、社会或文化事件。”之间；两者提示前后压力，并不把时间相邻写成单一因果。

本条首先回到《元史》；《高丽史》；《安南志略》及碑刻、地方文书（据事核）；本条只确认该年位于可追踪的共享年序中；正史、地方文书和外部材料的保存密度不一，不能由桥接标签推断行动、统属、疆界或社会影响。 因而元代区域治理、海域联通与地方社会的并行年序研究不是装饰性的书目，而是检验这一叙述边界的路径。

\subsection{1361 年（共享年序桥接）}

\noindent\eventanchor{SY-1361-CHRONOLOGY-BRIDGE}{【共享年序日期坐标】保留1361 年的多政权并行检索入口；本条不主张所列政权在当年发生直接互动，具体事件须另以单项材料入账。}\textbf{元与高丽、安南及地方政权（年序桥接）。}
作为共享年序的接点，“【共享年序日期坐标】保留1361 年的多政权并行检索入口；本条不主张所列政权在当年发生直接互动，具体事件须另以单项材料入账。”让元与高丽、安南及地方政权（年序桥接）与同年其他地区保持可比的日期位置。可确认的是文书使用这一时间框架；年内何事发生、何处落实，仍须由各自条目另证。

从账本的前后项看，“相邻已入账事件之间缺少该年度的共享年序入口，易使跨政权材料被单一政权叙事遮蔽。”提供了背景压力；“保留该年度的检索与补账位置；后续仅在获得可核单项材料时拆分为具体政治、财政、军事、社会或文化事件。”标出后来需要追踪的变化，二者之间仍存在未被材料填满的环节。

核心线索为《元史》；《高丽史》；《安南志略》及碑刻、地方文书（据事核）。本条只确认该年位于可追踪的共享年序中；正史、地方文书和外部材料的保存密度不一，不能由桥接标签推断行动、统属、疆界或社会影响。 现代研究可从元代区域治理、海域联通与地方社会的并行年序研究继续追问。

\subsection{1362 年（共享年序桥接）}

\noindent\eventanchor{SY-1362-CHRONOLOGY-BRIDGE}{【共享年序日期坐标】保留1362 年的多政权并行检索入口；本条不主张所列政权在当年发生直接互动，具体事件须另以单项材料入账。}\textbf{元与高丽、安南及地方政权（年序桥接）。}
元与高丽、安南及地方政权（年序桥接）在1362 年（共享年序桥接）的这条年序记录提供的是检索坐标：“【共享年序日期坐标】保留1362 年的多政权并行检索入口；本条不主张所列政权在当年发生直接互动，具体事件须另以单项材料入账。”。它保留材料稀薄处的空白，而不以一条王朝纪年替代地方、边疆或竞争政权的行动。

此前的条件是“相邻已入账事件之间缺少该年度的共享年序入口，易使跨政权材料被单一政权叙事遮蔽。”，此后的可见影响是“保留该年度的检索与补账位置；后续仅在获得可核单项材料时拆分为具体政治、财政、军事、社会或文化事件。”。这里的“影响”须按地点、机构和材料种类分别衡量。

这一叙述依据《元史》；《高丽史》；《安南志略》及碑刻、地方文书（据事核），但本条只确认该年位于可追踪的共享年序中；正史、地方文书和外部材料的保存密度不一，不能由桥接标签推断行动、统属、疆界或社会影响。。研究元代区域治理、海域联通与地方社会的并行年序研究可以补问材料未直接说出的空间与社会差异。

\subsection{1364 年（共享年序桥接）}

\noindent\eventanchor{SY-1364-CHRONOLOGY-BRIDGE}{【共享年序日期坐标】保留1364 年的多政权并行检索入口；本条不主张所列政权在当年发生直接互动，具体事件须另以单项材料入账。}\textbf{元与高丽、安南及地方政权（年序桥接）。}
元与高丽、安南及地方政权（年序桥接）在1364 年（共享年序桥接）的这条年序记录提供的是检索坐标：“【共享年序日期坐标】保留1364 年的多政权并行检索入口；本条不主张所列政权在当年发生直接互动，具体事件须另以单项材料入账。”。它保留材料稀薄处的空白，而不以一条王朝纪年替代地方、边疆或竞争政权的行动。

账本把这一节点放在“相邻已入账事件之间缺少该年度的共享年序入口，易使跨政权材料被单一政权叙事遮蔽。”与“保留该年度的检索与补账位置；后续仅在获得可核单项材料时拆分为具体政治、财政、军事、社会或文化事件。”之间；两者提示前后压力，并不把时间相邻写成单一因果。

本条首先回到《元史》；《高丽史》；《安南志略》及碑刻、地方文书（据事核）；本条只确认该年位于可追踪的共享年序中；正史、地方文书和外部材料的保存密度不一，不能由桥接标签推断行动、统属、疆界或社会影响。 因而元代区域治理、海域联通与地方社会的并行年序研究不是装饰性的书目，而是检验这一叙述边界的路径。

\subsection{1365 年（共享年序桥接）}

\noindent\eventanchor{SY-1365-CHRONOLOGY-BRIDGE}{【共享年序日期坐标】保留1365 年的多政权并行检索入口；本条不主张所列政权在当年发生直接互动，具体事件须另以单项材料入账。}\textbf{元与高丽、安南及地方政权（年序桥接）。}
1365 年（共享年序桥接）的元与高丽、安南及地方政权（年序桥接）材料只留下“【共享年序日期坐标】保留1365 年的多政权并行检索入口；本条不主张所列政权在当年发生直接互动，具体事件须另以单项材料入账。”这一纪年接点。它使同年的诏令、交涉或地方材料可以被放回并行时间轴，却不应被误写成一件已经具备完整情节的政治行动。

从账本的前后项看，“相邻已入账事件之间缺少该年度的共享年序入口，易使跨政权材料被单一政权叙事遮蔽。”提供了背景压力；“保留该年度的检索与补账位置；后续仅在获得可核单项材料时拆分为具体政治、财政、军事、社会或文化事件。”标出后来需要追踪的变化，二者之间仍存在未被材料填满的环节。

核心线索为《元史》；《高丽史》；《安南志略》及碑刻、地方文书（据事核）。本条只确认该年位于可追踪的共享年序中；正史、地方文书和外部材料的保存密度不一，不能由桥接标签推断行动、统属、疆界或社会影响。 现代研究可从元代区域治理、海域联通与地方社会的并行年序研究继续追问。

\subsection{1366 年（共享年序桥接）}

\noindent\eventanchor{SY-1366-CHRONOLOGY-BRIDGE}{【共享年序日期坐标】保留1366 年的多政权并行检索入口；本条不主张所列政权在当年发生直接互动，具体事件须另以单项材料入账。}\textbf{元与高丽、安南及地方政权（年序桥接）。}
作为共享年序的接点，“【共享年序日期坐标】保留1366 年的多政权并行检索入口；本条不主张所列政权在当年发生直接互动，具体事件须另以单项材料入账。”让元与高丽、安南及地方政权（年序桥接）与同年其他地区保持可比的日期位置。可确认的是文书使用这一时间框架；年内何事发生、何处落实，仍须由各自条目另证。

此前的条件是“相邻已入账事件之间缺少该年度的共享年序入口，易使跨政权材料被单一政权叙事遮蔽。”，此后的可见影响是“保留该年度的检索与补账位置；后续仅在获得可核单项材料时拆分为具体政治、财政、军事、社会或文化事件。”。这里的“影响”须按地点、机构和材料种类分别衡量。

这一叙述依据《元史》；《高丽史》；《安南志略》及碑刻、地方文书（据事核），但本条只确认该年位于可追踪的共享年序中；正史、地方文书和外部材料的保存密度不一，不能由桥接标签推断行动、统属、疆界或社会影响。。研究元代区域治理、海域联通与地方社会的并行年序研究可以补问材料未直接说出的空间与社会差异。
